% Complete cumulative LaTeX source for the named public edition.
% Build from releases/ inside the extracted source bundle; dependencies are under ../upstream/.
% XeLaTeX and the bundled guard workflow are documented in README.md.
\documentclass[12pt,a4paper,oneside]{report}
\usepackage[margin=22mm,headheight=15pt]{geometry}
\usepackage{fontspec}
\setmainfont[Script=Arabic]{Amiri}
\newfontfamily\latinfont{Latin Modern Roman}
\usepackage{amsmath,amssymb,amsfonts,amsthm,mathrsfs,nicefrac,stmaryrd}
\usepackage{xparse,etoolbox,xpunctuate,mfirstuc}
\usepackage{graphicx,tikz,xcolor,framed,comment,longtable,multirow,array}
\usepackage{bussproofs}
\input{../upstream/sty/bussproofs-extra.sty}
\usepackage[tableaux]{prooftrees}
\usepackage[unicode,colorlinks,linkcolor=blue,urlcolor=blue]{hyperref}
\usepackage{fancyhdr}
\usepackage{bidi}

\hypersetup{
  pdftitle={OpenLogic Pashto Pakistan - Propositional Syntax, Semantics, and Proof Systems},
  pdfauthor={Open Logic Project; independent Pashto (Pakistan) translation},
  pdfsubject={Partial edition: 70 source units; full 722-unit edition in progress}
}

% Load the upstream notation profile without its English token layer.  The
% builder expands every source token to the reviewed Pashto form before TeX.
\newcommand*{\DeclareDocumentMacro}[2]{\def#1{#2}}
\NewDocumentCommand{\settexttoken}{m s m m o o}{}
\makeatletter
\input{../upstream/sty/open-logic-formulas.sty}
\input{../upstream/sty/open-logic-selective.sty}
\input{../upstream/open-logic-config.sty}
\makeatother
\renewcommand{\MP}{\LR{\latinfont MP}}
\renewcommand{\QR}{\LR{\latinfont QR}}
\renewcommand{\Hyp}{\LR{\latinfont Hyp}}

% Source-relative file IDs and references are retained in the flattened
% reader.  A reference outside this partial volume prints its exact OpenLogic
% identifier instead of an unresolved question mark.
\newcommand{\olfileid}[3]{\gdef\pspart{#1}\gdef\pschapter{#2}\gdef\pssection{#3}}
\NewDocumentCommand{\olchapter}{o m m m}{%
  \gdef\pspart{#2}\gdef\pschapter{#3}\gdef\pssection{}%
  \IfNoValueTF{#1}{\chapter{#4}}{\chapter[#1]{#4}}%
  \label{#2:#3::chap}}
\NewDocumentCommand{\olsection}{o m}{%
  \IfNoValueTF{#1}{\section{#2}}{\section[#1]{#2}}%
  \label{\pspart:\pschapter:\pssection:sec}}
\newcommand{\ollabel}[1]{\label{\pspart:\pschapter:\pssection:#1}}
\NewDocumentCommand{\olref}{o o o m}{\begingroup
  \IfNoValueTF{#1}{\edef\pskey{\pspart:\pschapter:\pssection:#4}}{%
    \IfNoValueTF{#2}{\edef\pskey{\pspart:\pschapter:#1:#4}}{%
      \IfNoValueTF{#3}{\edef\pskey{\pspart:#1:#2:#4}}{\edef\pskey{#1:#2:#3:#4}}}}%
  \ifcsname r@\pskey\endcsname
    \LR{\ref{\pskey}}%
  \else
    \mbox{\LR{\latinfont\scriptsize OpenLogic \texttt{\pskey}}}%
  \fi
  \endgroup}
\let\Olref\olref
\NewDocumentCommand{\oliflabeldef}{m m m}{\ifcsname r@#1\endcsname #2\else #3\fi}
\providecommand{\OLEndChapterHook}{}

% Proof layouts are geometrically left-to-right even though the surrounding
% explanatory prose is right-to-left.
\renewenvironment{prooftree}
  {\begin{center}\begin{LTR}\bottomAlignProof}
  {\DisplayProof\end{LTR}\end{center}}
\newenvironment{oltableau}
  {\center\LTR\tableau{}}
  {\endtableau\endLTR\endcenter}

% The source environments remain visible in this review reader, including
% editorial notices that disclose source limitations and recorded repairs.
\newtheorem{thm}{قضيه}[chapter]
\newtheorem{lem}[thm]{لم}
\newtheorem{prop}[thm]{قضيه}
\newtheorem{cor}[thm]{نتيجه}
\theoremstyle{definition}
\newtheorem{defn}[thm]{تعريف}
\newtheorem{ex}[thm]{بېلګه}
\newtheorem{prob}[thm]{تمرين}
\newenvironment{defish}{\begin{framed}}{\end{framed}}
\newenvironment{editorial}{\begin{framed}\textbf{ادارتي يادونه۔}\enspace}{\end{framed}}
\newenvironment{explain}{}{}
\newenvironment{intro}{}{}
\newenvironment{digress}{\par\medskip}{\par\medskip}
\newenvironment{derivation}{%
  \par\smallskip\begin{LTR}\begin{tabular}{@{}rll@{}}}
  {\end{tabular}\end{LTR}\par\smallskip}

\renewcommand{\proofname}{ثبوت}
\renewcommand{\figurename}{شکل}
\renewcommand{\chaptername}{باب}
\renewcommand{\contentsname}{فهرست}
\renewcommand{\bibname}{\RL{\normalfont سرچينې}}
\makeatletter
\renewcommand{\@pnumwidth}{2.2em}
\renewcommand{\@tocrmarg}{4em}
\makeatother
\setlength{\parindent}{0pt}
\setlength{\parskip}{5.5pt}
\setlength{\emergencystretch}{6em}
% \sFmla intentionally lets a truth-sign subscript overhang a narrower
% formula nucleus by at most 3.3pt in three inline metavariable examples.
\hfuzz=4pt
\linespread{1.15}
\pagestyle{fancy}
\fancyhf{}
\fancyhead[RE,LO]{\small خلاص منطق — پښتو (پاکستان)}
\fancyhead[LE,RO]{\small بياني نحو، معناپوهنه او د ثبوت نظامونه}
\fancyfoot[C]{\thepage}

\begin{document}
\tracinglostchars=3
\setRTL
\begin{titlepage}
\thispagestyle{empty}
\begin{center}
\vspace*{24mm}
{\Huge خلاص منطق}\par\bigskip
{\Large پښتو — پاکستان}\par\bigskip
{\LARGE بياني نحو، معناپوهنه او د ثبوت نظامونه}\par\bigskip
د ماشيني ژباړې د روان کار څلورمه لوستيزه نسخه
\end{center}
\bigskip
په دې لوستيزه نسخه کښې د بياني منطق د نحو او معناپوهنې بشپړ باب،
د ثبوتي نظامونو ټولکتنه، او د لومړۍ درجې منطق د سېکوېنټ حساب، فطري
استنتاج، تابلو او بديهي اشتقاق بشپړ بابونه شامل دي۔ دا ټول
70 سرچينيز واحدونه دي؛ د ۷۲۲ واحدونو بشپړه ژباړه لا
روانه ده او تر اوسه 131 واحدونه په ژباړل شوي بنډل کښې
شته۔

په متن کښې چې کوم بهرني OpenLogic پېژندونه ښکاري، هغه د لومړۍ درجې
منطق هغو مخکينيو نحوي يا معنايي برخو ته کره اشارې دي چې لا په دې
جزوي لوستيزه نسخه کښې نۀ دي راغلي۔ دا پېژندونه ماتې حوالې نۀ دي او
په بشپړه نسخه کښې به شمېرل شويو حوالو ته واوړي۔

\begin{LTR}\latinfont\small
Independent Pashto (Pakistan) translation of Open Logic Project revision
\texttt{9620cc73\allowbreak f9c8e0ad\allowbreak 003c514a\allowbreak
5d3748f2\allowbreak 9611c4c0}. This reader contains
70 source units (OLP-0056 through OLP-0125). Pakistani usage is
primary; Afghan Pashto sources are labelled regional comparators. Specialist
terminology remains openly reviewable where exact native technical attestation
is sparse. No native-reader review or upstream endorsement is claimed.

Original text and adaptation: Creative Commons Attribution 4.0 International,
subject to the component notices preserved with the source snapshot.
\href{https://openlogicproject.org/}{Open Logic Project}.\quad
\href{https://github.com/KokunoYumeto/OpenLogic-translations}{Translation catalogue}.\quad
\href{https://doi.org/10.5281/zenodo.22307197}{Edition lineage}.
\end{LTR}
\end{titlepage}
\tableofcontents
\clearpage


\phantomsection\label{olpseg:OLP-0056-B004:start}
\olchapter{pl}{syn}{نحو او معنٰی پوهنه}
\label{olpseg:OLP-0056-B004:end}

\phantomsection\label{olpseg:OLP-0056-B005:start}
\begin{editorial}
  دا يوازې د تعريفونو ډېره لنډه لنډيزه ده۔ بايد پراخه شي، څو د
  فارمولونو پر جوړښت د استقرا د ثبوتونو اسانه پېژندنه او ډېرې نورې
  بېلګې پکې راشي۔
\end{editorial}
\label{olpseg:OLP-0056-B005:end}

\phantomsection\label{olpseg:OLP-0057-B004:start}
\olfileid{pl}{syn}{int}
\label{olpseg:OLP-0057-B004:end}

\phantomsection\label{olpseg:OLP-0057-B005:start}
\olsection{پېژندنه}
\label{olpseg:OLP-0057-B005:end}

\phantomsection\label{olpseg:OLP-0057-B006:start}
بياني منطق له هغو فارمولونه سره کار لري چې له
بياني متغيرونه څخه د \LR{$\lnot$}، \LR{$\land$}، \LR{$\lor$}، \LR{$\lif$} او
\LR{$\liff$} بياني نښلوونکو په کارولو جوړېږي۔ د تصور له مخې،
بياني متغير~\LR{$p$} د داسې يوې جملې يا بيان ځاے نيسي
چې رښتونے يا دروغژن وي۔ هر کله چې په فارمول کښې د
بياني متغير ``د رښتياوالي قيمت'' وټاکل شي، نو د
بياني نښلوونکو په کارولو ترې د جوړو شويو ټولو فارمولونه د
رښتياوالي قيمت هم ټاکل کېږي۔ وايو چې بياني منطق \emph{د رښتياوالي
  تابعوي} دے، ځکه چې معنٰی پوهنه ئې د رښتياوالي د قيمتونو د تابعو
له مخې ورکول کېږي۔ په تېره بيا، په بياني منطق کښې دا نۀ څېړو چې
رښتياوالے او دروغوالے نور څنګه ټاکل کېږي؛ مثلاً دا چې يو څه د يوازې
اتفاقي رښتياوالي پر ځاے بالضروره رښتيا دي، يا څوک پرې پوهېږي چې
رښتيا دي، يا اوس رښتيا دي نۀ دا چې پخوا رښتيا وو يا وروسته به رښتيا
وي۔ موږ يوازې د رښتيا (\LR{$\True$}) او دروغو (\LR{$\False$}) دوه قيمتونه په
پام کښې نيسو، نو دا امکان له بحثه وباسو چې يو بيان دې نۀ رښتونے وي
نۀ دروغژن، يا دې نيم رښتونے وي۔ موږ يوازې هغو نښلوونکو ته هم پام
کوو چې ترې د جوړ شوي فارمول د رښتياوالي قيمت د هغې د برخو د
رښتياوالي په قيمتونو پوره ټاکل کېږي، نۀ مثلاً د هغې په معنٰی۔ په
تېره بيا، دا لانجمنه ده چې په انګرېزۍ کښې د شرطيو رښتياوالي قيمت په
دې معنٰی تابعوي دے که نۀ۔ مادي شرطي~\LR{$\lif$} تابعوي دے؛ نور منطقونه
له هغو شرطيو سره کار کوي چې د رښتياوالي تابعوي نۀ وي۔
\label{olpseg:OLP-0057-B006:end}

\phantomsection\label{olpseg:OLP-0057-B007:start}
د رښتياوالي د تابعوي بياني منطق د نظريې او مابعدمنطق د ودې دپاره
بايد لومړے د هغه د عبارتونو نحو او معنٰی پوهنه تعريف کړو۔ موږ به له
بياني متغيرونه څخه د نښلوونکو په کارولو د فارمولونه
د جوړولو يوه لار بيان کړو۔ نور تعريفونه هم ممکن دي۔ نور نظامونه به
بېلې نښې غوره کړي، د بنسټيزو نښلوونکو بېل سټونه به وټاکي او قوسونه
به په بېله توګه وکاروي، يا به ئې هېڅ و نۀ کاروي، لکه په تش په نامه
پولنډۍ نښه‌اېښودنه کښې۔ خو په ټولو لارو کښې يو ګډ شے دا دے چې د
جوړېدو قاعدې د فارمولونه سټ \emph{په استقرايي ډول} تعريفوي۔ که
کار په سمه توګه وشي، هر عبارت د جوړېدو د قاعدو له مخې په اصل کښې
يوازې په يوه لار ترلاسه کېدے شي۔ هغه استقرايي تعريف چې
\emph{په يکتا ډول لوستل کېدونکي} عبارتونه ورکوي، موږ ته اجازه راکوي
چې په هماغه طريقه، يعنې په استقرايي تعريف، دغو عبارتونو ته معنٰی
ورکړو۔
\label{olpseg:OLP-0057-B007:end}

\phantomsection\label{olpseg:OLP-0057-B008:start}
عبارتونو ته معنٰی ورکول د معنٰی پوهنې ډګر دے۔ د بياني منطق په
معنٰی پوهنه کښې مرکزي مفهوم په ارزښت ټاکنه کښې صدق دے۔
ارزښت ټاکنه~\LR{$\pAssign{v}$} بياني متغيرونه ته د
رښتياوالي \LR{$\True$} او \LR{$\False$} قيمتونه ورکوي۔ هر ارزښت ټاکنه د هر
فارمول~\LR{$!A$} دپاره د رښتياوالي قيمت \LR{$\pValue{v}(!A)$} ټاکي۔
فارمول په ارزښت ټاکنه~\LR{$\pAssign{v}$} کښې هله او يوازې هله
صادق دے چې \LR{$\pValue{v}(!A) = \True$} وي؛ دا په \LR{$\pSat{v}{!A}$} ليکو۔
دا اړيکه د~\LR{$!A$} پر جوړښت د استقرا له مخې هم تعريفېدے شي: د منطقي
نښلوونکو د رښتياوالي تابعې کاروو، څو مثلاً د \LR{$!A \land !B$} صدق د
\LR{$!A$} او~\LR{$!B$} د صدق يا نۀ صدق له مخې تعريف کړو۔
\label{olpseg:OLP-0057-B008:end}

\phantomsection\label{olpseg:OLP-0057-B009:start}
د جملو د صدق د اړيکې \LR{$\pSat{v}{!A}$} پر بنسټ بيا د تاوتولوژۍ،
استلزام او د صدق وړتيا بنسټيز معنايي مفهومونه تعريفولے شو۔
فارمول هله تاوتولوژي، \LR{$\Entails !A$}، ده چې هر ارزښت ټاکنه
ئې صادق کړي، يعنې د هر~\LR{$\pAssign{v}$} دپاره
\LR{$\pValue{v}(!A) = \True$} وي۔ د فارمولونه يو سټ هله هغې ته
استلزام ورکوي، \LR{$\Gamma \Entails !A$}، چې هر هغه ارزښت ټاکنه چې د
~\LR{$\Gamma$} ټول فارمولونه صادق کړي، \LR{$!A$} هم صادق کړي۔ او د
فارمولونه يو سټ هله د صدق وړ دے چې کوم ارزښت ټاکنه ئې ټول
فارمولونه په يو وخت صادق کړي۔ ځکه چې فارمولونه په استقرايي
ډول تعريف شوي او صدق بيا د فارمولونه پر جوړښت د استقرا له مخې
تعريف شوے، موږ استقرا کارولے شو چې د خپلې معنٰی پوهنې خاصيتونه ثابت
او تعريف شوي معنايي مفهومونه يو له بل سره وتړو۔
\label{olpseg:OLP-0057-B009:end}

\phantomsection\label{olpseg:OLP-0058-B004:start}
\olfileid{pl}{syn}{fml}
\label{olpseg:OLP-0058-B004:end}

\phantomsection\label{olpseg:OLP-0058-B005:start}
\olsection{بياني فارمولونه}
\label{olpseg:OLP-0058-B005:end}

\phantomsection\label{olpseg:OLP-0058-B006:start}
د بياني منطق فارمولونه له \emph{بياني متغيرونه}، د دروغوالي له
  بياني ثابت~\LR{$\lfalse$}
    او د رښتياوالي له بياني ثابت~\LR{$\ltrue$} څخه د \emph{منطقي
  نښلوونکو} په کارولو جوړېږي۔
\label{olpseg:OLP-0058-B006:end}

\phantomsection\label{olpseg:OLP-0058-B007:start}
\begin{enumerate}
\item د بياني متغيرونه \LR{$\Obj p_0$}،
  \LR{$\Obj p_1$}، \dots يو شمېرېدونکے لامتناهي سټ~\LR{$\PVar$}
\item د دروغوالی بياني ثابت~\LR{$\lfalse$}۔
\item د رښتياوالی بياني ثابت~\LR{$\ltrue$}۔
\item منطقي نښلوونکي:
  \startycommalist
  \ycomma \LR{$\lnot$} (نفي)%
  \ycomma \LR{$\land$} (عطف)%
  \ycomma \LR{$\lor$} (فصل)%
  \ycomma \LR{$\lif$} (شرطي)%
  \ycomma \LR{$\liff$} (دوه اړخيز شرطي)%
\item د وقف نښې: (، ) او کامه۔
\end{enumerate}
\label{olpseg:OLP-0058-B007:end}

\phantomsection\label{olpseg:OLP-0058-B008:start}
د بياني منطق دا ژبه په \LR{$\Lang L_0$} ښيو۔
\label{olpseg:OLP-0058-B008:end}

\phantomsection\label{olpseg:OLP-0058-B009:start}
۔
\label{olpseg:OLP-0058-B009:end}



\phantomsection\label{olpseg:OLP-0058-B011:start}
% Alternate symbols
\label{olpseg:OLP-0058-B011:end}

\phantomsection\label{olpseg:OLP-0058-B012:start}
\begin{intro}
کېدے شي تاسو له هغو اصطلاحاتو او نښو سره بلد يئ چې زموږ له پورته
کارولو سره توپير لري۔ د منطق متنونه، او ښوونکي، عموماً د ``نفي''
دپاره \LR{$\sim$}، \LR{$\neg$} او~!، او د ``عطف'' دپاره \LR{$\wedge$}، \LR{$\cdot$} او
\LR{$\&$} کاروي۔ د ``شرطي'' يا ``استلزام'' عامې نښې \LR{$\rightarrow$}،
\LR{$\Rightarrow$} او \LR{$\supset$} دي۔
د ``دوه اړخيز شرطي''، ``دوه اړخيز استلزام''
  يا ``(مادي) هم ارزښت'' نښې \LR{$\leftrightarrow$}،
  \LR{$\Leftrightarrow$} او \LR{$\equiv$} دي۔
د \LR{$\lfalse$} نښې ته په بېلو ځايونو کښې
  ``دروغوالے''، ``فالسوم''، ``محال'' يا ``ښکته'' ويل کېږي۔
د \LR{$\ltrue$} نښې ته په بېلو ځايونو کښې
  ``رښتياوالے''، ``ويروم'' يا ``پورته'' ويل کېږي۔
\end{intro}
\label{olpseg:OLP-0058-B012:end}

\phantomsection\label{olpseg:OLP-0058-B013:start}
\begin{defn}[فارمول]
\ollabel{defn:formulas}
د بياني منطق د \emph{فارمولونه} سټ~\LR{$\Frm[L_0]$} په استقرايي ډول
داسې تعريفېږي:
\begin{enumerate}
\item \LR{$\lfalse$} يو اټومي فارمول دے۔
\label{olpseg:OLP-0058-B013:end}

\phantomsection\label{olpseg:OLP-0058-B014:start}
\item \LR{$\ltrue$} يو اټومي فارمول دے۔
\label{olpseg:OLP-0058-B014:end}

\phantomsection\label{olpseg:OLP-0058-B015:start}
\item هر بياني متغير~\LR{$\Obj p_i$} يو اټومي
  فارمول دے۔
\label{olpseg:OLP-0058-B015:end}

\phantomsection\label{olpseg:OLP-0058-B016:start}
\item که \LR{$!A$} يو فارمول وي، نو \LR{$\lnot !A$} هم
  يو فارمول دے۔
\label{olpseg:OLP-0058-B016:end}

\phantomsection\label{olpseg:OLP-0058-B017:start}
\item که \LR{$!A$} او \LR{$!B$} فارمولونه وي، نو \LR{$(!A \land
  !B)$} يو فارمول دے۔
\label{olpseg:OLP-0058-B017:end}

\phantomsection\label{olpseg:OLP-0058-B018:start}
\item که \LR{$!A$} او \LR{$!B$} فارمولونه وي، نو \LR{$(!A \lor !B)$}
  يو فارمول دے۔
\label{olpseg:OLP-0058-B018:end}

\phantomsection\label{olpseg:OLP-0058-B019:start}
\item که \LR{$!A$} او \LR{$!B$} فارمولونه وي، نو \LR{$(!A \lif !B)$}
  يو فارمول دے۔
\label{olpseg:OLP-0058-B019:end}

\phantomsection\label{olpseg:OLP-0058-B020:start}
\item که \LR{$!A$} او \LR{$!B$} فارمولونه وي، نو \LR{$(!A \liff !B)$}
  يو فارمول دے۔
\label{olpseg:OLP-0058-B020:end}

\phantomsection\label{olpseg:OLP-0058-B021:start}
\item بل هېڅ شے فارمول نۀ دے۔
\end{enumerate}
\end{defn}
\label{olpseg:OLP-0058-B021:end}

\phantomsection\label{olpseg:OLP-0058-B022:start}
\begin{explain}
د فارمولونه تعريف يو \emph{استقرايي تعريف} دے۔ په اصل کښې موږ
د فارمولونه سټ په نامتناهي ډېرو پړاوونو کښې جوړوو۔ په لومړني
پړاو کښې ټول اټومي فارمولونه فارمولونه ګڼو؛ دا د تعريف له لومړنيو
څو حالتونو سره سمون لري، يعنې د
\LR{$\ltrue$}، %
\LR{$\lfalse$}، %
\LR{$\Obj p_i$} حالتونه۔ په دې توګه ``اټومي فارمول'' د همدې بڼې هر
فارمول ته وايي۔
\label{olpseg:OLP-0058-B022:end}

\phantomsection\label{olpseg:OLP-0058-B023:start}
د تعريف نور حالتونه له هغو فارمولونه څخه د نويو فارمولونه د
جوړولو قاعدې راکوي چې له مخکښې جوړ شوي وي۔ په دويم پړاو کښې له دغو
قاعدو څخه په کار اخيستو له اټومي فارمولونه څخه فارمولونه
جوړولے شو۔ په درېيم پړاو کښې له اټومي فارمولونو او په دويم پړاو کښې
له ترلاسه شويو فارمولونو څخه نوي فارمولونه جوړوو، او همداسې مخکښې
ځو۔ فارمول هر هغه شے دے چې بالاخره په داسې کوم پړاو کښې جوړ
شي، او بل هېڅ شے نۀ دے۔
\end{explain}
\label{olpseg:OLP-0058-B023:end}

\phantomsection\label{olpseg:OLP-0058-B024:start}
کله چې د دوه ځايه نښلوونکي~\LR{$\ast$} په کارولو له \LR{$!B$} او \LR{$!C$} څخه
جوړ شوے فارمول \LR{$(!B \ast !C)$} ليکو، ډېر ځله د قوسونو تر ټولو بهرنی
جوړ پرېږدو او يوازې~\LR{$!B \ast !C$} ليکو۔
\label{olpseg:OLP-0058-B024:end}

\phantomsection\label{olpseg:OLP-0058-B034:start}
\begin{defn}[نحوي يوشانوالی]
د \LR{$\ident$} نښه د نښو د توريزو لړيو ترمنځ نحوي يوشانوالی څرګندوي؛
يعنې \LR{$!A \ident !B$} هله او يوازې هله چې \LR{$!A$} او \LR{$!B$} د نښو د يو
شان اوږدوالي توريزې لړۍ وي او په هر ځاے کښې يوه نښه ولري۔
\end{defn}
\label{olpseg:OLP-0058-B034:end}

\phantomsection\label{olpseg:OLP-0058-B035:start}
د \LR{$\ident$} نښې دواړو خواوو ته د يو ځاے کولو له لارې ترلاسه شوې
توريزې لړۍ هم راتلے شي۔ مثلاً \LR{$!A \ident (!B \lor !C)$} دا معنٰی لري:
د نښو توريزه لړۍ~\LR{$!A$} هماغه توريزه لړۍ ده چې په همدې ترتيب د پرانيستي
قوس، د \LR{$!B$} توريزې لړۍ، د \LR{$\lor$} نښې، د~\LR{$!C$} توريزې لړۍ او د تړلي
قوس په يو ځاے کولو ترلاسه کېږي۔ که داسې وي، نو پوهېږو چې د \LR{$!A$}
لومړۍ نښه پرانيستے قوس دے، \LR{$!A$} د \LR{$!B$} توريزه لړۍ د فرعي توريزې
لړۍ په توګه لري چې له دويمې نښې پيلېږي، ورپسې \LR{$\lor$} راځي، او داسې نور۔
\label{olpseg:OLP-0058-B035:end}

\phantomsection\label{olpseg:OLP-0059-B004:start}
\olfileid{pl}{syn}{pre}
\label{olpseg:OLP-0059-B004:end}

\phantomsection\label{olpseg:OLP-0059-B005:start}
\olsection{لومړنۍ پايلې}
\label{olpseg:OLP-0059-B005:end}

\phantomsection\label{olpseg:OLP-0059-B006:start}
\begin{thm}[\emph{د فارمولونه دپاره د استقرا اصل}]
  \ollabel{thm:induction}  
  که کوم خاصيت~\LR{$P$} د ټولو اټومي فارمولونه دپاره صدق کوي او داسې
  وي چې
  \begin{enumerate}
    \item هر کله چې د~\LR{$!A$} دپاره صدق کوي، د \LR{$\lnot !A$}
      دپاره هم صدق کوي؛
    \item هر کله چې د \LR{$!A$} او~\LR{$!B$} دپاره صدق کوي، د
      \LR{$(!A \land !B)$} دپاره هم صدق کوي؛
    \item هر کله چې د \LR{$!A$} او~\LR{$!B$} دپاره صدق کوي، د
      \LR{$(!A \lor !B)$} دپاره هم صدق کوي؛
    \item هر کله چې د \LR{$!A$} او~\LR{$!B$} دپاره صدق کوي، د
      \LR{$(!A \lif !B)$} دپاره هم صدق کوي؛
    \item هر کله چې د \LR{$!A$} او~\LR{$!B$} دپاره صدق کوي، د
      \LR{$(!A \liff !B)$} دپاره هم صدق کوي؛
  \end{enumerate}
  نو \LR{$P$} د ټولو فارمولونه دپاره صدق کوي۔
\end{thm}
\label{olpseg:OLP-0059-B006:end}

\phantomsection\label{olpseg:OLP-0059-B007:start}
\begin{proof}
  \LR{$S$} دې د هغو ټولو فارمولونه مجموعه وي چې خاصيت~\LR{$P$} لري۔ په
  څرګنده \LR{$S \subseteq \Frm[L_0]$}۔ \LR{$S$} د
  \olref[fml]{defn:formulas} ټول شرطونه پوره کوي: ټول اټومي
  فارمولونه لري او د منطقي عملګرونه له مخې بند دے۔ \LR{$\Frm[L_0]$} د
  دې ډول تر ټولو وړوکے ټولګے دے، نو \LR{$\Frm[L_0] \subseteq S$}۔ له دې
  امله \LR{$\Frm[L_0] = S$} او هر فارمول خاصيت~\LR{$P$} لري۔
\end{proof}
\label{olpseg:OLP-0059-B007:end}

\phantomsection\label{olpseg:OLP-0059-B008:start}
\begin{prop}\ollabel{prop:balanced}
   په~\LR{$\Frm[L_0]$} کښې هر فارمول \emph{متوازن} دے؛ يعنې د کيڼو
   قوسونو شمېر ئې د ښيو قوسونو له شمېر سره برابر دے۔
\end{prop}
\label{olpseg:OLP-0059-B008:end}

\phantomsection\label{olpseg:OLP-0059-B009:start}
\begin{prob}
د \olref[pl][syn][pre]{prop:balanced} ثبوت ورکړئ۔
\end{prob}
\label{olpseg:OLP-0059-B009:end}

\phantomsection\label{olpseg:OLP-0059-B010:start}
\begin{prop} \ollabel{prop:noinit}
د فارمول هېڅ خاصه پيلنۍ برخه فارمول نۀ ده۔
\end{prop}
\label{olpseg:OLP-0059-B010:end}

\phantomsection\label{olpseg:OLP-0059-B011:start}
\begin{prob}
  د \olref[pl][syn][pre]{prop:noinit} ثبوت ورکړئ۔
\end{prob}
\label{olpseg:OLP-0059-B011:end}

\phantomsection\label{olpseg:OLP-0059-B012:start}
\begin{prop}[يکتا لوستونتيا]
په \LR{$ \Frm[L_0]$} کښې هر فارمول~\LR{$!A$} په لاندې بڼو کښې د يوې په
توګه پوره يوه تجزيه لري:
\begin{enumerate}
\item \LR{$\lfalse$}۔
\label{olpseg:OLP-0059-B012:end}

\phantomsection\label{olpseg:OLP-0059-B013:start}
\item \LR{$\ltrue$}۔
\label{olpseg:OLP-0059-B013:end}

\phantomsection\label{olpseg:OLP-0059-B014:start}
\item د کوم \LR{$\Obj p_n \in  \PVar$} دپاره \LR{$\Obj p_n$}۔
\label{olpseg:OLP-0059-B014:end}

\phantomsection\label{olpseg:OLP-0059-B015:start}
\item د کوم فارمول~\LR{$!B$} دپاره \LR{$\lnot !B$}۔
\label{olpseg:OLP-0059-B015:end}

\phantomsection\label{olpseg:OLP-0059-B016:start}
\item د کومو فارمولونه \LR{$!B$} او~\LR{$!C$} دپاره
  \LR{$(!B \land !C)$}۔
\label{olpseg:OLP-0059-B016:end}

\phantomsection\label{olpseg:OLP-0059-B017:start}
\item د کومو فارمولونه \LR{$!B$} او~\LR{$!C$} دپاره
  \LR{$(!B \lor !C)$}۔
\label{olpseg:OLP-0059-B017:end}

\phantomsection\label{olpseg:OLP-0059-B018:start}
\item د کومو فارمولونه \LR{$!B$} او~\LR{$!C$} دپاره
  \LR{$(!B \lif !C)$}۔
\label{olpseg:OLP-0059-B018:end}

\phantomsection\label{olpseg:OLP-0059-B019:start}
\item د کومو فارمولونه \LR{$!B$} او~\LR{$!C$} دپاره
  \LR{$(!B \liff !C)$}۔
\end{enumerate}
سربېره پر دې، دا تجزيه \emph{يکتا} ده۔
\end{prop}
\label{olpseg:OLP-0059-B019:end}

\phantomsection\label{olpseg:OLP-0059-B020:start}
\begin{proof}
پر \LR{$!A$} استقرا کوو۔ مثلاً فرض کړئ چې \LR{$!A$} د \LR{$(!B \lif !C)$} او
\LR{$(!B' \lif !C')$} په توګه دوه بېلې لوستنې لري۔ نو \LR{$!B$} او \LR{$!B'$}
بايد يو شان وي، ګنې يو به د بل خاصه پيلنۍ برخه وي؛ نو که د \LR{$!A$}
دواړه لوستنې بېلې وي، علت ئې خامخا دا دے چې \LR{$!C$} او \LR{$!C'$} د نښو د
هماغې يوې لړۍ بېلې لوستنې دي، چې د استقرايي فرض له مخې ناشونې ده۔
\end{proof}
\label{olpseg:OLP-0059-B020:end}

\phantomsection\label{olpseg:OLP-0059-B021:start}
\begin{defn}[يو شان ځايناستي]
که \LR{$!A$} او \LR{$!B$} فارمولونه وي او \LR{$\Obj p_i$} يو بياني متغير وي، نو \LR{$\Subst{!A}{!B}{\Obj p_i}$} هغه پايله ښيي چې په
\LR{$!A$} کښې د \LR{$\Obj p_i$} هر راتګ د \LR{$!B$} په يوه راتګ بدل شي؛ همدارنګه،
د \LR{$\Obj p_1$}، \dots،~\LR{$\Obj p_n$} پر ځاے په يو وخت د فارمولونه
\LR{$!B_1$}، \dots،~\LR{$!B_n$} ځايناستي په
\LR{$\SSubst{!A}{\subst{!B_1}{\Obj p_1},\dots,\subst{!B_n}{\Obj p_n}}$}
ښودل کېږي۔
\end{defn}
\label{olpseg:OLP-0059-B021:end}

\phantomsection\label{olpseg:OLP-0059-B022:start}
\begin{prob} د لاندې پنځو فارمولونه په هر يوه کښې وټاکئ چې آيا
 فارمول د \( \Subst{!A}{!B}{\Obj p_i} \) د ځايناستي په توګه
 څرګندېدے شي، چې \( !A \) په دې کښې يو وي: (الف) \( \Obj p_0 \)؛
 (ب) \( ( \lnot \Obj p_0 \land \Obj p_1) \)؛ او (ج)
 \( ( ( \lnot \Obj p_0 \lif \Obj p_1 ) \land \Obj p_2 ) \)۔ په هر
 حالت کښې اړونده ځايناستي په ګوته کړئ۔
  \begin{enumerate}
    \item \( \Obj p_1 \) 
    \item \( ( \lnot \Obj p_0 \land \Obj p_0 ) \)
    \item \( ( ( \Obj p_0 \lor \Obj p_1 ) \land \Obj p_2 )  \)
    \item \( \lnot ( ( \Obj p_0 \lif \Obj p_1 ) \land \Obj p_2 ) \)
    \item \( (( \lnot ( \Obj p_0 \lif \Obj p_1 ) \lif ( \Obj p_0 \lor \Obj p_1 )) \land \lnot ( \Obj p_0 \land \Obj p_1 )) \)
  \end{enumerate}
\end{prob}
\label{olpseg:OLP-0059-B022:end}

\phantomsection\label{olpseg:OLP-0059-B023:start}
\begin{prob}
د \LR{$\Subst{!A}{!B}{p}$} يو له رياضي پلوه کره تعريف د استقرا له مخې
ورکړئ۔
\end{prob}
\label{olpseg:OLP-0059-B023:end}

\phantomsection\label{olpseg:OLP-0060-B004:start}
\olfileid{pl}{syn}{fseq}
\label{olpseg:OLP-0060-B004:end}

\phantomsection\label{olpseg:OLP-0060-B005:start}
\olsection{جوړښتي لړۍ}
\label{olpseg:OLP-0060-B005:end}

\phantomsection\label{olpseg:OLP-0060-B006:start}
په استقرايي تعريف سره د فارمولونه تعريفول، او د استقرا له لارې
د فارمولونه د خاصيتونو د ثابتولو بشپړوونکې طريقه، ښکلې او اغېزمنه
لار ده۔ خو کله کله دا هم ګټوره وي چې د فارمولونه جوړول له بېخه
پورته، ګام په ګام وڅېړو۔ دلته د \emph{جوړښتي لړۍ} په مفهوم همدا کار
کوو۔
\label{olpseg:OLP-0060-B006:end}

\phantomsection\label{olpseg:OLP-0060-B007:start}
\begin{defn}[د فارمولونو جوړښتي لړۍ]
\ollabel{defn:fseq-frm}
د ژبې~\LR{$\Lang L_0$} د نښو د توريزو لړيو يوه متناهي لړۍ
\LR{$\tuple{!A_0,\dotsc,!A_n}$} د \LR{$!A$} دپاره \emph{جوړښتي لړۍ} ده که
\LR{$!A \ident !A_n$} وي او د هر \LR{$i \leq n$} دپاره يا \LR{$!A_i$} اټومي
فارمول وي، يا داسې \LR{$j,k < i$} موجود وي چې له لاندې حالتونو څخه يو
صدق وکړي:
\begin{enumerate}
    \item \LR{$!A_i \ident \lnot !A_j$}۔%
    \item \LR{$!A_i \ident (!A_j \land !A_k)$}۔%
    \item \LR{$!A_i \ident (!A_j \lor !A_k)$}۔%
    \item \LR{$!A_i \ident (!A_j \lif !A_k)$}۔%
    \item \LR{$!A_i \ident (!A_j \liff !A_k)$}۔%
\end{enumerate}
\end{defn}
\label{olpseg:OLP-0060-B007:end}

\phantomsection\label{olpseg:OLP-0060-B008:start}
\begin{ex}
\[
\tuple{
    \Obj p_0,
    \Obj p_1,
    (\Obj p_1 \land \Obj p_0),
    \lnot (\Obj p_1 \land \Obj p_0)
}
\]
د \LR{$\lnot (\Obj p_1 \land \Obj p_0)$} يوه جوړښتي لړۍ ده، او دا هم ده:
\[
\tuple{
    \Obj p_0,
    \Obj p_1,
    \Obj p_0,
    (\Obj p_1 \land \Obj p_0),
    (\Obj p_0 \lif \Obj p_1),
    \lnot (\Obj p_1 \land \Obj p_0)
}.
\]
%
لکه چې له دويمې بېلګې ښکاري، په جوړښتي لړيو کښې ``بې‌ګټې توکي''
هم راتلے شي: داسې فارمولونه چې زياتي وي يا په جوړونه کښې ونډه نۀ
لري۔
\end{ex}
\label{olpseg:OLP-0060-B008:end}

\phantomsection\label{olpseg:OLP-0060-B009:start}
\begin{prop}
\ollabel{prop:formed}
په~\LR{$\Frm[L_0]$} کښې هر فارمول~\LR{$!A$} يوه جوړښتي لړۍ لري۔
\end{prop}
\label{olpseg:OLP-0060-B009:end}

\phantomsection\label{olpseg:OLP-0060-B010:start}
\begin{proof}
فرض کړئ \LR{$!A$} اټومي دے۔ نو لړۍ~\LR{$\tuple{!A}$} د~\LR{$!A$} جوړښتي لړۍ ده۔
%
اوس فرض کړئ چې \LR{$!B$} او~\LR{$!C$} په ترتيب سره جوړښتي لړۍ
\LR{$\tuple{!B_0,\dotsc,!B_n}$} او \LR{$\tuple{!C_0,\dotsc,!C_m}$} لري۔
%
\begin{enumerate}
  \item که \LR{$!A \ident \lnot !B$} وي، نو
      \LR{$\tuple{!B_0,\dotsc,!B_n,\lnot !B_n}$}
      د~\LR{$!A$} جوړښتي لړۍ ده۔
  \item که \LR{$!A \ident (!B \land !C)$} وي، نو
      \LR{$\tuple{!B_0,\dotsc,!B_n,!C_0,\dotsc,!C_m,(!B_n \land !C_m)}$}
      د~\LR{$!A$} جوړښتي لړۍ ده۔
  \item که \LR{$!A \ident (!B \lor !C)$} وي، نو
      \LR{$\tuple{!B_0,\dotsc,!B_n,!C_0,\dotsc,!C_m,(!B_n \lor !C_m)}$}
      د~\LR{$!A$} جوړښتي لړۍ ده۔
  \item که \LR{$!A \ident (!B \lif !C)$} وي، نو
      \LR{$\tuple{!B_0,\dotsc,!B_n,!C_0,\dotsc,!C_m,(!B_n \lif !C_m)}$}
      د~\LR{$!A$} جوړښتي لړۍ ده۔
  \item که \LR{$!A \ident (!B \liff !C)$} وي، نو
      \LR{$\tuple{!B_0,\dotsc,!B_n,!C_0,\dotsc,!C_m,(!B_n \liff !C_m)}$}
      د~\LR{$!A$} جوړښتي لړۍ ده۔
\end{enumerate}
د فارمولونه دپاره د استقرا د اصل له مخې، هر فارمول يوه
جوړښتي لړۍ لري۔
\end{proof}
\label{olpseg:OLP-0060-B010:end}

\phantomsection\label{olpseg:OLP-0060-B011:start}
د دې عکس هم ثابتولے شو۔ دا ځکه مهم دے چې ښيي د فارمولونو د تعريف
زموږ دواړه لارې همارزښته دي: دواړه يو شان پايلې ورکوي۔ دا هم ترې
معلومېږي چې د جوړښتي لړيو پر اوږدوالي عادي استقرا کارولے شو او د
فارمولونو په اړه قضيې ثابتولے شو۔
\label{olpseg:OLP-0060-B011:end}

\phantomsection\label{olpseg:OLP-0060-B012:start}
\begin{lem}
\ollabel{lem:fseq-init}
فرض کړئ \LR{$\tuple{!A_0,\dotsc,!A_n}$} د~\LR{$!A_n$} جوړښتي لړۍ ده او
\LR{$k \leq n$} دے۔ نو \LR{$\tuple{!A_0,\dotsc,!A_k}$} د~\LR{$!A_k$} جوړښتي لړۍ
ده۔
\end{lem}
\label{olpseg:OLP-0060-B012:end}

\phantomsection\label{olpseg:OLP-0060-B013:start}
\begin{proof}
تمرين۔
\end{proof}
\label{olpseg:OLP-0060-B013:end}

\phantomsection\label{olpseg:OLP-0060-B014:start}
\begin{thm}
\ollabel{thm:fseq-frm-equiv}
\LR{$\Frm[L_0]$} د ژبې~\LR{$\Lang L_0$} د نښو د ټولو هغو توريزو لړيو سټ دے
چې يوه جوړښتي لړۍ لري۔
\end{thm}
\label{olpseg:OLP-0060-B014:end}

\phantomsection\label{olpseg:OLP-0060-B015:start}
\begin{proof}
\LR{$F$} دې د ژبې~\LR{$\Lang L_0$} د نښو د هغو ټولو توريزو لړيو سټ وي چې
جوړښتي لړۍ لري۔ په \olref[pl][syn][fseq]{prop:formed} کښې مو وليدل
چې \LR{$\Frm[L_0] \subseteq F$}؛ اوس ئې عکس ثابتوو۔
\label{olpseg:OLP-0060-B015:end}

\phantomsection\label{olpseg:OLP-0060-B016:start}
فرض کړئ \LR{$!A$} يوه جوړښتي لړۍ \LR{$\tuple{!A_0,\dotsc,!A_n}$} لري۔ پر~\LR{$n$}
د قوي استقرا له لارې ثابتوو چې \LR{$!A \in \Frm[L_0]$}۔ زموږ استقرايي
فرض دا دے چې د نښو هره هغه توريزه لړۍ چې د \LR{$m < n$} اوږدوالي
جوړښتي لړۍ لري، په~\LR{$\Frm[L_0]$} کښې ده۔ د جوړښتي لړۍ د تعريف له مخې،
يا \LR{$!A_n$} اټومي دے، يا بايد داسې \LR{$j,k < n$} موجود وي چې له لاندې
حالتونو څخه يو صدق وکړي:
\begin{enumerate}
    \item \LR{$!A_n \ident \lnot !A_j$}۔%
    \item \LR{$!A_n \ident (!A_j \land !A_k)$}۔%
    \item \LR{$!A_n \ident (!A_j \lor !A_k)$}۔%
    \item \LR{$!A_n \ident (!A_j \lif !A_k)$}۔%
    \item \LR{$!A_n \ident (!A_j \liff !A_k)$}۔%
\end{enumerate}
اوس حالت په حالت استدلال کوو۔ که \LR{$!A_n$} اټومي وي، نو
\LR{$!A_n \in \Frm[L_0]$}۔ بل حالت کښې فرض کړئ چې \LR{$!A \ident
(!A_j \land !A_k)$}۔ د \olref[pl][syn][fseq]{lem:fseq-init} له مخې،
\LR{$\tuple{!A_0,\dotsc,!A_j}$} او \LR{$\tuple{!A_0,\dotsc,!A_k}$} په ترتيب
سره د \LR{$!A_j$} او~\LR{$!A_k$} جوړښتي لړۍ دي۔ ځکه چې دا د~\LR{$!A$} د جوړښتي
لړۍ خاصې پيلنۍ فرعي لړۍ دي، د دواړو اوږدوالے له~\LR{$n$} څخه لږ دے۔ نو
د استقرايي فرض له مخې \LR{$!A_j$} او~\LR{$!A_k$} په~\LR{$\Frm[L_0]$} کښې دي؛ او
بيا د فارمول د تعريف له مخې \LR{$(!A_j \land !A_k)$} هم پکې دے۔
نور حالتونه په ورته استدلال ثابتېږي۔
\textbf{د ژباړې د نسخې سمون (OLPL-002):} سرچينه په دې يوه کرښه کښې
د نحوي يوشانوالي د نښې پر ځاے د معنايي هم ارزښت نښه کاروي، حال دا
چې تعريف، ټول نور حالتونه او ثبوت د توريزو لړيو يوشانوالی غواړي۔
دلته د هماغه تعريف نښه بېرته کارول شوې ده۔
\end{proof}
\label{olpseg:OLP-0060-B016:end}

\phantomsection\label{olpseg:OLP-0061-B004:start}
\olfileid{pl}{syn}{val}
\label{olpseg:OLP-0061-B004:end}

\phantomsection\label{olpseg:OLP-0061-B005:start}
\olsection{ارزښت ټاکنې او صدق}
\label{olpseg:OLP-0061-B005:end}

\phantomsection\label{olpseg:OLP-0061-B006:start}
\begin{defn}[ارزښت ټاکنې]
\LR{$\{\True, \False\}$} دې د رښتياوالي د دوو قيمتونو، ``رښتيا'' او
``دروغ''، سټ وي۔ د \LR{$\Lang{L_0}$} دپاره \emph{ارزښت ټاکنه} يوه
تابع~\LR{$\pAssign{v}$} ده چې د ژبې بياني متغيرونه ته يا
\LR{$\True$} ورکوي يا \LR{$\False$}؛ يعنې
\LR{$\pAssign{v} \colon \PVar \to \{\True, \False \}$}۔
\end{defn}
\label{olpseg:OLP-0061-B006:end}

\phantomsection\label{olpseg:OLP-0061-B007:start}
\begin{defn}
  چې ارزښت ټاکنه~\LR{$\pAssign{v}$} راکړل شوے وي، د ارزونې تابع
  \LR{$\pValue{v} \colon \Frm[L_0] \to \{\True, \False \}$} په استقرايي
  ډول داسې تعريف کړئ:
  \begin{align*}
    \pValue{v}(\lfalse) & = \False; \\
    \pValue{v}(\ltrue) & = \True; \\
    \pValue{v}(\Obj p_n) & = \pAssign{v}(\Obj p_n); \\
    \pValue{v}(\lnot !A) & = \begin{cases}
        \True & \text{\RL{که }} \pValue{v}(!A) = \False;\\
        \False & \text{\RL{که داسې نۀ وي۔}}
      \end{cases} \\ 
    \pValue{v}(!A \land !B) & = \begin{cases}
        \True &
        \text{\RL{که \LR{$\pValue{v}(!A) = \True$} او \LR{$\pValue{v}(!B) = \True$} وي؛}}\\
        \False &
        \text{\RL{که \LR{$\pValue{v}(!A) = \False$} يا \LR{$\pValue{v}(!B) = \False$} وي}}.
      \end{cases}\\
    \pValue{v}(!A \lor !B) & = \begin{cases}
        \True &
        \text{\RL{که \LR{$\pValue{v}(!A) = \True$} يا \LR{$\pValue{v}(!B) = \True$} وي؛}}\\
        \False &
        \text{\RL{که \LR{$\pValue{v}(!A) = \False$} او \LR{$\pValue{v}(!B) = \False$} وي}}.
      \end{cases}\\
    \pValue{v}(!A \lif !B) & = \begin{cases}
        \True &
        \text{\RL{که \LR{$\pValue{v}(!A) = \False$} يا \LR{$\pValue{v}(!B) = \True$} وي؛}}\\
        \False &
        \text{\RL{که \LR{$\pValue{v}(!A) = \True$} او \LR{$\pValue{v}(!B) = \False$} وي}}.
      \end{cases}\\
    \pValue{v}(!A \liff !B) & = \begin{cases}
        \True &
        \text{\RL{که \LR{$\pValue{v}(!A) = \pValue{v}(!B)$} وي؛}}\\
        \False &
        \text{\RL{که \LR{$\pValue{v}(!A) \neq \pValue{v}(!B)$} وي}}.
      \end{cases}
\end{align*}
\end{defn}
\label{olpseg:OLP-0061-B007:end}

\phantomsection\label{olpseg:OLP-0061-B008:start}
\begin{explain}
دا مادې د رښتياوالي له لاندې جدولونو سره سمون لري:
\begin{center}

  \begin{tabular}{|c||c|} \hline
    \LR{$!A$} & \LR{$ \lnot !A$} \\
    \hline \hline
    \LR{$\True$} & \LR{$\False$} \\
    \LR{$\False$} & \LR{$\True$} \\
    \hline
  \end{tabular}


  \begin{tabular}{|cc||c|} \hline
    \LR{$!A$} & \LR{$!B$} & \LR{$!A \land !B$} \\
    \hline \hline
    \LR{$\True$} & \LR{$\True$} & \LR{$\True$} \\
    \LR{$\True$} & \LR{$\False$} & \LR{$\False$} \\
    \LR{$\False$} & \LR{$\True$} & \LR{$\False$} \\
    \LR{$\False$} & \LR{$\False$} & \LR{$\False$} \\
    \hline
  \end{tabular}


  \begin{tabular}{|cc||c|} \hline
    \LR{$!A$} & \LR{$!B$} & \LR{$!A \lor !B$} \\
    \hline \hline
    \LR{$\True$} & \LR{$\True$} & \LR{$\True$} \\
    \LR{$\True$} & \LR{$\False$} & \LR{$\True$} \\
    \LR{$\False$} & \LR{$\True$} & \LR{$\True$} \\
    \LR{$\False$} & \LR{$\False$} & \LR{$\False$} \\
    \hline
  \end{tabular}
%
\\[1em]

  \begin{tabular}{|cc||c|} \hline
    \LR{$!A$} & \LR{$!B$} & \LR{$!A \lif !B$} \\
    \hline \hline
    \LR{$\True$} & \LR{$\True$} & \LR{$\True$} \\
    \LR{$\True$} & \LR{$\False$} & \LR{$\False$} \\
    \LR{$\False$} & \LR{$\True$} & \LR{$\True$} \\
    \LR{$\False$} & \LR{$\False$} & \LR{$\True$} \\
    \hline
  \end{tabular}


  \begin{tabular}{|cc||c|} \hline
    \LR{$!A$} & \LR{$!B$} & \LR{$!A \liff !B$} \\
    \hline \hline
    \LR{$\True$} & \LR{$\True$} & \LR{$\True$} \\
    \LR{$\True$} & \LR{$\False$} & \LR{$\False$} \\
    \LR{$\False$} & \LR{$\True$} & \LR{$\False$} \\
    \LR{$\False$} & \LR{$\False$} & \LR{$\True$} \\
    \hline
  \end{tabular}

\end{center}
\end{explain}
\label{olpseg:OLP-0061-B008:end}

\phantomsection\label{olpseg:OLP-0061-B009:start}
\begin{prob}
فرض کړئ چې \LR{$\Lang{L_0}$} ته يو درې ځايه نښلوونکے \LR{$\diamondsuit$}
ورزياتوو چې ارزونه ئې داسې ورکول شوې ده:
\begin{gather*}
  \pValue{v}(\diamondsuit( !A , !B , !C ) ) = \begin{cases}
    \pValue{v}( !B ) &
    \text{\RL{که \LR{$\pValue{v}(!A) = \True$} وي؛}}\\
    \pValue{v}( !C ) &
    \text{\RL{که \LR{$\pValue{v}(!A) = \False $} وي}}۔
  \end{cases}
\end{gather*}
د دې نښلوونکي د رښتياوالي جدول وليکئ۔
\end{prob}
\label{olpseg:OLP-0061-B009:end}

\phantomsection\label{olpseg:OLP-0061-B010:start}
\begin{thm}[ځايي ټاکنه]
\ollabel{thm:LocalDetermination} فرض کړئ \LR{$\pAssign{v_1}$} او
\LR{$\pAssign{v_2}$} داسې ارزښت ټاکنې دي چې په \LR{$!A$} کښې په راتلونکو
بياني متغيرونه يو له بل سره موافق دي؛ يعنې هر کله چې
\LR{$\Obj p_n$} په فارمول~\LR{$!A$} کښې راشي،
\LR{$\pAssign{v_1}(\Obj p_n) = \pAssign{v_2}(\Obj p_n)$} وي۔ نو
\LR{$\pValue{v_1}$} او \LR{$\pValue{v_2}$} هم پر~\LR{$!A$} موافق دي؛ يعنې
\LR{$\pValue{v_1}(!A) = \pValue{v_2}(!A)$}۔
\end{thm}
\textbf{د ژباړې د نسخې سمون (OLPL-003):} سرچينه د فرض په دويمه
وينا کښې ``په کوم فارمول کښې'' ليکي، خو قضيه هماغه له مخکښې ټاکلے
شوے فارمول ارزوي۔ ژباړه د لومړۍ وينا سره سم د هماغه فارمول متغيرونه
اخلي؛ په دې سره د قضیې مقصود شرط څرګند پاتې کېږي۔
\label{olpseg:OLP-0061-B010:end}

\phantomsection\label{olpseg:OLP-0061-B011:start}
\begin{proof}
پر \LR{$!A$} استقرا کوو۔
\end{proof}
\label{olpseg:OLP-0061-B011:end}

\phantomsection\label{olpseg:OLP-0061-B012:start}
\begin{defn}[صدق]
\ollabel{defn:satisfaction} د \emph{فارمول~\LR{$!A$} صدق په
  ارزښت ټاکنه~\LR{$\pAssign{v}$} کښې}، يعنې \LR{$\pSat{v}{!A}$}، په
  استقرايي ډول داسې تعريفولے شو۔ (د ``\LR{$\pSat{v}{!A}$} نۀ دے'' د
  څرګندولو دپاره \LR{$\pSat/{v}{!A}$} ليکو۔)
\begin{enumerate}
\item %
  \indcase{!A}{\lfalse}{\LR{$\pSat/{v}{\indfrm}$}۔}
\label{olpseg:OLP-0061-B012:end}

\phantomsection\label{olpseg:OLP-0061-B013:start}
\item %
  \indcase{!A}{\ltrue}{\LR{$\pSat{v}{\indfrm}$}۔}
\label{olpseg:OLP-0061-B013:end}

\phantomsection\label{olpseg:OLP-0061-B014:start}
\item \indcase{!A}{\Obj p_i}{\LR{$\pSat{v}{\indfrm}$}
  هله او يوازې هله چې \LR{$\pAssign{v}(\Obj p_i) = \True$} وي۔}
\label{olpseg:OLP-0061-B014:end}

\phantomsection\label{olpseg:OLP-0061-B015:start}
\item %
  \indcase{!A}{\lnot !B}{\LR{$\pSat{v}{\indfrm}$} هله او يوازې هله چې
    \LR{$\pSat/{v}{!B}$} وي۔}
\label{olpseg:OLP-0061-B015:end}

\phantomsection\label{olpseg:OLP-0061-B016:start}
\item %
  \indcase{!A}{(!B \land !C)}{\LR{$\pSat{v}{\indfrm}$} هله او يوازې هله
    چې \LR{$\pSat{v}{!B}$} او \LR{$\pSat{v}{!C}$} دواړه وي۔}
\label{olpseg:OLP-0061-B016:end}

\phantomsection\label{olpseg:OLP-0061-B017:start}
\item %
  \indcase{!A}{(!B \lor !C)}{\LR{$\pSat{v}{\indfrm}$} هله او يوازې هله
    چې \LR{$\pSat{v}{!B}$} يا \LR{$\pSat{v}{!C}$}، يا دواړه، وي۔}
\label{olpseg:OLP-0061-B017:end}

\phantomsection\label{olpseg:OLP-0061-B018:start}
\item %
  \indcase{!A}{(!B \lif !C)}{\LR{$\pSat{v}{\indfrm}$} هله او يوازې هله
    چې \LR{$\pSat/{v}{!B}$} يا \LR{$\pSat{v}{!C}$}، يا دواړه، وي۔}
\label{olpseg:OLP-0061-B018:end}

\phantomsection\label{olpseg:OLP-0061-B019:start}
\item %
  \indcase{!A}{(!B \liff !C)}{\LR{$\pSat{v}{\indfrm}$} هله او يوازې هله
    چې يا \LR{$\pSat{v}{!B}$} او \LR{$\pSat{v}{!C}$} دواړه وي، يا نۀ
    \LR{$\pSat{v}{!B}$} وي او نۀ \LR{$\pSat{v}{!C}$}۔}
\end{enumerate}
که \LR{$\Gamma$} د فارمولونه يو سټ وي، نو \LR{$\pSat{v}{\Gamma}$} هله او
يوازې هله چې د هر~\LR{$!A \in \Gamma$} دپاره \LR{$\pSat{v}{!A}$} وي۔
\end{defn}
\label{olpseg:OLP-0061-B019:end}

\phantomsection\label{olpseg:OLP-0061-B020:start}
\begin{prop}\ollabel{prop:sat-value}
  \LR{$\pSat{v}{!A}$} هله او يوازې هله چې \LR{$\pValue{v}(!A) = \True$} وي۔
\end{prop}
\label{olpseg:OLP-0061-B020:end}

\phantomsection\label{olpseg:OLP-0061-B021:start}
\begin{proof}
  پر~\LR{$!A$} استقرا کوو۔
\end{proof}
\label{olpseg:OLP-0061-B021:end}

\phantomsection\label{olpseg:OLP-0061-B022:start}
\begin{prob}
  د \olref[pl][syn][val]{prop:sat-value} ثبوت ورکړئ۔
\end{prob}
\label{olpseg:OLP-0061-B022:end}

\phantomsection\label{olpseg:OLP-0062-B004:start}
\olfileid{pl}{syn}{sem}
\label{olpseg:OLP-0062-B004:end}

\phantomsection\label{olpseg:OLP-0062-B005:start}
\olsection{معنايي مفهومونه}
\label{olpseg:OLP-0062-B005:end}

\phantomsection\label{olpseg:OLP-0062-B006:start}
لاندې معنايي مفهومونه تعريفوو:
\label{olpseg:OLP-0062-B006:end}

\phantomsection\label{olpseg:OLP-0062-B007:start}
\begin{defn} 
\begin{enumerate}
\item فارمول~\LR{$!A$} هله \emph{د صدق وړ} دے چې د کوم
  \LR{$\pAssign{v}$} دپاره \LR{$\pSat{v}{!A}$} وي؛ او هله \emph{د صدق نۀ وړ}
  دے چې د هېڅ \LR{$\pAssign{v}$} دپاره \LR{$\pSat{v}{!A}$} نۀ وي؛
\item فارمول~\LR{$!A$} هله \emph{تاوتولوژي} ده چې د ټولو
  ارزښت ټاکنې~\LR{$\pAssign{v}$} دپاره \LR{$\pSat{v}{!A}$} وي؛
\item فارمول~\LR{$!A$} هله \emph{اتفاقي} دے چې د صدق وړ وي خو
  تاوتولوژي نۀ وي؛
\item که \LR{$\Gamma$} د فارمولونه يو سټ وي، نو \LR{$\Gamma \Entails !A$}
  (``\LR{$\Gamma$}، \LR{$!A$} ته استلزام ورکوي'') هله او يوازې هله چې د هر
  هغه ارزښت ټاکنه~\LR{$\pAssign{v}$} دپاره \LR{$\pSat{v}{!A}$} وي چې
  \LR{$\pSat{v}{\Gamma}$} وي۔
\item که \LR{$\Gamma$} د فارمولونه يو سټ وي، نو \LR{$\Gamma$} هله
  \emph{د صدق وړ} دے چې داسې ارزښت ټاکنه~\LR{$\pAssign{v}$} موجود وي
  چې \LR{$\pSat{v}{\Gamma}$} وي؛ او که داسې نۀ وي، \LR{$\Gamma$}
  \emph{د صدق نۀ وړ} دے۔
\end{enumerate} 
\end{defn}
\label{olpseg:OLP-0062-B007:end}

\phantomsection\label{olpseg:OLP-0062-B008:start}
\begin{prob} 
 د لاندې څلورو فارمولونه په هر يوه کښې وټاکئ چې آيا هغه
 (الف)~د صدق وړ، (ب)~تاوتولوژي او (ج)~اتفاقي دے۔
\begin{enumerate}
  \item \( ( \Obj p_0 \lif ( \lnot \Obj p_1 \lif \lnot \Obj p_0 ) ) \).
  \item \( ( ( \Obj p_0 \land \lnot \Obj p_1 ) \lif ( \lnot \Obj p_0 \land \Obj p_2 )) \liff ( ( \Obj p_2 \lif \Obj p_0 ) \lif ( \Obj p_0 \lif \Obj p_1 )) \).
  \item \( ( \Obj p_0 \liff \Obj p_1 ) \lif ( \Obj p_2 \liff \lnot \Obj p_1 ) \).
  \item \( (( \Obj p_0 \liff ( \lnot \Obj p_1 \land \Obj p_2 )) \lor ( \Obj p_2 \lif ( \Obj p_0 \liff \Obj p_1 ))) \).
\end{enumerate}
\end{prob}
\label{olpseg:OLP-0062-B008:end}

\phantomsection\label{olpseg:OLP-0062-B009:start}
\begin{prop}
\ollabel{prop:semanticalfacts} 
\begin{enumerate} 
\item \LR{$!A$} هله او يوازې هله تاوتولوژي ده چې
  \LR{$\emptyset \Entails !A$} وي؛ 
\item که \LR{$\Gamma \Entails !A$} او \LR{$\Gamma \Entails !A \lif !B$} وي،
  نو \LR{$\Gamma \Entails !B$}؛
\item که \LR{$\Gamma$} د صدق وړ وي، نو د \LR{$\Gamma$} هر متناهي فرعي سټ هم
  د صدق وړ دے؛ 
\item \ollabel{def:monotonicity} يو اړخيز زياتوالے: که
  \LR{$\Gamma \subseteq \Delta$} او \LR{$\Gamma \Entails !A$} وي، نو
  \LR{$\Delta \Entails !A$} هم؛
\item \ollabel{def:Cut} تعدي: که \LR{$\Gamma \Entails !A$} او
  \LR{$\Delta \cup \{ !A\} \Entails !B$} وي، نو
  \LR{$\Gamma \cup \Delta \Entails !B$}۔
\end{enumerate}
\end{prop}
\label{olpseg:OLP-0062-B009:end}

\phantomsection\label{olpseg:OLP-0062-B010:start}
\begin{proof}
تمرين۔
\end{proof}
\label{olpseg:OLP-0062-B010:end}

\phantomsection\label{olpseg:OLP-0062-B011:start}
\begin{prob}
د \olref[pl][syn][sem]{prop:semanticalfacts} ثبوت ورکړئ۔
\end{prob}
\label{olpseg:OLP-0062-B011:end}

\phantomsection\label{olpseg:OLP-0062-B012:start}
\begin{prop}\ollabel{prop:entails-unsat}
  \LR{$\Gamma \Entails !A$} هله او يوازې هله چې
  \LR{$\Gamma \cup \{\lnot !A\}$} د صدق نۀ وړ وي۔
\end{prop}
\label{olpseg:OLP-0062-B012:end}

\phantomsection\label{olpseg:OLP-0062-B013:start}
\begin{proof}
تمرين۔
\end{proof}
\label{olpseg:OLP-0062-B013:end}

\phantomsection\label{olpseg:OLP-0062-B014:start}
\begin{prob}
د \olref[pl][syn][sem]{prop:entails-unsat} ثبوت ورکړئ۔
\end{prob}
\label{olpseg:OLP-0062-B014:end}

\phantomsection\label{olpseg:OLP-0062-B015:start}
\begin{thm}[د معنايي استنتاج قضيه]
  \ollabel{thm:sem-deduction} \LR{$\Gamma \Entails !A \lif !B$} هله او
  يوازې هله چې \LR{$\Gamma \cup \{!A\} \Entails !B$} وي۔
\end{thm}
\label{olpseg:OLP-0062-B015:end}

\phantomsection\label{olpseg:OLP-0062-B016:start}
\begin{proof}
تمرين۔
\end{proof}
\label{olpseg:OLP-0062-B016:end}

\phantomsection\label{olpseg:OLP-0062-B017:start}
\begin{prob}
د \olref[pl][syn][sem]{thm:sem-deduction} ثبوت ورکړئ۔
\end{prob}
\label{olpseg:OLP-0062-B017:end}

\phantomsection\label{olpseg:OLP-0063-B004:start}
\olchapter{fol}{prf}{د اشتقاق نظامونه}
\label{olpseg:OLP-0063-B004:end}

\phantomsection\label{olpseg:OLP-0063-B005:start}
\begin{editorial}
دا څپرکے د اشتقاق د نظامونو عمومي مواد راغونډوي۔ کوم درسي
کتاب چې يو ځانګړے نظام کاروي، د هغه نظام د معرفي کوونکي څپرکي په
پيل کښې د پېژندنې برخه او د اړوندې ټولکتنې برخه ورګډولے شي۔
\end{editorial}
\label{olpseg:OLP-0063-B005:end}

\phantomsection\label{olpseg:OLP-0064-B004:start}
\olfileid{fol}{prf}{int}
\label{olpseg:OLP-0064-B004:end}

\phantomsection\label{olpseg:OLP-0064-B005:start}
\olsection{پېژندنه}
\label{olpseg:OLP-0064-B005:end}

\phantomsection\label{olpseg:OLP-0064-B006:start}
منطقونه عموماً معنٰی پوهنه او اشتقاق نظام دواړه لري۔ معنٰی
پوهنه د رښتياوالي، د صدق وړتيا، اعتبار او استلزام په څېر مفهومونو
سره کار لري۔ د اشتقاق د نظامونو موخه دا ده چې د استلزام او
اعتبار د ثابتولو يوه بشپړه نحوي طريقه ورکړي۔ دا نظامونه په دې معنٰی
بشپړ نحوي دي چې په داسې يوه نظام کښې اشتقاق يو متناهي نحوي
څيز وي، چې عموماً د تړلي فارمولونه يا فارمولونه يوه لړۍ يا بل
متناهي ترتيب وي۔ د اشتقاق ښۀ نظامونه دا خاصيت لري چې د
تړلي فارمولونه يا فارمولونه د هرې ورکړل شوې لړۍ يا ترتيب
``سموالے'' په ميخانيکي توګه کتل کېدے شي۔
\label{olpseg:OLP-0064-B006:end}

\phantomsection\label{olpseg:OLP-0064-B007:start}
د لومړۍ درجې منطق تر ټولو ساده، او له تاريخي پلوه لومړني،
اشتقاق نظامونه \emph{بديهي} وو۔ په داسې يوه نظام کښې د
فارمولونه يوه لړۍ هله اشتقاق ګڼل کېږي چې په هغې کښې هر
فارمول يا د ``بديهي اصولو'' په يوه ټاکلي سټ کښې وي، يا د ټاکلو
``استنتاجي قاعدو'' له يوې څخه په کار اخيستو په لړۍ کښې له مخکښو
فارمولونه څخه ترې راووځي؛ او په ميخانيکي توګه کتل کېدے شي چې آيا
فارمول بديهي اصل دے او آيا د استنتاج له کومې قاعدې سره سم له
نورو فارمولونه څخه راوتلے دے۔ د اشتقاق بديهي نظامونه
تشريح کول اسانه دي، او له مابعدمنطقي پلوه پرې کار کول هم اسانه دي،
خو په هغو کښې اشتقاقونه لوستل او پوهېدل ګران دي او جوړول ئې هم
ګران دي۔
\label{olpseg:OLP-0064-B007:end}

\phantomsection\label{olpseg:OLP-0064-B008:start}
د اشتقاق نور نظامونه د دې موخې دپاره جوړ شوي چې د
اشتقاقونه جوړول يا له بشپړېدو وروسته پر اشتقاقونه
پوهېدل اسانه کړي۔
بېلګې ئې طبيعي استنتاج، د رښتياوالي ونې چې د تابلو ثبوتونو په نامه
هم پېژندل کېږي، او د سېکوېنټ حساب دي۔ د اشتقاق ځينې نظامونه
په ځانګړې توګه د ميخانيکي کولو دپاره طرح شوي دي؛ مثلاً د حل طريقه
په سافټوير کښې په اسانۍ پلې کېږي، خو پر اشتقاقونه ئې پوهېدل
تقريباً ناشوني دي۔ د اشتقاق دا نور نظامونه زياتره
اشتقاقونه د لړيو پر ځاے د فارمولونه د ونو په توګه ښيي۔ په
دې سره په اسانۍ ليدل کېږي چې د اشتقاق کومې برخې پر کومو
نورو برخو تکيه لري۔
\label{olpseg:OLP-0064-B008:end}

\phantomsection\label{olpseg:OLP-0064-B009:start}
نو د يوه ورکړل شوي منطق، لکه د لومړۍ درجې منطق، د اشتقاق
بېل نظامونه به دا په بېلو ډولونو روښانه کړي چې تړلے فارمول کله
\emph{قضيه} ده او د نورو له ځينو څخه د تړلے فارمول د
د اشتقاق وړ کېدو معنٰی څه ده۔ دا که په بديهي اشتقاقونه،
طبيعي استنتاجونو، سېکوېنټي اشتقاقونه، د رښتياوالي په ونو يا د
حل په ابطالونو ترسره شي، غواړو چې دا اړيکې د اعتبار او استلزام له
معنايي مفهومونو سره سمون ولري۔ د ``\LR{$!A$}~قضيه ده'' دپاره
\LR{$\Proves !A$} او د ``\LR{$!A$} له~\LR{$\Gamma$} څخه د اشتقاق وړ دے'' دپاره
``\LR{$\Gamma \Proves !A$}'' ليکو۔ \LR{$\Proves$} چې هر ډول تعريف شي، غواړو
له \LR{$\Entails$} سره سمون ولري؛ يعنې:
\begin{enumerate}
\item \LR{$\Proves !A$} هله او يوازې هله چې \LR{$\Entails !A$}
\item \LR{$\Gamma \Proves !A$} هله او يوازې هله چې \LR{$\Gamma \Entails !A$}
\end{enumerate}
د پورته دواړو د ``يوازې هله'' لوري ته \emph{سموالے} ويل کېږي۔
اشتقاق نظام هله سم دے چې د اشتقاق وړتيا استلزام يا اعتبار
تضمين کړي۔ د اشتقاق هر د کار وړ نظام بايد سم وي؛ د
اشتقاق ناسم نظامونه هېڅ ګټه نۀ لري۔ په اصل کښې د
اشتقاق ټوله موخه دا ده چې د اعتبار يا استلزام نحوي تضمين
ورکړي۔ موږ به د وړاندې شويو اشتقاق نظامونو سموالے ثابت کړو۔
\label{olpseg:OLP-0064-B009:end}

\phantomsection\label{olpseg:OLP-0064-B010:start}
د دې عکس، يعنې ``هله'' لورے، هم مهم دے: دې ته \emph{بشپړتيا} ويل
کېږي۔ د اشتقاق بشپړ نظام دومره ځواکمن وي چې هر کله \LR{$!A$}
معتبر وي، وښيي چې \LR{$!A$} قضيه ده؛ او هر کله چې
\LR{$\Gamma \Entails !A$} وي، وښيي چې \LR{$\Gamma \Proves !A$}۔ بشپړتيا ثابتول
ګران دي، او ځينې منطقونه د اشتقاق هېڅ بشپړ نظام نۀ لري۔ د
لومړۍ درجې منطق ئې لري۔ کورت ګوډل لومړے کس و چې د لومړۍ درجې منطق د
اشتقاق د يوه نظام بشپړتيا ئې په خپله ۱۹۲۹ز کال څېړنيزه
رساله کښې ثابته کړه۔
\label{olpseg:OLP-0064-B010:end}

\phantomsection\label{olpseg:OLP-0064-B011:start}
له اشتقاق نظامونو سره يو بل تړلے مفهوم \emph{سازګاري} ده۔
د تړلي فارمولونه يو سټ هله ناسازګار بلل کېږي چې هر څه ترې
اشتقاق کېدے شي؛ او که داسې نۀ وي، سازګار دے۔ ناسازګاري د صدق
د نۀ وړتيا نحوي سياله ده: لکه د صدق نۀ وړ سټونو غوندې، د
تړلي فارمولونه ناسازګار سټونه ښې نظريې نۀ جوړوي او په بنسټيز ډول
نيمګړي وي۔ د تړلي فارمولونه سازګار سټونه کېدے شي رښتوني يا ګټور نۀ
وي، خو لږ تر لږه د منطقي ګټورتيا له دې تر ټولو ټيټ بريد څخه تېرېږي۔
د اشتقاق په بېلو نظامونو کښې د تړلي فارمولونه د سټونو د
سازګارۍ ځانګړے تعريف توپير کولے شي؛ خو لکه~\LR{$\Proves$}، غواړو
سازګاري د خپل معنايي سيال، د صدق وړتيا، سره سمه وي۔ غواړو تل داسې
وي چې \LR{$\Gamma$} هله او يوازې هله سازګار وي چې د صدق وړ وي۔ دلته د
``يوازې هله'' لورے له بشپړتيا سره برابر دے، ځکه سازګاري د صدق
وړتيا تضمينوي؛ او ``هله'' لورے له سموالي سره برابر دے، ځکه د صدق
وړتيا سازګاري تضمينوي۔ په حقيقت کښې د کلاسيکي لومړۍ درجې منطق
دپاره د سموالي او بشپړتيا دواړه بڼې همارزښته دي۔
\label{olpseg:OLP-0064-B011:end}

\phantomsection\label{olpseg:OLP-0065-B004:start}
\olfileid{fol}{prf}{seq}
\label{olpseg:OLP-0065-B004:end}

\phantomsection\label{olpseg:OLP-0065-B005:start}
\olsection{د سېکوېنټ حساب}
\label{olpseg:OLP-0065-B005:end}

\phantomsection\label{olpseg:OLP-0065-B006:start}
د اشتقاق ډېر نظامونه د تړلي فارمولونه له ترتيبونو سره کار
کوي، خو د سېکوېنټ حساب له \emph{سېکوېنټونو} سره کار کوي۔ سېکوېنټ
د دې بڼې يو عبارت دے:
\[
!A_1, \dots, !A_m \Sequent !B_1, \dots, !B_n,
\]
\textbf{د ژباړې د نسخې سمون (OLPF-001):} سرچينه د سېکوېنټ دواړو
خواوو ته يو شان وروستے اندېکس کاروي، چې دواړه لړۍ په ناپامۍ يو شان
اوږدوي؛ خو ورپسې څرګندوي چې هره لړۍ په خپلواکه توګه تشه کېدے شي۔
دلته د ښۍ لړۍ وروستے اندېکس بېل شوے، څو دواړه متناهي لړۍ خپلواکه
اوږدوالی ولري۔
يعنې دا د تړلي فارمولونه د لړيو يوه جوړه ده چې د سېکوېنټ په
نښه~\LR{$\Sequent$} بېله شوې ده۔ له دواړو لړيو هره يوه تشه کېدے شي۔ د
سېکوېنټ په حساب کښې اشتقاق د سېکوېنټونو يوه ونه ده۔ تر
ټولو پاسني سېکوېنټونه ځانګړې بڼه لري، او ``لومړني سېکوېنټونه'' يا
``بديهي اصول'' بلل کېږي؛ هر بل سېکوېنټ د استنتاج د کومې قاعدې له
لارې سم له پاسه سېکوېنټونو څخه راځي۔ د استنتاج قاعدې يا په
سېکوېنټونو کښې تړلي فارمولونه سمبالوي، يعنې کيڼې يا ښۍ خوا ته ئې
ورزياتوي، ترې لرې کوي يا ئې بيا ترتيبوي؛ او يا د قاعدې په پايله کښې
يو مرکب فارمول معرفي کوي۔ مثلاً د \LR{$\LeftR{\land}$} قاعده اجازه
راکوي چې له \LR{$!A, \Gamma \Sequent \Delta$} څخه \LR{$!A \land !B, \Gamma
\Sequent \Delta$} استنتاج کړو، او \LR{$\RightR{\lif}$} اجازه راکوي چې د
هر \LR{$\Gamma$}، \LR{$\Delta$}، \LR{$!A$} او~\LR{$!B$} دپاره له \LR{$!A, \Gamma \Sequent
\Delta, !B$} څخه \LR{$\Gamma \Sequent \Delta, !A \lif !B$} استنتاج کړو۔
(په تېره بيا، \LR{$\Gamma$} او~\LR{$\Delta$} تش کېدے شي۔)
\label{olpseg:OLP-0065-B006:end}

\phantomsection\label{olpseg:OLP-0065-B007:start}
د سېکوېنټ پر حساب ولاړه \LR{$\Proves$} اړيکه داسې تعريفېږي:
\LR{$\Gamma \Proves !A$} هله او يوازې هله چې داسې کومه لړۍ \LR{$\Gamma_0$}
موجوده وي چې په \LR{$\Gamma_0$} کښې هر \LR{$!A$} په~\LR{$\Gamma$} کښې وي، او داسې
اشتقاق موجود وي چې سېکوېنټ~\LR{$\Gamma_0 \Sequent !A$} ئې په
جرړه کښې وي۔ \LR{$!A$} هله د سېکوېنټ په حساب کښې قضيه ده چې سېکوېنټ
~\LR{$\Sequent !A$} کوم اشتقاق ولري۔ مثلاً دا اشتقاق
ښيي چې \LR{$\Proves (!A \land !B) \lif !A$}:
\begin{prooftree}
  \Axiom$!A \fCenter !A$
  \RightLabel{\LeftR{\land}}
  \UnaryInf$!A \land !B \fCenter !A$
  \RightLabel{\RightR{\lif}}
  \UnaryInf$\fCenter (!A \land !B) \lif !A$
\end{prooftree}
\label{olpseg:OLP-0065-B007:end}

\phantomsection\label{olpseg:OLP-0065-B008:start}
يو سټ \LR{$\Gamma$} د سېکوېنټ په حساب کښې هله ناسازګار دے چې د
\LR{$\Gamma_0 \Sequent$} کوم اشتقاق موجود وي، چې په کښې هر
\LR{$!A \in \Gamma_0$} په~\LR{$\Gamma$} کښې وي او د سېکوېنټ ښۍ خوا تشه وي۔
د \RightR{\Weakening} قاعدې په کارولو له ناسازګار سټ څخه هر
تړلے فارمول اشتقاق کېدے شي۔
\label{olpseg:OLP-0065-B008:end}

\phantomsection\label{olpseg:OLP-0065-B009:start}
د سېکوېنټ حساب په ۱۹۳۰مو کلونو کښې ګيرهارډ ګېنتڅن جوړ کړ۔ د منظمې
او متناظرې طرحې له امله دا د اشتقاقونه د نظريې د ودې دپاره
ډېر ګټور صوري نظام دے۔ د سېکوېنټ په حساب کښې د اشتقاقونه
موندل نسبتاً اسانه دي، خو دا اشتقاقونه ډېر ځله په سختۍ لوستل
کېږي او له ثبوتونو سره ئې تړاو کله کله په اسانۍ نۀ ښکاري۔ بيا هم،
دا د اشتقاق نظامونو ډېره ښکلې لار ثابته شوې، او ډېر منطقونه
د سېکوېنټ حساب نظامونه لري۔
\label{olpseg:OLP-0065-B009:end}

\phantomsection\label{olpseg:OLP-0066-B004:start}
\olfileid{fol}{prf}{ntd}
\label{olpseg:OLP-0066-B004:end}

\phantomsection\label{olpseg:OLP-0066-B005:start}
\olsection{طبيعي استنتاج}
\label{olpseg:OLP-0066-B005:end}

\phantomsection\label{olpseg:OLP-0066-B006:start}
طبيعي استنتاج د اشتقاق داسې نظام دے چې موخه ئې د رښتيني
استدلال، په تېره بيا د رياضي پوهانو د منظم استدلال، انځورول دي۔
رښتينے استدلال د يو شمېر ``طبيعي'' بڼو له مخې مخکښې ځي۔ مثلاً د
حالتونو ثبوت موږ ته اجازه راکوي چې د يوه فصلي فرض پر بنسټ پايله
ثابته کړو، په دې چې وښيو پايله د فصل له هرې برخې څخه راځي۔ نامستقيم
ثبوت اجازه راکوي چې پايله په دې سره ثابته کړو چې نفي ئې تناقض ته
رسېږي۔ شرطي ثبوت د ``که \dots نو \dots'' شرطي ادعا په دې سره ثابتوي
چې تالي له مقدم څخه راځي۔ طبيعي استنتاج د دغو طبيعي استنتاجونو د
ځينو صوري کول دي۔ هر منطقي نښلوونکے او هر سور دوه قاعدې لري، يوه د
معرفي کولو او بله د لرې کولو قاعده، او هره يوه ئې د داسې طبيعي
استنتاج له يوې بڼې سره سمه ده۔ مثلاً \LR{$\Intro{\lif}$} له شرطي ثبوت
سره او \LR{$\Elim{\lor}$} د حالتونو له ثبوت سره سمون لري۔ يوه په تېره
ساده قاعده \LR{$\Elim{\land}$} ده، چې اجازه راکوي له \LR{$!A \land !B$} څخه
~\LR{$!A$} يا \LR{$!B$} استنتاج کړو۔
\label{olpseg:OLP-0066-B006:end}

\phantomsection\label{olpseg:OLP-0066-B007:start}
يو خاصيت چې طبيعي استنتاج د اشتقاق له نورو نظامونو څخه
بېلوي، د فرضونو کارول دي۔ په طبيعي استنتاج کښې اشتقاق د
فارمولونه يوه ونه ده۔ يو فارمول د فارمولونه د ونې په جرړه
کښې وي، او د ونې ``پاڼې'' هغه فارمولونه وي چې پايله ترې راايستل
کېږي۔ په طبيعي استنتاج کښې ځينې پاڼه-فارمولونه د اشتقاق
دننه ونډه لري، خو تر هغې ``کارول شوي وي'' چې اشتقاق پايلې ته
رسېږي۔ دا په رښتيني استدلال کښې د هغو فرضيو د معرفي کولو له دود سره
سمون لري چې يوازې لږ وخت نافذې پاتې کېږي۔ مثلاً د حالتونو په ثبوت
کښې د فصل د هرې برخې رښتياوالے فرضوو؛ په شرطي ثبوت کښې د مقدم
رښتياوالے فرضوو؛ او په نامستقيم ثبوت کښې د پايلې د نفي رښتياوالے
فرضوو۔ د فرضي فرضونو د معرفي کولو او بيا د يوه منځني ګام د ثابتولو
دپاره د هغو ختمولو دا طريقه د طبيعي استنتاج څرګنده نښه ده۔ د طبيعي
استنتاج د يوه اشتقاق په پاڼو کښې فارمولونه فرضونه بلل کېږي،
او د استنتاج ځينې قاعدې کېدے شي هغه ``ايستل'' کړي۔ مثلاً که
مو له ځينو فرضونو څخه د~\LR{$!B$} کوم اشتقاق درلود چې~\LR{$!A$} هم
پکې شامل و، نو د \LR{$\Intro{\lif}$} قاعده اجازه راکوي چې~\LR{$!A \lif !B$}
استنتاج کړو او د~\LR{$!A$} بڼې هر فرض ختم کړو۔ (د دې ثبتولو
دپاره چې کوم فرضونه په کومو استنتاجونو کښې ختمېږي، استنتاج او هغه
فرضونه چې ختموي ئې په يوه شمېره نښه کوو۔) هغه فرضونه چې د
اشتقاق په پاے کښې نۀ ايستل شوې پاتې کېږي، په ګډه د پايلې
د رښتياوالي دپاره بس دي؛ نو اشتقاق ثابتوي چې د هغه
نۀ ايستل شوي فرضونه پايلې ته استلزام ورکوي۔
\label{olpseg:OLP-0066-B007:end}

\phantomsection\label{olpseg:OLP-0066-B008:start}
د طبيعي استنتاج پر بنسټ اړيکه \LR{$\Gamma \Proves !A$} هله صدق کوي چې
داسې اشتقاق موجود وي چې \LR{$!A$} د ونې وروستۍ تړلے فارمول وي،
او هره نۀ ايستل شوې پاڼه په~\LR{$\Gamma$} کښې وي۔ \LR{$!A$} په طبيعي
استنتاج کښې هله قضيه ده چې داسې اشتقاق موجود وي چې \LR{$!A$}
وروستۍ تړلے فارمول وي او ټول فرضونه ايستل شوي وي۔ مثلاً دا
اشتقاق ښيي چې \LR{$\Proves (!A \land !B) \lif !A$}:
\begin{prooftree}
  \AxiomC{\LR{$\Discharge{!A \land !B}{1}$}}
  \RightLabel{\Elim{\land}}
  \UnaryInfC{\LR{$!A$}}
  \DischargeRule{\Intro{\lif}}{1}
  \UnaryInfC{\LR{$(!A \land !B) \lif !A$}}
\end{prooftree}
نښه~\LR{$1$} څرګندوي چې فرض \LR{$!A \land !B$} د \Intro{\lif} په استنتاج کښې
ايستل کېږي۔
\label{olpseg:OLP-0066-B008:end}

\phantomsection\label{olpseg:OLP-0066-B009:start}
يو سټ~\LR{$\Gamma$} هله ناسازګار دے چې په طبيعي استنتاج کښې
\LR{$\Gamma \Proves \lfalse$} وي۔ د \FalseInt{} قاعده دا يقيني کوي چې له
ناسازګار سټ څخه هر تړلے فارمول اشتقاق کېدے شي۔
\label{olpseg:OLP-0066-B009:end}

\phantomsection\label{olpseg:OLP-0066-B010:start}
د طبيعي استنتاج نظامونه ګيرهارډ ګېنتڅن او ستانيسواف ياشکوفسکي په
۱۹۳۰مو کلونو کښې جوړ کړل، او وروسته ډاګ پراوېڅ او فريډرېک فېچ لا
پراخ کړل۔ ځکه چې استنتاجونه ئې د ثبوت طبيعي لارې انځوروي، فلسفيان
ئې خوښوي۔ د فېچ جوړې شوې بڼې زياتره د منطق په پېژندنيزو درسي
کتابونو کښې کارول کېږي۔ د منطق په فلسفه کښې کله کله د طبيعي استنتاج
قاعدې د منطقي عملګرو معنٰی ورکوي، چې دې ته ``ثبوت-نظري معنايي
پوهنه'' ويل کېږي۔
\label{olpseg:OLP-0066-B010:end}

\phantomsection\label{olpseg:OLP-0067-B004:start}
\olfileid{fol}{prf}{tab}
\label{olpseg:OLP-0067-B004:end}

\phantomsection\label{olpseg:OLP-0067-B005:start}
\olsection{تابلوګانې}
\label{olpseg:OLP-0067-B005:end}

\phantomsection\label{olpseg:OLP-0067-B006:start}
د اشتقاق ډېر نظامونه د تړلي فارمولونه له ترتيبونو سره کار
کوي، خو تابلوګانې له نښه لرونکي فارمولونه سره کار کوي۔
نښه لرونکے فارمول يوه جوړه ده چې د رښتياوالي له يوې نښې
(\LR{$\True$} يا \LR{$\False$}) او تړلے فارمول څخه جوړه ده:
\[
\sFmla{\True}{!A} \text{\RL{ يا }} \sFmla{\False}{!A}.
\]
تابلو له هغو نښه لرونکي فارمولونه څخه جوړ وي چې د لاندې خوا
ته څانګېدونکې ونه کښې ترتيب شوي وي۔ دا له يو شمېر \emph{فرضونو}
پيلېږي او په هغو نښه لرونکي فارمولونه دوام کوي چې د استنتاج د کومې
قاعدې په پلي کولو تر ځان له پاسه له يوه نښه لرونکي فارمولونه څخه
راوتلي وي۔ هره قاعده اجازه راکوي چې د يوې څانګې په پاے کښې يو يا
څو نښه لرونکي فارمولونه ورزيات کړو، يا دوه نښه لرونکي فارمولونه يو د
بل تر څنګ ورزيات کړو؛ په وروستي حالت کښې څانګه په دوو وېشل کېږي او
دوه ورزيات شوي نښه لرونکي فارمولونه د دواړو څانګو سرونه جوړوي۔
\label{olpseg:OLP-0067-B006:end}

\phantomsection\label{olpseg:OLP-0067-B007:start}
پر يوه مرکب نښه لرونکے فارمول د قاعدې له پلي کولو څخه داسې
نښه لرونکي فارمولونه ورزياتېږي چې سمدستي فرعي-فارمولونه وي۔ د هرې
قاعدې يوه جوړه راځي، د دواړو نښو هرې يوې دپاره يوه قاعده۔ مثلاً د
\LR{$\TRule{\True}{\land}$} قاعده پر \LR{$\sFmla{\True}{!A \land !B}$} پلې
کېږي او اجازه راکوي چې د هرې هغې څانګې په پاے کښې دواړه
نښه لرونکي فارمولونه \LR{$\sFmla{\True}{!A}$} او~\LR{$\sFmla{\True}{!B}$}
ورزيات کړو چې \LR{$\sFmla{\True}{!A \land !B}$} لري؛ او د
\LR{$\TRule{\False}{\land}$} قاعده اجازه راکوي چې يوه څانګه د
\LR{$\sFmla{\False}{!A}$} او \LR{$\sFmla{\False}{!B}$} يو د بل تر څنګ په
ورزياتولو ووېشو۔
\textbf{د ژباړې د نسخې سمون (OLPF-002):} په سرچينه کښې د دروغې
عطف قاعدې دويم ارګومېنټ ټول مرکب فارمول دے، خو د رښتوني عطف موازي
قاعده او د لاندې بېلګې ټولې قاعدې هلته يوازې نښلوونکے اخلي۔ دلته
هماغه د عطف نښلوونکے کارول شوے، څو د قاعدې نښه په يو شان بڼه جوړه
شي۔
تابلو هله تړلے دے چې د هغه هره څانګه د نښه لرونکي فارمولونه
سره برابره جوړه \LR{$\sFmla{\True}{!A}$} او \LR{$\sFmla{\False}{!A}$} ولري۔
\label{olpseg:OLP-0067-B007:end}

\phantomsection\label{olpseg:OLP-0067-B008:start}
پر تابلوګانې ولاړه \LR{$\Proves$} اړيکه داسې تعريفېږي:
\LR{$\Gamma \Proves !A$} هله او يوازې هله چې داسې کوم متناهي سټ
~\LR{$\Gamma_0 = \{!B_1, \dots, !B_n\} \subseteq \Gamma$} موجود وي چې د
لاندې فرضونو دپاره يو تړلے تابلو موجود وي:
\[
\{\sFmla{\False}{!A}, \sFmla{\True}{!B_1}, \dots, \sFmla{\True}{!B_n}\} 
\]
مثلاً دا تړلے تابلو ښيي چې \LR{$\Proves (!A \land !B) \lif !A$}:
\begin{oltableau}
  [\sFmla{\False}{(\formula{A} \land \formula{B}) \lif \formula{A}}, just = \TAss
    [\sFmla{\True}{\formula{A} \land \formula{B}}, just = {\TRule{\False}{\lif}[1]}
      [\sFmla{\False}{\formula{A}}, just = {\TRule{\False}{\lif}[1]}
        [\sFmla{\True}{\formula{A}}, just = {\TRule{\True}{\lif}[2]}
          [\sFmla{\True}{\formula{B}}, just = {\TRule{\True}{\lif}[2]}, close]
        ]
      ]
    ]
  ]
\end{oltableau}
\label{olpseg:OLP-0067-B008:end}

\phantomsection\label{olpseg:OLP-0067-B009:start}
يو سټ \LR{$\Gamma$} د تابلو په حساب کښې هله ناسازګار دے چې د
لاندې فرضونو دپاره يو تړلے تابلو موجود وي:
\[
\{\sFmla{\True}{!B_1}, \dots, \sFmla{\True}{!B_n}\} 
\]
د کوم \LR{$!B_i \in \Gamma$} دپاره۔
\label{olpseg:OLP-0067-B009:end}

\phantomsection\label{olpseg:OLP-0067-B010:start}
تابلوګانې په ۱۹۵۰مو کلونو کښې اېورټ بېت او ياکو هېنتيکا هر يوه
په خپلواکه توګه جوړ کړل، او رېمونډ سموليان ساده او عام کړل۔ دا کارول
ډېر اسانه دي، ځکه چې د تابلو جوړول ډېره منظمه طريقه ده۔ د
تابلوګانې د منظم طبيعت له امله په کمپيوټر کښې هم په اسانۍ پلي
کېږي۔ خو تابلو زياتره په سختۍ لوستل کېږي او له ثبوتونو سره ئې
تړاو کله کله په اسانۍ نۀ ښکاري۔ دا تګلاره هم ډېره عامه ده، او ډېر
بېل منطقونه د تابلو نظامونه لري۔ تابلوګانې له موږ سره د
هغو جوړښتونه په موندلو کښې هم مرسته کوي چې ورکړل شوي
تړلي فارمولونه، يا د هغو سټونه، صادقوي: که سټ د صدق وړ وي، تړلے
تابلو به و نۀ لري؛ يعنې هر تابلو به يوه پرانيستې څانګه
ولري۔ په دې شرط چې پر هغې څانګه هره ممکنه قاعده پلې شوې وي،
صادقوونکے جوړښت له يوې پرانيستې څانګې څخه ``لوستل کېدے''
شي۔ له سېکوېنټ حساب سره هم ډېر نږدې تړاو شته: په اصل کښې يو تړلے
تابلو د سېکوېنټ په حساب کښې يو لنډ شوے اشتقاق دے چې
سرچپه ليکل شوے وي۔
\label{olpseg:OLP-0067-B010:end}

\phantomsection\label{olpseg:OLP-0068-B004:start}
\olfileid{fol}{prf}{axd}
\label{olpseg:OLP-0068-B004:end}

\phantomsection\label{olpseg:OLP-0068-B005:start}
\olsection{بديهي اشتقاقونه}
\label{olpseg:OLP-0068-B005:end}

\phantomsection\label{olpseg:OLP-0068-B006:start}
بديهي اشتقاقونه د اشتقاق تر ټولو زاړۀ او ساده منطقي
نظامونه دي۔ د هغو اشتقاقونه يوازې د تړلي فارمولونه لړۍ دي۔ د
تړلي فارمولونه يوه لړۍ هله سم اشتقاق ګڼل کېږي چې په هغې کښې
هره تړلے فارمول~\LR{$!A$} له لاندې شرطونو څخه يو پوره کړي:
\begin{enumerate}
\item \LR{$!A$} بديهي اصل وي؛ يا
\item \LR{$!A$} د تړلي فارمولونه د ورکړل شوي سټ~\LR{$\Gamma$} يو غړے
  وي؛ يا
\item \LR{$!A$} د استنتاج په کومه قاعده توجيه شوے وي۔
\end{enumerate}
د بديهي اصل کېدو دپاره \LR{$!A$} بايد د يو شمېر ټاکلو تړلے فارمول
شېماوو له يوې سره بڼه ولري۔ د بديهي اصلونو د شېماوو ډېر داسې سټونه
شته چې د لومړۍ درجې منطق دپاره د اشتقاق د قناعت وړ، سم او
بشپړ نظام ورکوي۔ ځينې د هغو نښلوونکو له مخې ترتيب شوي وي چې واک پرې
چلوي؛ مثلاً دا شېماوې
\[
!A \lif (!B \lif !A) \qquad !B \lif (!B \lor !C) \qquad (!B \land !C) \lif !B
\]
هغه عام بديهي اصول دي چې پر \LR{$\lif$}، \LR{$\lor$} او~\LR{$\land$} واک چلوي۔ د
بديهي اصولو ځينې نظامونه غواړي د اصولو شمېر تر ټولو لږ وي۔ دا چې
کوم نښلوونکي بنسټيز وټاکل شي، ان داسې نظامونه هم موندل کېدے شي چې
يوازې له يوه بديهي اصل څخه جوړ وي۔
\label{olpseg:OLP-0068-B006:end}

\phantomsection\label{olpseg:OLP-0068-B007:start}
د استنتاج قاعده يوه شرطي وينا ده چې په اشتقاق کښې د
تړلے فارمول د توجيه کېدو دپاره بس شرط ورکوي۔ مودوس پوننس يوه ډېره
عامه داسې قاعده ده: وايي که \LR{$!A$} او \LR{$!A \lif !B$} له مخکښې توجيه
شوي وي، نو \LR{$!B$} توجيه شوے دے۔ دا معنٰی لري چې په اشتقاق
کښې هغه کرښه چې تړلے فارمول~\LR{$!B$} لري، په دې شرط توجيه شوې ده چې
\LR{$!A$} او \LR{$!A \lif !B$} دواړه، د کوم تړلے فارمول~\LR{$!A$} دپاره، په
اشتقاق کښې له~\LR{$!B$} څخه مخکښې راغلي وي۔
\label{olpseg:OLP-0068-B007:end}

\phantomsection\label{olpseg:OLP-0068-B008:start}
پر بديهي اشتقاقونه ولاړه \LR{$\Proves$} اړيکه داسې تعريفېږي:
\LR{$\Gamma \Proves !A$} هله او يوازې هله چې داسې اشتقاق موجود
وي چې تړلے فارمول~\LR{$!A$} ئې وروستے فارمول وي، او \LR{$\Gamma$} په هغه
اشتقاق کښې د هغو تړلي فارمولونه سټ ونيول شي چې د پورته (۲) له
مخې توجيه شوي وي۔ \LR{$!A$} هله قضيه ده چې \LR{$!A$} داسې اشتقاق
ولري چې~\LR{$\Gamma$} پکې تش وي؛ يعنې په اشتقاق کښې هره
تړلے فارمول يا د (۱) له مخې توجيه شوې وي يا د~(۳) له مخې۔ مثلاً دا
اشتقاق ښيي چې \LR{$\Proves !A  \lif (!B \lif (!B \lor !A))$}:
\begin{derivation}
  1. & \LR{$!B \lif (!B \lor !A)$} \\
  2. & \LR{$(!B \lif (!B \lor !A)) \lif (!A  \lif (!B \lif (!B \lor !A)))$}\\
  3. & \LR{$!A  \lif (!B \lif (!B \lor !A))$}
\end{derivation}
په ۱مه کرښه تړلے فارمول د بديهي اصل \LR{$!A \lif (!A \lor !B)$} بڼه
لري، چې د \LR{$!A$} او \LR{$!B$} ونډې پکې بدلې شوې دي۔ په ۲مه کرښه جمله د
بديهي اصل \LR{$!A \lif (!B \lif !A)$} بڼه لري۔ نو دواړه کرښې توجيه شوې
دي۔ ۳مه کرښه د مودوس پوننس له مخې توجيه شوې ده: که په لنډه ئې
\LR{$!D$} وليکو، نو ۲مه کرښه د \LR{$!C \lif !D$} بڼه لري، چې \LR{$!C$} هم
\LR{$!B \lif (!B \lor !A)$}، يعنې ۱مه کرښه، ده۔
\label{olpseg:OLP-0068-B008:end}

\phantomsection\label{olpseg:OLP-0068-B009:start}
يو سټ \LR{$\Gamma$} هله ناسازګار دے چې \LR{$\Gamma \Proves \lfalse$}۔ د
بديهي اصولو يو بشپړ نظام به دا هم ثابتوي چې د هر~\LR{$!A$} دپاره
\LR{$\lfalse \lif !A$}؛ نو که \LR{$\Gamma$} ناسازګار وي، د هر~\LR{$!A$} دپاره
\LR{$\Gamma \Proves !A$}۔
\label{olpseg:OLP-0068-B009:end}

\phantomsection\label{olpseg:OLP-0068-B010:start}
د منطق د بديهي اشتقاقونه نظامونه لومړے ګوټلوب فرېګه په خپل
۱۸۷۹ز اثر \emph{بېګريفشريفت} کښې وړاندې کړل، چې له همدې امله زياتره
د معاصر منطق لومړے اثر ګڼل کېږي۔ بيا الفرېډ نارت وايټ هېډ او برټرېنډ
رسل په \emph{پرېنسيپيا ماتيماتيکا} کښې، او ډيوېډ هيلبرټ او شاګردانو
ئې په ۱۹۲۰مو کلونو کښې بشپړ کړل۔ ځکه ورته زياتره ``د فرېګه نظامونه''
يا ``د هيلبرټ نظامونه'' ويل کېږي۔ دا ډېر انعطاف منونکي دي، ځکه چې د
يو منطق دپاره بديهي نظام زياتره په اسانۍ موندل کېږي۔ دا چې
اشتقاقونه ډېر ساده جوړښت او يوازې يوه يا دوه استنتاجي قاعدې
لري، د هغو \emph{په اړه} د څيزونو ثابتول هم نسبتاً اسانه دي۔ خو په
عمل کښې ئې کارول ډېر ګران دي؛ يعنې ثبوتونه موندل او ليکل ګران دي۔
\label{olpseg:OLP-0068-B010:end}

\phantomsection\label{olpseg:OLP-0069-B004:start}
\olchapter{fol}{seq}{د سېکوېنټ حساب}
\label{olpseg:OLP-0069-B004:end}

\phantomsection\label{olpseg:OLP-0069-B005:start}
\begin{editorial}
  دا څپرکے د کلاسیک د لومړۍ درجې منطق دپاره د ګېنتڅن معياري
  سېکوېنټ حساب \texttt{LK} وړاندې کوي۔ دې ته نور مثالونه او تمرينونه
  هم پکار دي۔ د ثبوت د نظام په توګه د سېکوېنټ له حساب سره د اړوندو
  موادو د شاملولو يا اېستلو دپاره د \texttt{prfLK} ټېګ وکاروئ۔
\end{editorial}
\label{olpseg:OLP-0069-B005:end}

\phantomsection\label{olpseg:OLP-0069-B008:start}
%


\label{olpseg:OLP-0069-B008:end}

\phantomsection\label{olpseg:OLP-0069-B012:start}
%


\label{olpseg:OLP-0069-B012:end}

\phantomsection\label{olpseg:OLP-0069-B016:start}
%


\label{olpseg:OLP-0069-B016:end}

\phantomsection\label{olpseg:OLP-0069-B018:start}
%



\label{olpseg:OLP-0069-B018:end}

\phantomsection\label{olpseg:OLP-0070-B004:start}
\olfileid{fol}{seq}{rul}
\label{olpseg:OLP-0070-B004:end}

\phantomsection\label{olpseg:OLP-0070-B005:start}
\olsection{قاعدې او اشتقاقونه}
\label{olpseg:OLP-0070-B005:end}

\phantomsection\label{olpseg:OLP-0070-B006:start}
د راتلونکو خبرو دپاره، پرېږدئ چې \LR{$\Gamma, \Delta, \Pi, \Lambda$} د
تړلي فارمولونه متناهي لړۍ وښيي۔
\label{olpseg:OLP-0070-B006:end}

\phantomsection\label{olpseg:OLP-0070-B007:start}
\begin{defn}[سېکوېنټ]
\emph{سېکوېنټ} د لاندې بڼې څرګندنه ده:
\[
\Gamma \Sequent \Delta
\]
دلته \LR{$\Gamma$} او \LR{$\Delta$} د ژبې \LR{$\Lang L$} د تړلي فارمولونه متناهي
(او کېدے شي تشې) لړۍ دي۔ \LR{$\Gamma$} ته \emph{مقدم} او \LR{$\Delta$} ته
\emph{تالي} ويل کېږي۔
\end{defn}
\label{olpseg:OLP-0070-B007:end}

\phantomsection\label{olpseg:OLP-0070-B008:start}
\begin{explain}
د يوه سېکوېنټ تر شا شهودي مفکوره دا ده: که د مقدم ټولې تړلي فارمولونه
صدق کوي، نو د تالي لږ تر لږه يوه تړلي فارمولونه صدق کوي۔ يعنې که
\LR{$\Gamma = \tuple{!A_1, \dots, !A_m}$} او
\LR{$\Delta = \tuple{!B_1, \dots, !B_n}$} وي، نو \LR{$\Gamma \Sequent \Delta$}
هله صدق کوي چې
\[
(!A_1 \land \cdots \land !A_m) \lif (!B_1 \lor \cdots \lor
!B_n)
\]
صدق کوي۔ دوه خاص حالتونه شته: يو هغه چې \LR{$\Gamma$} تش وي، او بل هغه
چې \LR{$\Delta$}~تش وي۔ کله چې \LR{$\Gamma$} تش وي، يعنې \LR{$m = 0$}، نو \LR{$\quad
\Sequent \Delta$} هله صدق کوي چې \LR{$!B_1 \lor \dots \lor !B_n$} صدق
کوي۔ کله چې \LR{$\Delta$} تش وي، يعنې \LR{$n = 0$}، نو \LR{$\Gamma \Sequent \quad$}
هله صدق کوي چې \LR{$\lnot(!A_1 \land \dots \land !A_m)$} صدق کوي۔
موږ يو سېکوېنټ هله معتبر بولو چې ورسره سمون لرونکې تړلے فارمول
معتبره وي۔
\end{explain}
\label{olpseg:OLP-0070-B008:end}

\phantomsection\label{olpseg:OLP-0070-B009:start}
که \LR{$\Gamma$} د تړلي فارمولونه يوه لړۍ وي، نو \LR{$\Gamma, !A$} د هغې
پايلې دپاره ليکو چې~\LR{$!A$} د~\LR{$\Gamma$} ښي پاے ته ورزيات شي (او \LR{$!A,
\Gamma$} د هغې پايلې دپاره چې~\LR{$!A$} د~\LR{$\Gamma$} کيڼ پاے ته ورزيات
شي)۔ که \LR{$\Delta$} هم د تړلي فارمولونه يوه لړۍ وي، نو \LR{$\Gamma,
\Delta$} د دواړو لړيو نښلول دي۔
\label{olpseg:OLP-0070-B009:end}

\phantomsection\label{olpseg:OLP-0070-B010:start}
\begin{defn}[لومړنے سېکوېنټ]
\emph{لومړنے سېکوېنټ} داسې سېکوېنټ دے
چې له لاندې بڼو څخه يوه لري:
  \begin{enumerate}
    \item \LR{$!A \Sequent !A$}
    \item \LR{$\quad \Sequent \ltrue$}
    \item \LR{$\lfalse \Sequent \quad$}
  \end{enumerate}، او \LR{$!A$} د ژبې هره تړلے فارمول
کېدے شي۔
\end{defn}
\label{olpseg:OLP-0070-B010:end}

\phantomsection\label{olpseg:OLP-0070-B011:start}
په سېکوېنټ حساب کښې اشتقاقونه د سېکوېنټونو داسې ټاکلې ونې
دي چې تر ټولو بره سېکوېنټونه ئې لومړني سېکوېنټونه وي؛ او که يو
سېکوېنټ د يوه يا دوو نورو سېکوېنټونو لاندې راشي، بايد د استنتاج د
يوې قاعدې له مخې ترې په سمه توګه راووځي۔ د~\LR{$\Log{LK}$} قاعدې پر دوو
اصلي ډولونو وېشل کېږي: \emph{منطقي} قاعدې او \emph{جوړښتي} قاعدې۔
منطقي قاعدې د لاندې سېکوېنټ د هغې تړلے فارمول د اصلي عملګر په نوم نومول کېږي چې \LR{$!A$} او/يا \LR{$!B$} لري۔ هره قاعده
دوه بڼې لري: يوه د داسې سېکوېنټ د استنتاج دپاره چې د
منطقي عملګر لرونکې تړلے فارمول ئې کيڼ لوري ته وي، او بله د هغه
سېکوېنټ دپاره چې دا تړلے فارمول ئې ښي لوري ته وي۔
\label{olpseg:OLP-0070-B011:end}

\phantomsection\label{olpseg:OLP-0071-B004:start}
\olfileid{fol}{seq}{prl}
\label{olpseg:OLP-0071-B004:end}

\phantomsection\label{olpseg:OLP-0071-B005:start}
\olsection{بياني قاعدې}
\label{olpseg:OLP-0071-B005:end}

\phantomsection\label{olpseg:OLP-0071-B006:start}
\subsection{د \LR{$\lnot$} دپاره قاعدې}
\label{olpseg:OLP-0071-B006:end}

\phantomsection\label{olpseg:OLP-0071-B007:start}
\begin{defish}
\Axiom$ \Gamma \fCenter \Delta, !A $
\RightLabel{\LeftR{\lnot}}
\UnaryInf$ \lnot !A, \Gamma \fCenter \Delta$
\DisplayProof
\hfill
\Axiom$!A, \Gamma \fCenter \Delta$
\RightLabel{\RightR{\lnot}}
\UnaryInf$ \Gamma \fCenter \Delta, \lnot !A $
\DisplayProof
\end{defish}
\label{olpseg:OLP-0071-B007:end}

\phantomsection\label{olpseg:OLP-0071-B008:start}
\subsection{د \LR{$\land$} دپاره قاعدې}
\label{olpseg:OLP-0071-B008:end}

\phantomsection\label{olpseg:OLP-0071-B009:start}
\begin{defish}\noindent
\begin{tabular}{l}
\Axiom$ !A, \Gamma \fCenter \Delta$
\RightLabel{\LeftR{\land}}
\UnaryInf$ !A \land !B, \Gamma \fCenter \Delta$
\DisplayProof
\\[3ex]
\Axiom$!B, \Gamma \fCenter \Delta$
\RightLabel{\LeftR{\land}}
\UnaryInf$!A \land !B, \Gamma \fCenter \Delta$
\DisplayProof
\end{tabular}
\hfill
\Axiom$\Gamma \fCenter \Delta, !A$
\Axiom$ \Gamma \fCenter \Delta, !B$
\RightLabel{\RightR{\land}}
\BinaryInf$ \Gamma \fCenter \Delta, !A \land !B $
\DisplayProof
\end{defish}
\label{olpseg:OLP-0071-B009:end}

\phantomsection\label{olpseg:OLP-0071-B010:start}
\subsection{د \LR{$\lor$} دپاره قاعدې}
\label{olpseg:OLP-0071-B010:end}

\phantomsection\label{olpseg:OLP-0071-B011:start}
\begin{defish}
\Axiom$!A, \Gamma \fCenter \Delta$
\Axiom$!B, \Gamma \fCenter \Delta$
\RightLabel{\LeftR{\lor}}
\BinaryInf$!A \lor !B, \Gamma \fCenter \Delta$
\DisplayProof
\hfill
\begin{tabular}{r}
\Axiom$\Gamma \fCenter \Delta, !A$
\RightLabel{\RightR{\lor}}
\UnaryInf$ \Gamma \fCenter \Delta, !A \lor !B$
\DisplayProof
\\[3ex]
\Axiom$ \Gamma \fCenter \Delta, !B$
\RightLabel{\RightR{\lor}}
\UnaryInf$ \Gamma \fCenter \Delta, !A \lor !B$
\DisplayProof
\end{tabular}
\end{defish}
\label{olpseg:OLP-0071-B011:end}

\phantomsection\label{olpseg:OLP-0071-B012:start}
\subsection{د \LR{$\lif$} دپاره قاعدې}
\label{olpseg:OLP-0071-B012:end}

\phantomsection\label{olpseg:OLP-0071-B013:start}
\begin{defish}
\Axiom$ \Gamma \fCenter \Delta, !A$
\Axiom$ !B, \Pi \fCenter \Lambda$
\RightLabel{\LeftR{\lif}}
\BinaryInf$ !A \lif !B, \Gamma, \Pi \fCenter \Delta, \Lambda$
\DisplayProof
\hfill
\Axiom$ !A, \Gamma \fCenter \Delta, !B$
\RightLabel{\RightR{\lif}}
\UnaryInf$ \Gamma \fCenter \Delta, !A \lif !B $
\DisplayProof
\end{defish}
\label{olpseg:OLP-0071-B013:end}

\phantomsection\label{olpseg:OLP-0072-B004:start}
\olfileid{fol}{seq}{qrl}
\label{olpseg:OLP-0072-B004:end}

\phantomsection\label{olpseg:OLP-0072-B005:start}
\olsection{د سورونو قاعدې}
\label{olpseg:OLP-0072-B005:end}

\phantomsection\label{olpseg:OLP-0072-B006:start}
\subsection{د \LR{$\lforall$} دپاره قاعدې}
\label{olpseg:OLP-0072-B006:end}

\phantomsection\label{olpseg:OLP-0072-B007:start}
\begin{defish}
\Axiom$ !A(t), \Gamma \fCenter \Delta$
\RightLabel{\LeftR{\lforall}}
\UnaryInf$ \lforall[x][!A(x)],\Gamma \fCenter \Delta$
\DisplayProof
\hfill
\Axiom$ \Gamma \fCenter \Delta, !A(a) $
\RightLabel{\RightR{\lforall}}
\UnaryInf$ \Gamma \fCenter \Delta, \lforall[x][!A(x)]$
\DisplayProof
\end{defish}
\label{olpseg:OLP-0072-B007:end}

\phantomsection\label{olpseg:OLP-0072-B008:start}
په \LeftR{\lforall} کښې \LR{$t$} يوه تړلې اصطلاح ده (يعنې داسې اصطلاح
چې متغيرونه نۀ لري)۔ په \RightR{\lforall} کښې \LR{$a$}~ثابت سمبول دے
چې بايد د \RightR{\lforall} قاعدې په لاندې سېکوېنټ کښې هېڅ ځاے نۀ
وي راغلے۔ موږ \LR{$a$} د \RightR{\forall} استنتاج \emph{ځانګړے متغير}
بولو۔\footnote{موږ د “ځانګړي متغير” اصطلاح کاروو، که څه هم په
پورتنۍ قاعده کښې \LR{$a$} ثابت سمبول دے۔ دا نوم تاريخي لاملونه لري۔}
\label{olpseg:OLP-0072-B008:end}

\phantomsection\label{olpseg:OLP-0072-B009:start}
\subsection{د \LR{$\lexists$} دپاره قاعدې}
\label{olpseg:OLP-0072-B009:end}

\phantomsection\label{olpseg:OLP-0072-B010:start}
\begin{defish}
\Axiom$ !A(a), \Gamma \fCenter \Delta $
\RightLabel{\LeftR{\lexists}}
\UnaryInf$ \lexists[x][!A(x)], \Gamma \fCenter \Delta$
\DisplayProof
\hfill
\Axiom$ \Gamma \fCenter \Delta, !A(t) $
\RightLabel{\RightR{\lexists}}
\UnaryInf$ \Gamma \fCenter \Delta, \lexists[x][!A(x)]$
\DisplayProof
\end{defish}
\label{olpseg:OLP-0072-B010:end}

\phantomsection\label{olpseg:OLP-0072-B011:start}
بيا هم، \LR{$t$}~يوه تړلې اصطلاح ده، او \LR{$a$}~داسې ثابت سمبول دے چې
د \LeftR{\lexists} قاعدې په لاندې سېکوېنټ کښې نۀ راځي۔ موږ \LR{$a$} د
\LeftR{\lexists} استنتاج \emph{ځانګړے متغير} بولو۔
\label{olpseg:OLP-0072-B011:end}

\phantomsection\label{olpseg:OLP-0072-B012:start}
دا شرط چې ځانګړے متغير د \RightR{\lforall} يا
\LeftR{\lexists} استنتاج په لاندې سېکوېنټ کښې نۀ راځي،
\emph{د ځانګړي متغير شرط} بلل کېږي۔
\label{olpseg:OLP-0072-B012:end}

\phantomsection\label{olpseg:OLP-0072-B013:start}
\begin{explain}
هغه دود راياد کړئ چې کله \LR{$!A$} داسې فارمول وي چې
متغير~\LR{$x$} پکې ازاد وي، نو دا په~\LR{$!A(x)$} ليکلو ښيو۔ په همدې
چوکاټ کښې بيا \LR{$!A(t)$} د~\LR{$\Subst{!A}{t}{x}$} لنډون دے۔ نو د
\RightR{\lexists} قاعده داسې هم ليکلے شو:
\begin{prooftree}
  \Axiom$\Gamma \fCenter \Delta, \Subst{!A}{t}{x}$
  \RightLabel{\RightR{\lexists}}
  \UnaryInf$\Gamma \fCenter \Delta, \lexists[x][!A]$
\end{prooftree}
پام وکړئ چې \LR{$t$} کېدے شي له مخکښې په~\LR{$!A$} کښې راغلے وي؛ مثلاً
\LR{$!A$}~کېدے شي~\LR{$\Atom{\Obj P}{t,x}$} وي۔ له دې امله، له
~\LR{$\Gamma \Sequent \Delta, \Atom{\Obj P}{t,t}$} څخه
\LR{$\Gamma \Sequent \Delta, \lexists[x][\Atom{\Obj P}{t,x}]$} استنتاجول
د~\RightR{\lexists} سمه کارونه ده---تاسو د مقدمې د~\LR{$t$} يوه يا
زياتې پېښې په تړلي متغير~\LR{$x$} “بدلولے” شئ، او اړينه نۀ ده
چې ټولې پېښې ئې بدلې کړئ۔ خو د \RightR{\lforall} او
~\LeftR{\lexists} د ځانګړي متغير شرطونه غواړي چې
ثابت سمبول~\LR{$a$} په~\LR{$!A$} کښې نۀ راځي۔ نو د~\LR{$\RightR{\lforall}$} په
کارولو سره له \LR{$\Gamma \Sequent \Delta, \Atom{\Obj P}{a,a}$} څخه
\LR{$\Gamma \Sequent \Delta, \lforall[x][\Atom{\Obj P}{a,x}]$} په سمه
توګه نۀ شئ استنتاجولے۔
\end{explain}
\label{olpseg:OLP-0072-B013:end}

\phantomsection\label{olpseg:OLP-0072-B014:start}
\begin{explain}
په \RightR{\lexists} او \LeftR{\lforall} کښې پر اصطلاح~\LR{$t$} هېڅ
بنديز نشته۔ بل لوري ته، د \LeftR{\lexists} او
\RightR{\lforall} په قاعدو کښې د ځانګړي متغير شرط غواړي چې
ثابت سمبول~\LR{$a$} په پورتني سېکوېنټ کښې د~\LR{$!A(a)$} څخه بهر هېڅ ځاے
نۀ راځي۔ دا شرط د نظام د سموالي د يقيني کولو دپاره اړين دے؛ يعنې
نظام يوازې معتبر سېکوېنټونه اشتقاق کړي۔ له دې شرط پرته لاندې
استنتاجونه روا کېدل:
\begin{prooftree}
  \Axiom$!A(a) \fCenter !A(a)$
  \RightLabel{*\LeftR{\lexists}}
  \UnaryInf$\lexists[x][!A(x)] \fCenter !A(a)$
  \RightLabel{\RightR{\lforall}}
  \UnaryInf$\lexists[x][!A(x)] \fCenter \lforall[x][!A(x)]$
  \DisplayProof\bottomAlignProof
  \qquad
  \Axiom$!A(a) \fCenter !A(a)$
  \RightLabel{*\RightR{\lforall}}
  \UnaryInf$!A(a) \fCenter \lforall[x][!A(x)]$
  \RightLabel{\LeftR{\lexists}}
  \UnaryInf$\lexists[x][!A(x)] \fCenter \lforall[x][!A(x)]$
\end{prooftree}
خو \LR{$\lexists[x][!A(x)] \Sequent \lforall[x][!A(x)]$} معتبر نۀ دے۔
\end{explain}
\label{olpseg:OLP-0072-B014:end}

\phantomsection\label{olpseg:OLP-0073-B004:start}
\olfileid{fol}{seq}{srl}
\label{olpseg:OLP-0073-B004:end}

\phantomsection\label{olpseg:OLP-0073-B005:start}
\olsection{جوړښتي قاعدې}
\label{olpseg:OLP-0073-B005:end}

\phantomsection\label{olpseg:OLP-0073-B006:start}
موږ څو داسې قاعدو ته هم اړتيا لرو چې د يوه سېکوېنټ په کيڼ او ښي
لوري کښې د تړلي فارمولونه د بيا ترتيبولو اجازه راکړي۔ ځکه چې منطقي
قاعدې غواړي هغه تړلي فارمولونه چې قاعده پرې کار کوي، يا په بشپړ کيڼ
او يا په بشپړ ښي پاے کښې وي، نو د “تبادلې” داسې قاعدې ته اړتيا
لرو چې تړلي فارمولونه سم ځاے ته يوسي۔ کله کله دا هم مهمه وي چې دوه
يو شان تړلي فارمولونه سره يو کړو او هر لوري ته تړلے فارمول
ورزياته کړو۔
\label{olpseg:OLP-0073-B006:end}

\phantomsection\label{olpseg:OLP-0073-B007:start}
\subsection{کمزوري کول}
\label{olpseg:OLP-0073-B007:end}

\phantomsection\label{olpseg:OLP-0073-B008:start}
\begin{defish}
\Axiom$ \Gamma \fCenter \Delta $
\RightLabel{\LeftR{\Weakening}}
\UnaryInf$ !A, \Gamma \fCenter \Delta$
\DisplayProof
\hfill
\Axiom$ \Gamma \fCenter \Delta$
\RightLabel{\RightR{\Weakening}}
\UnaryInf$ \Gamma \fCenter \Delta, !A$
\DisplayProof
\end{defish}
\label{olpseg:OLP-0073-B008:end}

\phantomsection\label{olpseg:OLP-0073-B009:start}
\subsection{انقباض}
\label{olpseg:OLP-0073-B009:end}

\phantomsection\label{olpseg:OLP-0073-B010:start}
\begin{defish}
\Axiom$ !A, !A, \Gamma \fCenter \Delta $
\RightLabel{\LeftR{\Contraction}}
\UnaryInf$ !A, \Gamma \fCenter \Delta$
\DisplayProof
\hfill
\Axiom$ \Gamma \fCenter \Delta, !A, !A$
\RightLabel{\RightR{\Contraction}}
\UnaryInf$ \Gamma \fCenter \Delta, !A$
\DisplayProof
\end{defish}
\label{olpseg:OLP-0073-B010:end}

\phantomsection\label{olpseg:OLP-0073-B011:start}
\subsection{تبادله}
\label{olpseg:OLP-0073-B011:end}

\phantomsection\label{olpseg:OLP-0073-B012:start}
\begin{defish}
\Axiom$ \Gamma, !A, !B, \Pi \fCenter \Delta $
\RightLabel{\LeftR{\Exchange}}
\UnaryInf$ \Gamma, !B, !A, \Pi \fCenter \Delta$
\DisplayProof
\hfill
\Axiom$ \Gamma \fCenter \Delta, !A, !B, \Lambda$
\RightLabel{\RightR{\Exchange}}
\UnaryInf$ \Gamma \fCenter \Delta, !B, !A, \Lambda$
\DisplayProof
\end{defish}
\label{olpseg:OLP-0073-B012:end}

\phantomsection\label{olpseg:OLP-0073-B013:start}
د کمزوري کولو، انقباض او تبادلې د استنتاجونو يوه لړۍ به زياتره د
استنتاج په دوو کرښو ښوول کېږي۔
\label{olpseg:OLP-0073-B013:end}

\phantomsection\label{olpseg:OLP-0073-B014:start}
لاندې قاعده چې “پرېکون” بلل کېږي، په کره توګه اړينه نۀ ده، خو د
اشتقاقونه بيا کارول او يو ځاے کول ډېر اسانوي۔
\begin{defish}
\[
\Axiom$ \Gamma \fCenter \Delta, !A$
\Axiom$ !A, \Pi \fCenter \Lambda $
\RightLabel{\Cut}
\BinaryInf$ \Gamma, \Pi \fCenter \Delta, \Lambda$
\DisplayProof
\]
\end{defish}
\label{olpseg:OLP-0073-B014:end}

\phantomsection\label{olpseg:OLP-0074-B004:start}
\olfileid{fol}{seq}{der}
\label{olpseg:OLP-0074-B004:end}

\phantomsection\label{olpseg:OLP-0074-B005:start}
\olsection{اشتقاقونه}
\label{olpseg:OLP-0074-B005:end}

\phantomsection\label{olpseg:OLP-0074-B006:start}
\begin{explain}
موږ وويل چې لومړنے سېکوېنټ څنګه وي او د استنتاج قاعدې مو ورکړې۔
په سېکوېنټ حساب کښې اشتقاقونه له دغو څخه په استقرايي توګه
جوړېږي: هر اشتقاق يا په يوازې ځان يو لومړنے سېکوېنټ وي،
او يا له يوه يا دوو اشتقاقونه او ورپسې يوه استنتاج څخه جوړ وي۔
\end{explain}
\label{olpseg:OLP-0074-B006:end}

\phantomsection\label{olpseg:OLP-0074-B007:start}
\begin{defn}[\LR{$\Log{LK}$} اشتقاق]
د يوه سېکوېنټ~\LR{$S$} \emph{\LR{$\Log{LK}$}-اشتقاق} د سېکوېنټونو
داسې متناهي ونه ده چې لاندې شرطونه پوره کوي:
\begin{enumerate}
\item د ونې تر ټولو بره سېکوېنټونه لومړني سېکوېنټونه دي۔
\item د ونې تر ټولو لاندې سېکوېنټ~\LR{$S$} دے۔
\item له~\LR{$S$} پرته د ونې هر سېکوېنټ د استنتاج د يوې قاعدې د سمې
  کارونې مقدمه ده، او د هغې کارونې پايله په ونه کښې د هماغه
  سېکوېنټ مخامخ لاندې ولاړه وي۔
\end{enumerate}
بيا وايو چې \LR{$S$} د اشتقاق \emph{وروستنے سېکوېنټ} دے او
\LR{$S$} \emph{په \LR{$\Log{LK}$} کښې د اشتقاق وړ} دے (يا
\LR{$\Log{LK}$}-د اشتقاق وړ دے)۔
\end{defn}
\label{olpseg:OLP-0074-B007:end}

\phantomsection\label{olpseg:OLP-0074-B008:start}
\begin{ex}
هر لومړنے سېکوېنټ، مثلاً \LR{$!C \Sequent !C$}، اشتقاق دے۔
له دې څخه د يوې بلې قاعدې، مثلاً د \LR{$\LeftR{\Weakening}$} قاعدې، په
کارولو نوے اشتقاق ترلاسه کولے شو:
\begin{prooftree}
\Axiom$ \Gamma \fCenter \Delta $
\RightLabel{\LeftR{\Weakening}}
\UnaryInf$ !A, \Gamma \fCenter \Delta$
\end{prooftree}
خو قاعده عمومي مراد لري: په قاعده کښې \LR{$!A$} د هرې تړلے فارمول،
مثلاً د~\LR{$!D$}، په ځاے راوستلے شو۔ که مقدمه زموږ له لومړني سېکوېنټ
\LR{$!C \Sequent !C$} سره سمون خوري، نو \LR{$\Gamma$} او \LR{$\Delta$} دواړه
يوازې~\LR{$!C$} دي، او پايله به بيا \LR{$!D, !C \Sequent !C$} وي۔ نو لاندې
اشتقاق دے:
\begin{prooftree}
\Axiom$ !C \fCenter !C $
\RightLabel{\LeftR{\Weakening}}
\UnaryInf$ !D, !C \fCenter !C$
\end{prooftree}
اوس يوه بله قاعده، مثلاً \LR{$\LeftR{\Exchange}$}، کارولے شو چې په کيڼ
لوري کښې د دوو تړلي فارمولونه د ځاے بدلولو اجازه راکوي۔ نو لاندې هم
سم اشتقاق دے:
\begin{prooftree}
\Axiom$ !C \fCenter !C $
\RightLabel{\LeftR{\Weakening}}
\UnaryInf$ !D, !C \fCenter !C$
\RightLabel{\LeftR{\Exchange}}
\UnaryInf$ !C, !D \fCenter !C$
\end{prooftree}
د دې قاعدې په دې کارونه کښې، چې داسې ورکړل شوې وه:
\begin{prooftree}
\Axiom$ \Gamma, !A, !B, \Pi \fCenter \Delta $
\RightLabel{\LeftR{\Exchange}}
\UnaryInf$ \Gamma, !B, !A, \Pi \fCenter \Delta,$
\end{prooftree}
\LR{$\Gamma$} او \LR{$\Pi$} دواړه تش وو، \LR{$\Delta$} هم \LR{$!C$} دے، او د \LR{$!A$} او
\LR{$!B$} ونډې په همدې ترتيب \LR{$!D$} او~\LR{$!C$} لوبوي۔ نژدې په همدې طريقه
وينو چې
\begin{prooftree}
\Axiom$ !D \fCenter !D $
\RightLabel{\LeftR{\Weakening}}
\UnaryInf$ !C, !D \fCenter !D$
\end{prooftree}
هم اشتقاق دے۔ اوس دا دواړه اشتقاقونه اخلو او د
\LR{$\RightR{\land}$} په کارولو ئې يو ځاے کوو۔ هغه قاعده داسې وه:
\begin{prooftree}
\Axiom$\Gamma \fCenter \Delta, !A$
\Axiom$ \Gamma \fCenter \Delta, !B$
\RightLabel{\RightR{\land}}
\BinaryInf$ \Gamma \fCenter \Delta, !A \land !B $
\end{prooftree}
زموږ په حالت کښې مقدمې بايد د هغو اشتقاقونه له وروستيو
سېکوېنټونو سره سمون وخوري چې په مقدمو پاے ته رسېږي۔ يعنې \LR{$\Gamma$}
د \LR{$!C, !D$}، \LR{$\Delta$} تشه، \LR{$!A$} د \LR{$!C$} او \LR{$!B$} د \LR{$!D$} سره برابره
ده۔ نو که استنتاج سم وي، پايله \LR{$!C, !D \Sequent !C
\land !D$} ده۔
\begin{prooftree}
\Axiom$ !C \fCenter !C $
\RightLabel{\LeftR{\Weakening}}
\UnaryInf$ !D, !C \fCenter !C$
\RightLabel{\LeftR{\Exchange}}
\UnaryInf$ !C, !D \fCenter !C$
\Axiom$ !D \fCenter !D $
\RightLabel{\LeftR{\Weakening}}
\UnaryInf$ !C, !D \fCenter !D$
\RightLabel{\RightR{\land}}
\BinaryInf$ !C, !D \fCenter !C \land !D $
\end{prooftree}
البته مقدمې سرچپه کولے هم شو؛ بيا \LR{$!A$} به \LR{$!D$} او \LR{$!B$} به~\LR{$!C$} وي۔
\begin{prooftree}
\Axiom$ !D \fCenter !D $
\RightLabel{\LeftR{\Weakening}}
\UnaryInf$ !C, !D \fCenter !D$
\Axiom$ !C \fCenter !C $
\RightLabel{\LeftR{\Weakening}}
\UnaryInf$ !D, !C \fCenter !C$
\RightLabel{\LeftR{\Exchange}}
\UnaryInf$ !C, !D \fCenter !C$
\RightLabel{\RightR{\land}}
\BinaryInf$ !C, !D \fCenter !D \land !C $
\end{prooftree}
\end{ex}
\label{olpseg:OLP-0074-B008:end}

\phantomsection\label{olpseg:OLP-0075-B004:start}
\olfileid{fol}{seq}{pro}
\label{olpseg:OLP-0075-B004:end}

\phantomsection\label{olpseg:OLP-0075-B005:start}
\olsection{د اشتقاقونه بېلګې}
\label{olpseg:OLP-0075-B005:end}

\phantomsection\label{olpseg:OLP-0075-B006:start}
\begin{ex}
د سېکوېنټ \LR{$!A \land !B \Sequent !A$} يو
\LR{$\Log{LK}$}-اشتقاق ورکړئ۔
\label{olpseg:OLP-0075-B006:end}

\phantomsection\label{olpseg:OLP-0075-B007:start}
د اشتقاق په لاندې برخه کښې د غوښتل شوي وروستي سېکوېنټ په
ليکلو پيل کوو۔
\begin{prooftree}
\AxiomC{}
\UnaryInf$!A\land !B \fCenter !A$
\end{prooftree}
بيا بايد معلومه کړو چې د کوم ډول استنتاج لاندې سېکوېنټ دا بڼه
لرلے شي۔ دا يوه جوړښتي قاعده کېدے شي، خو غوره ده چې لومړے کومه
منطقي قاعده ولټوؤ۔ په لاندې سېکوېنټ کښې يوازينے منطقي نښلوونکے
\LR{$\land$} دے؛ نو د \LR{$\land$} قاعده لټوؤ، او ځکه چې د \LR{$\land$} نښه په
مقدم کښې راځي، د \LeftR{\land} قاعده ګورو۔
\begin{prooftree}
\AxiomC{}
\RightLabel{\LeftR{\land}}
\UnaryInf$!A\land !B \fCenter !A$
\end{prooftree}
د \LeftR{\land} استنتاج د پورتني سېکوېنټ دپاره دوه امکانونه شته:
يا \LR{$!A \Sequent !A$} او يا \LR{$!B \Sequent !A$}۔ څرګنده ده چې
\LR{$!A \Sequent !A$} يو لومړنے سېکوېنټ دے، او دا ښه خبره ده؛ خو
\LR{$!B \Sequent !A$} په عمومي توګه د اشتقاق وړ نۀ دے۔ پورتنے سېکوېنټ
ورډکوو:
\begin{prooftree}
\Axiom$!A \fCenter !A$
\RightLabel{\LeftR{\land}}
\UnaryInf$!A\land !B \fCenter !A$
\end{prooftree}
اوس د سېکوېنټ \LR{$!A \land !B \Sequent !A$} سم
\LR{$\Log{LK}$}-اشتقاق لرو۔
\end{ex}
\label{olpseg:OLP-0075-B007:end}

\phantomsection\label{olpseg:OLP-0075-B008:start}
\begin{ex}
د سېکوېنټ \LR{$\lnot !A \lor !B \Sequent !A \lif !B$} يو
\LR{$\Log{LK}$}-اشتقاق ورکړئ۔
\label{olpseg:OLP-0075-B008:end}

\phantomsection\label{olpseg:OLP-0075-B009:start}
د اشتقاق په لاندې برخه کښې د غوښتل شوي وروستي سېکوېنټ په
ليکلو پيل وکړئ۔
\begin{prooftree}
\AxiomC{}
\UnaryInf$\lnot !A \lor !B \fCenter !A \lif !B$
\end{prooftree}
د داسې منطقي قاعدې د موندلو دپاره چې دا وروستنے سېکوېنټ راکړي، د
وروستي سېکوېنټ منطقي نښلوونکي ګورو: \LR{$\lnot$}، \LR{$\lor$} او \LR{$\lif$}۔
اوس مهال يوازې \LR{$\lor$} او \LR{$\lif$} مهم دي، ځکه چې د وروستي سېکوېنټ د
تړلي فارمولونه اصلي عملګرونه دي؛ خو \LR{$\lnot$} د بل نښلوونکي د
چاپېريال دننه دے، نو وروسته به ئې وڅېړو۔ له دې امله د وروستي
استنتاج د منطقي قاعدې امکانونه \LeftR{\lor} او \RightR{\lif} دي۔
په رښتيا هره يوه غوره کولے شو، خو \RightR{\lif} غوره کوو، يوازې
د دې دپاره چې په دوو څانګو وېشل لږ وروسته کړو۔ د هغو
\RightR{\lif} استنتاجونو د بڼې له مخې چې دا لاندې سېکوېنټ ورکولے
شي، بڼه بايد داسې وي:
\begin{prooftree}
\AxiomC{}
\UnaryInf$ !A, \lnot !A \lor !B \fCenter !B $
\RightLabel{\RightR{\lif}} \UnaryInf$ \lnot !A \lor !B \fCenter !A \lif !B $
\end{prooftree}
\label{olpseg:OLP-0075-B009:end}

\phantomsection\label{olpseg:OLP-0075-B010:start}
که \LR{$\lnot !A \lor !B$} د مقدم بهرني پاے ته يوسو، د
\LeftR{\lor} قاعده کارولے شو۔ د شېما له مخې دا بايد په لاندې ډول
پر دوو پورتنيو سېکوېنټونو ووېشل شي۔
\textbf{د ژباړې د نسخې سمون (OLSQ-001):} د سرچينې په هغو څلورو
قاعدوي ونو کښې چې د مقدم دوه فارمولونه سره بدلوي، د ښي لوري د
تبادلې نښه راغلې ده؛ دلته د قاعدې له واقعي لوري سره سمه د کيڼ
لوري د تبادلې نښه کارول کېږي۔
\begin{prooftree}
\AxiomC{}
\UnaryInf$\lnot !A, !A \fCenter !B$
\AxiomC{}
\UnaryInf$!B, !A \fCenter !B$
\RightLabel{\LeftR{\lor}}
\BinaryInf$ \lnot !A \lor !B, !A \fCenter !B $
\RightLabel{\LeftR{\Exchange}}
\UnaryInf$ !A, \lnot !A \lor !B \fCenter !B $
\RightLabel{\RightR{\lif}}
\UnaryInf$ \lnot !A \lor !B \fCenter !A \lif !B $
\end{prooftree}
راياد کړئ چې موږ هڅه کوو تر لومړنيو سېکوېنټونو پورې په پورته لوري
لار پيدا کړو؛ داسې ښکاري چې ډېر ورنژدې يو!{} ښۍ څانګه له لومړني
سېکوېنټ څخه يوازې يو کمزوري کول او يوه تبادله لرې ده، او بيا
بشپړېږي:
\begin{prooftree}
\AxiomC{}
\UnaryInf$\lnot !A, !A \fCenter !B$
\Axiom$!B \fCenter !B$
\RightLabel{\LeftR{\Weakening}}
\UnaryInf$!A, !B \fCenter !B$
\RightLabel{\LeftR{\Exchange}}
\UnaryInf$!B, !A \fCenter !B$
\RightLabel{\LeftR{\lor}}
\BinaryInf$\lnot !A \lor !B, !A \fCenter !B $
\RightLabel{\LeftR{\Exchange}}
\UnaryInf$ !A, \lnot !A \lor !B \fCenter !B $
\RightLabel{\RightR{\lif}}
\UnaryInf$ \lnot !A \lor !B \fCenter !A \lif !B $
\end{prooftree}
\label{olpseg:OLP-0075-B010:end}

\phantomsection\label{olpseg:OLP-0075-B011:start}
اوس کيڼه څانګه ګورو۔ په هره تړلے فارمول کښې يوازينے منطقي
نښلوونکے د مقدم د تړلي فارمولونه د \LR{$\lnot$} نښه ده، نو د
\LeftR{\lnot} قاعدې يوه کارونه ګورو۔
\begin{prooftree}
\AxiomC{}
\UnaryInf$ !A \fCenter !B, !A$
\RightLabel{\LeftR{\lnot}}
\UnaryInf$\lnot !A, !A \fCenter !B$
\Axiom$!B \fCenter !B$
\RightLabel{\LeftR{\Weakening}}
\UnaryInf$!A, !B \fCenter !B$
\RightLabel{\LeftR{\Exchange}}
\UnaryInf$!B, !A \fCenter !B$
\RightLabel{\LeftR{\lor}}
\BinaryInf$\lnot !A \lor !B, !A \fCenter !B $
\RightLabel{\LeftR{\Exchange}}
\UnaryInf$ !A, \lnot !A \lor !B \fCenter !B $
\RightLabel{\RightR{\lif}}
\UnaryInf$ \lnot !A \lor !B \fCenter !A \lif !B $
\end{prooftree}
لکه ښۍ څانګه چې مو بشپړه کړه، له دې کيڼې څانګې څخه هم يوازې يو
کمزوري کول او يوه تبادله لرې يو۔
\begin{prooftree}
\Axiom$!A \fCenter !A$
\RightLabel{\RightR{\Weakening}}
\UnaryInf$ !A \fCenter !A, !B$
\RightLabel{\RightR{\Exchange}}
\UnaryInf$ !A \fCenter !B, !A$
\RightLabel{\LeftR{\lnot}}
\UnaryInf$\lnot !A, !A \fCenter !B$
\Axiom$!B \fCenter !B$
\RightLabel{\LeftR{\Weakening}}
\UnaryInf$!A, !B \fCenter !B$
\RightLabel{\LeftR{\Exchange}}
\UnaryInf$!B, !A \fCenter !B$
\RightLabel{\LeftR{\lor}}
\BinaryInf$\lnot !A \lor !B, !A \fCenter !B $
\RightLabel{\LeftR{\Exchange}}
\UnaryInf$ !A, \lnot !A \lor !B \fCenter !B $
\RightLabel{\RightR{\lif}}
\UnaryInf$ \lnot !A \lor !B \fCenter !A \lif !B $
\end{prooftree}
\end{ex}
\label{olpseg:OLP-0075-B011:end}

\phantomsection\label{olpseg:OLP-0075-B012:start}
\begin{ex}
د سېکوېنټ \LR{$\lnot !A \lor \lnot !B
\Sequent \lnot (!A \land !B)$} يو \LR{$\Log{LK}$}-اشتقاق ورکړئ۔
\label{olpseg:OLP-0075-B012:end}

\phantomsection\label{olpseg:OLP-0075-B013:start}
د پورته طريقو په کارولو، غوښتل شوے وروستنے سېکوېنټ په لاندې برخه
کښې ليکو۔
\begin{prooftree}
\AxiomC{}
\UnaryInf$ \lnot !A \lor \lnot !B \fCenter \lnot (!A \land !B) $
\end{prooftree}
په وروستي سېکوېنټ کښې د تړلي فارمولونه شته اصلي نښلوونکي د \LR{$\lor$}
او \LR{$\lnot$} نښې دي۔ دلته هم \LeftR{\lor} او هم \RightR{\lnot}
کار ورکوي، خو په \RightR{\lnot} پيل کوو، ځکه چې د يوې شېبې دپاره
پر دوو څانګو وېشل ځنډوي:
\begin{prooftree}
\AxiomC{}
\UnaryInf$!A \land !B, \lnot !A \lor \lnot !B \fCenter $
\RightLabel{\RightR{\lnot}}
\UnaryInf$\lnot !A \lor \lnot !B \fCenter \lnot (!A \land !B)$
\end{prooftree}
اوس د \LeftR{\land} او \LeftR{\lor} تر منځ ټاکنه لرو۔ ګورو چې د
\LeftR{\land} قاعدې په کارولو څه کېږي: يا په سېکوېنټ \LR{$!A,
\lnot !A \lor \lnot !B \Sequent \quad$} او يا په سېکوېنټ \LR{$!B, \lnot !A
\lor \lnot !B \Sequent \quad$} پيل کولے شو۔ ځکه چې اشتقاق د
\LR{$!A$} او \LR{$!B$} له پلوه متناظر دے، لومړنے غوره کوو۔
\textbf{د ژباړې د نسخې سمون (OLSQ-002):} د سرچينې په دغو دوو
امکاني پورتنيو سېکوېنټونو کښې له دويمې جملې څخه نفي پاتې شوې وه؛
دلته د وروستي سېکوېنټ او ورپسې قاعدوي ونې له مخې په دواړو جملو
کښې نفي ساتل کېږي۔
\begin{prooftree}
\AxiomC{}
\UnaryInf$!A, \lnot !A \lor \lnot !B \fCenter $
\RightLabel{\LeftR{\land}}
\UnaryInf$!A \land !B, \lnot !A \lor \lnot !B \fCenter $
\RightLabel{\RightR{\lnot}}
\UnaryInf$\lnot !A \lor \lnot !B \fCenter \lnot (!A \land !B)$
\end{prooftree}
د اشتقاق د ډکولو په دوام وينو چې له يوې ستونزې سره مخ کېږو:
\begin{prooftree}
\Axiom$!A \fCenter !A$
\RightLabel{\LeftR{\lnot}}
\UnaryInf$ \lnot !A, !A \fCenter$
\AxiomC{}
\RightLabel{?}
\UnaryInf$!A \fCenter !B$
\RightLabel{\LeftR{\lnot}}
\UnaryInf$ \lnot !B, !A \fCenter$
\RightLabel{\LeftR{\lor}}
\BinaryInf$\lnot !A \lor \lnot !B, !A \fCenter $
\RightLabel{\LeftR{\Exchange}}
\UnaryInf$!A, \lnot !A \lor \lnot !B \fCenter $
\RightLabel{\LeftR{\land}}
\UnaryInf$!A \land !B, \lnot !A \lor \lnot !B \fCenter $
\RightLabel{\RightR{\lnot}}
\UnaryInf$\lnot !A \lor \lnot !B \fCenter \lnot (!A \land !B)$
\end{prooftree}
د ښۍ څانګې پورته برخه نوره نۀ شي راکمېدلے او د جوړښتي استنتاجونو
له لارې لومړني سېکوېنټ ته نۀ شي رسول کېدے؛ نو دا سمه لار نۀ ده۔
له دې امله پورته د \LeftR{\land} قاعدې کارول تېروتنه وه۔ بېرته
مخکيني حالت ته ځو او پر ځاے ئې \LeftR{\lor} قاعده پلې کوو، نو
ترلاسه کوو:
\begin{prooftree}
\AxiomC{}
\UnaryInf$\lnot !A, !A \land !B \fCenter $
\label{olpseg:OLP-0075-B013:end}

\phantomsection\label{olpseg:OLP-0075-B014:start}
\AxiomC{}
\UnaryInf$\lnot !B, !A \land !B \fCenter $
\label{olpseg:OLP-0075-B014:end}

\phantomsection\label{olpseg:OLP-0075-B015:start}
\RightLabel{\LeftR{\lor}}
\BinaryInf$\lnot !A \lor \lnot !B, !A \land !B \fCenter $
\RightLabel{\LeftR{\Exchange}}
\UnaryInf$!A \land !B, \lnot !A \lor \lnot !B \fCenter $
\RightLabel{\RightR{\lnot}}
\UnaryInf$\lnot !A \lor \lnot !B \fCenter \lnot (!A \land !B)$
\end{prooftree}
هره څانګه د مخکښې په شان بشپړه کړو، نو ترلاسه کوو:
\begin{prooftree}
\Axiom$ !A \fCenter!A$
\RightLabel{\LeftR{\land}}
\UnaryInf$!A \land !B \fCenter !A$
\RightLabel{\LeftR{\lnot}}
\UnaryInf$\lnot !A, !A \land !B \fCenter $
\label{olpseg:OLP-0075-B015:end}

\phantomsection\label{olpseg:OLP-0075-B016:start}
\Axiom$ !B \fCenter !B$
\RightLabel{\LeftR{\land}}
\UnaryInf$!A \land !B \fCenter !B$
\RightLabel{\LeftR{\lnot}}
\UnaryInf$\lnot !B, !A \land !B \fCenter $
\label{olpseg:OLP-0075-B016:end}

\phantomsection\label{olpseg:OLP-0075-B017:start}
\RightLabel{\LeftR{\lor}}
\BinaryInf$\lnot !A \lor \lnot !B, !A \land !B \fCenter $
\RightLabel{\LeftR{\Exchange}}
\UnaryInf$!A \land !B, \lnot !A \lor \lnot !B \fCenter $
\RightLabel{\RightR{\lnot}}
\UnaryInf$\lnot !A \lor \lnot !B \fCenter \lnot (!A \land !B)$
\end{prooftree}
(په دغو ګامونو کښې د \LR{$\land$} قاعدې د \LR{$\lnot$} له قاعدو څخه لاندې
هم پلې کېدے شوې او بيا به هم سم اشتقاق ترلاسه شوے و۔)
\end{ex}
\label{olpseg:OLP-0075-B017:end}

\phantomsection\label{olpseg:OLP-0075-B018:start}
\begin{ex}
تر اوسه مو د انقباض قاعده نۀ ده کارولې، خو کله کله اړينه وي۔ دا
ئې يوه بېلګه ده۔ فرض کړئ غواړو \LR{$\quad \Sequent !A \lor \lnot !A$}
ثابت کړو۔ د \LR{$\RightR{\lor}$} سرچپه پلي کول به له دغو دوو
اشتقاقونه څخه يو راکړي:
\begin{prooftree}
\AxiomC{}
\UnaryInf$ \fCenter !A$
\RightLabel{\RightR{\lor}}
\UnaryInf$ \fCenter !A \lor \lnot !A$
\DisplayProof\qquad\bottomAlignProof
\AxiomC{}
\UnaryInf$!A \fCenter $
\RightLabel{\RightR{\lnot}}
\UnaryInf$ \fCenter \lnot !A$
\RightLabel{\RightR{\lor}}
\UnaryInf$ \fCenter !A \lor \lnot !A$
\end{prooftree}
البته له دغو څخه هېڅ يو په لومړني سېکوېنټ نۀ پاے ته رسېږي۔ چل دا
دے چې پوه شو د انقباض قاعده د تړلے فارمول دوه نقلونه سره يو
کولے شي---او کله چې د ثبوت لټون کوو، يعنې له لاندې څخه پورته ځو،
د \LR{$!A \lor \lnot !A$} يوه کاپي په مقدمه کښې ساتلے شو، مثلاً:
\begin{prooftree}
\AxiomC{}
\UnaryInf$ \fCenter !A \lor \lnot !A, !A$
\RightLabel{\RightR{\lor}}
\UnaryInf$ \fCenter !A \lor \lnot !A, !A \lor \lnot !A$
\RightLabel{\RightR{\Contraction}}
\UnaryInf$ \fCenter !A \lor \lnot !A$
\end{prooftree}
اوس \LR{$\RightR{\lor}$} دويم ځل کارولے شو او~\LR{$\lnot !A$} هم ترلاسه کوو،
چې بشپړ اشتقاق ته رسېږي۔
\begin{prooftree}
\Axiom$!A \fCenter !A$
\RightLabel{\RightR{\lnot}}
\UnaryInf$\fCenter !A, \lnot !A$
\RightLabel{\RightR{\lor}}
\UnaryInf$\fCenter !A, !A \lor \lnot !A$
\RightLabel{\RightR{\Exchange}}
\UnaryInf$ \fCenter !A \lor \lnot !A, !A$
\RightLabel{\RightR{\lor}}
\UnaryInf$ \fCenter !A \lor \lnot !A, !A \lor \lnot !A$
\RightLabel{\RightR{\Contraction}}
\UnaryInf$ \fCenter !A \lor \lnot !A$
\end{prooftree}
\end{ex}
\label{olpseg:OLP-0075-B018:end}

\phantomsection\label{olpseg:OLP-0075-B019:start}
\begin{prob}
د لاندې سېکوېنټونو اشتقاقونه ورکړئ:
\begin{enumerate}
\item \LR{$!A \land (!B \land !C) \Sequent (!A \land !B) \land !C$}.
\item \LR{$!A \lor (!B \lor !C) \Sequent (!A \lor !B) \lor !C$}.
\item \LR{$!A \lif (!B \lif !C) \Sequent !B \lif (!A \lif !C)$}.
\item \LR{$!A \Sequent \lnot\lnot !A$}.
\end{enumerate}
\end{prob}
\label{olpseg:OLP-0075-B019:end}

\phantomsection\label{olpseg:OLP-0075-B020:start}
\begin{prob}
د لاندې سېکوېنټونو اشتقاقونه ورکړئ:
\begin{enumerate}
\item \LR{$(!A \lor !B) \lif !C \Sequent !A \lif !C$}.
\item \LR{$(!A \lif !C) \land (!B \lif !C) \Sequent (!A \lor !B) \lif !C$}.
\item \LR{$\Sequent \lnot(!A \land \lnot !A)$}.
\item \LR{$!B \lif !A \Sequent \lnot !A \lif \lnot !B$}.
\item \LR{$\Sequent (!A \lif \lnot !A) \lif \lnot !A$}.
\item \LR{$\Sequent \lnot(!A \lif !B) \lif \lnot !B$}.
\item \LR{$!A \lif !C \Sequent \lnot (!A \land \lnot !C)$}.
\item \LR{$!A \land \lnot !C \Sequent \lnot (!A \lif !C)$}.
\item \LR{$!A \lor !B, \lnot !B \Sequent !A$}.
\item \LR{$\lnot !A \lor \lnot !B \Sequent \lnot(!A \land !B)$}.
\item \LR{$\Sequent (\lnot !A \land \lnot !B) \lif\lnot(!A \lor !B)$}.
\item \LR{$\Sequent \lnot(!A \lor !B) \lif (\lnot !A \land \lnot !B)$}.
\end{enumerate}
\end{prob}
\label{olpseg:OLP-0075-B020:end}

\phantomsection\label{olpseg:OLP-0075-B021:start}
\begin{prob}
د لاندې سېکوېنټونو اشتقاقونه ورکړئ:
\begin{enumerate}
\item \LR{$\lnot(!A \lif !B) \Sequent !A$}.
\item \LR{$\lnot(!A \land !B) \Sequent \lnot !A \lor \lnot !B$}.
\item \LR{$!A \lif !B \Sequent \lnot !A \lor !B$}.
\item \LR{$\Sequent \lnot \lnot !A \lif !A$}.
\item \LR{$!A \lif !B, \lnot !A \lif !B \Sequent !B$}.
\item \LR{$(!A \land !B) \lif !C \Sequent (!A \lif !C) \lor (!B \lif !C)$}.
\item \LR{$(!A \lif !B) \lif !A \Sequent !A$}.
\item \LR{$\Sequent (!A \lif !B) \lor (!B \lif !C)$}.
\end{enumerate}
(دا ټول د~\LR{$\RightR{\Contraction}$} قاعدې ته اړتيا لري۔)
\end{prob}
\label{olpseg:OLP-0075-B021:end}

\phantomsection\label{olpseg:OLP-0076-B004:start}
\olfileid{fol}{seq}{prq}
\label{olpseg:OLP-0076-B004:end}

\phantomsection\label{olpseg:OLP-0076-B005:start}
\olsection{له سورونو سره اشتقاقونه}
\label{olpseg:OLP-0076-B005:end}

\phantomsection\label{olpseg:OLP-0076-B006:start}
\begin{ex}
د سېکوېنټ \LR{$\lexists[x][\lnot !A(x)]
\Sequent \lnot \lforall[x][!A(x)]$} يو
\LR{$\Log{LK}$}-اشتقاق ورکړئ۔
\label{olpseg:OLP-0076-B006:end}

\phantomsection\label{olpseg:OLP-0076-B007:start}
کله چې له سورونو سره کار کوو، بايد ډاډمن شو چې د ځانګړي متغير شرط
نۀ ماتوو، او کله کله د دې دپاره اړ يو چې د ځينو استنتاجونو د پلي
کولو ترتيب بدل کړو۔ په عمومي توګه دا ګټوره ده چې لومړے هغه قاعدې
وڅېړو چې د ځانګړي متغير د شرط تابع دي (په بشپړ شوي ثبوت کښې به
لاندې وي)۔ همدارنګه، غوره ده لږ وړاندې وګورو او اټکل وکړو چې
لومړنے سېکوېنټ به څنګه وي۔ زموږ په حالت کښې بايد د
\LR{$!A(a) \Sequent !A(a)$} په شان وي۔ يعنې کله چې د سورونو قاعدې
“سرچپه” کوو، بايد د \LR{$\lforall$} او \LR{$\lexists$} دواړو قاعدو دپاره
هماغه اصطلاح---چې \LR{$a$} به ئې وبولو---غوره کړو۔ که د هرې قاعدې
دپاره بېله اصطلاح غوره کړو، د \LR{$!A(a) \Sequent !A(b)$} په شان څه
ته رسېږو، چې البته د اشتقاق وړ نۀ دے۔
\label{olpseg:OLP-0076-B007:end}

\phantomsection\label{olpseg:OLP-0076-B008:start}
د معمول په شان په دې ليکلو پيل کوو:
\begin{prooftree}
\AxiomC{}
\UnaryInf$\lexists[x][\lnot !A(x)] \fCenter \lnot \lforall[x][!A(x)]$
\end{prooftree}
يا د \LeftR{\exists} قاعده پلې کولے شو او يا د
\RightR{\lnot} قاعده۔ ځکه چې \LeftR{\exists} د ځانګړي متغير د
شرط تابع ده، ښه ده چې دا ژر تر ژره وڅېړو؛ نو لومړے همدا پلې کوو۔
\begin{prooftree}
\AxiomC{}
\UnaryInf$ \lnot !A(a) \fCenter \lnot \lforall[x][!A(x)]$
\RightLabel{\LeftR{\lexists}}
\UnaryInf$ \lexists[x][\lnot !A(x)] \fCenter \lnot \lforall[x][!A(x)]$
\end{prooftree}
د \LeftR{\lnot} او \RightR{\lnot} قاعدو په سرچپه پلي کولو ترلاسه
کوو:
\begin{prooftree}
\AxiomC{}
\UnaryInf$\lforall[x][!A(x)] \fCenter !A(a)$
\RightLabel{\LeftR{\lnot}}
\UnaryInf$\lnot !A(a), \lforall[x][!A(x)] \fCenter $
\RightLabel{\LeftR{\Exchange}}
\UnaryInf$\lforall[x][!A(x)], \lnot !A(a) \fCenter $
\RightLabel{\RightR{\lnot}}
\UnaryInf$ \lnot !A(a) \fCenter \lnot \lforall[x] !A(x)$
\RightLabel{\LeftR{\lexists}}
\UnaryInf$ \lexists[x] \lnot !A(x) \fCenter \lnot \lforall[x] !A(x)$
\end{prooftree}
په دې ځاے کښې يوازينے امکان د \LeftR{\forall} قاعدې پلي کول دي۔
ځکه چې دا قاعده د ځانګړي متغير د بنديز تابع نۀ ده، نو کومه ستونزه
نشته۔ راياد کړئ چې غواړو يو لومړنے سېکوېنټ، يعنې د
\LR{$!A(a) \Sequent !A(a)$} په بڼه، ترلاسه کړو؛ نو د قاعدې د کارولو پر
مهال بايد \LR{$a$} د~\LR{$!A$} د ارګومېنټ په توګه غوره کړو۔
\begin{prooftree}
\Axiom$!A(a) \fCenter !A(a)$
\RightLabel{\LeftR{\lforall}}
\UnaryInf$\lforall[x][!A(x)] \fCenter !A(a)$
\RightLabel{\LeftR{\lnot}}
\UnaryInf$\lnot !A(a), \lforall[x][!A(x)] \fCenter $
\RightLabel{\LeftR{\Exchange}}
\UnaryInf$\lforall[x][!A(x)], \lnot !A(a) \fCenter $
\RightLabel{\RightR{\lnot}}
\UnaryInf$ \lnot !A(a) \fCenter \lnot \lforall[x][!A(x)]$
\RightLabel{\LeftR{\lexists}}
\UnaryInf$ \lexists[x][ \lnot !A(x)] \fCenter \lnot \lforall[x][!A(x)]$
\end{prooftree}
په تېره بيا له سورونو سره د کار پر مهال مهمه ده چې په دې ځاے کښې
بيا وګورو چې د ځانګړي متغير شرط خو نۀ دے مات شوے۔ ځکه چې زموږ له
کارول شوو قاعدو څخه يوازې \LeftR{\exists} د ځانګړي متغير د شرط
تابع وه، او ځانګړے متغير~\LR{$a$} د هغې په لاندې سېکوېنټ (وروستي
سېکوېنټ) کښې نۀ راځي، نو دا سم اشتقاق دے۔
\end{ex}
\label{olpseg:OLP-0076-B008:end}

\phantomsection\label{olpseg:OLP-0076-B009:start}
\begin{prob}
د لاندې سېکوېنټونو اشتقاقونه ورکړئ:
\begin{enumerate}
\item \LR{$\Sequent (\lforall[x][!A(x)] \land \lforall[y][!B(y)]) \lif
\lforall[z][(!A(z) \land !B(z))]$}.
\item \LR{$\Sequent (\lexists[x][!A(x)] \lor \lexists[y][!B(y)]) \lif
\lexists[z][(!A(z) \lor !B(z))]$}.
\item \LR{$\lforall[x][(!A(x) \lif !B)] \Sequent \lexists[y][!A(y)] \lif !B$}.
\item \LR{$\lforall[x][\lnot !A(x)] \Sequent \lnot\lexists[x][!A(x)]$}.
\item \LR{$\Sequent \lnot\lexists[x][!A(x)] \lif \lforall[x][\lnot !A(x)]$}.
\item \LR{$\Sequent \lnot\lexists[x][\lforall[y][((!A(x,y) \lif \lnot
!A(y,y)) \land (\lnot !A(y,y) \lif !A(x,y)))]]$}.
\end{enumerate}
\end{prob}
\label{olpseg:OLP-0076-B009:end}

\phantomsection\label{olpseg:OLP-0076-B010:start}
\begin{prob}
د لاندې سېکوېنټونو اشتقاقونه ورکړئ:
\begin{enumerate}
\item \LR{$\Sequent \lnot\lforall[x][!A(x)] \lif \lexists[x][\lnot!A(x)]$}.
\item \LR{$(\lforall[x][!A(x)] \lif !B) \Sequent \lexists[y][(!A(y) \lif !B)]$}.
\item \LR{$\Sequent \lexists[x][(!A(x) \lif \lforall[y][!A(y)])]$}.
\end{enumerate}
(دا ټول د~\LR{$\RightR{\Contraction}$} قاعدې ته اړتيا لري۔)
\end{prob}
\label{olpseg:OLP-0076-B010:end}

\phantomsection\label{olpseg:OLP-0077-B004:start}
\begin{editorial}
  دا برخه د سېکوېنټ حساب د اشتقاق وړتيا د اړيکې او سازګارۍ
  تعريفونه راغونډوي۔
  \textbf{د ژباړې د نسخې سمون (OLSQ-003):} په سرچينه کښې د طبيعي
  استنتاج نوم راغلے و؛ دلته د څپرکي او ورپسې تعريفونو له مخې د
  سېکوېنټ حساب نوم کارول کېږي۔
\end{editorial}
\label{olpseg:OLP-0077-B004:end}

\phantomsection\label{olpseg:OLP-0077-B005:start}
\olfileid{fol}{seq}{ptn}
\label{olpseg:OLP-0077-B005:end}

\phantomsection\label{olpseg:OLP-0077-B006:start}
\olsection{د ثبوت-تيوريکي مفاهيم}
\label{olpseg:OLP-0077-B006:end}

\phantomsection\label{olpseg:OLP-0077-B007:start}
\begin{explain}
لکه څنګه چې مو يو شمېر مهم معنايي مفاهيم (اعتبار، معنايي استلزام
او د صدق وړتيا) تعريف کړي دي، اوس ورسره سم \emph{ثبوت-تيوريکي
مفاهيم} تعريفوو۔ دا مفاهيم په جوړښتونه کښې د تړلي فارمولونه
د صدق پر بنسټ نۀ، بلکې د ځينو سېکوېنټونو د د اشتقاق وړتيا يا
د اشتقاق نۀ وړتيا پر بنسټ تعريفېږي۔ دا يوه مهمه موندنه وه چې دغه
مفاهيم سره سمون خوري۔ د دې سمون منځپانګه د \emph{سموالي} او
\emph{بشپړتيا قضيه} ده۔
\end{explain}
\label{olpseg:OLP-0077-B007:end}

\phantomsection\label{olpseg:OLP-0077-B008:start}
\begin{defn}[قضيې]
تړلے فارمول~\LR{$!A$} هله \emph{قضيه} ده چې په~\LR{$\Log{LK}$} کښې د
سېکوېنټ \LR{$\quad \Sequent !A$} اشتقاق موجود وي۔ که \LR{$!A$}
قضيه وي، \LR{$\Proves !A$} ليکو؛ او که نۀ وي، \LR{$\Proves/ !A$} ليکو۔
\end{defn}
\label{olpseg:OLP-0077-B008:end}

\phantomsection\label{olpseg:OLP-0077-B009:start}
\begin{defn}[د اشتقاق وړتيا]
تړلے فارمول~\LR{$!A$} د تړلي فارمولونه له سټ~\LR{$\Gamma$} څخه
\emph{د اشتقاق وړ} ده، يعنې \LR{$\Gamma \Proves !A$}، هله او يوازې
هله چې يو متناهي فرعي سټ~\LR{$\Gamma_0 \subseteq \Gamma$} او په
\LR{$\Gamma_0$} کښې د تړلي فارمولونه يو لړ~\LR{$\Gamma_0'$} موجود وي، داسې
چې \LR{$\Log{LK}$} د \LR{$\Gamma_0' \Sequent !A$} سېکوېنټ اشتقاق کړي۔
که \LR{$!A$} له \LR{$\Gamma$} څخه د اشتقاق وړ نۀ وي، \LR{$\Gamma \Proves/ !A$}
ليکو۔
\end{defn}
\label{olpseg:OLP-0077-B009:end}

\phantomsection\label{olpseg:OLP-0077-B010:start}
د انقباض، کمزوري کولو او تبادلې د قاعدو له امله، په~\LR{$\Gamma_0'$}
کښې د تړلي فارمولونه ترتيب او شمېر اهميت نۀ لري: که سېکوېنټ
\LR{$\Gamma_0' \Sequent !A$} د اشتقاق وړ وي، نو د هر داسې
\LR{$\Gamma_0''$} دپاره \LR{$\Gamma_0'' \Sequent !A$} هم د اشتقاق وړ دے
چې د~\LR{$\Gamma_0'$} هماغه تړلي فارمولونه لري۔ مثلاً که
\LR{$\Gamma_0 = \{!B, !C\}$} وي، نو هم
\LR{$\Gamma_0' = \tuple{!B, !B, !C}$} او هم
\LR{$\Gamma_0'' = \tuple{!C, !C, !B}$} داسې لړونه دي چې يوازې د
\LR{$\Gamma_0$} تړلي فارمولونه لري۔ که د يوه لړ لرونکے سېکوېنټ
د اشتقاق وړ وي، د بل لړ لرونکے هم دے، مثلاً:
\begin{prooftree}
  \AxiomC{}
  \Deduce$!B, !B, !C \fCenter !A$
  \RightLabel{\LeftR{\Contraction}}
  \UnaryInf$!B, !C \fCenter !A$
  \RightLabel{\LeftR{\Exchange}}
  \UnaryInf$!C, !B \fCenter !A$
  \RightLabel{\LeftR{\Weakening}}
  \UnaryInf$!C, !C, !B \fCenter !A$
\end{prooftree}
له دې وروسته به وايو چې که \LR{$\Gamma_0$} د تړلي فارمولونه يو متناهي
سټ وي، نو \LR{$\Gamma_0 \Sequent !A$} هر هغه سېکوېنټ دے چې مقدم ئې
په~\LR{$\Gamma_0$} کښې د تړلي فارمولونه يو لړ وي؛ او د اړتيا په وخت کښې
انقباضونه، تبادلې او کمزوري کول په ناويلې توګه ورګډوو۔
\label{olpseg:OLP-0077-B010:end}

\phantomsection\label{olpseg:OLP-0077-B011:start}
\begin{defn}[سازګاري]
د جملو سټ~\LR{$\Gamma$} هله او يوازې هله \emph{ناسازګار} دے چې يو
متناهي فرعي سټ~\LR{$\Gamma_0 \subseteq \Gamma$} موجود وي، داسې چې
\LR{$\Log{LK}$} د \LR{$\Gamma_0 \Sequent \quad$} سېکوېنټ اشتقاق کړي۔
که \LR{$\Gamma$} ناسازګار نۀ وي، يعنې که د هر متناهي
\LR{$\Gamma_0 \subseteq \Gamma$} دپاره \LR{$\Log{LK}$} د
\LR{$\Gamma_0 \Sequent \quad$} سېکوېنټ اشتقاق نۀ کړي، نو هغه
\emph{سازګار} بولو۔
\end{defn}
\label{olpseg:OLP-0077-B011:end}

\phantomsection\label{olpseg:OLP-0077-B012:start}
\begin{prop}[انعکاسيت]
\ollabel{prop:reflexivity}
که \LR{$!A \in \Gamma$} وي، نو \LR{$\Gamma \Proves !A$}۔
\end{prop}
\label{olpseg:OLP-0077-B012:end}

\phantomsection\label{olpseg:OLP-0077-B013:start}
\begin{proof}
لومړنے سېکوېنټ \LR{$!A \Sequent !A$} د اشتقاق وړ دے، او
\LR{$\{!A\} \subseteq \Gamma$}۔
\end{proof}
\label{olpseg:OLP-0077-B013:end}

\phantomsection\label{olpseg:OLP-0077-B014:start}
\begin{prop}[يکنواختي]
\ollabel{prop:monotonicity}
که \LR{$\Gamma \subseteq \Delta$} او \LR{$\Gamma \Proves !A$} وي، نو
\LR{$\Delta \Proves !A$}۔
\end{prop}
\label{olpseg:OLP-0077-B014:end}

\phantomsection\label{olpseg:OLP-0077-B015:start}
\begin{proof}
فرض کړئ \LR{$\Gamma \Proves !A$}، يعنې يو متناهي
\LR{$\Gamma_0 \subseteq \Gamma$} شته چې
\LR{$\Gamma_0 \Sequent !A$} د اشتقاق وړ دے۔ ځکه چې
\LR{$\Gamma \subseteq \Delta$}، نو \LR{$\Gamma_0$} د~\LR{$\Delta$} هم يو متناهي
فرعي سټ دے۔ له دې امله د \LR{$\Gamma_0 \Sequent !A$} اشتقاق
هم \LR{$\Delta \Proves !A$} ښيي۔
\end{proof}
\label{olpseg:OLP-0077-B015:end}

\phantomsection\label{olpseg:OLP-0077-B016:start}
\begin{prop}[انتقاليت]
\ollabel{prop:transitivity}
که \LR{$\Gamma \Proves !A$} او \LR{$\{!A\} \cup \Delta \Proves !B$} وي، نو
\LR{$\Gamma \cup \Delta \Proves !B$}۔
\end{prop}
\label{olpseg:OLP-0077-B016:end}

\phantomsection\label{olpseg:OLP-0077-B017:start}
\begin{proof}
که \LR{$\Gamma \Proves !A$} وي، يو متناهي
\LR{$\Gamma_0 \subseteq \Gamma$} او د
\LR{$\Gamma_0 \Sequent !A$} اشتقاق~\LR{$\pi_0$} شته۔ که
\LR{$\{!A\} \cup \Delta \Proves !B$} وي، نو د کوم متناهي فرعي سټ
\LR{$\Delta_0 \subseteq \Delta$} دپاره د
\LR{$!A, \Delta_0 \Sequent !B$} اشتقاق~\LR{$\pi_1$} شته۔ لاندې
اشتقاق وګورئ:
\begin{prooftree}
  \AxiomC{}
  \RightLabel{\LR{$\pi_0$}}
  \Deduce$\Gamma_0 \fCenter !A$
  \AxiomC{}
  \RightLabel{\LR{$\pi_1$}}
  \Deduce$!A, \Delta_0 \fCenter !B$
  \RightLabel{\Cut}
  \BinaryInf$\Gamma_0, \Delta_0 \fCenter !B$
\end{prooftree}
ځکه چې \LR{$\Gamma_0 \cup \Delta_0 \subseteq \Gamma \cup \Delta$}، نو
دا \LR{$\Gamma \cup \Delta \Proves !B$} ښيي۔
\end{proof}
\label{olpseg:OLP-0077-B017:end}

\phantomsection\label{olpseg:OLP-0077-B018:start}
پام وکړئ، په ځانګړې توګه د دې معنٰی دا ده چې که
\LR{$\Gamma \Proves !A$} او \LR{$!A \Proves !B$} وي، نو
\LR{$\Gamma \Proves !B$}۔ همدارنګه ترې دا هم راځي چې که
\LR{$!A_1, \dots, !A_n \Proves !B$} وي او د هر~\LR{$i$} دپاره
\LR{$\Gamma \Proves !A_i$} وي، نو \LR{$\Gamma \Proves !B$}۔
\label{olpseg:OLP-0077-B018:end}

\phantomsection\label{olpseg:OLP-0077-B019:start}
\begin{prop}
\ollabel{prop:incons}
\LR{$\Gamma$} هله او يوازې هله ناسازګار دے چې د هرې جملې~\LR{$!A$} دپاره
\LR{$\Gamma \Proves {!A}$} وي۔
\end{prop}
\label{olpseg:OLP-0077-B019:end}

\phantomsection\label{olpseg:OLP-0077-B020:start}
\begin{proof}
تمرين۔
\end{proof}
\label{olpseg:OLP-0077-B020:end}

\phantomsection\label{olpseg:OLP-0077-B021:start}
\begin{prob}
\olref[fol][seq][ptn]{prop:incons} ثابت کړئ۔
\end{prob}
\label{olpseg:OLP-0077-B021:end}

\phantomsection\label{olpseg:OLP-0077-B022:start}
\begin{prop}[متناهي شاهد]
  \ollabel{prop:proves-compact}
  \begin{enumerate}
  \item که \LR{$\Gamma \Proves !A$} وي، نو يو متناهي فرعي سټ
    \LR{$\Gamma_0 \subseteq \Gamma$} شته چې \LR{$\Gamma_0 \Proves !A$}۔
  \item که د~\LR{$\Gamma$} هر متناهي فرعي سټ سازګار وي، نو \LR{$\Gamma$}
    سازګار دے۔
  \end{enumerate}
\end{prop}
\label{olpseg:OLP-0077-B022:end}

\phantomsection\label{olpseg:OLP-0077-B023:start}
\begin{proof}
  \begin{enumerate}
    \item که \LR{$\Gamma \Proves !A$} وي، نو يو متناهي فرعي سټ
      \LR{$\Gamma_0 \subseteq \Gamma$} شته، داسې چې د
      \LR{$\Gamma_0 \Sequent !A$} سېکوېنټ اشتقاق لري۔ له دې امله
      \LR{$\Gamma_0 \Proves !A$}۔
    \item که \LR{$\Gamma$} ناسازګار وي، نو يو متناهي فرعي سټ
      \LR{$\Gamma_0 \subseteq \Gamma$} شته، داسې چې \LR{$\Log{LK}$} د
      \LR{$\Gamma_0 \Sequent \quad$} سېکوېنټ اشتقاق کړي۔ خو بيا
      \LR{$\Gamma_0$} د~\LR{$\Gamma$} يو ناسازګار متناهي فرعي سټ دے۔
  \end{enumerate}
\end{proof}
\label{olpseg:OLP-0077-B023:end}

\phantomsection\label{olpseg:OLP-0078-B004:start}
\olfileid{fol}{seq}{prv}
\label{olpseg:OLP-0078-B004:end}

\phantomsection\label{olpseg:OLP-0078-B005:start}
\olsection{د اشتقاق وړتيا او سازګاري}
\label{olpseg:OLP-0078-B005:end}

\phantomsection\label{olpseg:OLP-0078-B006:start}
اوس به د د اشتقاق وړتيا اړيکې يو شمېر خاصيتونه ثابت کړو۔ دا
خاصيتونه په خپلواک ډول هم په زړه پورې دي، خو هر يو به د بشپړتيا د
قضيې په ثبوت کښې ونډه ولري۔
\label{olpseg:OLP-0078-B006:end}

\phantomsection\label{olpseg:OLP-0078-B007:start}
\begin{prop}\ollabel{prop:provability-contr}
  که \LR{$\Gamma \Proves !A$} وي او \LR{$\Gamma \cup \{!A\}$} ناسازګار وي،
  نو \LR{$\Gamma$} ناسازګار دے۔
\end{prop}
\label{olpseg:OLP-0078-B007:end}

\phantomsection\label{olpseg:OLP-0078-B008:start}
\begin{proof}
داسې متناهي \LR{$\Gamma_0$} او \LR{$\Gamma_1 \subseteq \Gamma$} شته چې
\LR{$\Log{LK}$} د \LR{$\Gamma_0 \Sequent !A$} او
\LR{$!A, \Gamma_1 \Sequent \quad$} سېکوېنټونه اشتقاق کړي۔
د \LR{$\Gamma_0 \Sequent !A$} د \LR{$\Log{LK}$}-اشتقاق نوم~\LR{$\pi_0$}
او د \LR{$\Gamma_1, !A \Sequent \quad$} د
\LR{$\Log{LK}$}-اشتقاق نوم~\LR{$\pi_1$} کېږدو۔ بيا لاندې
اشتقاق کولے شو:
\begin{prooftree}
\AxiomC{}
\RightLabel{\LR{$\pi_0$}}
\Deduce$ \Gamma_0 \fCenter !A $
\AxiomC{}
\RightLabel{\LR{$\pi_1$}}
\Deduce$!A, \Gamma_1 \fCenter $
\RightLabel{\Cut}
\BinaryInf$ \Gamma_0,\Gamma_1 \fCenter $
\end{prooftree}
\label{olpseg:OLP-0078-B008:end}

\phantomsection\label{olpseg:OLP-0078-B009:start}
ځکه چې \LR{$\Gamma_0 \subseteq \Gamma$} او
\LR{$\Gamma_1 \subseteq \Gamma$}، نو
\LR{$\Gamma_0 \cup \Gamma_1 \subseteq \Gamma$}؛ له دې امله \LR{$\Gamma$}
ناسازګار دے۔
\end{proof}
\label{olpseg:OLP-0078-B009:end}

\phantomsection\label{olpseg:OLP-0078-B010:start}
\begin{prop}
\ollabel{prop:prov-incons}
\LR{$\Gamma \Proves !A$} هله او يوازې هله چې
\LR{$\Gamma \cup \{\lnot !A\}$} ناسازګار وي۔
\end{prop}
\label{olpseg:OLP-0078-B010:end}

\phantomsection\label{olpseg:OLP-0078-B011:start}
\begin{proof}
لومړے فرض کړئ \LR{$\Gamma \Proves !A$}، يعنې د
\LR{$\Gamma \Sequent !A$} اشتقاق~\LR{$\pi_0$} شته۔ د
\LR{$\LeftR{\lnot}$} قاعدې په ورزياتولو د
\LR{$\lnot !A, \Gamma \Sequent \quad$} اشتقاق ترلاسه کوو؛
يعنې \LR{$\Gamma \cup \{\lnot !A\}$} ناسازګار دے۔
\label{olpseg:OLP-0078-B011:end}

\phantomsection\label{olpseg:OLP-0078-B012:start}
که \LR{$\Gamma \cup \{\lnot !A\}$} ناسازګار وي، نو د
\LR{$\lnot !A, \Gamma \Sequent \quad$} اشتقاق~\LR{$\pi_1$} شته۔
لاندې د \LR{$\Gamma \Sequent !A$} اشتقاق دے:
\begin{prooftree}
  \Axiom$!A \fCenter !A$
  \RightLabel{\RightR{\lnot}}
  \UnaryInf$\fCenter !A, \lnot !A$
  \AxiomC{}
  \RightLabel{\LR{$\pi_1$}}
  \Deduce$\lnot !A, \Gamma \fCenter$
  \RightLabel{\Cut}
  \BinaryInf$\Gamma \fCenter !A$
\end{prooftree}
\end{proof}
\label{olpseg:OLP-0078-B012:end}

\phantomsection\label{olpseg:OLP-0078-B013:start}
\begin{prob}
ثابت کړئ چې \LR{$\Gamma \Proves \lnot !A$} هله او يوازې هله چې
\LR{$\Gamma \cup \{!A\}$} ناسازګار وي۔
\end{prob}
\label{olpseg:OLP-0078-B013:end}

\phantomsection\label{olpseg:OLP-0078-B014:start}
\begin{prop}\ollabel{prop:explicit-inc}
  که \LR{$\Gamma \Proves !A$} او \LR{$\lnot !A \in \Gamma$} وي، نو \LR{$\Gamma$}
  ناسازګار دے۔
\end{prop}
\label{olpseg:OLP-0078-B014:end}

\phantomsection\label{olpseg:OLP-0078-B015:start}
\begin{proof}
  فرض کړئ \LR{$\Gamma \Proves !A$} او \LR{$\lnot !A \in \Gamma$}۔ بيا د کوم
  سېکوېنټ \LR{$\Gamma_0 \Sequent !A$} اشتقاق~\LR{$\pi$} شته۔
  سېکوېنټ \LR{$\lnot !A, \Gamma_0 \Sequent \quad$} هم د اشتقاق وړ دے:
  \begin{prooftree}
    \AxiomC{}
    \RightLabel{\LR{$\pi$}}
    \Deduce$\Gamma_0 \fCenter !A$
    \Axiom$!A \fCenter !A$
    \RightLabel{\LeftR{\lnot}}
    \UnaryInf$\lnot !A, !A \fCenter$
    \RightLabel{\LeftR{\Exchange}}
    \UnaryInf$!A, \lnot !A \fCenter$
    \RightLabel{\Cut}
    \BinaryInf$\Gamma_0, \lnot !A \fCenter$
  \end{prooftree}
  ځکه چې \LR{$\lnot !A \in \Gamma$} او \LR{$\Gamma_0 \subseteq \Gamma$}،
  نو دا ښيي چې \LR{$\Gamma$} ناسازګار دے۔
\end{proof}
\label{olpseg:OLP-0078-B015:end}

\phantomsection\label{olpseg:OLP-0078-B016:start}
\begin{prop}\ollabel{prop:provability-exhaustive}
  که \LR{$\Gamma \cup \{!A\}$} او \LR{$\Gamma \cup \{\lnot !A\}$} دواړه
  ناسازګار وي، نو \LR{$\Gamma$} ناسازګار دے۔
\end{prop}
\label{olpseg:OLP-0078-B016:end}

\phantomsection\label{olpseg:OLP-0078-B017:start}
\begin{proof}
داسې متناهي سټونه \LR{$\Gamma_0 \subseteq \Gamma$} او
\LR{$\Gamma_1 \subseteq \Gamma$} او داسې
\LR{$\Log{LK}$}-اشتقاقونه \LR{$\pi_0$} او \LR{$\pi_1$} شته چې په همدې ترتيب
د \LR{$!A, \Gamma_0 \Sequent \quad$} او
\LR{$\lnot !A, \Gamma_1 \Sequent \quad$} اشتقاقونه دي۔ بيا لاندې
اشتقاق کولے شو:
\begin{prooftree}
\AxiomC{}
\RightLabel{\LR{$\pi_0$}}
\Deduce$ !A, \Gamma_0 \fCenter $
\RightLabel{\RightR{\lnot}}
\UnaryInf$ \Gamma_0 \fCenter \lnot !A$
\AxiomC{}
\RightLabel{\LR{$\pi_1$}}
\Deduce$\lnot !A, \Gamma_1 \fCenter  $
\RightLabel{\Cut}
\BinaryInf$ \Gamma_0, \Gamma_1 \fCenter $
\end{prooftree}
ځکه چې \LR{$\Gamma_0 \subseteq \Gamma$} او \LR{$\Gamma_1 \subseteq \Gamma$}،
نو \LR{$\Gamma_0 \cup \Gamma_1 \subseteq \Gamma$}۔ له دې امله \LR{$\Gamma$}
ناسازګار دے۔
\end{proof}
\label{olpseg:OLP-0078-B017:end}

\phantomsection\label{olpseg:OLP-0079-B005:start}
\olfileid{fol}{seq}{ppr}
\label{olpseg:OLP-0079-B005:end}

\phantomsection\label{olpseg:OLP-0079-B006:start}
\olsection{د اشتقاق وړتيا او قضييز نښلوونکي}
\label{olpseg:OLP-0079-B006:end}

\phantomsection\label{olpseg:OLP-0079-B007:start}
\begin{explain}
  موږ ثابتوو چې د سېکوېنټ حساب د د اشتقاق وړتيا اړيکه~\LR{$\Proves$}
  دومره پياوړې ده چې د قضييزو نښلوونکو په اړه ځينې بنسټيز حقيقتونه
  ثابت کړي، لکه \LR{$!A \land !B \Proves !A$} او
  \LR{$!A, !A \lif !B \Proves !B$} (مودوس پونېنس)۔ دغه حقيقتونه د
  بشپړتيا د قضيې د ثبوت دپاره اړين دي۔
\end{explain}
\label{olpseg:OLP-0079-B007:end}

\phantomsection\label{olpseg:OLP-0079-B008:start}
\begin{prop}\ollabel{prop:provability-land}
  \begin{enumerate}
  \item \ollabel{prop:provability-land-left} هم
    \LR{$!A \land !B \Proves !A$} او هم
    \LR{$!A \land !B \Proves !B$}۔
  \item \ollabel{prop:provability-land-right}
    \LR{$!A, !B \Proves !A \land !B$}۔
  \end{enumerate}
\end{prop}
\label{olpseg:OLP-0079-B008:end}

\phantomsection\label{olpseg:OLP-0079-B009:start}
\begin{proof}
  \begin{enumerate}
  \item دواړه سېکوېنټونه \LR{$!A \land !B \Sequent !A$} او
    \LR{$!A \land !B \Sequent !B$} د اشتقاق وړ دي:
    \begin{prooftree}
      \Axiom$!A \fCenter !A$
      \RightLabel{\LeftR{\land}}
      \UnaryInf$!A \land !B \fCenter !A$
      \DisplayProof\qquad\bottomAlignProof
      \Axiom$!B \fCenter !B$
      \RightLabel{\LeftR{\land}}
      \UnaryInf$!A \land !B \fCenter !B$
    \end{prooftree}
    \item دا د سېکوېنټ \LR{$!A, !B \Sequent !A \land !B$} يو
    اشتقاق دے۔
    \textbf{د ژباړې د نسخې سمون (OLSQ-004):} د سرچينې په ونه کښې
    د عطف-ښۍ قاعدې دواړو مقدمو يو شان مقدم نۀ درلود؛ دلته د
    جوړښتي استنتاجونو دوه ګونې کرښې دواړه مقدمې له پايلې سره سموي۔
    \begin{prooftree}
      \Axiom$!A \fCenter !A$
      \doubleLine
      \UnaryInf$!A, !B \fCenter !A$
      \Axiom$!B \fCenter !B$
      \doubleLine
      \UnaryInf$!A, !B \fCenter !B$
      \RightLabel{\RightR{\land}}
      \BinaryInf$!A, !B \fCenter !A \land !B$
    \end{prooftree}
  \end{enumerate}
\end{proof}
\label{olpseg:OLP-0079-B009:end}

\phantomsection\label{olpseg:OLP-0079-B010:start}
\begin{prop}\ollabel{prop:provability-lor}
  \begin{enumerate}
  \item \LR{$!A \lor !B, \lnot !A, \lnot !B$} ناسازګار دے۔
  \item هم \LR{$!A \Proves !A \lor !B$} او هم
    \LR{$!B \Proves !A \lor !B$}۔
  \end{enumerate}
\end{prop}
\label{olpseg:OLP-0079-B010:end}

\phantomsection\label{olpseg:OLP-0079-B011:start}
\begin{proof}
  \begin{enumerate}
  \item موږ د سېکوېنټ
    \LR{$!A \lor !B, \lnot !A, \lnot !B \Sequent$} يو
    اشتقاق ورکوو:
    \begin{prooftree}
      \Axiom$!A \fCenter !A$
      \RightLabel{\LeftR{\lnot}}
      \UnaryInf$\lnot !A, !A \fCenter$
      \doubleLine
      \UnaryInf$!A, \lnot !A, \lnot !B \fCenter$
      \Axiom$!B \fCenter !B$
      \RightLabel{\LeftR{\lnot}}
      \UnaryInf$\lnot !B, !B \fCenter$
      \doubleLine
      \UnaryInf$!B, \lnot !A, \lnot !B \fCenter$
      \RightLabel{\LeftR{\lor}}
      \BinaryInf$ !A \lor !B, \lnot !A, \lnot !B \fCenter $
    \end{prooftree}
    (راياد کړئ چې د استنتاج دوه ګونې کرښې څو کمزوري کول، انقباضونه
    او تبادلې ښيي۔)
  \item دواړه سېکوېنټونه \LR{$!A \Sequent !A \lor !B$} او
    \LR{$!B \Sequent !A \lor !B$} اشتقاقونه لري:
    \begin{prooftree}
      \Axiom$!A \fCenter !A$
      \RightLabel{\RightR{\lor}}
      \UnaryInf$!A \fCenter !A \lor !B$
      \DisplayProof\qquad\bottomAlignProof
      \Axiom$!B \fCenter !B$
      \RightLabel{\RightR{\lor}}
      \UnaryInf$!B \fCenter !A \lor !B$
    \end{prooftree}
  \end{enumerate}
\end{proof}
\label{olpseg:OLP-0079-B011:end}

\phantomsection\label{olpseg:OLP-0079-B012:start}
\begin{prop}\ollabel{prop:provability-lif}
  \begin{enumerate}
  \item \ollabel{prop:provability-lif-left}
    \LR{$!A, !A \lif !B \Proves !B$}۔
  \item \ollabel{prop:provability-lif-right}
    هم \LR{$\lnot !A \Proves !A \lif !B$} او هم
    \LR{$!B \Proves !A \lif !B$}۔
  \end{enumerate}
\end{prop}
\label{olpseg:OLP-0079-B012:end}

\phantomsection\label{olpseg:OLP-0079-B013:start}
\begin{proof}
  \begin{enumerate}
    \item سېکوېنټ \LR{$!A \lif !B, !A \Sequent !B$} د اشتقاق وړ دے:
      \begin{prooftree}
        \Axiom$!A \fCenter !A$
        \Axiom$!B \fCenter !B$
        \RightLabel{\LeftR{\lif}}
        \BinaryInf$!A \lif !B, !A  \fCenter !B$
      \end{prooftree}
    \item دواړه سېکوېنټونه
      \LR{$\lnot !A \Sequent !A \lif !B$} او
      \LR{$!B \Sequent !A \lif !B$} د اشتقاق وړ دي:
      \begin{prooftree}
        \Axiom$!A \fCenter !A$
        \RightLabel{\LeftR{\lnot}}
        \UnaryInf$\lnot !A, !A \fCenter$
        \RightLabel{\LeftR{\Exchange}}
        \UnaryInf$!A, \lnot !A \fCenter$
        \RightLabel{\RightR{\Weakening}}
        \UnaryInf$!A, \lnot !A \fCenter !B$
        \RightLabel{\RightR{\lif}}
        \UnaryInf$\lnot !A \fCenter !A \lif !B$
        \DisplayProof\qquad\bottomAlignProof
        \Axiom$!B \fCenter !B$
        \RightLabel{\LeftR{\Weakening}}
        \UnaryInf$!A, !B \fCenter !B$
        \RightLabel{\RightR{\lif}}
        \UnaryInf$!B \fCenter !A \lif !B$
      \end{prooftree}
  \end{enumerate}
\end{proof}
\label{olpseg:OLP-0079-B013:end}

\phantomsection\label{olpseg:OLP-0080-B004:start}
\olfileid{fol}{seq}{qpr}
\olsection{د اشتقاق وړتيا او سورونه}
\label{olpseg:OLP-0080-B004:end}

\phantomsection\label{olpseg:OLP-0080-B005:start}
\begin{explain}
  د بشپړتيا قضيه دا هم غواړي چې د سېکوېنټ حساب قاعدې په دې برخه
  کښې د~\LR{$\Proves$} په اړه ثابت شوي حقيقتونه راکړي۔
\end{explain}
\label{olpseg:OLP-0080-B005:end}

\phantomsection\label{olpseg:OLP-0080-B006:start}
\begin{thm}
\ollabel{thm:strong-generalization} که \LR{$c$} داسې ثابت وي چې په
\LR{$\Gamma$} يا \LR{$!A(x)$} کښې نۀ راځي، او \LR{$\Gamma \Proves !A(c)$} وي، نو
\LR{$\Gamma \Proves \lforall[x][!A(x)]$}۔
\end{thm}
\label{olpseg:OLP-0080-B006:end}

\phantomsection\label{olpseg:OLP-0080-B007:start}
\begin{proof}
د کوم متناهي \LR{$\Gamma_0 \subseteq \Gamma$} دپاره
\LR{$\Gamma_0 \Sequent !A(c)$} د \LR{$\Log{LK}$}-اشتقاق نوم~\LR{$\pi_0$}
کېږدئ۔ د \LR{$\RightR{\lforall}$} استنتاج په ورزياتولو د
\LR{$\Gamma_0 \Sequent \lforall[x][!A(x)]$} اشتقاق ترلاسه کوو؛
ځکه \LR{$c$} په \LR{$\Gamma$} يا \LR{$!A(x)$} کښې نۀ راځي، نو د ځانګړي متغير شرط
پوره کېږي۔
\end{proof}
\label{olpseg:OLP-0080-B007:end}

\phantomsection\label{olpseg:OLP-0080-B008:start}
\begin{prop}
\ollabel{prop:provability-quantifiers}
\begin{enumerate}
\item \LR{$!A(t) \Proves \lexists[x][!A(x)]$}.
\item \LR{$\lforall[x][!A(x)] \Proves !A(t)$}.
\end{enumerate}
\end{prop}
\label{olpseg:OLP-0080-B008:end}

\phantomsection\label{olpseg:OLP-0080-B009:start}
\begin{proof}
\begin{enumerate}
\item سېکوېنټ
  \LR{$!A(t) \Sequent \lexists[x][!A(x)]$} د اشتقاق وړ دے:
  \begin{prooftree}
    \Axiom$!A(t) \fCenter !A(t)$
    \RightLabel{\RightR{\lexists}}
    \UnaryInf$!A(t) \fCenter \lexists[x][!A(x)]$
    \end{prooftree}
\item سېکوېنټ
  \LR{$\lforall[x][!A(x)] \Sequent !A(t)$} د اشتقاق وړ دے:
  \begin{prooftree}
    \Axiom$!A(t) \fCenter !A(t)$
    \RightLabel{\LeftR{\lforall}}
    \UnaryInf$\lforall[x][!A(x)] \fCenter !A(t)$
    \end{prooftree}
\end{enumerate}
\end{proof}
\label{olpseg:OLP-0080-B009:end}

\phantomsection\label{olpseg:OLP-0081-B004:start}
\olfileid{fol}{seq}{sou}
\label{olpseg:OLP-0081-B004:end}

\phantomsection\label{olpseg:OLP-0081-B005:start}
\olsection{سموالے}
\label{olpseg:OLP-0081-B005:end}

\phantomsection\label{olpseg:OLP-0081-B006:start}
\begin{explain}
د سېکوېنټ حساب په شان اشتقاق نظام هله \emph{سم} دے چې
هغه څيزونه اشتقاق نۀ کړي چې په حقيقت کښې نۀ وي۔ نو سموالے د
اشتقاق نظامونو د تضمين شوې خونديتوب يو ډول خاصيت دے۔
د دې له مخې چې کوم ثبوت-تيوريکي خاصيت تر بحث لاندې وي، مثلاً غواړو
پوه شو چې:
\begin{enumerate}
\item هر د اشتقاق وړ~\LR{$!A$} معتبر ده؛
\item که تړلے فارمول له ځينو نورو څخه د اشتقاق وړ وي، د هغو
  معنايي پايله هم ده؛
\item که د تړلي فارمولونه يو سټ ناسازګار وي، د صدق وړ نۀ دے۔
\end{enumerate}
دا د اشتقاق نظام مهم خاصيتونه دي۔ که له دغو څخه کوم يو
سم نۀ وي، اشتقاق نظام نيمګړے دے---له حده زيات څيزونه به
اشتقاق کوي۔ له دې امله د اشتقاق نظام سموالے ثابتول
تر ټولو زيات اهميت لري۔
\label{olpseg:OLP-0081-B006:end}

\phantomsection\label{olpseg:OLP-0081-B007:start}
ځکه دغه ټول ثبوت-تيوريکي خاصيتونه په سېکوېنټ حساب کښې د ځينو
سېکوېنټونو د د اشتقاق وړتيا له لارې تعريف شوي دي، د پورته
(۱)--(۳) ثابتول غواړي چې د د اشتقاق وړ سېکوېنټونو د معنايي
خاصيتونو په اړه يو څه ثابت کړو۔ لومړے به تعريف کړو چې د سېکوېنټ
\emph{معتبر} کېدل څه معنٰی لري، او بيا به وښيو چې هر
د اشتقاق وړ سېکوېنټ معتبر دے۔ وروسته (۱)--(۳) د دې پايلې د
لازمو نتيجو په توګه راځي۔
\end{explain}
\label{olpseg:OLP-0081-B007:end}

\phantomsection\label{olpseg:OLP-0081-B008:start}
\begin{defn}
جوړښت~\LR{$\Struct
  M$} يو سېکوېنټ
\LR{$\Gamma \Sequent \Delta$} هله او يوازې هله \emph{پوره کوي} چې يا
د کوم \LR{$!A \in \Gamma$} دپاره
\LR{$\Sat/{M}{!A}$} وي، يا د کوم
\LR{$!A \in \Delta$} دپاره
\LR{$\Sat{M}{!A}$} وي۔
\label{olpseg:OLP-0081-B008:end}

\phantomsection\label{olpseg:OLP-0081-B009:start}
يو سېکوېنټ هله او يوازې هله \emph{معتبر} دے چې هر
جوړښت~\LR{$\Struct
  M$} هغه پوره کړي۔
\end{defn}
\label{olpseg:OLP-0081-B009:end}

\phantomsection\label{olpseg:OLP-0081-B010:start}
\begin{thm}[سموالے]
\ollabel{thm:sequent-soundness} که \LR{$\Log{LK}$} د
\LR{$\Theta \Sequent \Xi$} سېکوېنټ اشتقاق کړي، نو
\LR{$\Theta \Sequent \Xi$} معتبر دے۔
\end{thm}
\label{olpseg:OLP-0081-B010:end}

\phantomsection\label{olpseg:OLP-0081-B011:start}
\begin{proof}
فرض کړئ \LR{$\pi$} د \LR{$\Theta \Sequent \Xi$} اشتقاق دے۔ موږ
په~\LR{$\pi$} کښې د استنتاجونو پر شمېر~\LR{$n$} استقرا کوو۔
\label{olpseg:OLP-0081-B011:end}

\phantomsection\label{olpseg:OLP-0081-B012:start}
که د استنتاجونو شمېر~\LR{$0$} وي، نو \LR{$\pi$} يوازې له يوه لومړني
سېکوېنټ څخه جوړ دے۔ هر لومړنے سېکوېنټ
\LR{$!A \Sequent !A$} په څرګند ډول معتبر دے؛ ځکه د هر
\LR{$\Struct M$} دپاره يا
\LR{$\Sat/{M}{!A}$} وي يا
\LR{$\Sat{M}{!A}$}۔
\label{olpseg:OLP-0081-B012:end}

\phantomsection\label{olpseg:OLP-0081-B013:start}
که د استنتاجونو شمېر له~0 څخه زيات وي، د تر ټولو لاندې استنتاج
د ډول له مخې حالتونه بېلوو۔ د استقرا د فرضيې له مخې منلے شو چې د
دغه استنتاج مقدمې معتبرې دي؛ ځکه د هرې مقدمې په اشتقاق
کښې د استنتاجونو شمېر له~\LR{$n$} څخه کم دے۔
\label{olpseg:OLP-0081-B013:end}

\phantomsection\label{olpseg:OLP-0081-B014:start}
لومړے هغه ممکن استنتاجونه څېړو چې يوازې يوه مقدمه لري۔
\begin{enumerate}
\item وروستنے استنتاج کمزوري کول دي۔ بيا
  \LR{$\Theta \Sequent \Xi$} يا \LR{$!A, \Gamma \Sequent \Delta$} دے
  (که وروستنے استنتاج \LeftR{\Weakening} وي) يا
  \LR{$\Gamma \Sequent \Delta, !A$} دے
  (که \RightR{\Weakening} وي)، او اشتقاق له لاندې دواړو
  بڼو څخه په يوه پاے ته رسېږي:
  \begin{prooftree}
    \AxiomC{}
    \Deduce$\Gamma \fCenter \Delta$
    \RightLabel{\LeftR{\Weakening}}
    \UnaryInf$!A, \Gamma \fCenter \Delta$
    \DisplayProof\qquad\bottomAlignProof
    \AxiomC{}
    \Deduce$\Gamma \fCenter \Delta$
    \RightLabel{\RightR{\Weakening}}
    \UnaryInf$\Gamma \fCenter \Delta, !A$
  \end{prooftree}
  د استقرا د فرضيې له مخې \LR{$\Gamma \Sequent \Delta$} معتبر دے؛ يعنې
  د هر جوړښت~\LR{$\Struct{M}$} دپاره يا کوم
  \LR{$!C \in \Gamma$} شته چې
  \LR{$\Sat/{M}{!C}$} وي، يا کوم
  \LR{$!C \in \Delta$} شته چې
  \LR{$\Sat{M}{!C}$} وي۔
\label{olpseg:OLP-0081-B014:end}

\phantomsection\label{olpseg:OLP-0081-B015:start}
که د کوم \LR{$!C \in \Gamma$} دپاره
  \LR{$\Sat/{M}{!C}$} وي، نو
  \LR{$!C \in \Theta$} هم دے؛ ځکه \LR{$\Theta = !A, \Gamma$}، نو د کوم
  \LR{$!C \in \Theta$} دپاره
  \LR{$\Sat/{M}{!C}$} دے۔ همدارنګه، که د کوم
  \LR{$!C \in \Delta$} دپاره
  \LR{$\Sat{M}{!C}$} وي، ځکه
  \LR{$!C \in \Xi$}، نو د کوم \LR{$!C \in \Xi$} دپاره
  \LR{$\Sat{M}{!C}$} دے۔ له دې امله
  \LR{$\Theta \Sequent \Xi$} معتبر دے۔
\item وروستنے استنتاج \LeftR{\lnot} دے: بيا د وروستي استنتاج
  مقدمه \LR{$\Gamma \Sequent \Delta, !A$} او پايله
  \LR{$\lnot !A, \Gamma \Sequent \Delta$} ده، يعنې اشتقاق داسې
  پاے ته رسېږي:
  \begin{prooftree}
    \AxiomC{}
    \Deduce$\Gamma \fCenter \Delta, !A$
    \RightLabel{\LeftR{\lnot}}
    \UnaryInf$\lnot !A, \Gamma \fCenter \Delta$
  \end{prooftree}
  او \LR{$\Theta = \lnot !A, \Gamma$}، په داسې حال کښې چې
  \LR{$\Xi = \Delta$}۔
\label{olpseg:OLP-0081-B015:end}

\phantomsection\label{olpseg:OLP-0081-B016:start}
د استقرا فرضيه وايي چې \LR{$\Gamma \Sequent \Delta, !A$} معتبر دے؛
  يعنې د هر \LR{$\Struct{M}$} دپاره يا
  (الف) د کوم \LR{$!C \in \Gamma$} دپاره
  \LR{$\Sat/{M}{!C}$}، يا
  (ب) د کوم \LR{$!C \in \Delta$} دپاره
  \LR{$\Sat{M}{!C}$}، يا
  (ج) \LR{$\Sat{M}{!A}$} وي۔ غواړو وښيو چې
  \LR{$\Theta \Sequent \Xi$} هم معتبر دے۔ فرض کړئ
  \LR{$\Struct{M}$} جوړښت دے۔ که (الف) سم وي، نو داسې
  \LR{$!C \in \Gamma$} شته چې
  \LR{$\Sat/{M}{!C}$}، خو
  \LR{$!C \in \Theta$} هم دے۔ که (ب) سم وي، نو داسې
  \LR{$!C \in \Delta$} شته چې
  \LR{$\Sat{M}{!C}$}، خو
  \LR{$!C \in \Xi$} هم دے۔ په پاے کښې، که
  \LR{$\Sat{M}{!A}$} وي، نو
  \LR{$\Sat/{M}{\lnot !A}$} دے۔ ځکه
  \LR{$\lnot !A \in \Theta$}، نو داسې \LR{$!C \in \Theta$} شته چې
  \LR{$\Sat/{M}{!C}$} وي۔ له دې امله
  \LR{$\Theta \Sequent \Xi$} معتبر دے۔
\item وروستنے استنتاج \RightR{\lnot} دے: تمرين۔
\item وروستنے استنتاج \LeftR{\land} دے: دوه بڼې شته؛
  \LR{$!A \land !B$} په کيڼ لوري کښې د مقدمې له \LR{$!A$} يا \LR{$!B$} څخه
  استنتاج کېدے شي۔ په لومړي حالت کښې \LR{$\pi$} داسې پاے ته رسېږي:
  \begin{prooftree}
    \AxiomC{}
    \Deduce$!A, \Gamma \fCenter \Delta$
    \RightLabel{\LeftR{\land}}
    \UnaryInf$!A \land !B, \Gamma \fCenter \Delta$
  \end{prooftree}
  او \LR{$\Theta = !A \land !B, \Gamma$}، په داسې حال کښې چې
  \LR{$\Xi = \Delta$}۔ يو
  جوړښت~\LR{$\Struct
    M$} په پام کښې ونيسئ۔ ځکه د
  استقرا د فرضيې له مخې \LR{$!A, \Gamma \Sequent \Delta$} معتبر دے،
  يا (الف) \LR{$\Sat/{M}{!A}$}،
  يا (ب) د کوم \LR{$!C \in \Gamma$} دپاره
  \LR{$\Sat/{M}{!C}$}، يا
  (ج) د کوم \LR{$!C \in \Delta$} دپاره
  \LR{$\Sat{M}{!C}$} وي۔
  په (الف) حالت کښې
  \LR{$\Sat/{M}{!A \land !B}$}، نو داسې
  \LR{$!C \in \Theta$} (يعنې \LR{$!A \land !B$}) شته چې
  \LR{$\Sat/{M}{!C}$} وي۔ په (ب) حالت کښې داسې
  \LR{$!C \in \Gamma$} شته چې
  \LR{$\Sat/{M}{!C}$}، او
  \LR{$!C \in \Theta$} هم دے۔ په (ج) حالت کښې داسې
  \LR{$!C \in \Delta$} شته چې
  \LR{$\Sat{M}{!C}$}، او ځکه
  \LR{$\Xi = \Delta$}، \LR{$!C \in \Xi$} هم دے۔ نو په هر حالت کښې
  \LR{$\Struct M$} د
  \LR{$!A \land !B, \Gamma \Sequent \Delta$} سېکوېنټ پوره کوي۔
  \textbf{د ژباړې د نسخې سمون (OLSQ-005):} د سرچينې په ورپسې
  پايله کښې عطف لرونکې مقدمه پاتې شوې وه؛ دلته د همدغه ثابت شوي
  بشپړ سېکوېنټ اعتبار ليکل کېږي۔
  ځکه \LR{$\Struct{M}$} اختياري و،
  \LR{$!A \land !B, \Gamma \Sequent \Delta$} معتبر دے۔ هغه حالت چې
  \LR{$!A \land !B$} له \LR{$!B$} څخه استنتاج کېږي په همدې ډول، د \LR{$!A$} پر
  ځاے د \LR{$!B$} په راوړلو، حل کېږي۔
\item وروستنے استنتاج \RightR{\lor} دے: دوه بڼې شته؛
  \LR{$!A \lor !B$} په ښي لوري کښې د مقدمې له \LR{$!A$} يا \LR{$!B$} څخه
  استنتاج کېدے شي۔ په لومړي حالت کښې \LR{$\pi$} داسې پاے ته رسېږي:
  \begin{prooftree}
    \AxiomC{}
    \Deduce$\Gamma \fCenter \Delta, !A$
    \RightLabel{\RightR{\lor}}
    \UnaryInf$\Gamma \fCenter \Delta, !A \lor !B$
  \end{prooftree}
  اوس \LR{$\Theta = \Gamma$} او \LR{$\Xi = \Delta, !A \lor !B$} دي۔ يو
  جوړښت~\LR{$\Struct{M}$} په پام کښې ونيسئ۔ ځکه
  \LR{$\Gamma \Sequent \Delta, !A$} معتبر دے، يا
  (الف) \LR{$\Sat{M}{!A}$}،
  يا (ب) د کوم \LR{$!C \in \Gamma$} دپاره
  \LR{$\Sat/{M}{!C}$}، يا
  (ج) د کوم \LR{$!C \in \Delta$} دپاره
  \LR{$\Sat{M}{!C}$} وي۔
  په (الف) حالت کښې
  \LR{$\Sat{M}{!A \lor !B}$}۔ په (ب) حالت کښې
  داسې \LR{$!C \in \Gamma$} شته چې
  \LR{$\Sat/{M}{!C}$}۔ په (ج) حالت کښې داسې
  \LR{$!C \in \Delta$} شته چې
  \LR{$\Sat{M}{!C}$}۔ نو په هر حالت کښې
  \LR{$\Struct{M}$} د
  \LR{$\Gamma \Sequent \Delta, !A \lor !B$}، يعنې
  \LR{$\Theta \Sequent \Xi$}، پوره کوي۔ ځکه
  \LR{$\Struct{M}$} اختياري و،
  \LR{$\Theta \Sequent \Xi$} معتبر دے۔ هغه حالت چې \LR{$!A \lor !B$} له
  \LR{$!B$} څخه استنتاج کېږي په همدې ډول، د \LR{$!A$} پر ځاے د \LR{$!B$} په
  راوړلو، حل کېږي۔
\item وروستنے استنتاج \RightR{\lif} دے: بيا \LR{$\pi$} داسې پاے ته رسېږي:
  \begin{prooftree}
    \AxiomC{}
    \Deduce$!A, \Gamma \fCenter \Delta, !B$
    \RightLabel{\RightR{\lif}}
    \UnaryInf$\Gamma \fCenter \Delta, !A \lif !B$
  \end{prooftree}
  بيا هم د استقرا فرضيه وايي چې مقدمه معتبره ده؛ غواړو وښيو چې
  پايله هم معتبره ده۔ يو اختياري
  \LR{$\Struct{M}$} ونيسئ۔ ځکه
  \LR{$!A, \Gamma \Sequent \Delta, !B$} معتبر دے، لږ تر لږه يو له دغو
  حالتونو څخه سم دے: (الف)
  \LR{$\Sat/{M}{!A}$}، (ب)
  \LR{$\Sat{M}{!B}$}، (ج) د کوم
  \LR{$~!C \in \Gamma$} دپاره
  \LR{$\Sat/{M}{!C}$}، يا (د) د کوم
  \LR{$!C \in \Delta$} دپاره
  \LR{$\Sat{M}{!C}$}۔
  په (الف) او (ب) حالتونو کښې
  \LR{$\Sat{M}{!A \lif !B}$}، نو داسې
  \LR{$!C \in \Delta, !A \lif !B$} شته چې
  \LR{$\Sat{M}{!C}$} وي۔ په (ج) حالت کښې د کوم
  \LR{$!C \in \Gamma$} دپاره
  \LR{$\Sat/{M}{!C}$} دے۔ په (د) حالت کښې د
  کوم \LR{$!C \in \Delta$} دپاره
  \LR{$\Sat{M}{!C}$} دے۔ په هر حالت کښې
  \LR{$\Struct{M}$} د
  \LR{$\Gamma \Sequent \Delta, !A \lif !B$} سېکوېنټ پوره کوي۔ ځکه
  \LR{$\Struct{M}$} اختياري و،
  \LR{$\Gamma \Sequent \Delta, !A \lif !B$} معتبر دے۔
  %
\item وروستنے استنتاج \LeftR{\lforall} دے: بيا داسې
  فارمول~\LR{$!A(x)$} او تړلې اصطلاح~\LR{$t$} شته چې \LR{$\pi$} داسې
  پاے ته رسېږي:
  \begin{prooftree}
    \AxiomC{}
    \Deduce$!A(t), \Gamma \fCenter \Delta$
    \RightLabel{\LeftR{\lforall}}
    \UnaryInf$\lforall[x][!A(x)], \Gamma \fCenter \Delta$
  \end{prooftree}
  غواړو وښيو چې پايله
  \LR{$\lforall[x][!A(x)], \Gamma \Sequent \Delta$} معتبره ده۔ يو
  جوړښت~\LR{$\Struct M$} په پام کښې ونيسئ۔ ځکه چې مقدمه
  \LR{$!A(t), \Gamma \Sequent \Delta$} معتبره ده، يا
  (الف) \LR{$\Sat/{M}{!A(t)}$}، (ب) د کوم \LR{$!C \in \Gamma$} دپاره
  \LR{$\Sat/{M}{!C}$}، يا (ج) د کوم \LR{$!C \in \Delta$} دپاره
  \LR{$\Sat{M}{!C}$} دے۔ په (الف) حالت کښې، د
  \olref[syn][sem]{prop:quant-terms} له مخې، که
  \LR{$\Sat{M}{\lforall[x][!A(x)]}$} وي، نو \LR{$\Sat{M}{!A(t)}$} دے۔ ځکه
  \LR{$\Sat/{M}{!A(t)}$}، نو
  \LR{$\Sat/{M}{\lforall[x][!A(x)]}$}۔ په (ب) او (ج) حالتونو کښې
  \LR{$\Struct{M}$} هم \LR{$\lforall[x][!A(x)], \Gamma \Sequent \Delta$}
  پوره کوي۔ ځکه \LR{$\Struct M$} اختياري و،
  \LR{$\lforall[x][!A(x)], \Gamma \Sequent \Delta$} معتبر دے۔
\item وروستنے استنتاج \RightR{\lexists} دے: تمرين۔
\item وروستنے استنتاج \RightR{\lforall} دے: بيا داسې
  فارمول~\LR{$!A(x)$} او ثابت سمبول~\LR{$a$} شته چې \LR{$\pi$} داسې
  پاے ته رسېږي:
  \begin{prooftree}
    \AxiomC{}
    \Deduce$\Gamma \fCenter \Delta, !A(a)$
    \RightLabel{\RightR{\lforall}}
    \UnaryInf$\Gamma \fCenter \Delta, \lforall[x][!A(x)]$
  \end{prooftree}
  دلته د ځانګړي متغير شرط پوره کېږي؛ يعنې \LR{$a$} په \LR{$!A(x)$}،
  \LR{$\Gamma$} يا \LR{$\Delta$} کښې نۀ راځي۔ د استقرا د فرضيې له مخې د
  وروستي استنتاج مقدمه معتبره ده۔ بايد وښيو چې پايله هم معتبره ده؛
  يعنې د هر جوړښت~\LR{$\Struct M$} دپاره يا
  (الف) \LR{$\Sat{M}{\lforall[x][!A(x)]}$}، يا
  (ب) د کوم \LR{$!C \in \Gamma$} دپاره \LR{$\Sat/{M}{!C}$}، يا
  (ج) د کوم \LR{$!C \in \Delta$} دپاره \LR{$\Sat{M}{!C}$} وي۔
\label{olpseg:OLP-0081-B016:end}

\phantomsection\label{olpseg:OLP-0081-B017:start}
فرض کړئ \LR{$\Struct{M}$} يو اختياري جوړښت دے۔ که (ب) يا (ج)
  سم وي، کار بشپړ دے؛ نو فرض کړئ يو هم سم نۀ دے: د هر
  \LR{$!C \in \Gamma$} دپاره \LR{$\Sat{M}{!C}$}، او د هر
  \LR{$!C \in \Delta$} دپاره \LR{$\Sat/{M}{!C}$}۔ بايد وښيو چې (الف) سم دے،
  يعنې \LR{$\Sat{M}{\lforall[x][!A(x)]}$}۔ د
  \olref[syn][ass]{prop:sat-quant} له مخې کافي ده وښيو چې د ټولو
  متغير-ټاکنو~\LR{$s$} دپاره \LR{$\Sat{M}{!A(x)}[s]$}۔ نو يو اختياري
  متغير-ټاکنه~\LR{$s$} ونيسئ۔ داسې جوړښت~\LR{$\Struct{M'}$} په پام کښې
  ونيسئ چې د~\LR{$\Struct{M}$} په شان دے، پرته له دې چې
  \LR{$\Assign{a}{M'} = s(x)$}۔ د
  \olref[syn][ext]{cor:extensionality-sent} له مخې، د هر
  \LR{$!C \in \Gamma$} دپاره \LR{$\Sat{M'}{!C}$}؛ ځکه \LR{$a$} په~\LR{$\Gamma$} کښې
  نۀ راځي، او د هر \LR{$!C \in \Delta$} دپاره \LR{$\Sat/{M'}{!C}$}۔ خو
  مقدمه معتبره ده، نو \LR{$\Sat{M'}{!A(a)}$}۔ د
  \olref[syn][ass]{prop:sentence-sat-true} له مخې
  \LR{$\Sat{M'}{!A(a)}[s]$}، ځکه \LR{$!A(a)$} يوه جمله ده۔ اوس
  \LR{$\varAssign{s}{s}{x}$} او
  \LR{$s(x) = \Value{a}{M'}[s]$}، ځکه \LR{$\Struct{M'}$} مو همداسې تعريف
  کړے دے۔ نو \olref[syn][ext]{prop:ext-formulas} پلي کېږي او
  \LR{$\Sat{M'}{!A(x)}[s]$} ترلاسه کوو۔ ځکه \LR{$a$} په~\LR{$!A(x)$} کښې نۀ راځي،
  د \olref[syn][ext]{prop:extensionality} له مخې
  \LR{$\Sat{M}{!A(x)}[s]$}۔ ځکه \LR{$s$} اختياري و، دا مو ثابت کړل چې د ټولو
  متغير-ټاکنو دپاره \LR{$\Sat{M}{!A(x)}[s]$}۔
\item وروستنے استنتاج \LeftR{\lexists} دے: تمرين۔

\end{enumerate}
اوس هغه ممکن استنتاجونه څېړو چې دوه مقدمې لري۔
\begin{enumerate}
\item وروستنے استنتاج پرېکون دے: بيا \LR{$\pi$} داسې پاے ته رسېږي:
  \begin{prooftree}
    \AxiomC{}
    \Deduce$\Gamma \fCenter \Delta, !A$
    \AxiomC{}
    \Deduce$!A, \Pi \fCenter \Lambda$
    \RightLabel{\Cut}
    \BinaryInf$\Gamma, \Pi \fCenter \Delta, \Lambda$
  \end{prooftree}
  فرض کړئ \LR{$\Struct{M}$} جوړښت دے۔ د استقرا د فرضيې له مخې دواړه
  مقدمې معتبرې دي، نو
  \LR{$\Struct{M}$} دواړه پوره کوي۔
  دوه حالتونه بېلوو: (الف)
  \LR{$\Sat/{M}{!A}$} او (ب)
  \LR{$\Sat{M}{!A}$}۔ په (الف) حالت کښې، د دې
  دپاره چې \LR{$\Struct{M}$} کيڼه مقدمه
  پوره کړي، بايد \LR{$\Gamma \Sequent \Delta$} پوره کړي۔ بيا پايله هم
  پوره کوي۔ په (ب) حالت کښې، د دې دپاره چې
  \LR{$\Struct{M}$} ښۍ مقدمه پوره کړي،
  \textbf{د ژباړې د نسخې سمون (OLSQ-006):} سرچينې دلته د
  سېکوېنټ د نښې پر ځاے د سټ-توپير نښه لرله؛ د قاعدې له مقدمو سره
  سمه نښه راوستل کېږي۔
  بايد \LR{$\Pi \Sequent \Lambda$} پوره کړي۔ بيا هم
  \LR{$\Struct{M}$} پايله پوره کوي۔
\item وروستنے استنتاج \RightR{\land} دے۔ بيا \LR{$\pi$} داسې پاے ته رسېږي:
  \begin{prooftree}
    \AxiomC{}
    \Deduce$\Gamma \fCenter \Delta, !A$
    \AxiomC{}
    \Deduce$\Gamma \fCenter \Delta, !B$
    \RightLabel{\RightR{\land}}
    \BinaryInf$\Gamma \fCenter \Delta, !A \land !B$
  \end{prooftree}
  يو جوړښت~\LR{$\Struct
    M$} په پام کښې ونيسئ۔ که
  \LR{$\Struct{M}$} د
  \LR{$\Gamma \Sequent \Delta$} سېکوېنټ پوره کوي، کار بشپړ دے۔ نو فرض
  کړئ نۀ ئې پوره کوي۔ ځکه د استقرا د فرضيې له مخې
  \LR{$\Gamma \fCenter \Delta, !A$} معتبر دے،
  \LR{$\Sat{M}{!A}$}۔ همدارنګه، ځکه
  \LR{$\Gamma \Sequent \Delta, !B$} معتبر دے،
  \LR{$\Sat{M}{!B}$}۔ نو
  \LR{$\Sat{M}{!A \land !B}$}۔
\item وروستنے استنتاج \LeftR{\lor} دے: تمرين۔
\item وروستنے استنتاج \LeftR{\lif} دے۔ بيا \LR{$\pi$} داسې پاے ته رسېږي:
  \begin{prooftree}
    \AxiomC{}
    \Deduce$\Gamma \fCenter \Delta, !A$
    \AxiomC{}
    \Deduce$!B, \Pi \fCenter \Lambda$
    \RightLabel{\LeftR{\lif}}
    \BinaryInf$!A \lif !B, \Gamma, \Pi \fCenter \Delta, \Lambda$
  \end{prooftree}
  بيا هم يو
  جوړښت~\LR{$\Struct{M}$} په پام کښې ونيسئ او فرض کړئ چې
  \LR{$\Struct{M}$} د
  \LR{$\Gamma, \Pi \Sequent \Delta, \Lambda$} سېکوېنټ نۀ پوره کوي۔
  بايد وښيو چې
  \LR{$\Sat/{M}{!A \lif !B}$}۔ که
  \LR{$\Struct{M}$} د
  \LR{$\Gamma, \Pi \Sequent \Delta, \Lambda$} سېکوېنټ نۀ پوره کوي،
  نۀ \LR{$\Gamma \Sequent \Delta$} او نۀ \LR{$\Pi \Sequent \Lambda$} پوره
  کوي۔ ځکه \LR{$\Gamma \Sequent \Delta, !A$} معتبر دے،
  \LR{$\Sat{M}{!A}$}؛ او ځکه
  \LR{$!B, \Pi \Sequent \Lambda$} معتبر دے،
  \LR{$\Sat/{M}{!B}$}۔ خو بيا
  \LR{$\Sat/{M}{!A \lif !B}$}، او همدا مو
  غوښتل۔
\end{enumerate}
\end{proof}
\label{olpseg:OLP-0081-B017:end}

\phantomsection\label{olpseg:OLP-0081-B018:start}
\begin{prob}
د \olref[fol][seq][sou]{thm:sequent-soundness} ثبوت بشپړ کړئ۔
\end{prob}
\label{olpseg:OLP-0081-B018:end}



\phantomsection\label{olpseg:OLP-0081-B020:start}
\begin{cor}
\ollabel{cor:weak-soundness}
که \LR{$\Proves !A$} وي، نو \LR{$!A$} معتبر ده۔
\end{cor}
\label{olpseg:OLP-0081-B020:end}

\phantomsection\label{olpseg:OLP-0081-B021:start}
\begin{cor}
\ollabel{cor:entailment-soundness}
که \LR{$\Gamma \Proves !A$} وي، نو \LR{$\Gamma \Entails !A$}۔
\end{cor}
\label{olpseg:OLP-0081-B021:end}

\phantomsection\label{olpseg:OLP-0081-B022:start}
\begin{proof}
که \LR{$\Gamma \Proves !A$} وي، نو د کوم متناهي فرعي سټ
\LR{$\Gamma_0 \subseteq \Gamma$} دپاره د
\LR{$\Gamma_0 \Sequent !A$} اشتقاق شته۔ د
\olref{thm:sequent-soundness} له مخې، هر
جوړښت~\LR{$\Struct{M}$} يا کوم \LR{$!B \in \Gamma_0$}
ناسموي يا \LR{$!A$} رښتيا کوي۔ نو که
\LR{$\Sat{M}{\Gamma}$} وي، بيا
\LR{$\Sat{M}{!A}$} هم دے۔
\end{proof}
\label{olpseg:OLP-0081-B022:end}

\phantomsection\label{olpseg:OLP-0081-B023:start}
\begin{cor}
\ollabel{cor:consistency-soundness}
که \LR{$\Gamma$} د صدق وړ وي، نو سازګار دے۔
\end{cor}
\label{olpseg:OLP-0081-B023:end}

\phantomsection\label{olpseg:OLP-0081-B024:start}
\begin{proof}
موږ مقابل عکس ثابتوو۔ فرض کړئ \LR{$\Gamma$} سازګار نۀ دے۔ بيا يو
متناهي \LR{$\Gamma_0 \subseteq \Gamma$} او د
\LR{$\Gamma_0 \Sequent \quad$} اشتقاق شته۔ د
\olref{thm:sequent-soundness} له مخې
\LR{$\Gamma_0 \Sequent \quad$} معتبر دے۔ په بله وينا، د هر
جوړښت~\LR{$\Struct{M}$} دپاره داسې
\LR{$!C \in \Gamma_0$} شته چې
\LR{$\Sat/{M}{!C}$}، او ځکه
\LR{$\Gamma_0 \subseteq \Gamma$}، دغه \LR{$!C$} په~\LR{$\Gamma$} کښې هم دے۔
نو هېڅ \LR{$\Struct{M}$} د~\LR{$\Gamma$} سټ نۀ
پوره کوي، او \LR{$\Gamma$} د صدق وړ نۀ دے۔
\end{proof}
\label{olpseg:OLP-0081-B024:end}

\phantomsection\label{olpseg:OLP-0082-B004:start}
\olfileid{fol}{seq}{ide}
\label{olpseg:OLP-0082-B004:end}

\phantomsection\label{olpseg:OLP-0082-B005:start}
\olsection{له عينيت سره اشتقاقونه}
\label{olpseg:OLP-0082-B005:end}

\phantomsection\label{olpseg:OLP-0082-B006:start}
له عينيت سره اشتقاقونه اضافي لومړنيو سېکوېنټونو او
استنتاجي قاعدو ته اړتيا لري۔
\label{olpseg:OLP-0082-B006:end}

\phantomsection\label{olpseg:OLP-0082-B007:start}
\begin{defn}[د~\LR{$\eq$} دپاره لومړني سېکوېنټونه]
که \LR{$t$} تړلې اصطلاح وي، نو \LR{${} \Sequent \eq[t][t]$} يو لومړنے
سېکوېنټ دے۔
\end{defn}
\label{olpseg:OLP-0082-B007:end}

\phantomsection\label{olpseg:OLP-0082-B008:start}
د \LR{$\eq$} قاعدې دا دي (\LR{$t_1$} او \LR{$t_2$} تړلې اصطلاحګانې دي):
\label{olpseg:OLP-0082-B008:end}

\phantomsection\label{olpseg:OLP-0082-B009:start}
\begin{defish}
\Axiom$ \eq[t_1][t_2], \Gamma \fCenter \Delta, !A(t_1) $
\RightLabel{\LR{$\eq$}}
\UnaryInf$\eq[t_1][t_2], \Gamma \fCenter \Delta, !A(t_2)$
\DisplayProof
\hfill
\Axiom$\eq[t_1][t_2], \Gamma \fCenter \Delta, !A(t_2) $
\RightLabel{\LR{$\eq$}}
\UnaryInf$\eq[t_1][t_2], \Gamma  \fCenter \Delta, !A(t_1)$
\DisplayProof
\end{defish}
\label{olpseg:OLP-0082-B009:end}

\phantomsection\label{olpseg:OLP-0082-B010:start}
\begin{ex}
که \LR{$s$} او \LR{$t$} تړلې اصطلاحګانې وي، نو
\LR{$\eq[s][t], !A(s) \Proves !A(t)$}:
\begin{prooftree}
\Axiom$ !A(s) \fCenter !A(s)$
\RightLabel{\LeftR{\Weakening}}
\UnaryInf$\eq[s][t], !A(s)  \fCenter !A(s)$
\RightLabel{\LR{$\eq$}}
\UnaryInf$\eq[s][t], !A(s)  \fCenter !A(t)$
\end{prooftree}
دا به د يو شان څيزونو د تعويض د اصل يا د لايبنېڅ د قانون په نامه
اشنا وي۔
\label{olpseg:OLP-0082-B010:end}

\phantomsection\label{olpseg:OLP-0082-B011:start}
\LR{$\Log{LK}$} ثابتوي چې \LR{$\eq$} متناظره او انتقالي ده:
\begin{prooftree}
\Axiom$ \fCenter \eq[t_1][t_1] $
\RightLabel{\LeftR{\Weakening}}
\UnaryInf$ \eq[t_1][t_2] \fCenter \eq[t_1][t_1] $
\RightLabel{\LR{$\eq$}}
\UnaryInf$ \eq[t_1][t_2] \fCenter \eq[t_2][t_1]$
\DisplayProof\qquad\bottomAlignProof
\Axiom$ \eq[t_1][t_2] \fCenter \eq[t_1][t_2] $
\RightLabel{\LeftR{\Weakening}}
\UnaryInf$\eq[t_2][t_3], \eq[t_1][t_2]  \fCenter \eq[t_1][t_2] $
\RightLabel{\LR{$\eq$}}
\UnaryInf$\eq[t_2][t_3], \eq[t_1][t_2]  \fCenter \eq[t_1][t_3]$
\RightLabel{\LeftR{\Exchange}}
\UnaryInf$\eq[t_1][t_2], \eq[t_2][t_3]  \fCenter \eq[t_1][t_3]$
\end{prooftree}
په کيڼ اشتقاق کښې فارمول~\LR{$\eq[x][t_1]$} زموږ
\LR{$!A(x)$} دے۔ په ښي اړخ کښې \LR{$!A(x)$} د~\LR{$\eq[t_1][x]$} په توګه اخلو۔
\end{ex}
\label{olpseg:OLP-0082-B011:end}

\phantomsection\label{olpseg:OLP-0082-B012:start}
\begin{prob}
د لاندې سېکوېنټونو اشتقاقونه ورکړئ:
\begin{enumerate}
\item \LR{$\Sequent \lforall[x][\lforall[y][((x = y \land !A(x)) \lif !A(y))]]$}
\item \LR{$\lexists[x][!A(x)] \land \lforall[y][\lforall[z][((!A(y) \land
    !A(z)) \lif y = z)]] \Sequent 
\lexists[x][(!A(x) \land \lforall[y][(!A(y) \lif y = x)])]$}
\end{enumerate}
\end{prob}
\label{olpseg:OLP-0082-B012:end}

\phantomsection\label{olpseg:OLP-0083-B004:start}
\olfileid{fol}{seq}{sid}
\label{olpseg:OLP-0083-B004:end}

\phantomsection\label{olpseg:OLP-0083-B005:start}
\olsection{له عينيت سره سموالے}
\label{olpseg:OLP-0083-B005:end}

\phantomsection\label{olpseg:OLP-0083-B006:start}
\begin{prop}
\LR{$\Log{LK}$} د عينيت له لومړنيو سېکوېنټونو او قاعدو سره سم دے۔
\end{prop}
\label{olpseg:OLP-0083-B006:end}

\phantomsection\label{olpseg:OLP-0083-B007:start}
\begin{proof}
د \LR{${} \Sequent \eq[t][t]$} په بڼه لومړني سېکوېنټونه معتبر دي؛
ځکه د هر جوړښت~\LR{$\Struct M$} دپاره
\LR{$\Sat{M}{\eq[t][t]}$}۔ (پام وکړئ، فرض کوو چې اصطلاح~\LR{$t$} تړلې ده،
يعنې هېڅ متغير نۀ لري؛ نو متغير-ټاکنې اهميت نۀ لري۔)
\label{olpseg:OLP-0083-B007:end}

\phantomsection\label{olpseg:OLP-0083-B008:start}
فرض کړئ په اشتقاق کښې وروستنے استنتاج \LR{$=$} دے۔ بيا مقدمه
\LR{$\eq[t_1][t_2], \Gamma \Sequent \Delta, !A(t_1)$} او پايله
\LR{$\eq[t_1][t_2], \Gamma \Sequent \Delta, !A(t_2)$} ده۔ يو
جوړښت~\LR{$\Struct M$} په پام کښې ونيسئ۔ بايد وښيو چې پايله
معتبره ده؛ يعنې که \LR{$\Sat{M}{\eq[t_1][t_2]}$} او
\LR{$\Sat{M}{\Gamma}$} وي، نو يا د کوم \LR{$!C \in \Delta$} دپاره
\LR{$\Sat{M}{!C}$}، يا \LR{$\Sat{M}{!A(t_2)}$} دے۔
\label{olpseg:OLP-0083-B008:end}

\phantomsection\label{olpseg:OLP-0083-B009:start}
د استقرا د فرضيې له مخې مقدمه معتبره ده۔ د دې معنٰی دا ده چې که
\LR{$\Sat{M}{\eq[t_1][t_2]}$} او \LR{$\Sat{M}{\Gamma}$} وي، نو يا
(الف) د کوم \LR{$!C \in \Delta$} دپاره \LR{$\Sat{M}{!C}$}، يا
(ب) \LR{$\Sat{M}{!A(t_1)}$} دے۔ په (الف) حالت کښې کار بشپړ دے۔ (ب)
حالت وګورئ۔ داسې متغير-ټاکنه~\LR{$s$} ونيسئ چې
\LR{$s(x) = \Value{t_1}{M}$}۔ د
\olref[syn][ass]{prop:sentence-sat-true} له مخې
\LR{$\Sat{M}{!A(t_1)}[s]$}۔ ځکه \LR{$\varAssign{s}{s}{x}$}، د
\olref[syn][ext]{prop:ext-formulas} له مخې
\LR{$\Sat{M}{!A(x)}[s]$}۔ ځکه \LR{$\Sat{M}{\eq[t_1][t_2]}$}، نو
\LR{$\Value{t_1}{M} = \Value{t_2}{M}$}، او له دې امله
\LR{$s(x) = \Value{t_2}{M}$}۔ د
\olref[syn][ext]{prop:ext-formulas} په بيا پلي کولو
\LR{$\Sat{M}{!A(t_2)}[s]$} هم لرو۔ د
\olref[syn][ass]{prop:sentence-sat-true} له مخې
\LR{$\Sat{M}{!A(t_2)}$}۔
\end{proof}
\label{olpseg:OLP-0083-B009:end}

\phantomsection\label{olpseg:OLP-0084-B004:start}
\olchapter{fol}{ntd}{فطري استنتاج}
\label{olpseg:OLP-0084-B004:end}

\phantomsection\label{olpseg:OLP-0084-B005:start}
\begin{editorial}
دا څپرکے د ګېنتڅن/پراويڅ په طريقه د فطري استنتاج يو نظام وړاندې
کوي۔
\label{olpseg:OLP-0084-B005:end}

\phantomsection\label{olpseg:OLP-0084-B006:start}
د ثبوت د يوه نظام په توګه له فطري استنتاج سره اړوند مواد د شاملولو
يا ايستلو دپاره د~\texttt{prfND} ټګ وکاروئ۔
\end{editorial}
\label{olpseg:OLP-0084-B006:end}

\phantomsection\label{olpseg:OLP-0084-B009:start}
%


\label{olpseg:OLP-0084-B009:end}

\phantomsection\label{olpseg:OLP-0084-B012:start}
%


\label{olpseg:OLP-0084-B012:end}

\phantomsection\label{olpseg:OLP-0084-B016:start}
%


\label{olpseg:OLP-0084-B016:end}

\phantomsection\label{olpseg:OLP-0084-B018:start}
%



\label{olpseg:OLP-0084-B018:end}

\phantomsection\label{olpseg:OLP-0085-B004:start}
\olfileid{fol}{ntd}{rul}
\label{olpseg:OLP-0085-B004:end}

\phantomsection\label{olpseg:OLP-0085-B005:start}
\olsection{قاعدې او اشتقاقونه}
\label{olpseg:OLP-0085-B005:end}

\phantomsection\label{olpseg:OLP-0085-B006:start}
\begin{explain}
د فطري استنتاج له نظامونو څخه موخه دا ده چې په رياضيکي ثبوت کښې له
کارېدونکي غيررسمي استدلال سره نژدې موازي وي (همدا وجه ده چې تر يوه
حده “فطري” دے)۔ د فطري استنتاج ثبوتونه له فرضيو پيلېږي۔ بيا د
استنتاج قاعدې ورباندې کارول کېږي۔ فرضيې د~\Intro{\lnot}،
\Intro{\lif}، \Elim{\lor} او \Elim{\lexists} د استنتاج د قاعدو په وسيله “ايستل” کېږي،
او د روښانتيا دپاره د ايستل شوې فرضيې نښان د استنتاج تر
څنګ ليکل کېږي۔
\end{explain}
\label{olpseg:OLP-0085-B006:end}

\phantomsection\label{olpseg:OLP-0085-B007:start}
\begin{defn}[فرضيه]
يوه \emph{فرضيه} هره هغه تړلے فارمول ده چې د هرې څانګې په تر ټولو
پورتني ځاے کښې وي۔
\end{defn}
\label{olpseg:OLP-0085-B007:end}

\phantomsection\label{olpseg:OLP-0085-B008:start}
په فطري استنتاج کښې اشتقاقونه د تړلي فارمولونه ځانګړې ونې دي؛
تر ټولو پورتنۍ تړلي فارمولونه فرضيې وي، او که يوه تړلے فارمول تر
يوې، دوو يا دريو نورو جملو لاندې وي، بايد په سمه توګه د استنتاج له
يوې قاعدې څخه پايله شي۔ د استنتاج په سر کښې تړلي فارمولونه د هغې
\emph{مقدمې} او لاندې تړلے فارمول د استنتاج \emph{پايله} بلل کېږي۔
قاعدې په جوړو کښې راځي: د هر منطقي عملګر دپاره يوه د داخلولو او
يوه د ايستلو قاعده۔ دوئ په پايله کښې منطقي عملګر داخلوي يا
منطقي عملګر د قاعدې له يوې مقدمې څخه باسي۔ ځينې قاعدې اجازه
ورکوي چې د يوه ټاکلي
ډول فرضيه \emph{ايستل} شي۔ د دې د ښودلو دپاره چې کومه
فرضيه د کوم استنتاج په وسيله ايستل کېږي، فرضيې او استنتاج
دواړو ته نښانونه ورکوو۔ دا د فرضيې په “\LR{$\Discharge{!A}{n}$}” ليکلو
ښودل کېږي۔ \textbf{د ژباړې د نسخې سمون (OLND-001):} سرچينه په دې
ونه کښې لاندې ولاړ توکي “سېکوېنټونه” بولي؛ دلته “جملې” راوړل شوې،
ځکه خپله سرچينه ونه د جملو ونه تعريفوي او قاعدې پر جملو لګېږي۔
\label{olpseg:OLP-0085-B008:end}

\phantomsection\label{olpseg:OLP-0085-B009:start}
% Only include this sentence is on of \land, lor, lif or lnot is defined.
\label{olpseg:OLP-0085-B009:end}

\phantomsection\label{olpseg:OLP-0085-B010:start}
دا دود دے چې د ټولو منطقي عملګرونه~\LR{$\land$}، \LR{$\lor$}، \LR{$\lif$}،
    \LR{$\lnot$} او~\LR{$\lfalse$} دپاره قاعدې په پام کښې ونيول شي، که څه هم
    ځينې ئې تعريف شوي وي۔
\label{olpseg:OLP-0085-B010:end}

\phantomsection\label{olpseg:OLP-0086-B004:start}
\olfileid{fol}{ntd}{prl}
\label{olpseg:OLP-0086-B004:end}

\phantomsection\label{olpseg:OLP-0086-B005:start}
\olsection{قضيې قاعدې}
\label{olpseg:OLP-0086-B005:end}

\phantomsection\label{olpseg:OLP-0086-B006:start}
\subsection{د \LR{$\land$} دپاره قاعدې}
\label{olpseg:OLP-0086-B006:end}

\phantomsection\label{olpseg:OLP-0086-B007:start}
\begin{defish}
\AxiomC{\LR{$!A$}}
\AxiomC{\LR{$!B$}}
\RightLabel{\Intro{\land}}
\BinaryInfC{\LR{$!A \land !B$}}
\DisplayProof
\hfill
\begin{tabular}{r}
\AxiomC{\LR{$!A \land !B$}}
\RightLabel{\Elim{\land}}
\UnaryInfC{\LR{$!A$}}
\DisplayProof
\\[3ex]
\AxiomC{\LR{$!A \land !B$}}
\RightLabel{\Elim{\land}}
\UnaryInfC{\LR{$!B$}}
\DisplayProof
\end{tabular}
\end{defish}
\label{olpseg:OLP-0086-B007:end}

\phantomsection\label{olpseg:OLP-0086-B008:start}
\subsection{د \LR{$\lor$} دپاره قاعدې}
\label{olpseg:OLP-0086-B008:end}

\phantomsection\label{olpseg:OLP-0086-B009:start}
\begin{defish}
\begin{tabular}{r}
\AxiomC{\LR{$!A$}}
\RightLabel{\Intro{\lor}}
\UnaryInfC{\LR{$!A \lor !B$}}
\DisplayProof
\\[3ex]
\AxiomC{\LR{$!B$}}
\RightLabel{\Intro{\lor}}
\UnaryInfC{\LR{$!A \lor !B$}}
\DisplayProof
\end{tabular}
\hfill
\AxiomC{\LR{$!A \lor !B$}}
\AxiomC{\LR{$\Discharge{!A}{n}$}}
\DeduceC{\LR{$!C$}}
\AxiomC{\LR{$\Discharge{!B}{n}$}}
\DeduceC{\LR{$!C$}}
\DischargeRule{\Elim{\lor}}{n}
\TrinaryInfC{\LR{$!C$}}
\DisplayProof
\end{defish}
\label{olpseg:OLP-0086-B009:end}

\phantomsection\label{olpseg:OLP-0086-B010:start}
\subsection{د \LR{$\lif$} دپاره قاعدې}
\label{olpseg:OLP-0086-B010:end}

\phantomsection\label{olpseg:OLP-0086-B011:start}
\begin{defish}
\AxiomC{\LR{$\Discharge{!A}{n}$}}
\DeduceC{\LR{$!B$}}
\DischargeRule{\Intro{\lif}}{n}
\UnaryInfC{\LR{$!A \lif !B$}}
\DisplayProof
\hfill
\AxiomC{\LR{$!A \lif !B$}}
\AxiomC{\LR{$!A$}}
\RightLabel{\Elim{\lif}}
\BinaryInfC{\LR{$!B$}}
\DisplayProof
\end{defish}
\label{olpseg:OLP-0086-B011:end}

\phantomsection\label{olpseg:OLP-0086-B012:start}
\subsection{د \LR{$\lnot$} دپاره قاعدې}
\label{olpseg:OLP-0086-B012:end}

\phantomsection\label{olpseg:OLP-0086-B013:start}
\begin{defish}
\AxiomC{\LR{$\Discharge{!A}{n}$}}
\noLine
\DeduceC{\LR{$\lfalse$}}
\DischargeRule{\Intro{\lnot}}{n}
\UnaryInfC{\LR{$\lnot !A$}}
\DisplayProof
\hfill
\AxiomC{\LR{$\lnot !A$}}
\AxiomC{\LR{$!A$}}
\RightLabel{\Elim{\lnot}}
\BinaryInfC{\LR{$\lfalse$}}
\DisplayProof
\end{defish}
\label{olpseg:OLP-0086-B013:end}

\phantomsection\label{olpseg:OLP-0086-B014:start}
\subsection{د \LR{$\lfalse$} دپاره قاعدې}
\label{olpseg:OLP-0086-B014:end}

\phantomsection\label{olpseg:OLP-0086-B015:start}
\begin{defish}
\AxiomC{\LR{$\lfalse$}}
\RightLabel{\FalseInt}
\UnaryInfC{\LR{$!A$}}
\DisplayProof
\hfill
\AxiomC{\LR{$\Discharge{\lnot !A}{n}$}}
\DeduceC{\LR{$\lfalse$}}
\DischargeRule{\FalseCl}{n}
\UnaryInfC{\LR{$!A$}}
\DisplayProof
\end{defish}
\label{olpseg:OLP-0086-B015:end}

\phantomsection\label{olpseg:OLP-0086-B016:start}
پام وکړئ چې \LR{$\Intro{\lnot}$} او~\LR{$\FalseCl$} ډېر سره ورته دي: توپير
دا دے چې \LR{$\Intro{\lnot}$} يوه نفي شوې تړلے فارمول~\LR{$\lnot !A$} راوباسي،
خو \LR{$\FalseCl$} يوه مثبته تړلے فارمول~\LR{$!A$} راوباسي۔
\label{olpseg:OLP-0086-B016:end}

\phantomsection\label{olpseg:OLP-0086-B017:start}
کله چې يوه قاعده ښيي چې کومه فرضيه ايستل کېدے شي، موږ دا اجازه
بولو، شرط ئې نۀ بولو۔ مثلاً د~\LR{$\Intro{\lif}$} په قاعده کښې د
مقدمې~\LR{$!B$} په اشتقاق کښې د~\LR{$!A$} په بڼه هر شمېر فرضيې
ايستلے شو، چې صفر هم پکې شامل دے۔
\label{olpseg:OLP-0086-B017:end}

\phantomsection\label{olpseg:OLP-0087-B004:start}
\olfileid{fol}{ntd}{qrl}
\label{olpseg:OLP-0087-B004:end}

\phantomsection\label{olpseg:OLP-0087-B005:start}
\olsection{د سورونو قاعدې}
\label{olpseg:OLP-0087-B005:end}

\phantomsection\label{olpseg:OLP-0087-B006:start}
\subsection{د \LR{$\lforall$} دپاره قاعدې}
\label{olpseg:OLP-0087-B006:end}

\phantomsection\label{olpseg:OLP-0087-B007:start}
\begin{defish}
\AxiomC{\LR{$!A(a)$}}
\RightLabel{\Intro{\lforall}}
\UnaryInfC{\LR{$\lforall[x][\Atom{!A}{x}]$}}
\DisplayProof
\hfill
\AxiomC{\LR{$\lforall[x][\Atom{!A}{x}]$}}
\RightLabel{\Elim{\lforall}}
\UnaryInfC{\LR{$!A(t)$}}
\DisplayProof
\end{defish}
\label{olpseg:OLP-0087-B007:end}

\phantomsection\label{olpseg:OLP-0087-B008:start}
د~\LR{$\lforall$} په قاعدو کښې \LR{$t$} يوه تړلې اصطلاح ده (يعنې داسې اصطلاح
چې هېڅ متغير نۀ لري)، او \LR{$a$}~داسې ثابت سمبول دے چې په
پايله~\LR{$\lforall[x][!A(x)]$} يا په هرې هغې فرضيې کښې نۀ راځي چې د
مقدمې~\LR{$!A(a)$} په پای ته رسېدونکي اشتقاق کښې
نۀ ايستل شوې پاتې وي۔ موږ \LR{$a$} د~\Intro{\lforall} استنتاج
\emph{ځانګړے متغير} بولو۔\footnote{موږ د “ځانګړي متغير” اصطلاح
کاروو، که څه هم په پورتنۍ قاعده کښې \LR{$a$} يو ثابت دے۔ دا نوم تاريخي
لاملونه لري۔}
\label{olpseg:OLP-0087-B008:end}

\phantomsection\label{olpseg:OLP-0087-B009:start}
\subsection{د \LR{$\lexists$} دپاره قاعدې}
\label{olpseg:OLP-0087-B009:end}

\phantomsection\label{olpseg:OLP-0087-B010:start}
\begin{defish}
\AxiomC{\LR{$\Atom{!A}{t}$}}
\RightLabel{\Intro{\lexists}}
\UnaryInfC{\LR{$\lexists[x][\Atom{!A}{x}]$}}
\DisplayProof
\hfill
\AxiomC{\LR{$\lexists[x][\Atom{!A}{x}]$}}
\AxiomC{[\LR{$\Atom{!A}{a}$}]\LR{$^n$}}
\DeduceC{\LR{$!C$}}
\DischargeRule{\Elim{\lexists}}{n}
\BinaryInfC{\LR{$!C$}}
\DisplayProof
\end{defish}
\label{olpseg:OLP-0087-B010:end}

\phantomsection\label{olpseg:OLP-0087-B011:start}
بيا هم، \LR{$t$} يوه تړلې اصطلاح ده، او \LR{$a$} داسې ثابت سمبول دے چې
په مقدمه~\LR{$\lexists[x][!A(x)]$}، پايله~\LR{$!C$}، يا په هرې هغې فرضيې
کښې نۀ راځي چې د دواړو مقدمو په پای ته رسېدونکو اشتقاقونه کښې
نۀ ايستل شوې پاتې وي (د~\LR{$!A(a)$} فرضيو پرته)۔ موږ \LR{$a$} د
\Elim{\lexists} استنتاج \emph{ځانګړے متغير} بولو۔
\label{olpseg:OLP-0087-B011:end}

\phantomsection\label{olpseg:OLP-0087-B012:start}
هغه شرط چې په~\Intro{\lforall} يا~\Elim{\lexists} استنتاج کښې
ځانګړے متغير په پايله يا په هرې هغې فرضيې کښې نۀ راځي چې مقدمو ته
په رسېدونکو اشتقاقونه کښې نۀ ايستل شوې وي، او د
په وجودي مقدمه کښې هم نۀ راځي، \emph{د ځانګړي
متغير شرط} بلل کېږي۔ \textbf{د ژباړې د نسخې سمون (OLND-002):}
سرچينه په دې عمومي جمله کښې وايي چې ځانګړے متغير په “مقدمو” کښې
نۀ راځي، خو د دواړو قاعدو مرستندويه مقدمه همدا ځانګړے ثابت لري؛
دلته د سرچينې له مخکښې کره شرطونو سره سم پايله، نۀ ايستل شوې فرضيې
او د وجود ايستلو وجودي مقدمه ټاکل شوې دي۔
\label{olpseg:OLP-0087-B012:end}

\phantomsection\label{olpseg:OLP-0087-B013:start}
\begin{explain}
هغه دود راياد کړئ چې کله \LR{$!A$} داسې فارمول وي چې
متغير~\LR{$x$} پکې ازاد وي، نو دا په~\LR{$!A(x)$} ليکلو ښيو۔ په همدې
چوکاټ کښې بيا \LR{$!A(t)$} د~\LR{$\Subst{!A}{t}{x}$} لنډون دے۔ نو د
\LR{$\Intro\lexists$} قاعده داسې هم ليکلے شو:
\begin{prooftree}
\AxiomC{\LR{$\Subst{!A}{t}{x}$}}
\RightLabel{\Intro{\lexists}}
\UnaryInfC{\LR{$\lexists[x][!A]$}}
\end{prooftree}
پام وکړئ چې \LR{$t$} کېدے شي له مخکښې په~\LR{$!A$} کښې راغلے وي؛ مثلاً
\LR{$!A$}~کېدے شي~\LR{$\Atom{\Obj P}{t,x}$} وي۔ له دې امله، له
~\LR{$\Atom{\Obj P}{t,t}$} څخه \LR{$\lexists[x][\Atom{\Obj P}{t,x}]$}
استنتاجول د~\LR{$\Intro\lexists$} سمه کارونه ده---تاسو د مقدمې د~\LR{$t$}
يوه يا زياتې پېښې په تړلي متغير~\LR{$x$} “بدلولے” شئ، او اړينه
نۀ ده چې ټولې پېښې ئې بدلې کړئ۔ خو د~\LR{$\Intro\lforall$} او
~\LR{$\Elim\lexists$} د ځانګړي متغير شرطونه غواړي چې ثابت سمبول~\LR{$a$}
په~\LR{$!A$} کښې نۀ راځي۔ نو د~\LR{$\Intro\lforall$} په کارولو سره له
\LR{$\Atom{\Obj P}{a,a}$} څخه \LR{$\lforall[x][\Atom{\Obj P}{a,x}]$} په سمه
توګه نۀ شئ استنتاجولے۔
\end{explain}
\label{olpseg:OLP-0087-B013:end}

\phantomsection\label{olpseg:OLP-0087-B014:start}
\begin{explain}
په~\Intro{\lexists} او~\Elim{\lforall} کښې له تړلتيا پرته بل
بنديز نشته، او اصطلاح~\LR{$t$} هره تړلې اصطلاح کېدے شي؛ نو د بل شرط
اندېښنه نۀ شته۔ بل لوري ته، د~\Elim{\lexists} او
\Intro{\lforall} په قاعدو کښې د ځانګړي متغير شرط غواړي چې
ثابت سمبول~\LR{$a$} په پايله يا کومې نۀ ايستل شوې فرضيې کښې هېڅ
ځاے نۀ راځي۔ دا شرط د نظام د سموالي د يقيني کولو دپاره اړين دے؛
يعنې نظام له هغو نۀ ايستل شوې فرضيو څخه يوازې هغه تړلي فارمولونه
اشتقاق کړي چې ترې معنايي پايله کېږي۔ له دې شرط پرته لاندې به
روا و:
\begin{prooftree}
  \AxiomC{\LR{$\lexists[x][!A(x)]$}}
  \AxiomC{\LR{$\Discharge{!A(a)}{1}$}}
  \RightLabel{*\Intro{\lforall}}
  \UnaryInfC{\LR{$\lforall[x][!A(x)]$}}
  \RightLabel{\Elim{\lexists}}
  \BinaryInfC{\LR{$\lforall[x][!A(x)]$}}
\end{prooftree}
خو \LR{$\lexists[x][!A(x)] \Entails/ \lforall[x][!A(x)]$}۔
\label{olpseg:OLP-0087-B014:end}

\phantomsection\label{olpseg:OLP-0087-B015:start}
ځکه چې د سورونو د ايستلو قاعدې د متغيرونه پر ځاے يوازې د تړلو
اصطلاحاتو ايښوولو اجازه ورکوي، نو هره هغه فارمول چې د
تړلي فارمولونه له يوه سټ څخه وايستل شي، پخپله تړلے فارمول ده۔
\end{explain}
\label{olpseg:OLP-0087-B015:end}

\phantomsection\label{olpseg:OLP-0088-B004:start}
\olfileid{fol}{ntd}{der}
\label{olpseg:OLP-0088-B004:end}

\phantomsection\label{olpseg:OLP-0088-B005:start}
\olsection{اشتقاقونه}
\label{olpseg:OLP-0088-B005:end}

\phantomsection\label{olpseg:OLP-0088-B006:start}
\begin{explain}
موږ وويل چې فرضيه څه ده او د استنتاج قاعدې مو ورکړې۔ په فطري
استنتاج کښې اشتقاقونه له دغو څخه په استقرايي ډول جوړېږي: هر
اشتقاق يا يوازې يوه فرضيه وي، يا له يو، دوو يا دريو
اشتقاقونه او ورپسې له يوه سم استنتاج څخه جوړ وي۔
\end{explain}
\label{olpseg:OLP-0088-B006:end}

\phantomsection\label{olpseg:OLP-0088-B007:start}
\begin{defn}[اشتقاق]
 له فرضيو~\LR{$\Gamma$} څخه د تړلے فارمول~\LR{$!A$} يو \emph{اشتقاق}
د تړلي فارمولونه يوه متناهي ونه ده چې لاندې شرطونه پوره کوي:
\begin{enumerate}
\item د ونې تر ټولو پورتنۍ تړلي فارمولونه يا په~\LR{$\Gamma$} کښې وي يا
  په ونه کښې د يوه استنتاج په وسيله ايستل شوې وي۔
\item د ونې تر ټولو لاندې تړلے فارمول~\LR{$!A$} ده۔
\item په ونه کښې هره تړلے فارمول، پرته له لاندېنۍ جملې~\LR{$!A$}
  څخه، د استنتاج د يوې قاعدې د سمې کارونې مقدمه وي چې پايله ئې په
  ونه کښې نېغ د هغې تړلے فارمول لاندې ولاړه وي۔
\end{enumerate}
بيا وايو چې \LR{$!A$} د اشتقاق \emph{پايله} او~\LR{$\Gamma$} ئې
نۀ ايستل شوې فرضيې دي۔
\label{olpseg:OLP-0088-B007:end}

\phantomsection\label{olpseg:OLP-0088-B008:start}
که له~\LR{$\Gamma$} څخه د~\LR{$!A$} کوم اشتقاق وي، وايو چې \LR{$!A$}
له~\LR{$\Gamma$} څخه \emph{د اشتقاق وړ} دے، يا په نښو کښې:
\LR{$\Gamma \Proves !A$}۔ که د~\LR{$!A$} داسې اشتقاق وي چې هره
فرضيه پکې ايستل شوې وي، نو~\LR{$\Proves !A$} ليکو۔
\end{defn}
\label{olpseg:OLP-0088-B008:end}

\phantomsection\label{olpseg:OLP-0088-B009:start}
\begin{ex}
هره فرضيه په يوازې ځان اشتقاق ده۔ نو مثلاً \LR{$!A$} په يوازې
ځان اشتقاق دے او \LR{$!B$} هم په يوازې ځان همداسې دے۔
له دغو څخه د، مثلاً، \LR{$\Intro{\land}$} قاعدې په کارولو يو نوے
اشتقاق اخستلے شو:
\begin{prooftree}
\AxiomC{\LR{$!A$}}
\AxiomC{\LR{$!B$}}
\RightLabel{\Intro{\land}}
\BinaryInfC{\LR{$!A \land !B$}}
\end{prooftree}
دا قاعدې عامې دي: پکې \LR{$!A$} او~\LR{$!B$} په هر ډول تړلي فارمولونه، مثلاً
په~\LR{$!C$} او~\LR{$!D$}، بدلولے شو۔ بيا به پايله \LR{$!C \land !D$} وي؛ نو
\begin{prooftree}
\AxiomC{\LR{$!C$}}
\AxiomC{\LR{$!D$}}
\RightLabel{\Intro{\land}}
\BinaryInfC{\LR{$!C \land !D$}}
\end{prooftree}
يو سم اشتقاق دے۔ البته فرضيې سره اړولے هم شو، څو \LR{$!D$} د~\LR{$!A$}
او \LR{$!C$} د~\LR{$!B$} ونډه واخلي۔ له دې امله
\begin{prooftree}
\AxiomC{\LR{$!D$}}
\AxiomC{\LR{$!C$}}
\RightLabel{\Intro{\land}}
\BinaryInfC{\LR{$!D \land !C$}}
\end{prooftree}
هم يو سم اشتقاق دے۔
\label{olpseg:OLP-0088-B009:end}

\phantomsection\label{olpseg:OLP-0088-B010:start}
اوس بله قاعده، مثلاً \LR{$\Intro{\lif}$}، کارولے شو۔ دا اجازه راکوي چې
يو شرطي قضيه پايله کړو او هره هغه فرضيه ايستل کړو چې له
دغه شرطي قضيې له مقدم سره يو شان وي۔ نو لاندې دواړه سم
اشتقاقونه دي:
\begin{prooftree}
\AxiomC{\LR{$\Discharge{!C}{1}$}}
\AxiomC{\LR{$!D$}}
\RightLabel{\Intro{\land}}
\BinaryInfC{\LR{$!C \land !D$}}
\DischargeRule{\Intro{\lif}}{1}
\UnaryInfC{\LR{$!C \lif (!C \land !D)$}}
\DisplayProof\bottomAlignProof
\AxiomC{\LR{$!C$}}
\AxiomC{\LR{$\Discharge{!D}{1}$}}
\RightLabel{\Intro{\land}}
\BinaryInfC{\LR{$!C \land !D$}}
\DischargeRule{\Intro{\lif}}{1}
\UnaryInfC{\LR{$!D \lif (!C \land !D)$}}
\end{prooftree}
دوئ په ترتيب سره ښيي چې \LR{$!D \Proves !C \lif (!C \land !D)$} او
\LR{$!C \Proves !D \lif (!C \land !D)$}۔
\label{olpseg:OLP-0088-B010:end}

\phantomsection\label{olpseg:OLP-0088-B011:start}
راياد کړئ چې د فرضيو ايستل اجازه ده، شرط نۀ دے: لازم نۀ ده چې
فرضیې وباسو۔ په ځانګړي ډول، يوه قاعده هله هم کارولے شو چې فرضيې په
اشتقاق کښې موجودې نۀ وي۔ مثلاً لاندې اشتقاق روا دے، که څه
هم د ايستل کېدو دپاره هېڅ فرضيه~\LR{$!A$} نۀ شته:
\begin{prooftree}
  \AxiomC{\LR{$!B$}}
  \DischargeRule{\Intro{\lif}}{1}
  \UnaryInfC{\LR{$!A \lif !B$}}
\end{prooftree}
\end{ex}
\label{olpseg:OLP-0088-B011:end}

\phantomsection\label{olpseg:OLP-0089-B004:start}
\olfileid{fol}{ntd}{pro}
\label{olpseg:OLP-0089-B004:end}

\phantomsection\label{olpseg:OLP-0089-B005:start}
\olsection{د اشتقاقونه بېلګې}
\label{olpseg:OLP-0089-B005:end}

\phantomsection\label{olpseg:OLP-0089-B006:start}
\begin{ex}
راځئ چې د تړلے فارمول~\LR{$(!A \land !B) \lif !A$} يو اشتقاق
ورکړو۔
\label{olpseg:OLP-0089-B006:end}

\phantomsection\label{olpseg:OLP-0089-B007:start}
په پيل کښې مطلوبه پايله د اشتقاق په لاندې برخه کښې ليکو۔
\begin{prooftree}
\AxiomC{}
\UnaryInfC{\LR{$(!A\land !B) \lif !A$}}
\end{prooftree}
\label{olpseg:OLP-0089-B007:end}

\phantomsection\label{olpseg:OLP-0089-B008:start}
اوس بايد معلومه کړو چې د کوم ډول استنتاج پايله د دې بڼې
تړلے فارمول کېدے شي۔ د پايلې اصلي عملګر~\LR{$\lif$} دے، نو هڅه
کوو چې د~\Intro{\lif} قاعدې په کارولو پايلې ته ورسېږو۔ غوره ده چې
اړوندې فرضيې وليکو او د پرمختګ پر مهال د استنتاج قاعدو ته نښانونه
ورکړو، څو د ثبوت په پای کښې په اسانۍ ووينو چې ټولې فرضيې
ايستل شوې دي که نۀ۔
\begin{prooftree}
\AxiomC{\LR{$\Discharge{!A \land !B}{1}$}}
\DeduceC{\LR{$!A$}}
\DischargeRule{\Intro{\lif}}{1} 
\UnaryInfC{\LR{$(!A\land !B) \lif !A$}}
\end{prooftree}
\label{olpseg:OLP-0089-B008:end}

\phantomsection\label{olpseg:OLP-0089-B009:start}
اوس بايد له فرضيې~\LR{$!A \land !B$} څخه تر~\LR{$!A$} پورې پړاوونه ډک کړو۔
ځکه چې يوازې له يوه رابط~\LR{$\land$} سره کار لرو، بايد د~\LR{$\land$} د
ايستلو قاعده وکاروو۔ له دې څخه لاندې ثبوت لاس ته راځي:
\begin{prooftree}
\AxiomC{\LR{$\Discharge{!A \land !B}{1}$}}
\RightLabel{\Elim{\land}}
\UnaryInfC{\LR{$!A$}}
\DischargeRule{\Intro{\lif}}{1} 
\UnaryInfC{\LR{$(!A\land !B) \lif !A$}}
\end{prooftree}
اوس د~\LR{$(!A \land !B) \lif !A$} يو سم اشتقاق لرو۔
\end{ex}
\label{olpseg:OLP-0089-B009:end}

\phantomsection\label{olpseg:OLP-0089-B010:start}
\begin{ex}
اوس راځئ چې د~\LR{$(\lnot !A \lor !B) \lif (!A \lif !B)$} يو
اشتقاق ورکړو۔
\label{olpseg:OLP-0089-B010:end}

\phantomsection\label{olpseg:OLP-0089-B011:start}
په پيل کښې مطلوبه پايله د اشتقاق په لاندې برخه کښې ليکو۔
\begin{prooftree}
\AxiomC{}
\UnaryInfC{\LR{$(\lnot !A \lor !B) \lif (!A \lif !B)$}}
\end{prooftree}
د هغې منطقي قاعدې د موندلو دپاره چې دا پايله راکولے شي، په پايله
کښې منطقي رابطونو~\LR{$\lnot$}، \LR{$\lor$} او~\LR{$\lif$} ته ګورو۔ اوس يوازې د
\LR{$\lif$} لومړۍ پېښه راته مهمه ده، ځکه همدا د وروستۍ پايلې د
تړلے فارمول~اصلي عملګر دے؛ \LR{$\lnot$}، \LR{$\lor$} او د~\LR{$\lif$}
دويمه پېښه د بل رابط په ساحه کښې دي، نو وروسته به ئې وڅېړو۔ له دې
امله د~\Intro{\lif} له قاعدې پيل کوو۔ سمه کارونه ئې بايد داسې وي:
\textbf{د ژباړې د نسخې سمون (OLND-004):} سرچينه دلته د فطري
استنتاج وروستۍ جمله “په وروستي سېکوېنټ کښې” بولي؛ دلته “وروستۍ
پايله” راوړل شوې، ځکه دا ونه د جملو ثبوت دے او سېکوېنټ نۀ لري۔
\begin{prooftree}
\AxiomC{\LR{$\Discharge{\lnot !A \lor !B}{1}$}}
\DeduceC{\LR{$!A \lif !B$}}
\DischargeRule{\Intro{\lif}}{1}
\UnaryInfC{\LR{$(\lnot !A \lor !B) \lif (!A \lif !B)$}}
\end{prooftree}
له دې وروسته د دوام دوه لارې لرو۔ يا له لاندې څخه پورته کار روان
ساتلے او د~\Intro{\lif} بله کارونه لټولے شو، يا له پاسه لاندې کار
کولے او د~\Elim{\lor} قاعده کارولے شو۔ دويمه لار اخلو۔ فرضيه
\LR{$\lnot !A \lor !B$} د~\Elim{\lor} د تر ټولو کيڼې مقدمې په توګه
کاروو۔ د~\Elim{\lor} د سمې کارونې دپاره نورې دوه مقدمې بايد له
پايلې~\LR{$!A \lif !B$} سره يو شان وي، خو هره يوه په خپل وار له بلې
فرضیې، يعنې د~\LR{$\lnot !A \lor !B$} له دوو بېلندونو څخه له يوه،
رايستل کېدے شي۔ نو اشتقاق به داسې وي:
\begin{prooftree}
\AxiomC{\LR{$\Discharge{\lnot !A \lor !B}{1}$}}
\AxiomC{\LR{$\Discharge{\lnot !A}{2}$}}
\DeduceC{\LR{$!A \lif !B$}}
\AxiomC{\LR{$\Discharge{!B}{2}$}}
\DeduceC{\LR{$!A \lif !B$}}
\DischargeRule{\Elim{\lor}}{2}
\TrinaryInfC{\LR{$!A \lif !B$}}
\DischargeRule{\Intro{\lif}}{1} 
\UnaryInfC{\LR{$(\lnot !A \lor !B) \lif (!A \lif !B)$}}
\end{prooftree}
\label{olpseg:OLP-0089-B011:end}

\phantomsection\label{olpseg:OLP-0089-B012:start}
په ښي لوري دواړو څانګو کښې غواړو \LR{$!A \lif !B$}~اشتقاق کړو؛
غوره لار ئې د~\Intro{\lif} کارول دي۔
\begin{prooftree}
\AxiomC{\LR{$\Discharge{\lnot !A \lor !B}{1}$}}
\AxiomC{\LR{$\Discharge{\lnot !A}{2}, \Discharge{!A}{3}$}}
\DeduceC{\LR{$!B$}}
\DischargeRule{\Intro{\lif}}{3}
\UnaryInfC{\LR{$!A \lif !B$}}
\AxiomC{\LR{$\Discharge{!B}{2}, \Discharge{!A}{4}$}}
\DeduceC{\LR{$!B$}}
\DischargeRule{\Intro{\lif}}{4}
\UnaryInfC{\LR{$!A \lif !B$}}
\DischargeRule{\Elim{\lor}}{2}
\TrinaryInfC{\LR{$!A \lif !B$}}
\DischargeRule{\Intro{\lif}}{1} 
\UnaryInfC{\LR{$(\lnot !A \lor !B) \lif (!A \lif !B)$}}
\end{prooftree}
\label{olpseg:OLP-0089-B012:end}

\phantomsection\label{olpseg:OLP-0089-B013:start}
د اشتقاق په دوو تشو برخو کښې له~\LR{$\lnot !A$} او~\LR{$!A$} څخه،
او له~\LR{$!A$} او~\LR{$!B$} څخه، د~\LR{$!B$} اشتقاقونه غواړو۔ لومړۍ برخه
اخلو۔ \LR{$\lnot !A$} او~\LR{$!A$} د~\Elim{\lnot} دوه مقدمې دي:
\begin{prooftree}
\AxiomC{\LR{$\Discharge{\lnot !A}{2}$}}
\AxiomC{\LR{$\Discharge{!A}{3}$}}
\RightLabel{\Elim{\lnot}}
\BinaryInfC{\LR{$\lfalse$}}
\DeduceC{\LR{$!B$}}
\end{prooftree}
د~\FalseInt په کارولو \LR{$!B$} د پايلې په توګه اخستلے او څانګه بشپړولے
شو۔
\begin{prooftree}
\AxiomC{\LR{$\Discharge{\lnot !A \lor !B}{1}$}}
\AxiomC{\LR{$\Discharge{\lnot !A}{2}$}}
\AxiomC{\LR{$\Discharge{!A}{3}$}}
\RightLabel{\Elim{\lnot}}
\BinaryInfC{\LR{$\lfalse$}}
\RightLabel{\FalseInt}
\UnaryInfC{\LR{$!B$}}
\DischargeRule{\Intro{\lif}}{3}
\UnaryInfC{\LR{$!A \lif !B$}}
\AxiomC{\LR{$\Discharge{!B}{2}, \Discharge{!A}{4}$}}
\DeduceC{\LR{$!B$}}
\DischargeRule{\Intro{\lif}}{4}
\UnaryInfC{\LR{$!A \lif !B$}}
\DischargeRule{\Elim{\lor}}{2}
\TrinaryInfC{\LR{$!A \lif !B$}}
\DischargeRule{\Intro{\lif}}{1} 
\UnaryInfC{\LR{$(\lnot !A \lor !B) \lif (!A \lif !B)$}}
\end{prooftree}
\textbf{د ژباړې د نسخې سمون (OLND-003):} سرچينه په دې بشپړه ونه
کښې له نفي شوې مقدمې او مثبتې مقدمې څخه کاذبۍ ته استنتاج د کاذبۍ
د داخلولو په نښان ښيي؛ د مخکښې جزوي ونې او ټاکلې قاعدې سره سم دلته
د نفي ايستلو نښان راوړل شوے دے۔
\label{olpseg:OLP-0089-B013:end}

\phantomsection\label{olpseg:OLP-0089-B014:start}
اوس تر ټولو ښۍ څانګې ته ګورو۔ دلته بايد پوه شو چې د اشتقاق
تعريف \emph{د فرضيو ايستل روا بولي} خو \emph{لازمي ئې نۀ بولي}۔ په
بله وينا، که له فرضيو~\LR{$!A$} او~\LR{$!B$} څخه د يوې په کارولو \LR{$!B$} راوباسو
او بله و نۀ کاروو، دا روا ده۔ له~\LR{$!B$} څخه د~\LR{$!B$}~اشتقاق کول
څرګند دي: \LR{$!B$} په يوازې ځان داسې اشتقاق دے او هېڅ استنتاج
نۀ غواړي۔ نو فرضيه~\LR{$!A$} ساده ړنګولے شو۔
\begin{prooftree}
\AxiomC{\LR{$\Discharge{\lnot !A \lor !B}{1}$}}
\AxiomC{\LR{$\Discharge{\lnot !A}{2}$}}
\AxiomC{\LR{$\Discharge{!A}{3}$}}
\RightLabel{\Elim{\lnot}}
\BinaryInfC{\LR{$\lfalse$}}
\RightLabel{\FalseInt}
\UnaryInfC{\LR{$!B$}}
\DischargeRule{\Intro{\lif}}{3}
\UnaryInfC{\LR{$!A \lif !B$}}
\AxiomC{\LR{$\Discharge{!B}{2}$}}
\RightLabel{\Intro{\lif}}
\UnaryInfC{\LR{$!A \lif !B$}}
\DischargeRule{\Elim{\lor}}{2}
\TrinaryInfC{\LR{$!A \lif !B$}}
\DischargeRule{\Intro{\lif}}{1} 
\UnaryInfC{\LR{$(\lnot !A \lor !B) \lif (!A \lif !B)$}}
\end{prooftree}
پام وکړئ چې په بشپړ اشتقاق کښې د تر ټولو ښي اړخ
\Intro{\lif} استنتاج په حقيقت کښې هېڅ فرضيه نۀ باسي۔
\end{ex}
\label{olpseg:OLP-0089-B014:end}

\phantomsection\label{olpseg:OLP-0089-B015:start}
\begin{ex}
تر اوسه مو د~\FalseCl{} قاعدې ته اړتيا نۀ ده پېښه شوې۔ دا ځکه
ځانګړې ده چې داسې فرضيه ايستل راکوي چې د قاعدې د پايلې فرعي
فارمول نۀ وي۔ دا له~\FalseInt{} قاعدې سره نژدې تړلې ده۔ په
حقيقت کښې \FalseInt{} د~\FalseCl{} قاعدې ځانګړے حالت دے---د
“شهودي منطق” په نامه يو منطق شته چې يوازې \FalseInt{} پکې روا دے۔
کله چې بله هېڅ لار کار نۀ کوي، \FalseCl{} وروستۍ وسيله ده۔ مثلاً
فرض کړئ غواړو \LR{$!A \lor \lnot !A$}~اشتقاق کړو۔ زموږ عادي تګلاره
به دا وي چې د~\LR{$!A \lor \lnot !A$} د~\LR{$\Intro{\lor}$} په کارولو
اشتقاق کړو۔ خو دا
غواړي چې له هېڅ فرضیې پرته يا~\LR{$!A$} يا~\LR{$\lnot !A$}~اشتقاق کړو،
او دا نۀ کېږي۔ \FalseCl{} مو ژغوري!
\begin{prooftree}
  \AxiomC{\LR{$\Discharge{\lnot(!A \lor \lnot !A)}{1}$}}
  \DeduceC{\LR{$\lfalse$}}
  \DischargeRule{\FalseCl}{1}
  \UnaryInfC{\LR{$!A \lor \lnot !A$}}
\end{prooftree}
اوس له~\LR{$\lnot(!A \lor \lnot !A)$} څخه د~\LR{$\lfalse$} يو
اشتقاق لټوو۔ ځکه \LR{$\lfalse$} د~\LR{$\Elim{\lnot}$} پايله ده، نو
کېدے شي دا هڅه وکړو:
\begin{prooftree}
  \AxiomC{\LR{$\Discharge{\lnot(!A \lor \lnot !A)}{1}$}}
  \DeduceC{\LR{$\lnot !A$}}
  \AxiomC{\LR{$\Discharge{\lnot(!A \lor \lnot !A)}{1}$}}
  \DeduceC{\LR{$!A$}}
  \RightLabel{\Elim{\lnot}}
  \BinaryInfC{\LR{$\lfalse$}}
  \DischargeRule{\FalseCl}{1}
  \UnaryInfC{\LR{$!A \lor \lnot !A$}}
\end{prooftree}
د~\LR{$\lnot !A$} د اشتقاق موندلو تګلاره د~\LR{$\Intro{\lnot}$}
کارونه غواړي:
\begin{prooftree}
  \AxiomC{\LR{$\Discharge{\lnot(!A \lor \lnot !A)}{1}, \Discharge{!A}{2}$}}
  \DeduceC{\LR{$\lfalse$}}
  \DischargeRule{\Intro{\lnot}}{2}
  \UnaryInfC{\LR{$\lnot !A$}}
  \AxiomC{\LR{$\Discharge{\lnot(!A \lor \lnot !A)}{1}$}}
  \DeduceC{\LR{$!A$}}
  \RightLabel{\Elim{\lnot}}
  \BinaryInfC{\LR{$\lfalse$}}
  \DischargeRule{\FalseCl}{1}
  \UnaryInfC{\LR{$!A \lor \lnot !A$}}
\end{prooftree}
دلته \LR{$\lfalse$} په اسانۍ اخستلے شو: \LR{$\Elim{\lnot}$} پر
\LR{$\lnot(!A \lor \lnot !A)$} او~\LR{$!A \lor \lnot !A$} کاروو؛ دويمه
جمله زموږ له نوې فرضيې~\LR{$!A$} څخه د~\LR{$\Intro{\lor}$} په وسيله راځي:
\begin{prooftree}
  \AxiomC{\LR{$\Discharge{\lnot(!A \lor \lnot !A)}{1}$}}
  \AxiomC{\LR{$\Discharge{!A}{2}$}}
  \RightLabel{\Intro{\lor}}
  \UnaryInfC{\LR{$!A \lor \lnot !A$}}
  \RightLabel{\Elim{\lnot}}
  \BinaryInfC{\LR{$\lfalse$}}
  \DischargeRule{\Intro{\lnot}}{2}
  \UnaryInfC{\LR{$\lnot !A$}}
  \AxiomC{\LR{$\Discharge{\lnot(!A \lor \lnot !A)}{1}$}}
  \DeduceC{\LR{$!A$}}
  \RightLabel{\Elim{\lnot}}
  \BinaryInfC{\LR{$\lfalse$}}
  \DischargeRule{\FalseCl}{1}
  \UnaryInfC{\LR{$!A \lor \lnot !A$}}
\end{prooftree}
په ښي لوري کښې همدا تګلاره کاروو، خو \LR{$!A$} د~\FalseCl په وسيله
اخلو:
\begin{prooftree}
  \AxiomC{\LR{$\Discharge{\lnot(!A \lor \lnot !A)}{1}$}}
  \AxiomC{\LR{$\Discharge{!A}{2}$}}
  \RightLabel{\Intro{\lor}}
  \UnaryInfC{\LR{$!A \lor \lnot !A$}}
  \RightLabel{\Elim{\lnot}}
  \BinaryInfC{\LR{$\lfalse$}}
  \DischargeRule{\Intro{\lnot}}{2}
  \UnaryInfC{\LR{$\lnot !A$}}
  \AxiomC{\LR{$\Discharge{\lnot(!A \lor \lnot !A)}{1}$}}
  \AxiomC{\LR{$\Discharge{\lnot !A}{3}$}}
  \RightLabel{\Intro{\lor}}
  \UnaryInfC{\LR{$!A \lor \lnot !A$}}
  \RightLabel{\Elim{\lnot}}
  \BinaryInfC{\LR{$\lfalse$}}
  \DischargeRule{\FalseCl}{3}
  \UnaryInfC{\LR{$!A$}}
  \RightLabel{\Elim{\lnot}}
  \BinaryInfC{\LR{$\lfalse$}}
  \DischargeRule{\FalseCl}{1}
  \UnaryInfC{\LR{$!A \lor \lnot !A$}}
\end{prooftree}
\end{ex}
\label{olpseg:OLP-0089-B015:end}

\phantomsection\label{olpseg:OLP-0089-B016:start}
\begin{prob}
داسې اشتقاقونه ورکړئ چې لاندې وښيي:
\begin{enumerate}
\item \LR{$!A \land (!B \land !C) \Proves (!A \land !B) \land !C$}.
\item \LR{$!A \lor (!B \lor !C) \Proves (!A \lor !B) \lor !C$}.
\item \LR{$!A \lif (!B \lif !C) \Proves !B \lif (!A \lif !C)$}.
\item \LR{$!A \Proves \lnot\lnot !A$}.
\end{enumerate}
\end{prob}
\label{olpseg:OLP-0089-B016:end}

\phantomsection\label{olpseg:OLP-0089-B017:start}
\begin{prob}
داسې اشتقاقونه ورکړئ چې لاندې وښيي:
\begin{enumerate}
\item \LR{$(!A \lor !B) \lif !C \Proves !A \lif !C$}.
\item \LR{$(!A \lif !C) \land (!B \lif !C) \Proves (!A \lor !B) \lif !C$}.
\item \LR{$\Proves \lnot(!A \land \lnot !A)$}.
\item \LR{$!B \lif !A \Proves \lnot !A \lif \lnot !B$}.
\item \LR{$\Proves (!A \lif \lnot !A) \lif \lnot !A$}.
\item \LR{$\Proves \lnot(!A \lif !B) \lif \lnot !B$}.
\item \LR{$!A \lif !C \Proves \lnot (!A \land \lnot !C)$}.
\item \LR{$!A \land \lnot !C \Proves \lnot (!A \lif !C)$}.
\item \LR{$!A \lor !B, \lnot !B \Proves !A$}.
\item \LR{$\lnot !A \lor \lnot !B \Proves \lnot(!A \land !B)$}.
\item \LR{$\Proves (\lnot !A \land \lnot !B) \lif\lnot(!A \lor !B)$}.
\item \LR{$\Proves \lnot(!A \lor !B) \lif (\lnot !A \land \lnot !B)$}.
\end{enumerate}
\end{prob}
\label{olpseg:OLP-0089-B017:end}

\phantomsection\label{olpseg:OLP-0089-B018:start}
\begin{prob}
داسې اشتقاقونه ورکړئ چې لاندې وښيي:
\begin{enumerate}
\item \LR{$\lnot(!A \lif !B) \Proves !A$}.
\item \LR{$\lnot(!A \land !B) \Proves \lnot !A \lor \lnot !B$}.
\item \LR{$!A \lif !B \Proves \lnot !A \lor !B$}.
\item \LR{$\Proves \lnot \lnot !A \lif !A$}.
\item \LR{$!A \lif !B, \lnot !A \lif !B \Proves !B$}.
\item \LR{$(!A \land !B) \lif !C \Proves (!A \lif !C) \lor (!B \lif !C)$}.
\item \LR{$(!A \lif !B) \lif !A \Proves !A$}.
\item \LR{$\Proves (!A \lif !B) \lor (!B \lif !C)$}.
\end{enumerate}
(دا ټول د~\LR{$\FalseCl$} قاعده غواړي۔)
\end{prob}
\label{olpseg:OLP-0089-B018:end}

\phantomsection\label{olpseg:OLP-0090-B004:start}
\olfileid{fol}{ntd}{prq}
\label{olpseg:OLP-0090-B004:end}

\phantomsection\label{olpseg:OLP-0090-B005:start}
\olsection{له سورونو سره اشتقاقونه}
\label{olpseg:OLP-0090-B005:end}

\phantomsection\label{olpseg:OLP-0090-B006:start}
\begin{ex}
له سورونو سره د کار پر مهال بايد ډاډ واخلو چې د ځانګړي متغير شرط
نۀ ماتوو، او کله کله د دې دپاره د ځينو استنتاجونو د عملي کولو ترتيب
بدلول را بدلول غواړي۔ په عمومي توګه ګټه کوي چې لومړے د هغو قاعدو
د حل هڅه وکړو چې د ځانګړي متغير شرط ورباندې لګېږي (په بشپړ ثبوت
کښې به دوئ لاندې راځي)۔
\label{olpseg:OLP-0090-B006:end}

\phantomsection\label{olpseg:OLP-0090-B007:start}
راځئ وګورو چې د فارمول~\LR{$\lexists[x][\lnot !A(x)] \lif \lnot
\lforall[x][!A(x)]$} يو اشتقاق څنګه ورکوو۔ د معمول په شان
داسې پيل کوو:
\begin{prooftree}
\AxiomC{}
\UnaryInfC{\LR{$\lexists[x][\lnot !A(x)]\lif \lnot \lforall[x][!A(x)]$}}
\end{prooftree}
لومړے ليکو چې د~\Intro{\lif} د قاعدې په وسيله د وروستي پړاو د
توجيه کولو دپاره څه پکار دي۔
\begin{prooftree}
\AxiomC{\LR{$\Discharge{\lexists[x][\lnot !A(x)]}{1}$}}
\DeduceC{\LR{$\lnot \lforall[x][!A(x)]$}}
\DischargeRule{\Intro{\lif}}{1}
\UnaryInfC{\LR{$\lexists[x][\lnot !A(x)]\lif \lnot \lforall[x][!A(x)]$}}
\end{prooftree}
ځکه پر~\LR{$\lnot \lforall[x][!A(x)]$} د کارولو دپاره څرګنده قاعده نۀ
شته، اشتقاق داسې برابروو چې د~\Elim{\lexists} قاعده
وکارولے شو۔ دلته بايد د ځانګړي متغير شرط ته پام وکړو او داسې ثابت
وټاکو چې په~\LR{$\lexists[x][!A(x)]$} يا د هغې په هرې اړوندې فرضيې کښې
نۀ راځي۔ (خو ځکه چې هېڅ ثابت سمبولونه نۀ دي راغلي، هر انتخاب سم
دے۔)
\begin{prooftree}
\AxiomC{\LR{$\Discharge{\lexists[x][\lnot !A(x)]}{1}$}}
\AxiomC{\LR{$\Discharge{\lnot !A(a)}{2}$}}
\DeduceC{\LR{$\lnot \lforall[x][!A(x)]$}}
\DischargeRule{\Elim{\lexists}}{2}
\BinaryInfC{\LR{$\lnot \lforall[x][!A(x)]$}}
\DischargeRule{\Intro{\lif}}{1}
\UnaryInfC{\LR{$\lexists[x][\lnot !A(x)] \lif \lnot \lforall[x][!A(x)]$}}
\end{prooftree}
د~\LR{$\lnot \lforall[x][!A(x)]$} د راايستلو دپاره د~\Intro{\lnot}
قاعده کارول غواړو: دا غواړي چې يو تناقض راوباسو، او کېدے شي
\LR{$\lforall[x][!A(x)]$} د يوې اضافي فرضيې په توګه وکاروو۔ البته په دې
تناقض کښې هغه فرضيه~\LR{$\lnot !A(a)$} هم راتلے شي چې د
\Elim{\lexists} استنتاج به ئې وباسي۔ داسې ئې برابروو:
\begin{prooftree}
\AxiomC{\LR{$\Discharge{\lexists[x][\lnot !A(x)]}{1}$}}
\AxiomC{\LR{$\Discharge{\lnot !A(a)}{2}, \Discharge{\lforall[x][!A(x)]}{3}$}}
\DeduceC{\LR{$\lfalse$}}
\DischargeRule{\Intro{\lnot}}{3}
\UnaryInfC{\LR{$\lnot \lforall[x][!A(x)]$}}
\DischargeRule{\Elim{\lexists}}{2}
\BinaryInfC{\LR{$\lnot \lforall[x][!A(x)]$}}
\DischargeRule{\Intro{\lif}}{1}
\UnaryInfC{\LR{$\lexists[x][\lnot !A(x)]\lif \lnot \lforall[x][!A(x)]$}}
\end{prooftree}
داسې ښکاري چې تناقض ته نژدې يو۔ تر ټولو اسانه قاعده
\Elim{\lforall} ده چې د ځانګړي متغير هېڅ شرط نۀ لري۔ ځکه د عمومي
سور لرونکي~\LR{$x$} پر ځاے هره تړلې اصطلاح کارولے شو، ښه ده چې هماغه
~\LR{$a$} وکاروو څو تناقض ته ورسېږو۔
\begin{prooftree}
\AxiomC{\LR{$\Discharge{\lexists[x][\lnot !A(x)]}{1}$}}
\AxiomC{\LR{$\Discharge{\lnot !A(a)}{2}$}}
\AxiomC{\LR{$\Discharge{\lforall[x][!A(x)]}{3}$}}
\RightLabel{\Elim{\lforall}}
\UnaryInfC{\LR{$!A(a)$}}
\RightLabel{\Elim{\lnot}}
\BinaryInfC{\LR{$\lfalse$}}
\DischargeRule{\Intro{\lnot}}{3}
\UnaryInfC{\LR{$\lnot \lforall[x][!A(x)]$}}
\DischargeRule{\Elim{\lexists}}{2}
\BinaryInfC{\LR{$\lnot \lforall[x][!A(x)]$}}
\DischargeRule{\Intro{\lif}}{1}
\UnaryInfC{\LR{$\lexists[x][\lnot !A(x)]\lif \lnot \lforall[x][!A(x)]$}}
\end{prooftree}
\label{olpseg:OLP-0090-B007:end}

\phantomsection\label{olpseg:OLP-0090-B008:start}
په ځانګړي ډول له سورونو سره د کار پر مهال اوس دا بيا کتل مهم دي
چې د ځانګړي متغير شرط نۀ دے مات شوے۔ دلته دغه شرط يوازې پر
\Elim{\exists} لګېدۀ، او ځانګړے متغير~\LR{$a$} په هغو فرضيو کښې نۀ
راځي چې استنتاج ورپورې تړلے دے؛ نو دا سم اشتقاق دے۔
\end{ex}
\label{olpseg:OLP-0090-B008:end}

\phantomsection\label{olpseg:OLP-0090-B009:start}
\begin{ex}
کله کله له نورو فارمولونه څخه فارمول راايستلے شو۔ په دغو
حالتونو کښې کېدے شي نۀ ايستل شوې فرضيې ولرو۔ د خپلې وروستۍ موخې
تر څنګ د فرضيو حساب ساتل هم مهم دي۔
\label{olpseg:OLP-0090-B009:end}

\phantomsection\label{olpseg:OLP-0090-B010:start}
راځئ وګورو چې له فرضيو~\LR{$\lexists[x][(!A(x) \land !B(x))]$} او
\LR{$\lforall[x][(!B(x) \lif !C(x,b))]$} څخه د فارمول
\LR{$\lexists[x][!C(x,b)]$} يو اشتقاق څنګه ورکوو۔ د معمول په
شان پايله لاندې ليکو۔
\begin{prooftree}
\AxiomC{}
\UnaryInfC{\LR{$\lexists[x][!C(x,b)]$}}
\end{prooftree}
\label{olpseg:OLP-0090-B010:end}

\phantomsection\label{olpseg:OLP-0090-B011:start}
د کار دپاره دوه مقدمې لرو۔ د لومړۍ د کارولو، يعنې له
\LR{$\lexists[x][(!A(x) \land !B(x))]$} څخه د
\LR{$\lexists[x][!C(x, b)]$} د اشتقاق موندلو دپاره به د
\Elim{\lexists} قاعده کاروو۔ ځکه د ځانګړي متغير شرط لري، دا قاعده
لومړۍ کاروو۔ لاندې لاس ته راځي:
\begin{prooftree}
\AxiomC{\LR{$\lexists[x][(!A(x) \land !B(x))]$}}
\AxiomC{\LR{$\Discharge{!A(a) \land !B(a)}{1}$}}
\DeduceC{\LR{$\lexists[x][!C(x,b)]$}}
\DischargeRule{\Elim{\lexists}}{1}
\BinaryInfC{\LR{$\lexists[x][!C(x,b)]$}}
\end{prooftree}
هغه دوه فرضيې چې ورسره کار کوو، \LR{$!B$} شريکوي۔ ښايي اوس د
\Elim{\land} په کارولو د~\LR{$!B(a)$} بېلول ګټور وي۔
\begin{prooftree}
\AxiomC{\LR{$\lexists[x][(!A(x) \land !B(x)])$}}
\AxiomC{\LR{$\Discharge{!A(a) 
\land !B(a)}{1}$}}
\RightLabel{\Elim{\land}}
\UnaryInfC{\LR{$!B(a)$}}
\DeduceC{\LR{$\lexists[x][!C(x,b)]$}}
\DischargeRule{\Elim{\lexists}}{1}
\BinaryInfC{\LR{$\lexists[x][!C(x,b)]$}}
\end{prooftree}
\label{olpseg:OLP-0090-B011:end}

\phantomsection\label{olpseg:OLP-0090-B012:start}
د کار دپاره دويمه فرضيه~\LR{$\lforall[x][(!B(x) \lif !C(x,b))]$} ده۔
ځکه د ځانګړي متغير شرط نۀ لري، د~\Elim{\lforall} په کارولو د
ثابت سمبول~\LR{$a$} په وسيله \LR{$x$} بېلګه کولے او~\LR{$!B(a) \lif !C(a, b)$}
اخستلے شو۔ اوس \LR{$!B(a) \lif !C(a,b)$} او~\LR{$!B(a)$} دواړه لرو۔ بل ګام
بايد د~\Elim{\lif} نېغه کارونه وي۔
\begin{prooftree}
\AxiomC{\LR{$\lexists[x][(!A(x) \land !B(x))]$}}
\AxiomC{\LR{$\lforall[x][(!B(x) \lif !C(x,b))]$}}
\RightLabel{\Elim{\lforall}}
\UnaryInfC{\LR{$!B(a) \lif !C(a,b)$}}
\AxiomC{\LR{$\Discharge{!A(a) 
\land !B(a)}{1}$}}
\RightLabel{\Elim{\land}}
\UnaryInfC{\LR{$!B(a)$}}
\RightLabel{\Elim{\lif}}
\BinaryInfC{\LR{$!C(a,b)$}}
\DeduceC{\LR{$\lexists[x][!C(x,b)]$}}
\DischargeRule{\Elim{\lexists}}{1}
\insertBetweenHyps{\hspace{-5em}}
\BinaryInfC{\LR{$\lexists[x][!C(x,b)]$}}
\end{prooftree}
موخې ته ډېر نژدې يو!{} د~\Intro{\lexists} يوه کارونه موخې ته مو
رسوي۔
\begin{prooftree}
\AxiomC{\LR{$\lexists[x][(!A(x) \land !B(x))]$}}
\AxiomC{\LR{$\lforall[x][(!B(x) \lif !C(x,b))]$}}
\RightLabel{\Elim{\lforall}}
\UnaryInfC{\LR{$!B(a) \lif !C(a,b)$}}
\AxiomC{\LR{$\Discharge{!A(a) 
\land !B(a)}{1}$}}
\RightLabel{\Elim{\land}}
\UnaryInfC{\LR{$!B(a)$}}
\RightLabel{\Elim{\lif}}
\BinaryInfC{\LR{$!C(a,b)$}}
\RightLabel{\Intro{\lexists}}
\UnaryInfC{\LR{$\lexists[x][!C(x,b)]$}}
\DischargeRule{\Elim{\lexists}}{1}
\insertBetweenHyps{\hspace{-5em}}
\BinaryInfC{\LR{$\lexists[x][!C(x,b)]$}}
\end{prooftree}
ځکه په هر پړاو کښې مو ډاډ اخستے چې د ځانګړي متغير شرطونه نۀ دي
مات شوي، باور کولے شو چې دا سم اشتقاق دے۔
\end{ex}
\label{olpseg:OLP-0090-B012:end}

\phantomsection\label{olpseg:OLP-0090-B013:start}
\begin{ex}
له فرضيو~\LR{$\lforall[x][!A(x)] \lif \lexists[y][!B(y)]$} او
\LR{$\lnot\lexists[y][!B(y)]$} څخه د فارمول
\LR{$\lnot\lforall[x][!A(x)]$} يو اشتقاق ورکړئ۔ د معمول په شان
موخې فارمول لاندې ليکو۔
\begin{prooftree}
\AxiomC{}
\UnaryInfC{\LR{$\lnot\lforall[x][!A(x)]$}}
\end{prooftree}
د اشتقاق وروستۍ کرښه نفي ده، نو \Intro{\lnot} کارول غواړو۔
دا غواړي چې معلومه کړو تناقض څنګه اشتقاق کړو۔
\begin{prooftree}
\AxiomC{\LR{$\Discharge{\lforall[x][!A(x)]}{1}$}}
\DeduceC{\LR{$\lfalse$}}
\DischargeRule{\Intro{\lnot}}{1}
\UnaryInfC{\LR{$\lnot\lforall[x][!A(x)]$}}
\end{prooftree}
تر دې ځايه ښه يو۔ \Elim{\lforall} کارولے شو، خو څرګنده نۀ ده چې
موخې ته به مو ورسوي۔ پر ځاے ئې خپله يوه فرضيه کاروو۔
\LR{$\lforall[x][!A(x)] \lif \lexists[y][!B(y)]$} له
\LR{$\lforall[x][!A(x)]$} سره د~\Elim{\lif} قاعدې کارول راکوي۔
\begin{prooftree}
\AxiomC{\LR{$\lforall[x][!A(x)] \lif \lexists[y][!B(y)]$}}
\AxiomC{\LR{$\Discharge{\lforall[x][!A(x)]}{1}$}}
\RightLabel{\Elim{\lif}}
\BinaryInfC{\LR{$\lexists[y][!B(y)]$}}
\DeduceC{\LR{$\lfalse$}}
\DischargeRule{\Intro{\lnot}}{1}
\UnaryInfC{\LR{$\lnot\lforall[x][!A(x)]$}}
\end{prooftree}
اوس د کار دپاره يوه وروستۍ فرضيه لرو، او ښکاري چې د~\Elim{\lnot}
په کارولو تناقض ته مو رسولے شي۔
\begin{prooftree}
\AxiomC{\LR{$\lnot\lexists[y][!B(y)]$}}
\AxiomC{\LR{$\lforall[x][!A(x)] \lif \lexists[y][!B(y)]$}}
\AxiomC{\LR{$\Discharge{\lforall[x][!A(x)]}{1}$}}
\RightLabel{\Elim{\lif}}
\BinaryInfC{\LR{$\lexists[y][!B(y)]$}}
\RightLabel{\Elim{\lnot}}
\BinaryInfC{\LR{$\lfalse$}}
\DischargeRule{\Intro{\lnot}}{1}
\UnaryInfC{\LR{$\lnot\lforall[x][!A(x)]$}}
\end{prooftree}
\end{ex}
\label{olpseg:OLP-0090-B013:end}

\phantomsection\label{olpseg:OLP-0090-B014:start}
\begin{prob}
داسې اشتقاقونه ورکړئ چې لاندې وښيي:
\begin{enumerate}
\item \LR{$\Proves (\lforall[x][!A(x)] \land \lforall[y][!B(y)]) \lif
\lforall[z][(!A(z) \land !B(z))]$}.
\item \LR{$\Proves (\lexists[x][!A(x)] \lor \lexists[y][!B(y)]) \lif
\lexists[z][(!A(z) \lor !B(z))]$}.
\item \LR{$\lforall[x][(!A(x) \lif !B)] \Proves \lexists[y][!A(y)] \lif !B$}.
\item \LR{$\lforall[x][\lnot !A(x)] \Proves \lnot\lexists[x][!A(x)]$}.
\item \LR{$\Proves \lnot\lexists[x][!A(x)] \lif \lforall[x][\lnot !A(x)]$}.
\item \LR{$\Proves \lnot\lexists[x][\lforall[y][((!A(x,y) \lif \lnot
!A(y,y)) \land (\lnot !A(y,y) \lif !A(x,y)))]]$}.
\end{enumerate}
\end{prob}
\label{olpseg:OLP-0090-B014:end}

\phantomsection\label{olpseg:OLP-0090-B015:start}
\begin{prob}
داسې اشتقاقونه ورکړئ چې لاندې وښيي:
\begin{enumerate}
\item \LR{$\Proves \lnot\lforall[x][!A(x)] \lif \lexists[x][\lnot!A(x)]$}.
\item \LR{$(\lforall[x][!A(x)] \lif !B) \Proves \lexists[y][(!A(y) \lif !B)]$}.
\item \LR{$\Proves \lexists[x][(!A(x) \lif \lforall[y][!A(y)])]$}.
\end{enumerate}
(دا ټول د~\LR{$\FalseCl$} قاعده غواړي۔)
\end{prob}
\label{olpseg:OLP-0090-B015:end}

\phantomsection\label{olpseg:OLP-0091-B004:start}
\olfileid{fol}{ntd}{ptn}
\label{olpseg:OLP-0091-B004:end}

\phantomsection\label{olpseg:OLP-0091-B005:start}
\olsection{ثبوت-تيوريکي مفاهيم}
\label{olpseg:OLP-0091-B005:end}

\phantomsection\label{olpseg:OLP-0091-B006:start}
\begin{editorial}
دا برخه د فطري استنتاج د اثبات‌پذيرۍ اړيکې او سازګارۍ تعريفونه
راټولوي۔
\end{editorial}
\label{olpseg:OLP-0091-B006:end}

\phantomsection\label{olpseg:OLP-0091-B007:start}
\begin{explain}
لکه څنګه چې مو يو شمېر مهم معنايي مفاهيم (اعتبار، معنايي استلزام،
د صدق وړتيا) تعريف کړل، اوس ورسره سم \emph{ثبوت-تيوريکي مفاهيم}
تعريفوو۔ دا مفاهيم په جوړښتونه کښې د تړلي فارمولونه د پوره
کېدو له لارې نۀ، بلکې له نورو څخه د ځينو تړلي فارمولونه د
د اشتقاق وړتيا يا د اشتقاق نۀ وړتيا له لارې تعريفېږي۔ دا مهمه
موندنه وه چې دغه مفاهيم سره سمون خوري۔ د دې سمون منځپانګه د
\emph{سموالي} او \emph{بشپړتيا قضيې} دي۔
\end{explain}
\label{olpseg:OLP-0091-B007:end}

\phantomsection\label{olpseg:OLP-0091-B008:start}
\begin{defn}[قضيې]
يوه تړلے فارمول~\LR{$!A$} هله \emph{قضيه} ده چې په فطري استنتاج کښې
د~\LR{$!A$} داسې اشتقاق وي چې ټولې فرضيې پکې ايستل
شوې وي۔ که \LR{$!A$} قضيه وي، \LR{$\Proves !A$} ليکو، او که نۀ وي
\LR{$\Proves/ !A$} ليکو۔
\end{defn}
\label{olpseg:OLP-0091-B008:end}

\phantomsection\label{olpseg:OLP-0091-B009:start}
\begin{defn}[د اشتقاق وړتيا]
تړلے فارمول~\LR{$!A$} د تړلي فارمولونه له سټ~\LR{$\Gamma$} څخه،
\LR{$\Gamma \Proves !A$}، هله \emph{د اشتقاق وړ} ده چې داسې
اشتقاق وي چې پايله ئې~\LR{$!A$} وي او هره فرضيه پکې يا
ايستل شوې وي يا په~\LR{$\Gamma$} کښې وي۔ که \LR{$!A$} له~\LR{$\Gamma$}
څخه د اشتقاق وړ نۀ وي، \LR{$\Gamma \Proves/ !A$} ليکو۔
\end{defn}
\label{olpseg:OLP-0091-B009:end}

\phantomsection\label{olpseg:OLP-0091-B010:start}
\begin{defn}[سازګاري]
د تړلي فارمولونه يو سټ~\LR{$\Gamma$} هله او يوازې هله \emph{ناسازګار}
دے چې \LR{$\Gamma \Proves \lfalse$} وي۔ که \LR{$\Gamma$} ناسازګار نۀ وي،
يعنې \LR{$\Gamma \Proves/ \lfalse$} وي، نو \emph{سازګار} ئې بولو۔
\end{defn}
\label{olpseg:OLP-0091-B010:end}

\phantomsection\label{olpseg:OLP-0091-B011:start}
\begin{prop}[انعکاسيت]
\ollabel{prop:reflexivity}
که \LR{$!A \in \Gamma$} وي، نو \LR{$\Gamma \Proves !A$}۔
\end{prop}
\label{olpseg:OLP-0091-B011:end}

\phantomsection\label{olpseg:OLP-0091-B012:start}
\begin{proof}
فرضيه~\LR{$!A$} په يوازې ځان د~\LR{$!A$} داسې اشتقاق دے چې هره
نۀ ايستل شوې فرضيه ئې (يعنې \LR{$!A$}) په~\LR{$\Gamma$} کښې ده۔
\end{proof}
\label{olpseg:OLP-0091-B012:end}

\phantomsection\label{olpseg:OLP-0091-B013:start}
\begin{prop}[يکنواختي]
\ollabel{prop:monotonicity}
که \LR{$\Gamma \subseteq \Delta$} او \LR{$\Gamma \Proves !A$} وي، نو
\LR{$\Delta \Proves !A$}۔
\end{prop}
\label{olpseg:OLP-0091-B013:end}

\phantomsection\label{olpseg:OLP-0091-B014:start}
\begin{proof}
له~\LR{$\Gamma$} څخه د~\LR{$!A$} هر اشتقاق له~\LR{$\Delta$} څخه هم د~\LR{$!A$}
اشتقاق دے۔
\end{proof}
\label{olpseg:OLP-0091-B014:end}

\phantomsection\label{olpseg:OLP-0091-B015:start}
\begin{prop}[انتقاليت]
\ollabel{prop:transitivity}
که \LR{$\Gamma \Proves !A$} او \LR{$\{!A\} \cup \Delta \Proves !B$} وي، نو
\LR{$\Gamma \cup \Delta \Proves !B$}۔
\end{prop}
\label{olpseg:OLP-0091-B015:end}

\phantomsection\label{olpseg:OLP-0091-B016:start}
\begin{proof}
که \LR{$\Gamma \Proves !A$} وي، د~\LR{$!A$} داسې اشتقاق~\LR{$\delta_0$}
شته چې ټولې نۀ ايستل شوې فرضيې ئې په~\LR{$\Gamma$} کښې دي۔ که
\LR{$\{!A\} \cup \Delta \Proves !B$} وي، د~\LR{$!B$} داسې
اشتقاق~\LR{$\delta_1$} شته چې ټولې نۀ ايستل شوې فرضيې ئې
په~\LR{$\{!A\} \cup \Delta$} کښې دي۔ اوس لاندې وګورئ:
\begin{prooftree}
  \AxiomC{\LR{$\Delta, \Discharge{!A}{1}$}}
  \RightLabel{\LR{$\delta_1$}}
  \DeduceC{\LR{$!B$}}
  \DischargeRule{\Intro{\lif}}{1}
  \UnaryInfC{\LR{$!A \lif !B$}}
  \AxiomC{\LR{$\Gamma$}}
  \RightLabel{\LR{$\delta_0$}}
  \DeduceC{\LR{$!A$}}
  \RightLabel{\Elim{\lif}}
  \BinaryInfC{\LR{$!B$}}
\end{prooftree}
اوس ټولې نۀ ايستل شوې فرضيې په~\LR{$\Gamma \cup \Delta$} کښې دي؛
نو دا \LR{$\Gamma \cup \Delta \Proves !B$} ښيي۔
\end{proof}
\label{olpseg:OLP-0091-B016:end}

\phantomsection\label{olpseg:OLP-0091-B017:start}
کله چې \LR{$\Gamma = \{!A_1, !A_2, \ldots, !A_k\}$} يو متناهي سټ وي،
د~\LR{$\Gamma \Proves !B$} دپاره لنډه نښه
\LR{$!A_1,!A_2,\ldots,!A_k \Proves !B$} کارولے شو؛ په ځانګړي ډول
\LR{$!A \Proves !B$} دا معنٰی لري چې \LR{$\{!A\} \Proves !B$}۔
\label{olpseg:OLP-0091-B017:end}

\phantomsection\label{olpseg:OLP-0091-B018:start}
پام وکړئ که \LR{$\Gamma \Proves !A$} او \LR{$!A \Proves !B$} وي، نو
\LR{$\Gamma \Proves !B$}۔ همدارنګه، که
\LR{$!A_1, \dots, !A_n \Proves !B$} او د هر~\LR{$i$} دپاره
\LR{$\Gamma \Proves !A_i$} وي، نو \LR{$\Gamma \Proves !B$}۔
\label{olpseg:OLP-0091-B018:end}

\phantomsection\label{olpseg:OLP-0091-B019:start}
\begin{prop}
\ollabel{prop:incons}
لاندې شرطونه معادل دي۔
    \begin{enumerate}
      \item \( \Gamma \) ناسازګار دے۔
      \item د هرې تړلے فارمول~\( {!A} \) دپاره \( \Gamma \Proves {!A} \)۔
      \item د کومې تړلے فارمول~\( {!A} \) دپاره \( \Gamma \Proves {!A} \) او \( \Gamma \Proves \lnot {!A} \)۔
    \end{enumerate}
\end{prop}
\label{olpseg:OLP-0091-B019:end}

\phantomsection\label{olpseg:OLP-0091-B020:start}
\begin{proof}
تمرين۔
\end{proof}
\label{olpseg:OLP-0091-B020:end}

\phantomsection\label{olpseg:OLP-0091-B021:start}
\begin{prob}
\olref[fol][ntd][ptn]{prop:incons} ثابت کړئ۔
\end{prob}
\label{olpseg:OLP-0091-B021:end}



\phantomsection\label{olpseg:OLP-0091-B023:start}
\begin{prop}[تراکم]
  \ollabel{prop:proves-compact}
  \begin{enumerate}
  \item که \LR{$\Gamma \Proves !A$} وي، نو داسې متناهي فرعي سټ
    \LR{$\Gamma_0 \subseteq \Gamma$} شته چې \LR{$\Gamma_0 \Proves !A$} وي۔
  \item که د~\LR{$\Gamma$} هر متناهي فرعي سټ سازګار وي، نو \LR{$\Gamma$}
    سازګار دے۔
  \end{enumerate}
\end{prop}
\label{olpseg:OLP-0091-B023:end}

\phantomsection\label{olpseg:OLP-0091-B024:start}
\begin{proof}
  \begin{enumerate}
    \item که \LR{$\Gamma \Proves !A$} وي، نو له~\LR{$\Gamma$} څخه د~\LR{$!A$}
      کوم اشتقاق~\LR{$\delta$} شته۔ \LR{$\Gamma_0$} د~\LR{$\delta$} د
      نۀ ايستل شوې فرضيو سټ ونيسئ۔ ځکه هر اشتقاق متناهي
      دے، \LR{$\Gamma_0$} يوازې متناهي ډېرې تړلي فارمولونه لرلے شي۔ نو
      \LR{$\delta$} له متناهي~\LR{$\Gamma_0 \subseteq \Gamma$} څخه د~\LR{$!A$}
      اشتقاق دے۔
    \item دا د~(۱) د هغه ځانګړي حالت نقيض عکس دے چې
      \LR{$!A \ident \lfalse$} وي۔
  \end{enumerate}
\end{proof}
\label{olpseg:OLP-0091-B024:end}

\phantomsection\label{olpseg:OLP-0092-B004:start}
\olfileid{fol}{ntd}{prv}
\label{olpseg:OLP-0092-B004:end}

\phantomsection\label{olpseg:OLP-0092-B005:start}
\olsection{د اشتقاق وړتيا او سازګاري}
\label{olpseg:OLP-0092-B005:end}

\phantomsection\label{olpseg:OLP-0092-B006:start}
اوس به د د اشتقاق وړتيا اړيکې يو شمېر خاصيتونه ثابت کړو۔ دوئ په
خپلواکه توګه په زړه پورې دي، خو هر يو به د بشپړتيا قضيې په ثبوت
کښې ونډه ولري۔
\label{olpseg:OLP-0092-B006:end}

\phantomsection\label{olpseg:OLP-0092-B007:start}
\begin{prop}\ollabel{prop:provability-contr}
  که \LR{$\Gamma \Proves !A$} وي او \LR{$\Gamma \cup \{!A\}$} ناسازګار وي،
  نو \LR{$\Gamma$} ناسازګار دے۔
\end{prop}
\label{olpseg:OLP-0092-B007:end}

\phantomsection\label{olpseg:OLP-0092-B008:start}
\begin{proof}
له~\LR{$\Gamma$} څخه د~\LR{$!A$} اشتقاق~\LR{$\delta_1$} او له
\LR{$\Gamma \cup \{!A\}$} څخه د~\LR{$\lfalse$} اشتقاق~\LR{$\delta_2$}
ونوموئ۔ بيا داسې اشتقاق کولے شو:
\begin{prooftree}
\AxiomC{\LR{$\Gamma, \Discharge{!A}{1}$}}
\RightLabel{\LR{$\delta_2$}}
\DeduceC{\LR{$\lfalse$}}
\DischargeRule{\Intro{\lnot}}{1}
\UnaryInfC{\LR{$\lnot !A$}}
\AxiomC{\LR{$\Gamma$}}
\RightLabel{\LR{$\delta_1$}}
\DeduceC{\LR{$!A$}}
\RightLabel{\Elim{\lnot}}
\BinaryInfC{\LR{$\lfalse$}}
\end{prooftree}
په نوي اشتقاق کښې فرضيه~\LR{$!A$}~ايستل شوې ده، نو دا
له~\LR{$\Gamma$} څخه اشتقاق دے۔
\end{proof}
\label{olpseg:OLP-0092-B008:end}

\phantomsection\label{olpseg:OLP-0092-B009:start}
\begin{prop}
\ollabel{prop:prov-incons}
\LR{$\Gamma \Proves !A$} هله او يوازې هله چې
\LR{$\Gamma \cup \{\lnot !A\}$} ناسازګار وي۔
\end{prop}
\label{olpseg:OLP-0092-B009:end}

\phantomsection\label{olpseg:OLP-0092-B010:start}
\begin{proof}
لومړے فرض کړئ \LR{$\Gamma \Proves !A$}؛ يعنې د~\LR{$!A$} داسې
اشتقاق~\LR{$\delta_0$} شته چې نۀ ايستل شوې فرضيې ئې
\LR{$\Gamma$} دي۔ له~\LR{$\Gamma \cup \{\lnot !A\}$} څخه د~\LR{$\lfalse$} يو
اشتقاق داسې اخلو:
\begin{prooftree}
  \AxiomC{\LR{$\lnot !A$}}
  \AxiomC{\LR{$\Gamma$}}
  \RightLabel{\LR{$\delta_0$}}
  \DeduceC{\LR{$!A$}}
  \RightLabel{\Elim{\lnot}}
  \BinaryInfC{\LR{$\lfalse$}}
\end{prooftree}
\label{olpseg:OLP-0092-B010:end}

\phantomsection\label{olpseg:OLP-0092-B011:start}
اوس فرض کړئ \LR{$\Gamma \cup \{\lnot !A\}$} ناسازګار دے، او
\LR{$\delta_1$} له~\LR{$\Gamma \cup \{\lnot !A\}$} کښې له نۀ ايستل شوې
فرضيو څخه د~\LR{$\lfalse$} اړوند اشتقاق دے۔ يوازې له~\LR{$\Gamma$}
څخه د~\LR{$!A$} يو اشتقاق د~\LR{$\FalseCl$} په کارولو اخلو:
\begin{prooftree}
  \AxiomC{\LR{$\Gamma, \Discharge{\lnot !A}{1}$}}
  \RightLabel{\LR{$\delta_1$}}
  \DeduceC{\LR{$\lfalse$}}
  \DischargeRule{\FalseCl}{1}
  \UnaryInfC{\LR{$!A$}}
\end{prooftree}
\textbf{د ژباړې د نسخې سمون (OLND-005):} سرچينه په دې ونه کښې
له ايستلو نښان سره د همدې قاعدې يو دويم ښي نښان هم ږدي؛ زيات نښان
لرې شوے او د فرضيې ايستل او د قاعدې نوم په ټاکلي ګډ نښان کښې ساتل
شوي دي۔
\end{proof}
\label{olpseg:OLP-0092-B011:end}

\phantomsection\label{olpseg:OLP-0092-B012:start}
\begin{prob}
ثابت کړئ چې \LR{$\Gamma \Proves \lnot !A$} هله او يوازې هله چې
\LR{$\Gamma \cup \{!A\}$} ناسازګار وي۔
\end{prob}
\label{olpseg:OLP-0092-B012:end}

\phantomsection\label{olpseg:OLP-0092-B013:start}
\begin{prop}\ollabel{prop:explicit-inc}
  که \LR{$\Gamma \Proves !A$} او \LR{$\lnot !A \in \Gamma$} وي، نو \LR{$\Gamma$}
  ناسازګار دے۔
\end{prop}
\label{olpseg:OLP-0092-B013:end}

\phantomsection\label{olpseg:OLP-0092-B014:start}
\begin{proof}
  فرض کړئ \LR{$\Gamma \Proves !A$} او \LR{$\lnot !A \in \Gamma$}۔ بيا
  له~\LR{$\Gamma$} څخه د~\LR{$!A$} کوم اشتقاق~\LR{$\delta$} شته۔ د
  \LR{$\Elim{\lnot}$} قاعدې دا ساده کارونه وګورئ:
  \begin{prooftree}
    \AxiomC{\LR{$\lnot !A$}}
    \AxiomC{\LR{$\Gamma$}}
    \RightLabel{\LR{$\delta$}}
    \DeduceC{\LR{$!A$}}
    \RightLabel{\Elim{\lnot}}
    \BinaryInfC{\LR{$\lfalse$}}
  \end{prooftree}
  ځکه \LR{$\lnot !A \in \Gamma$} او ټولې نۀ ايستل شوې فرضيې
  په~\LR{$\Gamma$} کښې دي، دا \LR{$\Gamma \Proves \lfalse$} ښيي۔
\end{proof}
\label{olpseg:OLP-0092-B014:end}

\phantomsection\label{olpseg:OLP-0092-B015:start}
\begin{prop}\ollabel{prop:provability-exhaustive}
  که \LR{$\Gamma \cup \{!A\}$} او \LR{$\Gamma \cup \{\lnot !A\}$} دواړه
  ناسازګار وي، نو \LR{$\Gamma$} ناسازګار دے۔
\end{prop}
\label{olpseg:OLP-0092-B015:end}

\phantomsection\label{olpseg:OLP-0092-B016:start}
\begin{proof}
په ترتيب سره له~\LR{$\Gamma \cup \{ !A \}$} څخه د~\LR{$\lfalse$} او له
\LR{$\Gamma \cup \{ \lnot !A \}$} څخه د~\LR{$\lfalse$}~اشتقاقونه
\LR{$\delta_1$} او~\LR{$\delta_2$} شته۔ بيا داسې اشتقاق کولے شو:
\begin{prooftree}
\AxiomC{\LR{$\Gamma, \Discharge{\lnot !A}{2}$}}
\RightLabel{\LR{$\delta_2$}}
\DeduceC{\LR{$\lfalse$}}
\DischargeRule{\Intro{\lnot}}{2}
\UnaryInfC{\LR{$\lnot \lnot !A$}}
\AxiomC{\LR{$\Gamma, \Discharge{!A}{1}$}}
\RightLabel{\LR{$\delta_1$}}
\DeduceC{\LR{$\lfalse$}}
\DischargeRule{\Intro{\lnot}}{1}
\UnaryInfC{\LR{$\lnot !A$}}
\RightLabel{\Elim{\lnot}}
\BinaryInfC{\LR{$\lfalse$}}
\end{prooftree}
ځکه فرضيې~\LR{$!A$} او~\LR{$\lnot !A$}~ايستل شوې دي، دا يوازې
له~\LR{$\Gamma$} څخه د~\LR{$\lfalse$} يو اشتقاق دے۔ نو \LR{$\Gamma$}
ناسازګار دے۔
\end{proof}
\label{olpseg:OLP-0092-B016:end}

\phantomsection\label{olpseg:OLP-0093-B005:start}
\olfileid{fol}{ntd}{ppr}
\label{olpseg:OLP-0093-B005:end}

\phantomsection\label{olpseg:OLP-0093-B006:start}
\olsection{د اشتقاق وړتيا او قضيې رابطونه}
\label{olpseg:OLP-0093-B006:end}

\phantomsection\label{olpseg:OLP-0093-B007:start}
\begin{explain}
  موږ ثابتوو چې د فطري استنتاج د د اشتقاق وړتيا اړيکه~\LR{$\Proves$}
  دومره پياوړې ده چې د قضيې رابطونو په اړه ځينې بنسټيز حقيقتونه،
  لکه \LR{$!A \land !B \Proves !A$} او
  \LR{$!A, !A \lif !B \Proves !B$} (قياسي استنتاج)، ثابت کړي۔ دا
  حقيقتونه د بشپړتيا قضيې د ثبوت دپاره پکار دي۔
\end{explain}
\label{olpseg:OLP-0093-B007:end}

\phantomsection\label{olpseg:OLP-0093-B008:start}
\begin{prop}\ollabel{prop:provability-land}
  \begin{enumerate}
  \item \ollabel{prop:provability-land-left} دواړه
    \LR{$!A \land !B \Proves !A$} او \LR{$!A \land !B \Proves !B$}۔
  \item \ollabel{prop:provability-land-right}
    \LR{$!A, !B \Proves !A \land !B$}۔
  \end{enumerate}
\end{prop}
\label{olpseg:OLP-0093-B008:end}

\phantomsection\label{olpseg:OLP-0093-B009:start}
\begin{proof}
  \begin{enumerate}
  \item دواړه داسې اشتقاق کولے شو:
    \begin{prooftree}
      \AxiomC{\LR{$!A \land !B$}}
      \RightLabel{\Elim{\land}}
      \UnaryInfC{\LR{$!A$}}
      \DisplayProof\qquad\bottomAlignProof
      \AxiomC{\LR{$!A \land !B$}}
      \RightLabel{\Elim{\land}}
      \UnaryInfC{\LR{$!B$}}
    \end{prooftree}
  \item داسې ئې اشتقاق کولے شو:
    \begin{prooftree}
      \AxiomC{\LR{$!A$}}
      \AxiomC{\LR{$!B$}}
      \RightLabel{\Intro{\land}}
      \BinaryInfC{\LR{$!A \land !B$}}
    \end{prooftree}
  \end{enumerate}
\end{proof}
\label{olpseg:OLP-0093-B009:end}

\phantomsection\label{olpseg:OLP-0093-B010:start}
\begin{prop}\ollabel{prop:provability-lor}
  \begin{enumerate}
  \item \LR{$!A \lor !B, \lnot !A, \lnot !B$} ناسازګار دے۔
  \item دواړه \LR{$!A \Proves !A \lor !B$} او \LR{$!B \Proves !A \lor !B$}۔
  \end{enumerate}
\end{prop}
\label{olpseg:OLP-0093-B010:end}

\phantomsection\label{olpseg:OLP-0093-B011:start}
\begin{proof}
  \begin{enumerate}
  \item لاندې اشتقاق وګورئ:
    \begin{prooftree}
      \AxiomC{\LR{$!A \lor !B$}}
      \AxiomC{\LR{$\lnot !A$}}
      \AxiomC{\LR{$\Discharge{!A}{1}$}}
      \RightLabel{\Elim{\lnot}}
      \BinaryInfC{\LR{$\lfalse$}}
      \AxiomC{\LR{$\lnot !B$}}
      \AxiomC{\LR{$\Discharge{!B}{1}$}}
      \RightLabel{\Elim{\lnot}}
      \BinaryInfC{\LR{$\lfalse$}}
      \DischargeRule{\Elim{\lor}}{1}
      \TrinaryInfC{\LR{$\lfalse$}}
    \end{prooftree}
    دا له نۀ ايستل شوې فرضيو~\LR{$!A \lor !B$}، \LR{$\lnot !A$} او
    \LR{$\lnot !B$} څخه د~\LR{$\lfalse$} يو اشتقاق دے۔
  \item دواړه داسې اشتقاق کولے شو:
    \begin{prooftree}
      \AxiomC{\LR{$!A$}}
      \RightLabel{\Intro{\lor}}
      \UnaryInfC{\LR{$!A \lor !B$}}
      \DisplayProof\qquad\bottomAlignProof
      \AxiomC{\LR{$!B$}}
      \RightLabel{\Intro{\lor}}
      \UnaryInfC{\LR{$!A \lor !B$}}
    \end{prooftree}
  \end{enumerate}
\end{proof}
\label{olpseg:OLP-0093-B011:end}

\phantomsection\label{olpseg:OLP-0093-B012:start}
\begin{prop}\ollabel{prop:provability-lif}
  \begin{enumerate}
  \item \ollabel{prop:provability-lif-left}
    \LR{$!A, !A \lif !B \Proves !B$}۔
  \item \ollabel{prop:provability-lif-right}
    دواړه \LR{$\lnot !A \Proves !A \lif !B$} او
    \LR{$!B \Proves !A \lif !B$}۔
  \end{enumerate}
\end{prop}
\label{olpseg:OLP-0093-B012:end}

\phantomsection\label{olpseg:OLP-0093-B013:start}
\begin{proof}
  \begin{enumerate}
  \item داسې ئې اشتقاق کولے شو:
    \begin{prooftree}
      \AxiomC{\LR{$!A \lif !B$}}
      \AxiomC{\LR{$!A$}}
      \RightLabel{\Elim{\lif}}
      \BinaryInfC{\LR{$!B$}}
    \end{prooftree}
\label{olpseg:OLP-0093-B013:end}

\phantomsection\label{olpseg:OLP-0093-B014:start}
\item دا د لاندې دوو اشتقاقونه په وسيله ښودل کېږي:
    \begin{prooftree}
      \AxiomC{\LR{$\lnot !A$}}
      \AxiomC{\LR{$\Discharge{!A}{1}$}}
      \RightLabel{\Elim{\lnot}}
      \BinaryInfC{\LR{$\lfalse$}}
      \RightLabel{\FalseInt}
      \UnaryInfC{\LR{$!B$}}
      \DischargeRule{\Intro{\lif}}{1}
      \UnaryInfC{\LR{$!A \lif !B$}}
      \DisplayProof\qquad\bottomAlignProof
      \AxiomC{\LR{$!B$}}
      \RightLabel{\Intro{\lif}}
      \UnaryInfC{\LR{$!A \lif !B$}}
    \end{prooftree}
    پام وکړئ چې \LR{$\Intro{\lif}$} د فرضيې~\LR{$!A$}~ايستل کولے
    شي، خو لازمه نۀ ده چې و ئې باسي۔
  \end{enumerate}
\end{proof}
\label{olpseg:OLP-0093-B014:end}

\phantomsection\label{olpseg:OLP-0094-B005:start}
\olfileid{fol}{ntd}{qpr}
\label{olpseg:OLP-0094-B005:end}

\phantomsection\label{olpseg:OLP-0094-B006:start}
\olsection{د اشتقاق وړتيا او سورونه}
\label{olpseg:OLP-0094-B006:end}

\phantomsection\label{olpseg:OLP-0094-B007:start}
\begin{explain}
  د بشپړتيا قضيه دا هم غواړي چې د فطري استنتاج قاعدې د~\LR{$\Proves$}
  په اړه هغه حقيقتونه ورکړي چې په دې برخه کښې ثابتېږي۔
\end{explain}
\label{olpseg:OLP-0094-B007:end}

\phantomsection\label{olpseg:OLP-0094-B008:start}
\begin{thm}
\ollabel{thm:strong-generalization} که \LR{$c$} داسې ثابت وي چې په
\LR{$\Gamma$} يا~\LR{$!A(x)$} کښې نۀ راځي، او \LR{$\Gamma \Proves !A(c)$} وي، نو
\LR{$\Gamma \Proves \lforall[x][!A(x)]$}۔
\end{thm}
\label{olpseg:OLP-0094-B008:end}

\phantomsection\label{olpseg:OLP-0094-B009:start}
\begin{proof}
\LR{$\delta$} له~\LR{$\Gamma$} څخه د~\LR{$!A(c)$} يو اشتقاق ونيسئ۔ د
\Intro{\lforall} استنتاج په زياتولو د~\LR{$\lforall[x][!A(x)]$} يو
اشتقاق اخلو۔ ځکه \LR{$c$} په~\LR{$\Gamma$} يا~\LR{$!A(x)$} کښې نۀ راځي،
د ځانګړي متغير شرط پوره کېږي۔
\end{proof}
\label{olpseg:OLP-0094-B009:end}

\phantomsection\label{olpseg:OLP-0094-B010:start}
\begin{prop}
\ollabel{prop:provability-quantifiers}
\begin{enumerate}
\item \LR{$!A(t) \Proves \lexists[x][!A(x)]$}.
\label{olpseg:OLP-0094-B010:end}

\phantomsection\label{olpseg:OLP-0094-B011:start}
\item \LR{$\lforall[x][!A(x)] \Proves !A(t)$}.
\end{enumerate}
\end{prop}
\label{olpseg:OLP-0094-B011:end}

\phantomsection\label{olpseg:OLP-0094-B012:start}
\begin{proof}
\begin{enumerate}
\item لاندې له~\LR{$!A(t)$} څخه د~\LR{$\lexists[x][!A(x)]$} يو
اشتقاق دے:
\begin{prooftree}
\AxiomC{\LR{$!A(t)$}}
\RightLabel{\Intro{\lexists}}
\UnaryInfC{\LR{$\lexists[x][!A(x)]$}}
\end{prooftree}
\label{olpseg:OLP-0094-B012:end}

\phantomsection\label{olpseg:OLP-0094-B013:start}
\item لاندې له~\LR{$\lforall[x][!A(x)]$} څخه د~\LR{$!A(t)$} يو
اشتقاق دے:
\begin{prooftree}
\AxiomC{\LR{$\lforall[x][!A(x)]$}}
\RightLabel{\Elim{\lforall}}
\UnaryInfC{\LR{$!A(t)$}}
\end{prooftree}
\end{enumerate}
\end{proof}
\label{olpseg:OLP-0094-B013:end}

\phantomsection\label{olpseg:OLP-0095-B004:start}
\olfileid{fol}{ntd}{sou}
\label{olpseg:OLP-0095-B004:end}

\phantomsection\label{olpseg:OLP-0095-B005:start}
\olsection{سموالے}
\label{olpseg:OLP-0095-B005:end}

\phantomsection\label{olpseg:OLP-0095-B006:start}
\begin{explain}
د فطري استنتاج په شان اشتقاق نظام هله \emph{سم} دے چې
هغه څيزونه اشتقاق نۀ کړي چې په حقيقت کښې ترې نۀ لازمېږي۔
نو سموالے د اشتقاق نظامونو د تضمين شوې خونديتوب يو ډول
خاصيت دے۔ د دې له مخې چې کوم ثبوت-تيوريکي خاصيت تر بحث لاندې وي،
مثلاً غواړو پوه شو چې:
\begin{enumerate}
\item هره د اشتقاق وړ~تړلے فارمول معتبره ده؛
\item که تړلے فارمول له ځينو نورو څخه د اشتقاق وړ وي، د هغو
  معنايي پايله هم ده؛
\item که د تړلي فارمولونه يو سټ ناسازګار وي، د صدق وړ نۀ دے۔
\end{enumerate}
دا د اشتقاق نظام مهم خاصيتونه دي۔ که له دغو څخه کوم يو
سم نۀ وي، اشتقاق نظام نيمګړے دے---له حده زيات څيزونه به
اشتقاق کوي۔ له دې امله د اشتقاق نظام سموالے ثابتول
تر ټولو زيات اهميت لري۔
\end{explain}
\label{olpseg:OLP-0095-B006:end}

\phantomsection\label{olpseg:OLP-0095-B007:start}
\begin{thm}[سموالے]
\ollabel{thm:soundness}
که \LR{$!A$} له نۀ ايستل شوې فرضيو~\LR{$\Gamma$} څخه د اشتقاق وړ وي،
نو \LR{$\Gamma \Entails !A$}۔
\end{thm}
\label{olpseg:OLP-0095-B007:end}

\phantomsection\label{olpseg:OLP-0095-B008:start}
\begin{proof}
فرض کړئ \LR{$\delta$} د~\LR{$!A$} يو اشتقاق دے۔ موږ په~\LR{$\delta$}
کښې د استنتاجونو پر شمېر استقرا کوو۔
\label{olpseg:OLP-0095-B008:end}

\phantomsection\label{olpseg:OLP-0095-B009:start}
د استقرا د بنسټ دپاره هغه حالت ثابتوو چې د استنتاجونو شمېر~\LR{$0$}
وي۔ په دې حالت کښې \LR{$\delta$} يوازې له يوې تړلے فارمول~\LR{$!A$} څخه
جوړ دے، يعنې يوه فرضيه ده۔ دغه فرضيه نۀ ايستل شوې ده؛ ځکه
فرضیې يوازې د استنتاجونو په وسيله ايستل کېدے شي، او دلته
هېڅ استنتاج نشته۔ نو هر
جوړښت~\LR{$\Struct{M}$}
چې د ثبوت ټولې نۀ ايستل شوې فرضيې پوره کوي، \LR{$!A$} هم پوره کوي۔
\label{olpseg:OLP-0095-B009:end}

\phantomsection\label{olpseg:OLP-0095-B010:start}
اوس د استقرا پړاو اخلو۔ فرض کړئ \LR{$\delta$}، \LR{$n$} استنتاجونه لري۔
د تر ټولو لاندې استنتاج مقدمه يا مقدمې د فرعي-اشتقاقونه په
وسيله اشتقاق شوې دي، او په هر يوه کښې تر~\LR{$n$} کم استنتاجونه
دي۔ د استقرا فرضيه منو: د تر ټولو لاندې استنتاج مقدمې د هغو
فرعي-اشتقاقونه له نۀ ايستل شوې فرضيو څخه لازمېږي چې په
دغو مقدمو پاے ته رسېږي۔ بايد وښيو چې پايله~\LR{$!A$} د ټول ثبوت له
نۀ ايستل شوې فرضيو څخه لازمېږي۔
\label{olpseg:OLP-0095-B010:end}

\phantomsection\label{olpseg:OLP-0095-B011:start}
د تر ټولو لاندې استنتاج د ډول له مخې حالتونه بېلوو۔ لومړے هغه
ممکن استنتاجونه څېړو چې يوازې يوه مقدمه لري۔
\label{olpseg:OLP-0095-B011:end}

\phantomsection\label{olpseg:OLP-0095-B012:start}
\begin{enumerate}
\item فرض کړئ وروستنے استنتاج \Intro{\lnot} دے: اشتقاق
  دا بڼه لري:
  \begin{prooftree}
    \AxiomC{\LR{$\Gamma, \Discharge{!A}{n}$}}
    \RightLabel{\LR{$\delta_1$}}
    \DeduceC{\LR{$\lfalse$}}
    \DischargeRule{\Intro{\lnot}}{n}
    \UnaryInfC{\LR{$\lnot !A$}}
  \end{prooftree}
  د استقرا د فرضيې له مخې \LR{$\lfalse$} د~\LR{$\delta_1$} له
  نۀ ايستل شوې فرضيو~\LR{$\Gamma \cup \{!A\}$} څخه لازمېږي۔ يو
  جوړښت~\LR{$\Struct{M}$}
  ونيسئ۔ بايد وښيو چې که
  \LR{$\Sat{M}{\Gamma}$} وي، نو
  \LR{$\Sat{M}{\lnot !A}$} دے۔ د خلف په توګه
  فرض کړئ \LR{$\Sat{M}{\Gamma}$}، خو
  \LR{$\Sat/{M}{\lnot !A}$}، يعنې
  \LR{$\Sat{M}{!A}$}۔ له دې به
  \LR{$\Sat{M}{\Gamma \cup \{!A\}}$} لازم شي۔
  دا زموږ د استقرا د فرضيې خلاف دے۔ نو
  \LR{$\Sat{M}{\lnot !A}$}۔
\label{olpseg:OLP-0095-B012:end}

\phantomsection\label{olpseg:OLP-0095-B013:start}
\item وروستنے استنتاج \Elim{\land} دے: دوه بڼې لري؛ له مقدمې
  \LR{$!A \land !B$} څخه يا~\LR{$!A$} يا~\LR{$!B$} راايستل کېدے شي۔ لومړے حالت
  ونيسئ۔ اشتقاق~\LR{$\delta$} داسې ښکاري:
  \begin{prooftree}
    \AxiomC{\LR{$\Gamma$}}
    \RightLabel{\LR{$\delta_1$}}
    \DeduceC{\LR{$!A \land !B$}}
    \RightLabel{\Elim{\land}}
    \UnaryInfC{\LR{$!A$}}
  \end{prooftree}
  د استقرا د فرضيې له مخې \LR{$!A \land !B$} د~\LR{$\delta_1$} له
  نۀ ايستل شوې فرضيو~\LR{$\Gamma$} څخه لازمېږي۔
  يو جوړښت~\LR{$\Struct{M}$} ونيسئ۔
  بايد وښيو چې که
  \LR{$\Sat{M}{\Gamma}$} وي، نو
  \LR{$\Sat{M}{!A}$} دے۔ فرض کړئ
  \LR{$\Sat{M}{\Gamma}$}۔ زموږ د استقرا د
  فرضيې (\LR{$\Gamma \Entails !A \land !B$}) له مخې پوهېږو چې
  \LR{$\Sat{M}{!A \land !B}$}۔ د تعريف له مخې
  \LR{$\Sat{M}{!A \land !B}$} هله او يوازې هله
  چې \LR{$\Sat{M}{!A}$} او
  \LR{$\Sat{M}{!B}$} دواړه وي۔ (هغه حالت چې له
  \LR{$!A \land !B$} څخه~\LR{$!B$} راايستل کېږي، په ورته ډول څېړل کېږي۔)
\label{olpseg:OLP-0095-B013:end}

\phantomsection\label{olpseg:OLP-0095-B014:start}
\item وروستنے استنتاج \Intro{\lor} دے: دوه بڼې لري؛
  \LR{$!A \lor !B$} يا له مقدمې~\LR{$!A$} يا له مقدمې~\LR{$!B$} څخه راايستل
  کېدے شي۔ لومړے حالت ونيسئ۔ اشتقاق دا بڼه لري:
  \begin{prooftree}
    \AxiomC{\LR{$\Gamma$}}
    \RightLabel{\LR{$\delta_1$}}
    \DeduceC{\LR{$!A$}}
    \RightLabel{\Intro{\lor}}
    \UnaryInfC{\LR{$!A \lor !B$}}
  \end{prooftree}
  د استقرا د فرضيې له مخې \LR{$!A$} د~\LR{$\delta_1$} له
  نۀ ايستل شوې فرضيو~\LR{$\Gamma$} څخه لازمېږي۔ يو
  جوړښت~\LR{$\Struct{M}$}
  ونيسئ۔ بايد وښيو چې که
  \LR{$\Sat{M}{\Gamma}$} وي، نو
  \LR{$\Sat{M}{!A \lor !B}$} دے۔ فرض کړئ
  \LR{$\Sat{M}{\Gamma}$}؛ بيا
  \LR{$\Sat{M}{!A}$}، ځکه
  \LR{$\Gamma \Entails !A$} (د استقرا فرضيه)۔ نو
  \LR{$\Sat{M}{!A \lor !B}$} هم بايد وي۔
  (هغه حالت چې \LR{$!A \lor !B$} له~\LR{$!B$} څخه راايستل کېږي، په ورته ډول
  څېړل کېږي۔)
\label{olpseg:OLP-0095-B014:end}

\phantomsection\label{olpseg:OLP-0095-B015:start}
\item وروستنے استنتاج \Intro{\lif} دے: \LR{$!A \lif !B$} له داسې فرعي
  ثبوت څخه راايستل کېږي چې فرضيه ئې~\LR{$!A$} او پايله ئې~\LR{$!B$} ده،
  يعنې:
  \begin{prooftree}
    \AxiomC{\LR{$\Gamma, \Discharge{!A}{n}$}}
    \RightLabel{\LR{$\delta_1$}}
    \DeduceC{\LR{$!B$}}
    \DischargeRule{\Intro{\lif}}{n}
    \UnaryInfC{\LR{$!A \lif !B$}}
  \end{prooftree}
  د استقرا د فرضيې له مخې \LR{$!B$} د~\LR{$\delta_1$} له
  نۀ ايستل شوې فرضيو څخه لازمېږي، يعنې
  \LR{$\Gamma \cup \{!A\} \Entails !B$}۔ يو
  جوړښت~\LR{$\Struct{M}$}
  ونيسئ۔ د~\LR{$\delta$}~نۀ ايستل شوې فرضيې يوازې~\LR{$\Gamma$} دي؛ ځکه
  \LR{$!A$} په وروستي استنتاج کښې ايستل شوې ده۔ نو بايد وښيو
  چې \LR{$\Gamma \Entails !A \lif !B$}۔ د خلف په توګه فرض کړئ چې د کوم
  جوړښت~\LR{$\Struct{M}$}
  دپاره \LR{$\Sat{M}{\Gamma}$}، خو
  \LR{$\Sat/{M}{!A \lif !B}$} وي۔ نو
  \LR{$\Sat{M}{!A}$} او
  \LR{$\Sat/{M}{!B}$}۔ خو د فرضيې له مخې
  \LR{$!B$} د~\LR{$\Gamma \cup \{!A\}$} معنايي پايله ده، يعنې
  \LR{$\Sat{M}{!B}$}، او دا تناقض دے۔ نو
  \LR{$\Gamma \Entails !A \lif !B$}۔
\label{olpseg:OLP-0095-B015:end}

\phantomsection\label{olpseg:OLP-0095-B016:start}
\item وروستنے استنتاج \FalseInt دے: دلته \LR{$\delta$} داسې پاے ته رسېږي:
  \begin{prooftree}
    \AxiomC{\LR{$\Gamma$}}
    \RightLabel{\LR{$\delta_1$}}
    \DeduceC{\LR{$\lfalse$}}
    \RightLabel{\FalseInt}
    \UnaryInfC{\LR{$!A$}}
  \end{prooftree}
  د استقرا د فرضيې له مخې \LR{$\Gamma \Entails \lfalse$}۔ بايد وښيو چې
  \LR{$\Gamma \Entails !A$}۔ فرض کړئ داسې نۀ وي؛ بيا د کوم
  \LR{$\Struct{M}$} دپاره
  \LR{$\Sat{M}{\Gamma}$} او
  \LR{$\Sat/{M}{!A}$} دي۔ خو
  کاذبي هېڅ جوړښت يا ارزونه نۀ پوره کوي؛ يعنې تل
  \LR{$\Sat/{M}{\lfalse}$} وي؛ نو له دې به
  \LR{$\Gamma \Entails/ \lfalse$} لازم شي، چې د استقرا د فرضيې خلاف دے۔
\label{olpseg:OLP-0095-B016:end}

\phantomsection\label{olpseg:OLP-0095-B017:start}
\item وروستنے استنتاج \FalseCl دے: تمرين۔
\label{olpseg:OLP-0095-B017:end}

\phantomsection\label{olpseg:OLP-0095-B018:start}

\item وروستنے استنتاج \Intro{\lforall} دے: بيا \LR{$\delta$} دا بڼه لري:
  \begin{prooftree}
    \AxiomC{\LR{$\Gamma$}}
    \RightLabel{\LR{$\delta_1$}}
    \DeduceC{\LR{$!A(a)$}}
    \RightLabel{\Intro{\lforall}}
    \UnaryInfC{\LR{$\lforall[x][!A(x)]$}}
  \end{prooftree}
  د استقرا د فرضيې له مخې مقدمه~\LR{$!A(a)$} د~نۀ ايستل شوې
  فرضيو~\LR{$\Gamma$} معنايي پايله ده۔ داسې جوړښت~\LR{$\Struct{M}$} ونيسئ
  چې \LR{$\Sat{M}{\Gamma}$} وي۔ بايد وښيو چې
  \LR{$\Sat{M}{\lforall[x][!A(x)]}$}۔ ځکه
  \LR{$\lforall[x][!A(x)]$} يوه تړلے فارمول ده، د دې معنٰی دا ده چې
  د هرې متغير-ګومارنې~\LR{$s$} دپاره بايد
  \LR{$\Sat{M}{!A(x)}[s]$} وښيو
  (\olref[syn][ass]{prop:sat-quant})۔ ځکه \LR{$\Gamma$} ټول له جملو جوړ
  دے، د هر~\LR{$!B \in \Gamma$} دپاره
  \LR{$\Sat{M}{!B}[s]$} د~\olref[syn][sat]{defn:satisfaction} له مخې
  دے۔ \LR{$\Struct{M'}$} داسې جوړښت ونيسئ چې د~\LR{$\Struct{M}$} په شان وي،
  خو \LR{$\Assign{a}{M'} = s(x)$} وي۔ ځکه \LR{$a$} په~\LR{$\Gamma$} کښې نۀ راځي،
  د~\olref[syn][ext]{cor:extensionality-sent} له مخې
  \LR{$\Sat{M'}{\Gamma}$}۔ ځکه \LR{$\Gamma \Entails !A(a)$}،
  \LR{$\Sat{M'}{!A(a)}$}۔ ځکه \LR{$!A(a)$} يوه تړلے فارمول ده، د
  \olref[syn][ass]{prop:sentence-sat-true} له مخې
  \LR{$\Sat{M'}{!A(a)}[s]$}۔ د
  \olref[syn][ext]{prop:ext-formulas} له مخې
  \LR{$\Sat{M'}{!A(x)}[s]$} هله او يوازې هله چې
  \LR{$\Sat{M'}{!A(a)}$} (ياد ولرئ \LR{$!A(a)$} هماغه
  \LR{$\Subst{!A(x)}{a}{x}$} دے)۔ نو \LR{$\Sat{M'}{!A(x)}[s]$}۔ ځکه \LR{$a$}
  په~\LR{$!A(x)$} کښې نۀ راځي، د
  \olref[syn][ext]{prop:extensionality} له مخې
  \LR{$\Sat{M}{!A(x)}[s]$}۔ خو \LR{$s$} خپل‌سری متغير-ګومارنه وه؛ نو
  \LR{$\Sat{M}{\lforall[x][!A(x)]}$}۔
\label{olpseg:OLP-0095-B018:end}

\phantomsection\label{olpseg:OLP-0095-B019:start}
\item وروستنے استنتاج \Intro{\lexists} دے: تمرين۔
\label{olpseg:OLP-0095-B019:end}

\phantomsection\label{olpseg:OLP-0095-B020:start}
\item وروستنے استنتاج \Elim{\forall} دے: تمرين۔

\end{enumerate}
\label{olpseg:OLP-0095-B020:end}

\phantomsection\label{olpseg:OLP-0095-B021:start}
اوس هغه ممکن استنتاجونه څېړو چې څو مقدمې لري:
\Elim{\lor}، \Intro{\land}، \Elim{\lnot}،
\Elim{\lif} او~\Elim{\lexists}۔
\textbf{د ژباړې د نسخې سمون (OLND-006):} سرچينې د څو-مقدمې
قاعدو له دې بشپړ لړ څخه د نفي ايستلو قاعده غورځولې ده، سره له دې
چې لاندې تمرين ئې هم راوړي؛ دلته بېرته شامله شوه۔
\begin{enumerate}
\label{olpseg:OLP-0095-B021:end}

\phantomsection\label{olpseg:OLP-0095-B022:start}
\item وروستنے استنتاج \Intro{\land} دے۔ \LR{$!A \land !B$} له
  مقدمو~\LR{$!A$} او~\LR{$!B$} څخه راايستل کېږي او \LR{$\delta$} دا بڼه لري:
  \begin{prooftree}
    \AxiomC{\LR{$\Gamma_1$}}
    \RightLabel{\LR{$\delta_1$}}
    \DeduceC{\LR{$!A$}}
    \AxiomC{\LR{$\Gamma_2$}}
    \RightLabel{\LR{$\delta_2$}}
    \DeduceC{\LR{$!B$}}
    \RightLabel{\Intro{\land}}
    \BinaryInfC{\LR{$!A \land !B$}}
  \end{prooftree}
  د استقرا د فرضيې له مخې \LR{$!A$} د~\LR{$\delta_1$} له
  نۀ ايستل شوې فرضيو~\LR{$\Gamma_1$} څخه، او~\LR{$!B$} د~\LR{$\delta_2$} له
  نۀ ايستل شوې فرضيو~\LR{$\Gamma_2$} څخه لازمېږي۔ د~\LR{$\delta$}
  نۀ ايستل شوې فرضيې \LR{$\Gamma_1 \cup \Gamma_2$} دي؛ نو بايد
  \LR{$\Gamma_1 \cup \Gamma_2 \Entails !A \land !B$} وښيو۔ داسې
  جوړښت~\LR{$\Struct{M}$}
  ونيسئ چې
  \LR{$\Sat{M}{\Gamma_1 \cup \Gamma_2}$} وي۔
  ځکه \LR{$\Sat{M}{\Gamma_1}$}، او
  \LR{$\Gamma_1 \Entails !A$}، نو
  \LR{$\Sat{M}{!A}$}۔ همدارنګه، ځکه
  \LR{$\Sat{M}{\Gamma_2}$} او
  \LR{$\Gamma_2 \Entails !B$}، نو
  \LR{$\Sat{M}{!B}$}۔ دواړه يو ځاے
  \LR{$\Sat{M}{!A \land !B}$} راکوي۔
\label{olpseg:OLP-0095-B022:end}

\phantomsection\label{olpseg:OLP-0095-B023:start}
\item وروستنے استنتاج \Elim{\lor} دے: تمرين۔
\label{olpseg:OLP-0095-B023:end}

\phantomsection\label{olpseg:OLP-0095-B024:start}
\item وروستنے استنتاج \Elim{\lif} دے۔ \LR{$!B$} له مقدمو
  \LR{$!A \lif !B$} او~\LR{$!A$} څخه راايستل کېږي۔ اشتقاق~\LR{$\delta$}
  داسې ښکاري:
  \begin{prooftree}
    \AxiomC{\LR{$\Gamma_1$}}
    \RightLabel{\LR{$\delta_1$}}
    \DeduceC{\LR{$!A \lif !B$}}
    \AxiomC{\LR{$\Gamma_2$}}
    \RightLabel{\LR{$\delta_2$}}
    \DeduceC{\LR{$!A$}}
    \RightLabel{\Elim{\lif}}
    \BinaryInfC{\LR{$!B$}}
  \end{prooftree}
  د استقرا د فرضيې له مخې \LR{$!A \lif !B$} د~\LR{$\delta_1$} له
  نۀ ايستل شوې فرضيو~\LR{$\Gamma_1$} څخه، او~\LR{$!A$} د~\LR{$\delta_2$} له
  نۀ ايستل شوې فرضيو~\LR{$\Gamma_2$} څخه لازمېږي۔ يو
  جوړښت~\LR{$\Struct{M}$}
  ونيسئ۔ بايد وښيو چې که
  \LR{$\Sat{M}{\Gamma_1 \cup \Gamma_2}$} وي، نو
  \LR{$\Sat{M}{!B}$} دے۔ فرض کړئ
  \LR{$\Sat{M}{\Gamma_1 \cup \Gamma_2}$}۔
  ځکه \LR{$\Gamma_1 \Entails !A \lif !B$}، نو
  \LR{$\Sat{M}{!A \lif !B}$}۔ ځکه
  \LR{$\Gamma_2 \Entails !A$}، نو
  \LR{$\Sat{M}{!A}$}۔ له دې څخه
  \LR{$\Sat{M}{!B}$} لازمېږي (ځکه که
  \LR{$\Sat/{M}{!B}$} وے، نو له
  \LR{$\Sat{M}{!A}$} سره به
  \LR{$\Sat/{M}{!A \lif !B}$} وے، چې د
  \LR{$\Sat{M}{!A \lif !B}$} خلاف دے)۔
\label{olpseg:OLP-0095-B024:end}

\phantomsection\label{olpseg:OLP-0095-B025:start}
\item وروستنے استنتاج \Elim{\lnot} دے: تمرين۔
\label{olpseg:OLP-0095-B025:end}

\phantomsection\label{olpseg:OLP-0095-B026:start}
\item وروستنے استنتاج \Elim{\lexists} دے: تمرين۔
\end{enumerate}
\end{proof}
\label{olpseg:OLP-0095-B026:end}

\phantomsection\label{olpseg:OLP-0095-B027:start}
\begin{prob}
د~\olref[fol][ntd][sou]{thm:soundness} ثبوت بشپړ کړئ۔
\end{prob}
\label{olpseg:OLP-0095-B027:end}



\phantomsection\label{olpseg:OLP-0095-B029:start}
\begin{cor}
\ollabel{cor:weak-soundness}
که \LR{$\Proves !A$} وي، نو \LR{$!A$} معتبره ده۔
\end{cor}
\label{olpseg:OLP-0095-B029:end}

\phantomsection\label{olpseg:OLP-0095-B030:start}
\begin{cor}
\ollabel{cor:consistency-soundness}
که \LR{$\Gamma$} د صدق وړ وي، نو سازګار دے۔
\end{cor}
\label{olpseg:OLP-0095-B030:end}

\phantomsection\label{olpseg:OLP-0095-B031:start}
\begin{proof}
نقيض عکس ثابتوو۔ فرض کړئ \LR{$\Gamma$} سازګار نۀ دے۔ بيا
\LR{$\Gamma \Proves \lfalse$}، يعنې له~\LR{$\Gamma$} څخه د
نۀ ايستل شوې فرضيو پر بنسټ د~\LR{$\lfalse$} يو اشتقاق شته۔
د~\olref{thm:soundness} له مخې هر
جوړښت~\LR{$\Struct{M}$}
چې \LR{$\Gamma$} پوره کوي، بايد \LR{$\lfalse$} هم پوره کړي۔ ځکه د هر
جوړښت~\LR{$\Struct{M}$}
دپاره \LR{$\Sat/{M}{\lfalse}$}، نو هېڅ
\LR{$\Struct{M}$}، \LR{$\Gamma$} نۀ شي پوره
کولے؛ يعنې \LR{$\Gamma$} د صدق وړ نۀ دے۔
\end{proof}
\label{olpseg:OLP-0095-B031:end}

\phantomsection\label{olpseg:OLP-0096-B004:start}
\olfileid{fol}{ntd}{ide}
\label{olpseg:OLP-0096-B004:end}

\phantomsection\label{olpseg:OLP-0096-B005:start}
\olsection{له عينيت سره اشتقاقونه}
\label{olpseg:OLP-0096-B005:end}

\phantomsection\label{olpseg:OLP-0096-B006:start}
له عينيت سره اشتقاقونه اضافي استنتاجي قاعدو ته اړتيا لري۔
\label{olpseg:OLP-0096-B006:end}

\phantomsection\label{olpseg:OLP-0096-B007:start}
\begin{defish}
\AxiomC{}
\RightLabel{\Intro{\eq}}
\UnaryInfC{\LR{$\eq[t][t]$}}
\DisplayProof
\hfill
\begin{tabular}{r}
\AxiomC{\LR{$\eq[t_1][t_2]$}}
\AxiomC{\LR{$!A(t_1)$}}
\RightLabel{\Elim{\eq}}
\BinaryInfC{\LR{$!A(t_2)$}}
\DisplayProof
\\[3ex]
\AxiomC{\LR{$\eq[t_1][t_2]$}}
\AxiomC{\LR{$!A(t_2)$}}
\RightLabel{\Elim{\eq}}
\BinaryInfC{\LR{$!A(t_1)$}}
\DisplayProof
\end{tabular}
\end{defish}
\label{olpseg:OLP-0096-B007:end}

\phantomsection\label{olpseg:OLP-0096-B008:start}
په پورته قاعدو کښې \LR{$t$}، \LR{$t_1$} او~\LR{$t_2$} تړلې اصطلاحګانې دي۔
د~\Intro{\eq} قاعده راکوي چې د~\LR{$\eq[t][t]$} په بڼه هره د يووالي
جمله نېغ په نېغه، له هېڅ فرضيې پرته، اشتقاق کړو۔
\label{olpseg:OLP-0096-B008:end}

\phantomsection\label{olpseg:OLP-0096-B009:start}
\begin{ex}
که \LR{$s$} او~\LR{$t$} تړلې اصطلاحګانې وي، نو
\LR{$!A(s), \eq[s][t] \Proves !A(t)$}:
\begin{prooftree}
\AxiomC{\LR{$\eq[s][t]$}}
\AxiomC{\LR{$!A(s)$}}
\RightLabel{\LR{$\Elim{\eq}$}}
\BinaryInfC{\LR{$!A(t)$}}
\end{prooftree}
دا به د «يو شان څيزونو د تعويض د اصل» يا د لايبنېڅ د قانون په نامه
اشنا وي۔
\end{ex}
\label{olpseg:OLP-0096-B009:end}

\phantomsection\label{olpseg:OLP-0096-B010:start}
\begin{prob}
ثابت کړئ چې \LR{$=$} متناظره او انتقالي دواړه ده؛ يعنې د
\LR{$\lforall[x][\lforall[y][(\eq[x][y] \lif
    \eq[y][x])]]$} او~\LR{$\lforall[x][\lforall[y][\lforall[z]((\eq[x][y]
    \land \eq[y][z]) \lif \eq[x][z])]]$}
اشتقاقونه ورکړئ۔
\end{prob}
\label{olpseg:OLP-0096-B010:end}

\phantomsection\label{olpseg:OLP-0096-B011:start}
\begin{ex}
موږ دا تړلے فارمول~اشتقاق کوو:
\begin{align*}
& \lforall[x][\lforall[y][((!A(x) \land !A(y)) \lif \eq[x][y])]]
\intertext{\RL{له دې تړلے فارمول څخه:}}
& \lexists[x][\lforall[y][(!A(y) \lif \eq[y][x])]]
\end{align*}
اشتقاق له لاندې څخه پورته جوړوو:
\begin{prooftree}
\AxiomC{\LR{$\lexists[x][\lforall[y][(!A(y) \lif \eq[y][x])]]
\quad \Discharge{!A(a) \land !A(b)}{1}$}}
\DeduceC{\LR{$\eq[a][b]$}}
\DischargeRule{\Intro{\lif}}{1}
\UnaryInfC{\LR{$((!A(a) \land !A(b)) \lif \eq[a][b])$}}
\RightLabel{\Intro{\lforall}}
\UnaryInfC{\LR{$\lforall[y][((!A(a) \land !A(y)) \lif \eq[a][y])]$}}
\RightLabel{\Intro{\lforall}}
\UnaryInfC{\LR{$\lforall[x][\lforall[y][((!A(x) \land !A(y)) \lif \eq[x][y])]]$}}
\end{prooftree}
اوس بايد اصلي فرضيه وکاروو: ځکه دا وجودي فارمول ده، د
منځنۍ پايلې~\LR{$\eq[a][b]$} د~اشتقاق کولو دپاره
\Elim{\lexists} کاروو۔
\begin{prooftree}
\AxiomC{\LR{$\lexists[x][\lforall[y][(!A(y) \lif \eq[y][x])]]$}}
\AxiomC{\LR{$\Discharge{\lforall[y][(!A(y) \lif \eq[y][c]])}{2}$}}
\noLine
\UnaryInfC{\LR{$\Discharge{!A(a) \land !A(b)}{1}$}}
\DeduceC{\LR{$\eq[a][b]$}}
\DischargeRule{\Elim{\lexists}}{2}
\BinaryInfC{\LR{$\eq[a][b]$}}
\DischargeRule{\Intro{\lif}}{1}
\UnaryInfC{\LR{$((!A(a) \land !A(b)) \lif \eq[a][b])$}}
\RightLabel{\Intro{\lforall}}
\UnaryInfC{\LR{$\lforall[y][((!A(a) \land !A(y)) \lif \eq[a][y])]$}}
\RightLabel{\Intro{\lforall}}
\UnaryInfC{\LR{$\lforall[x][\lforall[y][((!A(x) \land !A(y)) \lif \eq[x][y])]]$}}
\end{prooftree}
د پاس په ښي لوري کښې فرعي-اشتقاق د خپلو فرضيو په کارولو
بشپړېږي، څو \LR{$\eq[a][c]$} او~\LR{$\eq[b][c]$} وښيي۔ دا دوه جلا
اشتقاقونه غواړي۔ د~\LR{$\eq[a][c]$}~اشتقاق داسې دے:
\begin{prooftree}
\AxiomC{\LR{$\Discharge{\lforall[y][(!A(y) \lif \eq[y][c]])}{2}$}}
\RightLabel{\Elim{\lforall}}
\UnaryInfC{\LR{$!A(a) \lif \eq[a][c]$}}
\AxiomC{\LR{$\Discharge{!A(a) \land !A(b)}{1}$}}
\RightLabel{\Elim{\land}}
\UnaryInfC{\LR{$!A(a)$}}
\RightLabel{\Elim{\lif}}
\BinaryInfC{\LR{$\eq[a][c]$}}
\end{prooftree}
له~\LR{$\eq[a][c]$} او~\LR{$\eq[b][c]$} څخه
\LR{$\eq[a][b]$} د~\Elim{\eq} په وسيله اشتقاق کوو۔
\end{ex}
\label{olpseg:OLP-0096-B011:end}

\phantomsection\label{olpseg:OLP-0096-B012:start}
\begin{prob}
د لاندې فارمولونه~اشتقاقونه ورکړئ:
\begin{enumerate}
\item \LR{$\lforall[x][\lforall[y][((\eq[x][y] \land !A(x)) \lif !A(y))]]$}
\item \LR{$\lexists[x][!A(x)] \land \lforall[y][\lforall[z][((!A(y) \land
    !A(z)) \lif \eq[y][z])]] \lif \lexists[x][(!A(x) \land
  \lforall[y][(!A(y) \lif \eq[y][x])])]$}
\end{enumerate}
\end{prob}
\label{olpseg:OLP-0096-B012:end}

\phantomsection\label{olpseg:OLP-0097-B004:start}
\olfileid{fol}{ntd}{sid}
\label{olpseg:OLP-0097-B004:end}

\phantomsection\label{olpseg:OLP-0097-B005:start}
\olsection{له عينيت سره سموالے}
\label{olpseg:OLP-0097-B005:end}

\phantomsection\label{olpseg:OLP-0097-B006:start}
\begin{prop}
د~\LR{$\eq$} له قاعدو سره فطري استنتاج سم دے۔
\end{prop}
\label{olpseg:OLP-0097-B006:end}

\phantomsection\label{olpseg:OLP-0097-B007:start}
\begin{proof}
د~\LR{$\eq[t][t]$} په بڼه هره فارمول معتبره ده؛ ځکه د هر
جوړښت~\LR{$\Struct M$} دپاره \LR{$\Sat{M}{\eq[t][t]}$}۔ (پام وکړئ
چې اصطلاح~\LR{$t$} تړلې فرضوو، يعنې هېڅ متغير نۀ لري؛ نو
متغير-ګومارنې اغېز نۀ لري۔)
\label{olpseg:OLP-0097-B007:end}

\phantomsection\label{olpseg:OLP-0097-B008:start}
فرض کړئ په کوم اشتقاق کښې وروستنے استنتاج \Elim{\eq} دے،
يعنې اشتقاق لاندې بڼه لري:
\begin{prooftree}
  \AxiomC{\LR{$\Gamma_1$}}
  \RightLabel{\LR{$\delta_1$}}
  \DeduceC{\LR{$\eq[t_1][t_2]$}}
  \AxiomC{\LR{$\Gamma_2$}}
  \RightLabel{\LR{$\delta_2$}}
  \DeduceC{\LR{$!A(t_1)$}}
  \RightLabel{\Elim{\eq}}
  \BinaryInfC{\LR{$!A(t_2)$}}
\end{prooftree}
مقدمې~\LR{$\eq[t_1][t_2]$} او~\LR{$!A(t_1)$} په خپل وار له
نۀ ايستل شوې فرضيو~\LR{$\Gamma_1$} او~\LR{$\Gamma_2$} څخه
اشتقاق شوې دي۔ غواړو وښيو چې \LR{$!A(t_2)$} له
\LR{$\Gamma_1 \cup \Gamma_2$} څخه لازمېږي۔ داسې
جوړښت~\LR{$\Struct{M}$} ونيسئ چې
\LR{$\Sat{M}{\Gamma_1 \cup \Gamma_2}$} وي۔ د استقرا د فرضيې له مخې
\LR{$\Sat{M}{!A(t_1)}$} او~\LR{$\Sat{M}{\eq[t_1][t_2]}$}۔ نو
\LR{$\Value{t_1}{M} = \Value{t_2}{M}$}۔ \LR{$s$} هره متغير-ګومارنه ونيسئ،
او \LR{$m = \Value{t_1}{M} = \Value{t_2}{M}$}۔ د
\olref[fol][syn][ext]{prop:ext-formulas} له مخې
\LR{$\Sat{M}{!A(t_1)}[s]$} هله او يوازې هله چې
\LR{$\Sat{M}{!A(x)}[\Subst{s}{m}{x}]$}، او دا هله او يوازې هله چې
\LR{$\Sat{M}{!A(t_2)}[s]$}۔ ځکه \LR{$\Sat{M}{!A(t_1)}$}، نو
\LR{$\Sat{M}{!A(t_2)}$}۔
\end{proof}
\label{olpseg:OLP-0097-B008:end}

\phantomsection\label{olpseg:OLP-0098-B004:start}
\olchapter{fol}{tab}{تابلوګانې}
\label{olpseg:OLP-0098-B004:end}

\phantomsection\label{olpseg:OLP-0098-B005:start}
\begin{editorial}
دا څپرکے يو نښه لرونکے تحليلي تابلو نظام وړاندې کوي۔
\label{olpseg:OLP-0098-B005:end}

\phantomsection\label{olpseg:OLP-0098-B006:start}
د ثبوت د يوه نظام په توګه له تابلو سره اړوند مواد د شاملولو يا
ايستلو دپاره د~\texttt{prfTab} ټګ وکاروئ۔
\textbf{د ژباړې د نسخې سمون (OLTAB-001):} سرچينه دلته په تېروتنه
فطري استنتاج يادوي؛ د څپرکي، ټګ او ټول وارد شوي موادو له مخې مطلب
تابلو نظام دے، نو نوم ئې دلته سم شوے دے۔
\end{editorial}
\label{olpseg:OLP-0098-B006:end}

\phantomsection\label{olpseg:OLP-0098-B009:start}
%


\label{olpseg:OLP-0098-B009:end}

\phantomsection\label{olpseg:OLP-0098-B012:start}
%


\label{olpseg:OLP-0098-B012:end}

\phantomsection\label{olpseg:OLP-0098-B016:start}
%


\label{olpseg:OLP-0098-B016:end}

\phantomsection\label{olpseg:OLP-0098-B018:start}
%



\label{olpseg:OLP-0098-B018:end}

\phantomsection\label{olpseg:OLP-0099-B004:start}
\olfileid{fol}{tab}{rul}
\label{olpseg:OLP-0099-B004:end}

\phantomsection\label{olpseg:OLP-0099-B005:start}
\olsection{قاعدې او تابلوګانې}
\label{olpseg:OLP-0099-B005:end}

\phantomsection\label{olpseg:OLP-0099-B006:start}
تابلو په جوړښت کښې د~تړلے فارمول د رښتيا يا
دروغ کېدو د ممکنو لارو منظمه پلټنه ده۔ د تابلو بنسټيز توکي
نښه لرونکي فارمولونه دي: تړلي فارمولونه له يوه رښتياوالي ارزښت
«نښې» سره، چې يا~\LR{$\True$} وي يا~\LR{$\False$}۔ دا نښه لرونکي فارمولونه
په داسې ونه کښې ترتيبېږي چې لاندې خوا ته غځېږي۔
\label{olpseg:OLP-0099-B006:end}

\phantomsection\label{olpseg:OLP-0099-B007:start}
\begin{defn}
  \emph{نښه لرونکے فارمول} له يوه رښتياوالي ارزښت او~تړلے فارمول
  څخه جوړه ده، يعنې له دې دواړو څخه يوه:
  \[
  \sFmla{\True}{!A} \text{\RL{ يا }} \sFmla{\False}{!A}.
  \]
\end{defn}
\label{olpseg:OLP-0099-B007:end}

\phantomsection\label{olpseg:OLP-0099-B008:start}
په شهودي ډول \LR{$\sFmla{\True}{!A}$} داسې لوستلے شو چې «کېدے شي
\LR{$!A$} رښتيا وي»، او~\LR{$\sFmla{\False}{!A}$} داسې چې «کېدے شي
\LR{$!A$} دروغ وي» (په کوم جوړښت کښې)۔
\label{olpseg:OLP-0099-B008:end}

\phantomsection\label{olpseg:OLP-0099-B009:start}
په ونه کښې هر نښه لرونکے فارمول يا \emph{فرض} دے (چې د ونې په
تر ټولو پاسنۍ برخه کښې ليکل کېږي)، يا تر ځان له پاسه له کوم
نښه لرونکے فارمول څخه د يوې استنتاجي قاعدې په وسيله اخستل شوے
دے۔ د مخکښې فارمول د هر ممکن اصلي عملګر دپاره دوه
قاعدې شته: يوه د هغه حالت دپاره چې نښه~\LR{$\True$} وي، او بله د هغه
حالت دپاره چې نښه~\LR{$\False$} وي۔ ځينې قاعدې ونې ته څانګې ورکوي او
ځينې يوازې څانګې ته نښه لرونکي فارمولونه ورزياتوي۔ يوه قاعده
کېدے شي (او ډېر ځله بايد) سمدستي پر مخکښ نښه لرونکے فارمول نۀ،
بلکې له ريښې څخه د قاعدې د پلي کېدو تر ځاے پورې په څانګه کښې پر هر
نښه لرونکے فارمول پلې شي۔
\label{olpseg:OLP-0099-B009:end}

\phantomsection\label{olpseg:OLP-0099-B010:start}
يوه څانګه هله \emph{تړلې} ده چې
\LR{$\sFmla{\True}{!A}$} او~\LR{$\sFmla{\False}{!A}$} دواړه ولري۔
تړلے تابلو هغه دے چې هره څانګه ئې تړلې وي۔ د شهودي تعبير له
مخې هره څانګه يو ګډ امکان بيانوي، خو
\LR{$\sFmla{\True}{!A}$} او~\LR{$\sFmla{\False}{!A}$} يو ځاے ممکن نۀ دي۔
په بله وينا، که يوه څانګه تړلې وي، هغه امکان رد شوے دے چې څانګه
ئې بيانوي۔ په ځانګړي ډول، د دې معنٰی دا ده چې تړلے تابلو هغه
ټول امکانات ردوي چې د~\LR{$\sFmla{\True}{!A}$} په بڼه هر فرض په يو وخت
رښتيا او د~\LR{$\sFmla{\False}{!A}$} په بڼه هر فرض دروغ کړي۔
\label{olpseg:OLP-0099-B010:end}

\phantomsection\label{olpseg:OLP-0099-B011:start}
\emph{د~\LR{$!A$} دپاره} تړلے تابلو داسې تړلے تابلو دے چې
ريښه ئې~\LR{$\sFmla{\False}{!A}$} وي۔ که داسې تړلے تابلو موجود
وي، د~\LR{$!A$} د دروغ کېدو ټول امکانات رد شوي؛ يعنې \LR{$!A$} بايد په هر
جوړښت کښې رښتيا وي۔
\label{olpseg:OLP-0099-B011:end}

\phantomsection\label{olpseg:OLP-0100-B004:start}
\olfileid{fol}{tab}{prl}
\label{olpseg:OLP-0100-B004:end}

\phantomsection\label{olpseg:OLP-0100-B005:start}
\olsection{قضيوي قاعدې}
\label{olpseg:OLP-0100-B005:end}

\phantomsection\label{olpseg:OLP-0100-B006:start}
\subsection{د~\LR{$\lnot$} قاعدې}
\label{olpseg:OLP-0100-B006:end}

\phantomsection\label{olpseg:OLP-0100-B007:start}
\begin{defish}
\AxiomC{\sFmla{\True}{\lnot !A}}
\RightLabel{\TRule{\True}{\lnot}}
\UnaryInfC{\sFmla{\False}{!A}}
\DisplayProof
\hfill
\AxiomC{\sFmla{\False}{\lnot !A}}
\RightLabel{\TRule{\False}{\lnot}}
\UnaryInfC{\sFmla{\True}{!A}}
\DisplayProof
\end{defish}
\label{olpseg:OLP-0100-B007:end}

\phantomsection\label{olpseg:OLP-0100-B008:start}
\subsection{د~\LR{$\land$} قاعدې}
\label{olpseg:OLP-0100-B008:end}

\phantomsection\label{olpseg:OLP-0100-B009:start}
\begin{defish}\noindent
\AxiomC{\sFmla{\True}{!A \land !B}}
\RightLabel{\TRule{\True}{\land}}
\UnaryInfC{\sFmla{\True}{!A}}
\noLine
\UnaryInfC{\sFmla{\True}{!B}}
\DisplayProof
\hfill
\AxiomC{\sFmla{\False}{!A \land !B}}
\RightLabel{\TRule{\False}{\land}}
\UnaryInfC{\LR{$\sFmla{\False}{!A} \quad \mid \quad \sFmla{\False}{!B}$}}
\DisplayProof
\end{defish}
\label{olpseg:OLP-0100-B009:end}

\phantomsection\label{olpseg:OLP-0100-B010:start}
\subsection{د~\LR{$\lor$} قاعدې}
\label{olpseg:OLP-0100-B010:end}

\phantomsection\label{olpseg:OLP-0100-B011:start}
\begin{defish}
\AxiomC{\sFmla{\True}{!A \lor !B}}
\RightLabel{\TRule{\True}{\lor}}
\UnaryInfC{\LR{$\sFmla{\True}{!A} \quad \mid \quad \sFmla{\True}{!B}$}}
\DisplayProof
\hfill
\AxiomC{\sFmla{\False}{!A \lor !B}}
\RightLabel{\TRule{\False}{\lor}}
\UnaryInfC{\sFmla{\False}{!A}}
\noLine
\UnaryInfC{\sFmla{\False}{!B}}
\DisplayProof
\end{defish}
\label{olpseg:OLP-0100-B011:end}

\phantomsection\label{olpseg:OLP-0100-B012:start}
\subsection{د~\LR{$\lif$} قاعدې}
\label{olpseg:OLP-0100-B012:end}

\phantomsection\label{olpseg:OLP-0100-B013:start}
\begin{defish}
\AxiomC{\sFmla{\True}{!A \lif !B}}
\RightLabel{\TRule{\True}{\lif}}
\UnaryInfC{\LR{$\sFmla{\False}{!A} \quad \mid \quad \sFmla{\True}{!B}$}}
\DisplayProof
\hfill
\AxiomC{\sFmla{\False}{!A \lif !B}}
\RightLabel{\TRule{\False}{\lif}}
\UnaryInfC{\sFmla{\True}{!A}}
\noLine
\UnaryInfC{\sFmla{\False}{!B}}
\DisplayProof
\end{defish}
\label{olpseg:OLP-0100-B013:end}

\phantomsection\label{olpseg:OLP-0100-B014:start}
\subsection{د پرېکون قاعده}
\label{olpseg:OLP-0100-B014:end}

\phantomsection\label{olpseg:OLP-0100-B015:start}
\begin{defish}
\AxiomC{}
\RightLabel{\Cut}
\UnaryInfC{\LR{$\sFmla{\True}{!A} \quad \mid \quad \sFmla{\False}{!A}$}}
\DisplayProof
\end{defish}
\label{olpseg:OLP-0100-B015:end}

\phantomsection\label{olpseg:OLP-0100-B016:start}
د~\Cut{} قاعده پر مخکښ نښه لرونکے فارمول نۀ «پلې کېږي»؛ بلکې
راکوي چې په تابلو کښې هره څانګه په دوو ووېشو: يوه څانګه
\LR{$\sFmla{\True}{!A}$} او بله~\LR{$\sFmla{\False}{!A}$} لري۔ دا اړينه نۀ
ده---د نښه لرونکي فارمولونه هر هغه سټ چې تړلے تابلو ولري،
له~\Cut پرته هم يو تړلے تابلو لري---خو راکوي چې تابلوګانې په
اسانه طريقه سره يو ځاے کړو۔
\label{olpseg:OLP-0100-B016:end}

\phantomsection\label{olpseg:OLP-0101-B004:start}
\olfileid{fol}{tab}{qrl}
\label{olpseg:OLP-0101-B004:end}

\phantomsection\label{olpseg:OLP-0101-B005:start}
\olsection{د سورونو قاعدې}
\label{olpseg:OLP-0101-B005:end}

\phantomsection\label{olpseg:OLP-0101-B006:start}
\subsection{د~\LR{$\lforall$} قاعدې}
\label{olpseg:OLP-0101-B006:end}

\phantomsection\label{olpseg:OLP-0101-B007:start}
\begin{defish}
\AxiomC{\sFmla{\True}{\lforall[x][!A(x)]}}
\RightLabel{\TRule{\True}{\forall}}
\UnaryInfC{\sFmla{\True}{!A(t)}}
\DisplayProof
\hfill
\AxiomC{\sFmla{\False}{\lforall[x][!A(x)]}}
\RightLabel{\TRule{\False}{\lforall}}
\UnaryInfC{\sFmla{\False}{!A(a)}}
\DisplayProof
\end{defish}
\label{olpseg:OLP-0101-B007:end}

\phantomsection\label{olpseg:OLP-0101-B008:start}
په~\TRule{\True}{\lforall} کښې \LR{$t$} تړلې اصطلاح ده (يعنې متغيرونه
نۀ لري)۔ په~\TRule{\False}{\lforall} کښې \LR{$a$} داسې ثابت سمبول
دے چې بايد د~\TRule{\False}{\lforall} تر قاعدې پاس په څانګه کښې
هېڅ ځاے نۀ وي راغلے۔ \LR{$a$} د~\TRule{\False}{\forall} استنتاج
\emph{ځانګړے متغير} بولو۔\footnote{د «ځانګړي متغير» اصطلاح کاروو،
سره له دې چې په پورته قاعده کښې \LR{$a$} يو ثابت سمبول دے۔ دا
تاريخي لاملونه لري۔}
\label{olpseg:OLP-0101-B008:end}

\phantomsection\label{olpseg:OLP-0101-B009:start}
\subsection{د~\LR{$\lexists$} قاعدې}
\label{olpseg:OLP-0101-B009:end}

\phantomsection\label{olpseg:OLP-0101-B010:start}
\begin{defish}
\AxiomC{\sFmla{\True}{\lexists[x][!A(x)]}}
\RightLabel{\TRule{\True}{\lexists}}
\UnaryInfC{\sFmla{\True}{!A(a)}}
\DisplayProof
\hfill
\AxiomC{\sFmla{\False}{\lexists[x][!A(x)]}}
\RightLabel{\TRule{\False}{\lexists}}
\UnaryInfC{\sFmla{\False}{!A(t)}}
\DisplayProof
\end{defish}
\label{olpseg:OLP-0101-B010:end}

\phantomsection\label{olpseg:OLP-0101-B011:start}
بيا هم \LR{$t$} تړلې اصطلاح ده، او~\LR{$a$} داسې ثابت سمبول دے چې د
\TRule{\True}{\lexists} تر قاعدې پاس په څانګه کښې نۀ راځي۔
\LR{$a$} د~\TRule{\True}{\lexists} استنتاج \emph{ځانګړے متغير} بولو۔
\label{olpseg:OLP-0101-B011:end}

\phantomsection\label{olpseg:OLP-0101-B012:start}
هغه شرط چې ځانګړے متغير د~\TRule{\False}{\lforall} يا
\TRule{\True}{\lexists} استنتاج تر پاسه په څانګه کښې نۀ راځي،
\emph{د ځانګړي متغير شرط} بلل کېږي۔
\label{olpseg:OLP-0101-B012:end}

\phantomsection\label{olpseg:OLP-0101-B013:start}
\begin{explain}
هغه قرارداد را ياد کړئ چې کله \LR{$!A$} داسې فارمول وي چې
متغير~\LR{$x$} پکې ازاد وي، دا د~\LR{$!A(x)$} په ليکلو ښيو۔ په همدې
سياق کښې \LR{$!A(t)$} د~\LR{$\Subst{!A}{t}{x}$} لنډه بڼه ده۔ نو د
\TRule{\False}{\lexists} قاعده داسې هم ليکلے شو:
\begin{prooftree}
  \AxiomC{\sFmla{\False}{\lexists[x][!A]}}
  \RightLabel{\TRule{\False}{\lexists}}
  \UnaryInfC{\sFmla{\False}{\Subst{!A}{t}{x}}}
\end{prooftree}
پام وکړئ چې \LR{$t$} له مخکښې په~\LR{$!A$} کښې راتلے شي؛ مثلاً \LR{$!A$} کېدے
شي \LR{$\Atom{\Obj P}{t,x}$} وي۔ نو له
\LR{$ \sFmla{\False}{\lexists[x][\Atom{\Obj P}{t,x}]}$} څخه
\LR{$\sFmla{\False}{\Atom{\Obj P}{t,t}}$} استنتاجول د
\TRule{\False}{\lexists} سمه کارونه ده۔ خو په
\TRule{\False}{\lforall} او~\TRule{\True}{\lexists} کښې د ځانګړي
متغير شرطونه غواړي چې ثابت سمبول~\LR{$a$} په~\LR{$!A$} کښې نۀ راځي۔ نو له
\LR{$\sFmla{\False}{\lforall[x][\Atom{\Obj P}{a,x}]}$}
څخه د~\LR{$\TRule{\False}{\lforall}$} په کارولو
\LR{$\sFmla{\False}{\Atom{\Obj P}{a,a}}$} په سمه توګه نۀ شئ استنتاجولے۔
\end{explain}
\label{olpseg:OLP-0101-B013:end}

\phantomsection\label{olpseg:OLP-0101-B014:start}
\begin{explain}
په~\TRule{\True}{\lforall} او~\TRule{\False}{\lexists} کښې د
تړلې اصطلاح~\LR{$t$} پر انتخاب نور بنديز نشته۔ بل لوري ته، د
\TRule{\True}{\lexists} او~\TRule{\False}{\lforall} په قاعدو کښې
د ځانګړي متغير شرط غواړي چې ثابت سمبول~\LR{$a$} د اړوند استنتاج تر
پاسه په څانګو کښې هېڅ ځاے نۀ راځي۔ دا شرط د نظام د سموالي د
يقيني کولو دپاره اړين دے۔ له دې شرط پرته لاندې به د
\LR{$\lexists[x][\formula{A}(x)] \lif \lforall[x][\formula{A}(x)]$}
دپاره يو تړلے تابلو و:
\begin{center}
\begin{tableau}{}
  [\sFmla{\False}{\lexists[x][\formula{A}(x)] \lif \lforall[x][\formula{A}(x)]}, just=\TAss
    [\sFmla{\True}{\lexists[x][\formula{A}(x)]},
      just={\TRule{\False}{\lif}[1]}
      [\sFmla{\False}{\lforall[x][\formula{A}(x)]},
        just={\TRule{\False}{\lif}[1]}
        [\sFmla{\True}{\formula{A}(a)},
          just={\TRule{\True}{\lexists}[2]}
          [\sFmla{\False}{\formula{A}(a)},
            just={\TRule{\False}{\lforall}[3]}, close]
        ]
      ]
    ]
  ]
\end{tableau}
\end{center}
خو \LR{$\lexists[x][\formula{A}(x)] \lif
\lforall[x][\formula{A}(x)]$} معتبر نۀ دے۔
\end{explain}
\label{olpseg:OLP-0101-B014:end}

\phantomsection\label{olpseg:OLP-0102-B004:start}
\olfileid{fol}{tab}{der}
\label{olpseg:OLP-0102-B004:end}

\phantomsection\label{olpseg:OLP-0102-B005:start}
\olsection{تابلوګانې}
\label{olpseg:OLP-0102-B005:end}

\phantomsection\label{olpseg:OLP-0102-B006:start}
\begin{explain}
و مو ويل چې فرض څه دے، او د استنتاج قاعدې مو ورکړې۔ له دغو څخه
تابلوګانې په استقرايي ډول جوړېږي: هر تابلو يا يوه يوازېنۍ
څانګه ده چې له يوه يا زياتو فرضونو جوړه وي، يا پر يوه څانګه د
استنتاج د کومې قاعدې په پلي کولو له بل تابلو څخه اخستل کېږي۔
\end{explain}
\label{olpseg:OLP-0102-B006:end}

\phantomsection\label{olpseg:OLP-0102-B007:start}
\begin{defn}[تابلو]
د فرضونو~\LR{$\sFmla{S_1}{!A_1}$}، \dots،
\LR{$\sFmla{S_n}{!A_n}$} دپاره (چې هر~\LR{$S_i$} يا~\LR{$\True$} وي
يا~\LR{$\False$}) تابلو د~نښه لرونکي فارمولونه يوه متناهي ونه
ده چې لاندې شرطونه پوره کوي:
\begin{enumerate}
\item د ونې تر ټولو پاسني \LR{$n$}~نښه لرونکي فارمولونه
  \LR{$\sFmla{S_i}{!A_i}$} دي، يو تر بل لاندې۔
\item په ونه کښې هر نښه لرونکے فارمول چې له فرضونو څخه نۀ وي، تر
  ځان پاس په څانګه کښې پر نښه لرونکے فارمول د يوې استنتاجي قاعدې
  له سمې کارونې څخه راوتلے وي۔
\end{enumerate}
د~تابلو يوه څانګه هله او يوازې هله \emph{تړلې} ده چې
\LR{$\sFmla{\True}{!A}$} او~\LR{$\sFmla{\False}{!A}$} دواړه ولري؛ که نۀ وي
\emph{پرانيستې} ده۔ هغه تابلو چې هره څانګه ئې تړلې وي، د
خپلو فرضونو د سټ \emph{تړلے تابلو} دے۔ که تابلو تړلے
نۀ وي، يعنې لږ تر لږه يوه پرانيستې څانګه ولري، نو \emph{پرانيستے}
دے۔
\end{defn}
\label{olpseg:OLP-0102-B007:end}

\phantomsection\label{olpseg:OLP-0102-B008:start}
\begin{ex}
  د فرضونو هر سټ په يوازې ځان تابلو دے، خو په عام ډول به
  تړلے نۀ وي۔ (په څرګند ډول يوازې هله تړلے دے چې فرضونه له مخکښې
  د~نښه لرونکي فارمولونه يوه جوړه
  \LR{$\sFmla{\True}{!A}$} او~\LR{$\sFmla{\False}{!A}$} ولري۔)
\label{olpseg:OLP-0102-B008:end}

\phantomsection\label{olpseg:OLP-0102-B009:start}
له يوه تابلو (پرانيستي يا تړلي) څخه د کومې استنتاجي قاعدې
  په پلي کولو يو نوے، لوے تابلو اخستلے شو۔ قاعده په هغه کښې پر
  نښه لرونکے فارمول~\LR{$!A$} پلې کوو۔ دا قاعده به د هرې هغې څانګې په
  پاے کښې يو يا زيات نښه لرونکي فارمولونه ورزيات کړي چې د~\LR{$!A$} هغه
  پېښه لري چې قاعده مو پرې پلې کړې ده۔
\label{olpseg:OLP-0102-B009:end}

\phantomsection\label{olpseg:OLP-0102-B010:start}
مثلاً فرض~\LR{$\sFmla{\True}{!A \land \lnot !A}$} ونيسئ۔ لاندې هغه
  (پرانيستے) تابلو دے چې يوازې له همدې فرض څخه جوړ دے:
  \begin{center}
    \begin{tableau}{}
      [\sFmla{\True}{\formula{A} \land \lnot \formula{A}}, just=\TAss]
    \end{tableau}{}
  \end{center}
  د~\LR{$\TRule{\True}{\land}$} قاعده پر فرض پلې کوو او ترې نوے
  تابلو اخلو۔ دا قاعده راکوي چې تابلو ته دوه نوې کرښې،
  \LR{$\sFmla{\True}{!A}$} او~\LR{$\sFmla{\True}{\lnot !A}$}، ورزياتې کړو:
  \begin{center}
    \begin{tableau}{}
      [\sFmla{\True}{\formula{A} \land \lnot \formula{A}}, just=\TAss
        [\sFmla{\True}{\formula{A}}, just={\TRule{\True}{\land}[1]},
          [\sFmla{\True}{\lnot \formula{A}}, just={\TRule{\True}{\land}[1]}
          ]
        ]
      ]
    \end{tableau}{}
  \end{center}
  کله چې تابلوګانې ليکو، پلې شوې قاعدې په ښي لوري کښې ثبتوو
  (مثلاً \LR{$\TRule{\True}{\land} 1$} دا معنٰی لري چې پر دې کرښه
  نښه لرونکے فارمول د~\LR{$\TRule{\True}{\land}$} قاعدې د هغې کارونې
  پايله ده چې د~\LR{$1$} پر کرښه نښه لرونکے فارمول ته شوې ده)۔ دا نوے
  تابلو اوس نور نښه لرونکي فارمولونه لري، خو يوازې پر يوه
  (\LR{$\sFmla{\True}{\lnot !A}$}) قاعده پلې کولے شو (په دې حالت کښې
  \LR{$\TRule{\True}{\lnot}$} قاعده)۔ له دې څخه لاندې تړلے
  تابلو راځي:
  \begin{center}
    \begin{tableau}{}
      [\sFmla{\True}{\formula{A} \land \lnot \formula{A}}, just=\TAss
        [\sFmla{\True}{\formula{A}}, just={\TRule{\True}{\land}[1]}
          [\sFmla{\True}{\lnot \formula{A}}, just={\TRule{\True}{\land}[1]}
            [\sFmla{\False}{\formula{A}}, just={\TRule{\True}{\lnot}[3]}, close]
          ]
        ]
      ]
    \end{tableau}{}
  \end{center}
\end{ex}
\label{olpseg:OLP-0102-B010:end}

\phantomsection\label{olpseg:OLP-0103-B004:start}
\olfileid{fol}{tab}{pro}
\label{olpseg:OLP-0103-B004:end}

\phantomsection\label{olpseg:OLP-0103-B005:start}
\olsection{د~تابلوګانې بېلګې}
\label{olpseg:OLP-0103-B005:end}

\phantomsection\label{olpseg:OLP-0103-B006:start}
\begin{ex}
راځئ د~تړلے فارمول~\LR{$(!A \land !B) \lif !A$} دپاره يو تړلے
تابلو ومومو۔
\label{olpseg:OLP-0103-B006:end}

\phantomsection\label{olpseg:OLP-0103-B007:start}
لومړے اړوند فرض د~تابلو په سر کښې ليکو۔
\begin{oltableau}
  [\sFmla{\False}{(\formula{A} \land \formula{B}) \lif \formula{A}},
    just = \TAss]
\end{oltableau}
\label{olpseg:OLP-0103-B007:end}

\phantomsection\label{olpseg:OLP-0103-B008:start}
يوازې يو فرض شته، نو يوازې يو نښه لرونکے فارمول دے چې قاعده پرې
پلې کولے شو۔ (د هر نښه لرونکے فارمول دپاره تل تر ټولو زياته يوه
قاعده پلې کېدے شي: دا د~تړلے فارمول د اړوندې نښې او
اصلي عملګر قاعده ده۔) په دې حالت کښې بايد
\LR{$\TRule{\False}{\lif}$} پلې کړو۔
\begin{oltableau}
  [\sFmla{\False}{(\formula{A} \land \formula{B}) \lif \formula{A}},
    checked, just = \TAss
    [\sFmla{\True}{\formula{A} \land \formula{B}},
      just={\TRule{\False}{\lif}[1]}
      [\sFmla{\False}{\formula{A}}, just={\TRule{\False}{\lif}[1]}]
    ]
  ]
\end{oltableau}
د دې حساب ساتلو دپاره چې پر کومو نښه لرونکي فارمولونه مو اړوندې
قاعدې پلې کړې دي، د جملې تر څنګ د سم نښان ليکو۔ خو د سم نښان
\emph{يوازې} هله وليکئ چې قاعده پر ټولو پرانيستو څانګو پلې شوې
وي۔ کله چې پر کوم نښه لرونکے فارمول په هره پرانيستې څانګه کښې
اړونده قاعده پلې شي، بيا ورته د بېرته ورګرځېدو او د قاعدې د بيا
پلي کولو اړتيا نۀ لرو۔ په دې حالت کښې يوازې يوه څانګه ده، نو قاعده
يوازې يو ځل پلې کول غواړي۔ (پام وکړئ چې د سم نښانونه يوازې د
تابلوګانو د جوړولو اسانتيا ده او په رسمي ډول د تابلوګانو د نحوې
برخه نۀ ده۔)
\label{olpseg:OLP-0103-B008:end}

\phantomsection\label{olpseg:OLP-0103-B009:start}
يو نوے نښه لرونکے فارمول شته چې قاعده پرې پلې کولے شو: پر
\LR{$2$}~کرښه \LR{$\sFmla{\True}{!A \land !B}$}۔ د
\LR{$\TRule{\True}{\land}$} قاعدې له پلي کولو څخه دا راځي:
\begin{oltableau}
  [\sFmla{\False}{(\formula{A} \land \formula{B}) \lif \formula{A}},
    checked, just = \TAss
    [\sFmla{\True}{\formula{A} \land \formula{B}},
      just={\TRule{\False}{\lif}[1]}, checked
      [\sFmla{\False}{\formula{A}}, just={\TRule{\False}{\lif}[1]}
        [\sFmla{\True}{\formula{A}}, just={\TRule{\True}{\land}[2]}
          [\sFmla{\True}{\formula{B}}, just={\TRule{\True}{\land}[2]}, close
          ]
        ]
      ]
    ]
  ]
\end{oltableau}
ځکه اوس څانګه \LR{$\sFmla{\True}{!A}$} (پر~\LR{$4$} کرښه) او
\LR{$\sFmla{\False}{!A}$} (پر~\LR{$3$} کرښه) دواړه لري، څانګه تړلې ده۔
ځکه همدا يوازېنۍ څانګه ده، تابلو تړلے دے۔ د
\LR{$(!A \land !B) \lif !A$} دپاره مو يو تړلے تابلو وموند۔
\end{ex}
\label{olpseg:OLP-0103-B009:end}

\phantomsection\label{olpseg:OLP-0103-B010:start}
\begin{ex}
اوس راځئ د~\LR{$(\lnot !A \lor !B) \lif (!A
\lif !B)$} دپاره يو تړلے تابلو ومومو۔
\label{olpseg:OLP-0103-B010:end}

\phantomsection\label{olpseg:OLP-0103-B011:start}
له اړوند فرض څخه پيل کوو:
\begin{oltableau}
  [\sFmla{\False}{(\lnot \formula{A} \lor \formula{B}) \lif
      (\formula{A} \lif \formula{B})}, just=\TAss]
\end{oltableau}
په دې تابلو کښې د يوازېني نښه لرونکے فارمول
اصلي عملګر~\LR{$\lif$} او نښه~\LR{$\False$} ده، نو د
\LR{$\TRule{\False}{\lif}$} قاعده پرې پلې کوو او دا اخلو:
\begin{oltableau}
  [\sFmla{\False}{(\lnot \formula{A} \lor \formula{B})
      \lif (\formula{A} \lif \formula{B})}, just=\TAss, checked
    [\sFmla{\True}{\lnot \formula{A} \lor \formula{B}},
      just={\TRule{\False}{\lif}[1]}
      [\sFmla{\False}{(\formula{A} \lif \formula{B})},
        just={\TRule{\False}{\lif}[1]}
      ]
    ]
  ]
\end{oltableau}
اوس دا ټاکنه لرو چې پر~\LR{$2$} کرښه~\LR{$\TRule{\True}{\lor}$} پلې کړو
که پر~\LR{$3$} کرښه \LR{$\TRule{\False}{\lif}$}۔ په حقيقت کښې د ترتيب
ټاکنه توپير نۀ کوي، تر هغې چې د هر نښه لرونکے فارمول اړونده قاعده
په هره څانګه کښې پلې شي۔ نو لومړۍ اخلو۔ د
\LR{$\TRule{\True}{\lor}$} قاعده تابلو ته څانګې ورکوي، او د
قاعدې دوه پايلې به هغه نوي نښه لرونکي فارمولونه وي چې دوو نويو
څانګو ته ورزياتېږي۔ له دې څخه دا راځي:
\begin{oltableau}
  [\sFmla{\False}{(\lnot \formula{A} \lor \formula{B}) \lif
      (\formula{A} \lif \formula{B})}, just=\TAss, checked
    [\sFmla{\True}{\lnot \formula{A} \lor \formula{B}},
      just={\TRule{\False}{\lif}[1]}, checked
      [\sFmla{\False}{(\formula{A} \lif \formula{B})},
        just={\TRule{\False}{\lif}[1]}
        [\sFmla{\True}{\lnot \formula{A}}, just={\TRule{\True}{\lor}[2]}]
        [\sFmla{\True}{\formula{B}}, just={\TRule{\True}{\lor}[2]}]
      ]
    ]
  ]
\end{oltableau}
لا مو پر~\LR{$3$} کرښه د~\LR{$\TRule{\False}{\lif}$} قاعده نۀ ده پلې
کړې؛ اوس ئې کوو۔ د وخت د سپما دپاره ئې پر دواړو څانګو پلې کوو۔
را ياد کړئ چې د کوم نښه لرونکے فارمول تر څنګ يوازې هله د سم نښان
ليکو چې اړونده قاعده مو په هره پرانيستې څانګه کښې پلې کړې وي۔ نو
غوره ده چې قاعده د هرې هغې څانګې په پاے کښې پلې کړو چې اړوند
نښه لرونکے فارمول لري۔ په دې توګه به په بېلابېلو څانګو کښې لاندې
هماغه نښه لرونکے فارمول ته بېرته نۀ ورګرځو۔
\begin{oltableau}
  [\sFmla{\False}{(\lnot \formula{A} \lor \formula{B}) \lif
      (\formula{A} \lif \formula{B})}, just=\TAss, checked
    [\sFmla{\True}{\lnot \formula{A} \lor \formula{B}},
      just={\TRule{\False}{\lif}[1]}, checked
      [\sFmla{\False}{(\formula{A} \lif \formula{B})},
        just={\TRule{\False}{\lif}[1]}, checked
        [\sFmla{\True}{\lnot \formula{A}}, just={\TRule{\True}{\lor}[2]}
          [\sFmla{\True}{\formula{A}}, just={\TRule{\False}{\lif}[3]}
            [\sFmla{\False}{\formula{B}}, just={\TRule{\False}{\lif}[3]}]
          ]
        ]
        [\sFmla{\True}{\formula{B}}, just={\TRule{\True}{\lor}[2]}
          [\sFmla{\True}{\formula{A}}, just={\TRule{\False}{\lif}[3]}
            [\sFmla{\False}{\formula{B}}, just={\TRule{\False}{\lif}[3]}, close]
          ]
        ]
      ]
    ]
  ]
\end{oltableau}
اوس ښۍ څانګه تړلې ده۔ په کيڼه څانګه کښې لا هم پر~\LR{$4$} کرښه د
\LR{$\TRule{\True}{\lnot}$} قاعده پلې کولے شو۔ له دې څخه
\LR{$\sFmla{\False}{!A}$} راځي او کيڼه څانګه تړل کېږي:
\begin{oltableau}
  [\sFmla{\False}{(\lnot \formula{A} \lor \formula{B}) \lif
      (\formula{A} \lif \formula{B})}, just=\TAss, checked
    [\sFmla{\True}{\lnot \formula{A} \lor \formula{B}},
      just={\TRule{\False}{\lif}[1]}, checked
      [\sFmla{\False}{(\formula{A} \lif \formula{B})},
        just={\TRule{\False}{\lif}[1]}, checked
        [\sFmla{\True}{\lnot \formula{A}}, just={\TRule{\True}{\lor}[2]}
          [\sFmla{\True}{\formula{A}}, just={\TRule{\False}{\lif}[3]}
            [\sFmla{\False}{\formula{B}}, just={\TRule{\False}{\lif}[3]}
              [\sFmla{\False}{\formula{A}},
                just={\TRule{\True}{\lnot}[4]}, close]
            ]
          ]
        ]
        [\sFmla{\True}{\formula{B}}, just={\TRule{\True}{\lor}[2]}
          [\sFmla{\True}{\formula{A}}, just={\TRule{\False}{\lif}[3]}
            [\sFmla{\False}{\formula{B}},
              just={\TRule{\False}{\lif}[3]}, close]
          ]
        ]
      ]
    ]
  ]
\end{oltableau}
\end{ex}
\label{olpseg:OLP-0103-B011:end}

\phantomsection\label{olpseg:OLP-0103-B012:start}
\begin{ex}
د فرضونو په توګه د هر شمېر نښه لرونکي فارمولونه دپاره
تابلوګانې ورکولے شو۔ ډېر ځله له يوې زياتې داسې قاعدې پلي کول هم
اړين وي چې څانګې ورکوي؛ او په عام ډول تابلو هر شمېر څانګې
لرلے شي۔ مثلاً د
\LR{$\{\sFmla{\True}{!A \lor (!B \land !C)}, \sFmla{\False}{(!A \lor !B) \land (!A
  \lor !C)}\}$} دپاره تابلو ونيسئ۔ لومړے پر لومړي فرض د
\LR{$\TRule{\True}{\lor}$} قاعده پلې کوو:
\begin{oltableau}
  [\sFmla{\True}{\formula{A} \lor (\formula{B} \land \formula{C})},
    just=\TAss, checked
    [\sFmla{\False}{(\formula{A} \lor \formula{B}) \land
        (\formula{A} \lor \formula{C})}, just=\TAss
      [\sFmla{\True}{\formula{A}}, just={\TRule{\True}{\lor}[1]}]
      [\sFmla{\True}{\formula{B} \land \formula{C}},
        just={\TRule{\True}{\lor}[1]}]
    ]
  ]
\end{oltableau}
اوس پر~\LR{$2$} کرښه د~\LR{$\TRule{\False}{\land}$} قاعده پلې کولے
شو۔ دا په دواړو څانګو کښې په يو وخت کوو، نو پر~\LR{$2$} کرښه د سم نښان
وهلے شو:
\begin{oltableau}
  [\sFmla{\True}{\formula{A} \lor (\formula{B} \land \formula{C})},
    just=\TAss, checked
    [\sFmla{\False}{(\formula{A} \lor \formula{B}) \land
        (\formula{A} \lor \formula{C})}, just=\TAss, checked
      [\sFmla{\True}{\formula{A}}, just={\TRule{\True}{\lor}[1]}
        [\sFmla{\False}{\formula{A} \lor \formula{B}},
          just={\TRule{\False}{\land}[2]}]
        [\sFmla{\False}{\formula{A} \lor \formula{C}},
          just={\TRule{\False}{\land}[2]}]
      ]
      [\sFmla{\True}{\formula{B} \land \formula{C}},
        just={\TRule{\True}{\lor}[1]}
        [\sFmla{\False}{\formula{A} \lor \formula{B}},
          just={\TRule{\False}{\land}[2]}]
        [\sFmla{\False}{\formula{A} \lor \formula{C}},
          just={\TRule{\False}{\land}[2]}]
      ]
    ]
  ]
\end{oltableau}
اوس \LR{$\TRule{\False}{\lor}$} پر ټولو هغو څانګو پلې کولے شو چې
\LR{$!A \lor !B$} لري:
\begin{oltableau}
  [\sFmla{\True}{\formula{A} \lor (\formula{B} \land \formula{C})},
    just=\TAss, checked
    [\sFmla{\False}{(\formula{A} \lor \formula{B}) \land
        (\formula{A} \lor \formula{C})}, just=\TAss, checked
      [\sFmla{\True}{\formula{A}}, just={\TRule{\True}{\lor}[1]}
        [\sFmla{\False}{\formula{A} \lor \formula{B}},
          just={\TRule{\False}{\land}[2]}, checked
          [\sFmla{\False}{\formula{A}}, just={\TRule{\False}{\lor}[4]}
            [\sFmla{\False}{\formula{B}},
              just={\TRule{\False}{\lor}[4]}, close]
          ]
        ]
        [\sFmla{\False}{\formula{A} \lor \formula{C}},
          just={\TRule{\False}{\land}[2]}]
      ]
      [\sFmla{\True}{\formula{B} \land \formula{C}},
        just={\TRule{\True}{\lor}[1]}
        [\sFmla{\False}{\formula{A} \lor \formula{B}},
          just={\TRule{\False}{\land}[2]}, checked
          [\sFmla{\False}{\formula{A}}, just={\TRule{\False}{\lor}[4]}
            [\sFmla{\False}{\formula{B}},
              just={\TRule{\False}{\lor}[4]}]
          ]
        ]
        [\sFmla{\False}{\formula{A} \lor \formula{C}},
          just={\TRule{\False}{\land}[2]}]
      ]
    ]
  ]
\end{oltableau}
اوس تر ټولو کيڼه څانګه تړلې ده۔ راځئ
\LR{$\TRule{\False}{\lor}$} پر~\LR{$!A \lor !C$} پلې کړو:
\begin{oltableau}
  [\sFmla{\True}{\formula{A} \lor (\formula{B} \land \formula{C})},
    just=\TAss, checked
    [\sFmla{\False}{(\formula{A} \lor \formula{B}) \land
        (\formula{A} \lor \formula{C})}, just=\TAss, checked
      [\sFmla{\True}{\formula{A}}, just={\TRule{\True}{\lor}[1]}
        [\sFmla{\False}{\formula{A} \lor \formula{B}},
          just={\TRule{\False}{\land}[2]}, checked
          [\sFmla{\False}{\formula{A}},
            just={\TRule{\False}{\lor}[4]}
            [\sFmla{\False}{\formula{B}},
              just={\TRule{\False}{\lor}[4]}, close]
          ]
        ]
        [\sFmla{\False}{\formula{A} \lor \formula{C}},
          just={\TRule{\False}{\land}[2]}, checked
          [\sFmla{\False}{\formula{A}},
            just={\TRule{\False}{\lor}[4]}, move by=2
            [\sFmla{\False}{\formula{C}},
              just={\TRule{\False}{\lor}[4]}, close]
          ]
        ]
      ]
      [\sFmla{\True}{\formula{B} \land \formula{C}},
        just={\TRule{\True}{\lor}[1]}
        [\sFmla{\False}{\formula{A} \lor \formula{B}},
          just={\TRule{\False}{\land}[2]}, checked
          [\sFmla{\False}{\formula{A}},
            just={\TRule{\False}{\lor}[4]}
            [\sFmla{\False}{\formula{B}},
              just={\TRule{\False}{\lor}[4]}
            ]
          ]
        ]
        [\sFmla{\False}{\formula{A} \lor \formula{C}},
          just={\TRule{\False}{\land}[2]}, checked
          [\sFmla{\False}{\formula{A}},
            just={\TRule{\False}{\lor}[4]}, move by=2
            [\sFmla{\False}{\formula{C}},
              just={\TRule{\False}{\lor}[4]}]
          ]
        ]
      ]
    ]
  ]
\end{oltableau}
پام وکړئ چې د روښانتيا دپاره مو د
\LR{$\TRule{\False}{\lor}$} د دويم ځل کارونې پايله لاندې يوړه۔ په دې
حالت کښې دې ته اړتيا نۀ وه؛ ځکه توجيهونه به يو شان وو۔
\label{olpseg:OLP-0103-B012:end}

\phantomsection\label{olpseg:OLP-0103-B013:start}
دوه څانګې پرانيستې پاتې دي، او پر~\LR{$3$} کرښه
\LR{$\sFmla{\True}{!B \land !C}$} لا بې د سم نښانه دے۔ پر هغې
\LR{$\TRule{\True}{\land}$} پلې کوو او تړلے تابلو اخلو:
\begin{oltableau}
  [\sFmla{\True}{\formula{A} \lor (\formula{B} \land \formula{C})},
    just=\TAss, checked
    [\sFmla{\False}{(\formula{A} \lor \formula{B}) \land
        (\formula{A} \lor \formula{C})}, just=\TAss, checked
      [\sFmla{\True}{\formula{A}}, just={\TRule{\True}{\lor}[1]}
        [\sFmla{\False}{\formula{A} \lor \formula{B}},
          just={\TRule{\False}{\land}[2]}, checked
          [\sFmla{\False}{\formula{A}},
            just={\TRule{\False}{\lor}[4]}
            [\sFmla{\False}{\formula{B}},
              just={\TRule{\False}{\lor}[4]}, close]
          ]
        ]
        [\sFmla{\False}{\formula{A} \lor \formula{C}},
          just={\TRule{\False}{\land}[2]}, checked
          [\sFmla{\False}{\formula{A}},
            just={\TRule{\False}{\lor}[4]}
            [\sFmla{\False}{\formula{C}},
              just={\TRule{\False}{\lor}[4]}, close]
          ]
        ]
      ]
      [\sFmla{\True}{\formula{B} \land \formula{C}},
        just={\TRule{\True}{\lor}[1]}, checked
        [\sFmla{\False}{\formula{A} \lor \formula{B}},
          just={\TRule{\False}{\land}[2]}, checked
          [\sFmla{\False}{\formula{A}},
            just={\TRule{\False}{\lor}[4]}
            [\sFmla{\False}{\formula{B}},
              just={\TRule{\False}{\lor}[4]}
              [\sFmla{\True}{\formula{B}},
                just={\TRule{\True}{\land}[3]}
                [\sFmla{\True}{\formula{C}},
                  just={\TRule{\True}{\land}[3]},close
                ]
              ]
            ]
          ]
        ]
        [\sFmla{\False}{\formula{A} \lor \formula{C}},
          just={\TRule{\False}{\land}[2]}, checked
          [\sFmla{\False}{\formula{A}},
            just={\TRule{\False}{\lor}[4]}
            [\sFmla{\False}{\formula{C}},
              just={\TRule{\False}{\lor}[4]}
              [\sFmla{\True}{\formula{B}},
                just={\TRule{\True}{\land}[3]}
                [\sFmla{\True}{\formula{C}},
                  just={\TRule{\True}{\land}[3]},close
                ]
              ]
            ]
          ]
        ]
      ]
    ]
  ]
\end{oltableau}
\label{olpseg:OLP-0103-B013:end}

\phantomsection\label{olpseg:OLP-0103-B014:start}
د پرتلې دپاره، لاندې د فرضونو د همدې سټ يو تړلے تابلو دے
چې قاعدې پکې په بل ترتيب پلې شوې دي:
\begin{oltableau}
  [\sFmla{\True}{\formula{A} \lor (\formula{B} \land \formula{C})},
    just=\TAss, checked
    [\sFmla{\False}{(\formula{A} \lor \formula{B}) \land
        (\formula{A} \lor \formula{C})}, just=\TAss, checked
      [\sFmla{\False}{\formula{A} \lor \formula{B}},
        just={\TRule{\False}{\land}[2]}, checked
        [\sFmla{\False}{\formula{A}},
          just={\TRule{\False}{\lor}[3]}
          [\sFmla{\False}{\formula{B}},
            just={\TRule{\False}{\lor}[3]}
            [\sFmla{\True}{\formula{A}},
              just={\TRule{\True}{\lor}[1]},close
            ]
            [\sFmla{\True}{\formula{B} \land \formula{C}},
              just={\TRule{\True}{\lor}[1]}, checked
              [\sFmla{\True}{\formula{B}},
                just={\TRule{\True}{\land}[6]}
                [\sFmla{\True}{\formula{C}},
                  just={\TRule{\True}{\land}[6]},close
                ]
              ]
            ]
          ]
        ]
      ]
      [\sFmla{\False}{\formula{A} \lor \formula{C}},
        just={\TRule{\False}{\land}[2]}, checked
        [\sFmla{\False}{\formula{A}},
          just={\TRule{\False}{\lor}[3]}
          [\sFmla{\False}{\formula{C}},
            just={\TRule{\False}{\lor}[3]}
            [\sFmla{\True}{\formula{A}},
              just={\TRule{\True}{\lor}[1]},close
            ]
            [\sFmla{\True}{\formula{B} \land \formula{C}},
              just={\TRule{\True}{\lor}[1]}, checked
              [\sFmla{\True}{\formula{B}},
                just={\TRule{\True}{\land}[6]}
                [\sFmla{\True}{\formula{C}},
                  just={\TRule{\True}{\land}[6]},close
                ]
              ]
            ]
          ]
        ]
      ]
    ]
  ]
\end{oltableau}
\end{ex}
\label{olpseg:OLP-0103-B014:end}

\phantomsection\label{olpseg:OLP-0103-B015:start}
\begin{prob}
د لاندې دپاره تړلي تابلوګانې ورکړئ:
\begin{enumerate}
\item \LR{$\sFmla{\True}{!A \land (!B \land !C)}, \sFmla{\False}{(!A \land !B) \land !C}$}.
\item \LR{$\sFmla{\True}{!A \lor (!B \lor !C)}, \sFmla{\False}{(!A \lor !B) \lor !C}$}.
\item \LR{$\sFmla{\True}{!A \lif (!B \lif !C)}, \sFmla{\False}{!B \lif (!A \lif !C)}$}.
\item \LR{$\sFmla{\True}{!A}, \sFmla{\False}{\lnot\lnot !A}$}.
\end{enumerate}
\end{prob}
\label{olpseg:OLP-0103-B015:end}

\phantomsection\label{olpseg:OLP-0103-B016:start}
\begin{prob}
د لاندې دپاره تړلي تابلوګانې ورکړئ:
\begin{enumerate}
\item \LR{$\sFmla{\True}{(!A \lor !B) \lif !C}, \sFmla{\False}{!A \lif !C}$}.
\item \LR{$\sFmla{\True}{(!A \lif !C) \land (!B \lif !C)}, \sFmla{\False}{(!A \lor !B) \lif !C}$}.
\item \LR{$\sFmla{\False}{\lnot(!A \land \lnot !A)}$}.
\item \LR{$\sFmla{\True}{!B \lif !A}, \sFmla{\False}{\lnot !A \lif \lnot !B}$}.
\item \LR{$\sFmla{\False}{(!A \lif \lnot !A) \lif \lnot !A}$}.
\item \LR{$\sFmla{\False}{\lnot(!A \lif !B) \lif \lnot !B}$}.
\item \LR{$\sFmla{\True}{!A \lif !C}, \sFmla{\False}{\lnot (!A \land \lnot !C)}$}.
\item \LR{$\sFmla{\True}{!A \land \lnot !C}, \sFmla{\False}{\lnot (!A \lif !C)}$}.
\item \LR{$\sFmla{\True}{!A \lor !B}, \sFmla{\True}{\lnot !B}, \sFmla{\False}{!A}$}.
\textbf{د ژباړې د نسخې سمون (OLTAB-002):} سرچينې دوه رښتوني فرضونه
د يو نښه لرونکي فارمول په دننه کښې سره يو ځاے کړي وو؛ دلته هر فرض
خپله رښتيا نښه لري، لکه د مسئلې په اړوند استلزام کښې۔
\item \LR{$\sFmla{\True}{\lnot !A \lor \lnot !B}, \sFmla{\False}{\lnot(!A \land !B)}$}.
\item \LR{$\sFmla{\False}{(\lnot !A \land \lnot !B) \lif\lnot(!A \lor !B)}$}.
\item \LR{$\sFmla{\False}{\lnot(!A \lor !B) \lif (\lnot !A \land \lnot !B)}$}.
\end{enumerate}
\end{prob}
\label{olpseg:OLP-0103-B016:end}

\phantomsection\label{olpseg:OLP-0103-B017:start}
\begin{prob}
د لاندې دپاره تړلي تابلوګانې ورکړئ:
\begin{enumerate}
\item \LR{$\sFmla{\True}{\lnot(!A \lif !B)}, \sFmla{\False}{!A}$}.
\item \LR{$\sFmla{\True}{\lnot(!A \land !B)}, \sFmla{\False}{\lnot !A \lor \lnot !B}$}.
\item \LR{$\sFmla{\True}{!A \lif !B}, \sFmla{\False}{\lnot !A \lor !B}$}.
\item \LR{$\sFmla{\False}{\lnot \lnot !A \lif !A}$}.
\item \LR{$\sFmla{\True}{!A \lif !B}, \sFmla{\True}{\lnot !A \lif !B}, \sFmla{\False}{!B}$}.
\item \LR{$\sFmla{\True}{(!A \land !B) \lif !C}, \sFmla{\False}{(!A \lif !C) \lor (!B \lif !C)}$}.
\item \LR{$\sFmla{\True}{(!A \lif !B) \lif !A}, \sFmla{\False}{!A}$}.
\item \LR{$\sFmla{\False}{(!A \lif !B) \lor (!B \lif !C)}$}.
\end{enumerate}
\end{prob}
\label{olpseg:OLP-0103-B017:end}

\phantomsection\label{olpseg:OLP-0104-B004:start}
\olfileid{fol}{tab}{prq}
\label{olpseg:OLP-0104-B004:end}

\phantomsection\label{olpseg:OLP-0104-B005:start}
\olsection{له سورونو سره تابلوګانې}
\label{olpseg:OLP-0104-B005:end}

\phantomsection\label{olpseg:OLP-0104-B006:start}
\begin{ex}
له سورونو سره د کار پر مهال بايد ډاډ واخلو چې د ځانګړي متغير
شرط نۀ ماتوو، او کله کله د دې دپاره د ځينو استنتاجونو د پلي کولو
ترتيب بدلول را بدلول غواړي۔ په عام ډول ګټه کوي چې لومړے د هغو
قاعدو د حل هڅه وکړو چې د ځانګړي متغير شرط ورباندې لګېږي (په بشپړ
تابلو کښې به دوئ پاس راځي)۔
\label{olpseg:OLP-0104-B006:end}

\phantomsection\label{olpseg:OLP-0104-B007:start}
راځئ وګورو چې د~تړلے فارمول
\LR{$\lexists[x][\lnot !A(x)] \lif \lnot \lforall[x][!A(x)]$}
دپاره تابلو څنګه ورکوو۔ د معمول په شان د فرض له ثبتولو
پيل کوو:
\begin{oltableau}
[\sFmla{\False}{\lexists[x][\lnot \formula{A}(x)]\lif \lnot
    \lforall[x][\formula{A}(x)]}, just=\TAss]
\end{oltableau}
ځکه اصلي عملګر~\LR{$\lif$} دے، نو
\LR{$\TRule{\False}{\lif}$} پلې کوو:
\begin{oltableau}
[\sFmla{\False}{\lexists[x][\lnot \formula{A}(x)]\lif \lnot
    \lforall[x][\formula{A}(x)]}, just=\TAss, checked
  [\sFmla{\True}{\lexists[x][\lnot \formula{A}(x)]},
    just={\TRule{\False}{\lif}[1]}
    [\sFmla{\False}{\lnot\lforall[x][\formula{A}(x)]},
      just={\TRule{\False}{\lif}[1]}]
  ]
]
\end{oltableau}
راتلونکې د کار کرښه~\LR{$2$} ده۔ \LR{$\TRule{\True}{\lexists}$} کاروو۔
دا يو نوے ثابت سمبول غواړي؛ ځکه لا هېڅ ثابت سمبولونه نۀ دي
راغلي، هر يو ټاکلے شو، مثلاً~\LR{$a$}۔
\begin{oltableau}
[\sFmla{\False}{\lexists[x][\lnot \formula{A}(x)]\lif \lnot
    \lforall[x][\formula{A}(x)]}, just=\TAss, checked
  [\sFmla{\True}{\lexists[x][\lnot \formula{A}(x)]},
    just={\TRule{\False}{\lif}[1]}, checked
    [\sFmla{\False}{\lnot\lforall[x][\formula{A}(x)]},
      just={\TRule{\False}{\lif}[1]}
      [\sFmla{\True}{\lnot\formula{A}(a)}, just={\TRule{\True}{\lexists}[2]}]
    ]
  ]
]
\end{oltableau}
اوس پر~\LR{$3$} کرښه \LR{$\TRule{\False}{\lnot}$} پلې کوو:
\begin{oltableau}
[\sFmla{\False}{\lexists[x][\lnot \formula{A}(x)]\lif \lnot
    \lforall[x][\formula{A}(x)]}, just=\TAss, checked
  [\sFmla{\True}{\lexists[x][\lnot \formula{A}(x)]},
    just={\TRule{\False}{\lif}[1]}, checked
    [\sFmla{\False}{\lnot\lforall[x][\formula{A}(x)]},
      just={\TRule{\False}{\lif}[1]}, checked
      [\sFmla{\True}{\lnot\formula{A}(a)}, just={\TRule{\True}{\lexists}[2]}
        [\sFmla{\True}{\lforall[x][\formula{A}(x)]},
          just = {\TRule{\False}{\lnot}[3]}
        ]
      ]
    ]
  ]
]
\end{oltableau}
پر~\LR{$4$} کرښه د~\LR{$\TRule{\True}{\lnot}$} او ورپسې پر~\LR{$5$} کرښه د
\LR{$\TRule{\True}{\lforall}$} په پلي کولو تړلے تابلو اخلو۔
\begin{oltableau}
[\sFmla{\False}{\lexists[x][\lnot \formula{A}(x)]\lif \lnot
    \lforall[x][\formula{A}(x)]}, just=\TAss, checked
  [\sFmla{\True}{\lexists[x][\lnot \formula{A}(x)]},
    just={\TRule{\False}{\lif}[1]}, checked
    [\sFmla{\False}{\lnot\lforall[x][\formula{A}(x)]},
      just={\TRule{\False}{\lif}[1]}, checked
      [\sFmla{\True}{\lnot\formula{A}(a)}, just={\TRule{\True}{\lexists}[2]}
        [\sFmla{\True}{\lforall[x][\formula{A}(x)]},
          just = {\TRule{\False}{\lnot}[3]}
          [\sFmla{\False}{\formula{A}(a)}, just={\TRule{\True}{\lnot}[4]}
            [\sFmla{\True}{\formula{A}(a)},
              just={\TRule{\True}{\lforall}[5]}, close
            ]
          ]
        ]
      ]
    ]
  ]
]
\end{oltableau}
\end{ex}
\label{olpseg:OLP-0104-B007:end}

\phantomsection\label{olpseg:OLP-0104-B008:start}
\begin{ex}
راځئ وګورو چې د لاندې سټ دپاره تابلو څنګه ورکوو:
\[
\sFmla{\False}{\lexists[x][!C(x,b)]},
\sFmla{\True}{\lexists[x][(!A(x) \land !B(x))]},
\sFmla{\True}{\lforall[x][(!B(x) \lif !C(x,b))]}.
\]
د معمول په شان له فرضونو پيل کوو:
\begin{oltableau}
  [\sFmla{\False}{\lexists[x][\formula{C}(x,b)]}, just=\TAss
    [\sFmla{\True}{\lexists[x][(\formula{A}(x) \land \formula{B}(x))]},
      just=\TAss
      [\sFmla{\True}{\lforall[x][(\formula{B}(x) \lif \formula{C}(x,b))]},
        just=\TAss]
    ]
  ]
\end{oltableau}
تل بايد لومړے هغه قاعده پلې کړو چې د ځانګړي متغير شرط لري؛ په
دې حالت کښې هغه پر~\LR{$2$} کرښه \LR{$\TRule{\True}{\lexists}$} ده۔ ځکه
فرضونه ثابت سمبول~\LR{$b$} لري، بايد بل يو وکاروو؛ بيا~\LR{$a$} اخلو۔
\begin{oltableau}
  [\sFmla{\False}{\lexists[x][\formula{C}(x,b)]}, just=\TAss
    [\sFmla{\True}{\lexists[x][(\formula{A}(x) \land \formula{B}(x))]},
      just=\TAss, checked
      [\sFmla{\True}{\lforall[x][(\formula{B}(x) \lif \formula{C}(x,b))]},
        just=\TAss
        [\sFmla{\True}{\formula{A}(a) \land \formula{B}(a)},
          just={\TRule{\True}{\lexists}[2]}]
      ]
    ]
  ]
\end{oltableau}
که اوس پر~\LR{$1$} کرښه \LR{$\TRule{\False}{\lexists}$} يا پر~\LR{$3$} کرښه
\LR{$\TRule{\True}{\lforall}$} پلې کړو، بايد وټاکو چې د~\LR{$x$} پر ځاے
کومه اصطلاح~\LR{$t$} کښېږدو۔ ځکه دا قاعدې د ځانګړي متغير شرط نۀ لري،
هره خوښه اصطلاح ټاکلے شو۔ په ځينو حالتونو کښې ښايي قاعده څو ځله
له بېلابېلو~\LR{$t$}s سره پلې کول هم لازم شي۔ خو د عمومي قاعدې په توګه
ګټه کوي چې له هغو اصطلاحګانو څخه يوه واخلو چې له مخکښې په
تابلو کښې راغلې وي---دلته \LR{$a$} او~\LR{$b$}---او په دې حالت کښې
اټکل کولے شو چې له~\LR{$a$} سره د تړلې څانګې لاس ته راتلل ډېر ممکن دي۔
\begin{oltableau}
  [\sFmla{\False}{\lexists[x][\formula{C}(x,b)]}, just=\TAss
    [\sFmla{\True}{\lexists[x][(\formula{A}(x) \land \formula{B}(x))]},
      just=\TAss, checked
      [\sFmla{\True}{\lforall[x][(\formula{B}(x) \lif \formula{C}(x,b))]}, just=\TAss
        [\sFmla{\True}{\formula{A}(a) \land \formula{B}(a)}, just={\TRule{\True}{\lexists}[2]}
          [\sFmla{\False}{\formula{C}(a,b)}, just={\TRule{\False}{\lexists}[1]}
            [\sFmla{\True}{\formula{B}(a) \lif \formula{C}(a,b)},
              just={\TRule{\True}{\lforall}[3]}
            ]
          ]
        ]
      ]
    ]
  ]
\end{oltableau}
پر~\LR{$1$} او~\LR{$3$} کرښو نښه لرونکي فارمولونه ته د سم نښان نۀ ورکوو؛
ځکه ښايي بيا ئې وکاروو۔ اوس پر~\LR{$4$} کرښه
\LR{$\TRule{\True}{\land}$} پلې کړئ:
\begin{oltableau}
  [\sFmla{\False}{\lexists[x][\formula{C}(x,b)]}, just=\TAss
    [\sFmla{\True}{\lexists[x][(\formula{A}(x) \land \formula{B}(x))]},
      just=\TAss, checked
      [\sFmla{\True}{\lforall[x][(\formula{B}(x) \lif \formula{C}(x,b))]}, just=\TAss
        [\sFmla{\True}{\formula{A}(a) \land \formula{B}(a)},
          just={\TRule{\True}{\lexists}[2]}, checked
          [\sFmla{\False}{\formula{C}(a,b)}, just={\TRule{\False}{\lexists}[1]}
            [\sFmla{\True}{\formula{B}(a) \lif \formula{C}(a,b)},
              just={\TRule{\True}{\lforall}[3]}
              [\sFmla{\True}{\formula{A}(a)}, just={\TRule{\True}{\land}[4]}
                [\sFmla{\True}{\formula{B}(a)}, just={\TRule{\True}{\land}[4]}
                ]
              ]
            ]
          ]
        ]
      ]
    ]
  ]
\end{oltableau}
که اوس پر~\LR{$6$} کرښه \LR{$\TRule{\True}{\lif}$} پلې کړو،
تابلو تړل کېږي:
\begin{oltableau}
  [\sFmla{\False}{\lexists[x][\formula{C}(x,b)]}, just=\TAss
    [\sFmla{\True}{\lexists[x][(\formula{A}(x) \land \formula{B}(x))]},
      just=\TAss, checked
      [\sFmla{\True}{\lforall[x][(\formula{B}(x) \lif \formula{C}(x,b))]},
        just=\TAss
        [\sFmla{\True}{\formula{A}(a) \land \formula{B}(a)},
          just={\TRule{\True}{\lexists}[2]}, checked
          [\sFmla{\False}{\formula{C}(a,b)},
            just={\TRule{\False}{\lexists}[1]}
            [\sFmla{\True}{\formula{B}(a) \lif \formula{C}(a,b)},
              just={\TRule{\True}{\lforall}[3]}, checked
              [\sFmla{\True}{\formula{A}(a)},
                just={\TRule{\True}{\land}[4]}
                [\sFmla{\True}{\formula{B}(a)},
                  just={\TRule{\True}{\land}[4]}
                  [\sFmla{\False}{\formula{B}(a)},
                    just={\TRule{\True}{\lif}[6]},close
                  ]
                  [\sFmla{\True}{\formula{C}(a,b)},
                    just={\TRule{\True}{\lif}[6]},close
                  ]
                ]
              ]
            ]
          ]
        ]
      ]
    ]
  ]
\end{oltableau}
\label{olpseg:OLP-0104-B008:end}

\phantomsection\label{olpseg:OLP-0104-B009:start}
\end{ex}
\label{olpseg:OLP-0104-B009:end}

\phantomsection\label{olpseg:OLP-0104-B010:start}
\begin{ex}
د لاندې سټ دپاره تابلو جوړوو:
\[
\sFmla{\True}{\lforall[x][!A(x)]}, \sFmla{\True}{\lforall[x][!A(x)] 
  \lif \lexists[y][!B(y)]}, \sFmla{\True}{\lnot\lexists[y][!B(y)]}.
\]
د معمول په شان فرضونه ليکو:
\begin{oltableau}
  [\sFmla{\True}{\lforall[x][\formula{A}(x)]}, just=\TAss
    [\sFmla{\True}{\lforall[x][\formula{A}(x)] \lif
        \lexists[y][\formula{B}(y)]}, just=\TAss
      [\sFmla{\True}{\lnot\lexists[y][\formula{B}(y)]}, just=\TAss
      ]
    ]
  ]
\end{oltableau}
لومړے پر~\LR{$3$} کرښه د~\LR{$\TRule{\True}{\lnot}$} قاعده پلې کوو۔
د «تل لومړے د ځانګړي متغير شرط لرونکې قاعدې پلې کړئ» قاعدې يوه
پايله دا ده: «بې د ځانګړي متغير شرطه د سورونو قاعدې تر اړتيا پورې
وځنډوئ۔» همدارنګه هغه قاعدې هم وځنډوئ چې وېش رامنځته کوي۔
\begin{oltableau}
  [\sFmla{\True}{\lforall[x][\formula{A}(x)]}, just=\TAss
    [\sFmla{\True}{\lforall[x][\formula{A}(x)] \lif
        \lexists[y][\formula{B}(y)]}, just=\TAss
      [\sFmla{\True}{\lnot\lexists[y][\formula{B}(y)]}, just=\TAss, checked
        [\sFmla{\False}{\lexists[y][\formula{B}(y)]}, just={\TRule{\True}{\lnot}[3]}]
      ]
    ]
  ]
\end{oltableau}
نوې \LR{$4$}~کرښه \LR{$\TRule{\False}{\lexists}$} غواړي، چې د ځانګړي
متغير له شرط پرته د سور قاعده ده۔ نو دا ځنډوو او پر~\LR{$2$} کرښه
\LR{$\TRule{\True}{\lif}$} کاروو۔
\begin{oltableau}
  [\sFmla{\True}{\lforall[x][\formula{A}(x)]}, just=\TAss
    [\sFmla{\True}{\lforall[x][\formula{A}(x)] \lif
        \lexists[y][\formula{B}(y)]}, just=\TAss, checked
      [\sFmla{\True}{\lnot\lexists[y][\formula{B}(y)]}, just=\TAss, checked
        [\sFmla{\False}{\lexists[y][\formula{B}(y)]},
          just={\TRule{\True}{\lnot}[3]},
          [\sFmla{\False}{\lforall[x][\formula{A}(x)]}, just={\TRule{\True}{\lif}[2]}]
          [\sFmla{\True}{\lexists[y][\formula{B}(y)]}, just={\TRule{\True}{\lif}[2]}]
        ]
      ]
    ]
  ]
\end{oltableau}
دواړه نوي نښه لرونکي فارمولونه د ځانګړي متغير د شرط لرونکې
قاعدې غواړي، نو راتلونکے ګام بايد همدا وي:
\begin{oltableau}
  [\sFmla{\True}{\lforall[x][\formula{A}(x)]}, just=\TAss
    [\sFmla{\True}{\lforall[x][\formula{A}(x)] \lif
        \lexists[y][\formula{B}(y)]}, just=\TAss, checked
      [\sFmla{\True}{\lnot\lexists[y][\formula{B}(y)]}, just=\TAss, checked
        [\sFmla{\False}{\lexists[y][\formula{B}(y)]},
          just={\TRule{\True}{\lnot}[3]}
          [\sFmla{\False}{\lforall[x][\formula{A}(x)]}, just={\TRule{\True}{\lif}[2]},checked
            [\sFmla{\False}{\formula{A}(b)}, just={\TRule{\False}{\lforall}[5]}]
          ]
          [\sFmla{\True}{\lexists[y][\formula{B}(y)]}, just={\TRule{\True}{\lif}[2]},checked
            [\sFmla{\True}{\formula{B}(c)}, just={\TRule{\True}{\lexists}[5]}]
          ]
        ]
      ]
    ]
  ]
\end{oltableau}
د څانګو د تړلو دپاره بايد پر~\LR{$1$} او~\LR{$4$} کرښو
نښه لرونکي فارمولونه وکاروو۔ د هغو اړوندې قاعدې
(\TRule{\True}{\lforall} او~\TRule{\False}{\lexists}) د ځانګړي
متغير شرط نۀ لري، نو هرې مناسبې اصطلاحګانې ټاکلے شو۔ په دې حالت
کښې په خپل وار همدا \LR{$b$} او~\LR{$c$} دي۔
\textbf{د ژباړې د نسخې سمون (OLTAB-003):} سرچينه دلته د
دروغ-وجودي نښه لرونکي فارمول د کرښې شمېر درې ښيي، خو درېيمه کرښه
د نفي مقدمه ده؛ اړونده دروغ-وجودي قاعده او وروستۍ ونه دواړه څلورمه
کرښه کاروي، نو اشاره سمه شوه۔
\begin{oltableau}
  [\sFmla{\True}{\lforall[x][\formula{A}(x)]}, just=\TAss
    [\sFmla{\True}{\lforall[x][\formula{A}(x)] \lif
        \lexists[y][\formula{B}(y)]}, just=\TAss, checked
      [\sFmla{\True}{\lnot\lexists[y][\formula{B}(y)]}, just=\TAss, checked
        [\sFmla{\False}{\lexists[y][\formula{B}(y)]},
          just={\TRule{\True}{\lnot}[3]}
          [\sFmla{\False}{\lforall[x][\formula{A}(x)]},
            just={\TRule{\True}{\lif}[2]},checked
            [\sFmla{\False}{\formula{A}(b)},
              just={\TRule{\False}{\lforall}[5]}
              [\sFmla{\True}{\formula{A}(b)},
                just={\TRule{\True}{\lforall}[1]},close
              ]
            ]
          ]
          [\sFmla{\True}{\lexists[y][\formula{B}(y)]},
            just={\TRule{\True}{\lif}[2]},checked
            [\sFmla{\True}{\formula{B}(c)},
              just={\TRule{\True}{\lexists}[5]}
              [\sFmla{\False}{\formula{B}(c)},
                just={\TRule{\False}{\lexists}[4]},close
              ]
            ]
          ]
        ]
      ]
    ]
  ]
\end{oltableau}
\end{ex}
\label{olpseg:OLP-0104-B010:end}

\phantomsection\label{olpseg:OLP-0104-B011:start}
\begin{prob}
د لاندې دپاره تړلي تابلوګانې ورکړئ:
\begin{enumerate}
\item \LR{$\sFmla{\False}{(\lforall[x][!A(x)] \land \lforall[y][!B(y)])
\lif \lforall[z][(!A(z) \land !B(z))]}$}.
\item \LR{$\sFmla{\False}{(\lexists[x][!A(x)] \lor \lexists[y][!B(y)])
\lif \lexists[z][(!A(z) \lor !B(z))]}$}.
\item \LR{$\sFmla{\True}{\lforall[x][(!A(x) \lif !B)]},
\sFmla{\False}{\lexists[y][!A(y)] \lif !B}$}.
\item \LR{$\sFmla{\True}{\lforall[x][\lnot !A(x)]},
\sFmla{\False}{\lnot\lexists[x][!A(x)]}$}.
\item \LR{$\sFmla{\False}{\lnot\lexists[x][!A(x)] \lif \lforall[x][\lnot
!A(x)]}$}.
\item \LR{$\sFmla{\False}{\lnot\lexists[x][\lforall[y][((!A(x,y) \lif
\lnot !A(y,y)) \land (\lnot !A(y,y) \lif !A(x,y)))]]}$}.
\end{enumerate}
\end{prob}
\label{olpseg:OLP-0104-B011:end}

\phantomsection\label{olpseg:OLP-0104-B012:start}
\begin{prob}
د لاندې دپاره تړلي تابلوګانې ورکړئ:
\begin{enumerate}
\item \LR{$\sFmla{\False}{\lnot\lforall[x][!A(x)] \lif \lexists[x][\lnot!A(x)]}$}.
\item \LR{$\sFmla{\True}{(\lforall[x][!A(x)] \lif !B)}, \sFmla{\False}{\lexists[y][(!A(y) \lif !B)]}$}.
\item \LR{$\sFmla{\False}{\lexists[x][(!A(x) \lif \lforall[y][!A(y)])]}$}.
\end{enumerate}
\end{prob}
\label{olpseg:OLP-0104-B012:end}

\phantomsection\label{olpseg:OLP-0105-B004:start}
\olfileid{fol}{tab}{ptn}
\label{olpseg:OLP-0105-B004:end}

\phantomsection\label{olpseg:OLP-0105-B005:start}
\olsection{ثبوت-تيوريکي مفاهيم}
\label{olpseg:OLP-0105-B005:end}

\phantomsection\label{olpseg:OLP-0105-B006:start}
\begin{editorial}
دا برخه د تابلوګانو د اثبات‌پذيرۍ اړيکې او سازګارۍ تعريفونه
راټولوي۔
\end{editorial}
\label{olpseg:OLP-0105-B006:end}

\phantomsection\label{olpseg:OLP-0105-B007:start}
\begin{explain}
لکه څنګه چې مو يو شمېر مهم معنايي مفاهيم (اعتبار، استلزام او د
صدق وړتيا) تعريف کړل، اوس ورته \emph{ثبوت-تيوريکي مفاهيم}
تعريفوو۔ دا مفاهيم په جوړښتونه کښې د تړلي فارمولونه د پوره
کېدو له مخې نۀ، بلکې د ځينو تړلو تابلوګانې د موجوديت له مخې
تعريفېږي۔ دا مهمه موندنه وه چې دغه مفاهيم سره سمون خوري۔ د دې
سمون منځپانګه د \emph{سموالي} او \emph{بشپړتيا قضيې} دي۔
\end{explain}
\label{olpseg:OLP-0105-B007:end}

\phantomsection\label{olpseg:OLP-0105-B008:start}
\begin{defn}[قضيې]
يوه تړلے فارمول~\LR{$!A$} هله \emph{قضيه} ده چې د
\LR{$\sFmla{\False}{!A}$} دپاره تړلے تابلو موجود وي۔ که \LR{$!A$}
قضيه وي، \LR{$\Proves !A$} ليکو، او که نۀ وي \LR{$\Proves/ !A$} ليکو۔
\end{defn}
\label{olpseg:OLP-0105-B008:end}

\phantomsection\label{olpseg:OLP-0105-B009:start}
\begin{defn}[د اشتقاق وړتيا]
تړلے فارمول~\LR{$!A$} د تړلي فارمولونه له سټ~\LR{$\Gamma$} څخه،
\LR{$\Gamma \Proves !A$}، هله \emph{د اشتقاق وړ} ده چې داسې متناهي
سټ~\LR{$\{!B_1, \dots, !B_n\} \subseteq \Gamma$} او د لاندې سټ دپاره
تړلے تابلو موجود وي:
\[
\{
  \sFmla{\False}{!A},
  \sFmla{\True}{!B_1}, \dots,
  \sFmla{\True}{!B_n}
  \}.
\]
که \LR{$!A$} له~\LR{$\Gamma$} څخه د اشتقاق وړ نۀ وي،
\LR{$\Gamma \Proves/ !A$} ليکو۔
\end{defn}
\label{olpseg:OLP-0105-B009:end}

\phantomsection\label{olpseg:OLP-0105-B010:start}
\begin{defn}[سازګاري]
د تړلي فارمولونه يو سټ~\LR{$\Gamma$} هله او يوازې هله \emph{ناسازګار}
دے چې داسې متناهي سټ~\LR{$\{!B_1, \dots, !B_n\} \subseteq \Gamma$}
او د لاندې سټ دپاره تړلے تابلو موجود وي:
\[
\{
  \sFmla{\True}{!B_1}, \dots,
  \sFmla{\True}{!B_n}  
  \}.
\]
که \LR{$\Gamma$} ناسازګار نۀ وي، \emph{سازګار} ئې بولو۔
\end{defn}
\label{olpseg:OLP-0105-B010:end}

\phantomsection\label{olpseg:OLP-0105-B011:start}
\begin{prop}[انعکاسيت]
\ollabel{prop:reflexivity}
که \LR{$!A \in \Gamma$} وي، نو \LR{$\Gamma \Proves !A$}۔
\end{prop}
\label{olpseg:OLP-0105-B011:end}

\phantomsection\label{olpseg:OLP-0105-B012:start}
\begin{proof}
که \LR{$!A \in \Gamma$} وي، نو \LR{$\{!A\}$} د~\LR{$\Gamma$} يو متناهي فرعي سټ
دے، او لاندې تابلو تړلے دے:
\begin{oltableau}
  [\sFmla{\False}{\formula{A}}, just = \TAss
    [\sFmla{\True}{\formula{A}}, just = \TAss,close]
  ]
\end{oltableau}
\end{proof}
\label{olpseg:OLP-0105-B012:end}

\phantomsection\label{olpseg:OLP-0105-B013:start}
\begin{prop}[يکنواختي]
\ollabel{prop:monotonicity}
که \LR{$\Gamma \subseteq \Delta$} او \LR{$\Gamma \Proves !A$} وي، نو
\LR{$\Delta \Proves !A$}۔
\end{prop}
\label{olpseg:OLP-0105-B013:end}

\phantomsection\label{olpseg:OLP-0105-B014:start}
\begin{proof}
د~\LR{$\Gamma$} هر متناهي فرعي سټ د~\LR{$\Delta$} هم متناهي فرعي سټ دے۔
\end{proof}
\label{olpseg:OLP-0105-B014:end}

\phantomsection\label{olpseg:OLP-0105-B015:start}
\begin{prop}[انتقاليت]
\ollabel{prop:transitivity}
که \LR{$\Gamma \Proves !A$} او
\LR{$\{!A\} \cup \Delta \Proves !B$} وي، نو
\LR{$\Gamma \cup \Delta \Proves !B$}۔
\end{prop}
\label{olpseg:OLP-0105-B015:end}

\phantomsection\label{olpseg:OLP-0105-B016:start}
\begin{proof}
که \LR{$\{!A\} \cup \Delta \Proves !B$} وي، نو داسې متناهي فرعي سټ
\LR{$\Delta_0 = \{!C_1, \dots, !C_n\} \subseteq \Delta$} شته چې
\begin{align*}
\{\sFmla{\False}{!B}, & \sFmla{\True}{!A}, \sFmla{\True}{!C_1},
\dots, \sFmla{\True}{!C_n}\}
\intertext{\RL{تړلے تابلو لري۔ که \LR{$\Gamma \Proves !A$} وي، نو
  داسې \LR{$!D_1$}, \dots, \LR{$!D_m \in \Gamma$} شته چې}}
\{\sFmla{\False}{!A}, & \sFmla{\True}{!D_1},
\dots, \sFmla{\True}{!D_m}\}
\end{align*}
تړلے تابلو لري۔
\textbf{د ژباړې د نسخې سمون (OLTAB-004):} سرچينه په منځنۍ جمله
کښې د ډي-لړ يوازې وروستۍ جمله د ګاما فرعي سټ بولي؛ دلته د لړ
ټولې جملې د ګاما غړي ښودل شوې، لکه د انتقاليت تعريف چې غواړي۔
\label{olpseg:OLP-0105-B016:end}

\phantomsection\label{olpseg:OLP-0105-B017:start}
اوس د لاندې فرضونو تابلو ونيسئ:
\[
\sFmla{\False}{!B},
\sFmla{\True}{!C_1}, \dots, \sFmla{\True}{!C_n},
\sFmla{\True}{!D_1}, \dots, \sFmla{\True}{!D_m}.
\]
پر~\LR{$!A$} د~\Cut{} قاعده پلې کړئ۔ دا دوه څانګې جوړوي: يوه پکې
\LR{$\sFmla{\True}{!A}$}، او بله پکې \LR{$\sFmla{\False}{!A}$} لري۔ نو په
يوه څانګه کښې دا ټول موجود دي:
\[
\{\sFmla{\False}{!B}, \sFmla{\True}{!A},
\sFmla{\True}{!C_1}, \dots, \sFmla{\True}{!C_n}\}
\]
ځکه د دغو فرضونو دپاره تړلے تابلو شته، هغه له دې څانګې سره
نښلولے شو؛ هره هغه څانګه تړل کېږي چې له
\LR{$\sFmla{\True}{!A}$} څخه تېرېږي۔ په بله څانګه کښې دا ټول موجود دي:
\[
\{\sFmla{\False}{!A}, \sFmla{\True}{!D_1}, \dots,
\sFmla{\True}{!D_m}\}
\]
نو بل لورے هم بشپړولے او تړلے تابلو اخستلے شو۔ دا
\LR{$\Gamma \cup \Delta \Proves !B$} ښيي۔
\end{proof}
\label{olpseg:OLP-0105-B017:end}

\phantomsection\label{olpseg:OLP-0105-B018:start}
پام وکړئ د دې يوه ځانګړې پايله دا ده: که \LR{$\Gamma \Proves !A$} او
\LR{$!A \Proves !B$} وي، نو \LR{$\Gamma \Proves !B$}۔ له دې څخه دا هم
لازمېږي چې که \LR{$!A_1,\dots,!A_n \Proves !B$} او د هر~\LR{$i$} دپاره
\LR{$\Gamma \Proves !A_i$} وي، نو \LR{$\Gamma \Proves !B$}۔
\label{olpseg:OLP-0105-B018:end}

\phantomsection\label{olpseg:OLP-0105-B019:start}
\begin{prop}
\ollabel{prop:incons}
\LR{$\Gamma$} هله او يوازې هله ناسازګار دے چې د هرې
تړلے فارمول~\LR{$!A$} دپاره \LR{$\Gamma \Proves !A$} وي۔
\end{prop}
\label{olpseg:OLP-0105-B019:end}

\phantomsection\label{olpseg:OLP-0105-B020:start}
\begin{proof}
تمرين۔
\end{proof}
\label{olpseg:OLP-0105-B020:end}

\phantomsection\label{olpseg:OLP-0105-B021:start}
\begin{prob}
\olref[fol][tab][ptn]{prop:incons} ثابت کړئ۔
\end{prob}
\label{olpseg:OLP-0105-B021:end}



\phantomsection\label{olpseg:OLP-0105-B023:start}
\begin{prop}[تراکم]
  \ollabel{prop:proves-compact}
  \begin{enumerate}
  \item که \LR{$\Gamma \Proves !A$} وي، نو داسې متناهي فرعي سټ
    \LR{$\Gamma_0 \subseteq \Gamma$} شته چې \LR{$\Gamma_0 \Proves !A$} وي۔
  \item که د~\LR{$\Gamma$} هر متناهي فرعي سټ سازګار وي، نو \LR{$\Gamma$}
    سازګار دے۔
  \end{enumerate}
\end{prop}
\label{olpseg:OLP-0105-B023:end}

\phantomsection\label{olpseg:OLP-0105-B024:start}
\begin{proof}
  \begin{enumerate}
    \item که \LR{$\Gamma \Proves !A$} وي، نو داسې متناهي فرعي سټ
      \LR{$\Gamma_0 = \{!B_1, \dots, !B_n\}$} او د لاندې دپاره تړلے
      تابلو شته:
      \[
      \{\sFmla{\False}{!A}, \sFmla{\True}{!B_1}, \dots, \sFmla{\True}{!B_n}\}
      \]
      همدا تابلو، \LR{$\Gamma_0 \Proves !A$} هم ښيي۔
    \item که \LR{$\Gamma$} ناسازګار وي، نو د کوم متناهي فرعي سټ
      \LR{$\Gamma_0 = \{!B_1, \dots, !B_n\}$} دپاره د لاندې يو تړلے
      تابلو شته:
      \[
      \{\sFmla{\True}{!B_1}, \dots, \sFmla{\True}{!B_n}\}
      \]
      دا تړلے تابلو ښيي چې \LR{$\Gamma_0$} ناسازګار دے۔
  \end{enumerate}
\end{proof}
\label{olpseg:OLP-0105-B024:end}

\phantomsection\label{olpseg:OLP-0106-B004:start}
\olfileid{fol}{tab}{prv}
\label{olpseg:OLP-0106-B004:end}

\phantomsection\label{olpseg:OLP-0106-B005:start}
\olsection{د اشتقاق وړتيا او سازګاري}
\label{olpseg:OLP-0106-B005:end}

\phantomsection\label{olpseg:OLP-0106-B006:start}
اوس د~د اشتقاق وړتيا اړيکې يو شمېر خاصيتونه ثابتوو۔ دوئ په
خپلواک ډول هم په زړه پورې دي، خو هر يو به د بشپړتيا قضیې په ثبوت
کښې ونډه ولري۔
\label{olpseg:OLP-0106-B006:end}

\phantomsection\label{olpseg:OLP-0106-B007:start}
\begin{prop}\ollabel{prop:provability-contr}
  که \LR{$\Gamma \Proves !A$} او \LR{$\Gamma \cup \{!A\}$} ناسازګار وي،
  نو \LR{$\Gamma$} ناسازګار دے۔
\end{prop}
\label{olpseg:OLP-0106-B007:end}

\phantomsection\label{olpseg:OLP-0106-B008:start}
\begin{proof}
داسې متناهي
\LR{$\Gamma_0 = \{!B_1, \dots, !B_n\}$} او
\LR{$\Gamma_1 =\{!C_1, \dots, !C_m\} \subseteq \Gamma$} شته چې
  \begin{align*}
    \{\sFmla{\False}{!A}, &
    \sFmla{\True}{!B_1}, \dots, \sFmla{\True}{!B_n}\} \\
    \{\sFmla{\True}{!A}, &
    \sFmla{\True}{!C_1}, \dots, \sFmla{\True}{!C_m}\}
  \end{align*}
  تړلي تابلوګانې لري۔ پر~\LR{$!A$} د~\Cut{} قاعدې په کارولو دا په
  يوه تړلي تابلو کښې يو ځاے کولے شو چې د
  \LR{$\Gamma_0 \cup \Gamma_1$} ناسازګاري ښيي۔ ځکه
  \LR{$\Gamma_0 \subseteq \Gamma$} او~\LR{$\Gamma_1 \subseteq \Gamma$}،
  نو \LR{$\Gamma_0 \cup \Gamma_1 \subseteq \Gamma$}؛ له دې امله
  \LR{$\Gamma$} ناسازګار دے۔
\textbf{د ژباړې د نسخې سمون (OLTAB-005):} سرچينه د دويم متناهي
سټ په تعريف کښې لړ په اېن شاخص پاے ته رسوي، خو په تابلو کښې ئې
وروستے غړے د اېم شاخص لري؛ دلته د لړ وروستے شاخص يو شان شوے دے۔
\end{proof}
\label{olpseg:OLP-0106-B008:end}

\phantomsection\label{olpseg:OLP-0106-B009:start}
\begin{prop}
\ollabel{prop:prov-incons}
\LR{$\Gamma \Proves !A$} هله او يوازې هله چې
\LR{$\Gamma \cup \{\lnot !A\}$} ناسازګار وي۔
\end{prop}
\label{olpseg:OLP-0106-B009:end}

\phantomsection\label{olpseg:OLP-0106-B010:start}
\begin{proof}
لومړے فرض کړئ \LR{$\Gamma \Proves !A$}، يعنې د لاندې دپاره تړلے
تابلو شته:
\[
\{\sFmla{\False}{!A},
\sFmla{\True}{!B_1}, \dots, \sFmla{\True}{!B_n}\}
\]
د~\LR{$\TRule{\True}{\lnot}$} قاعدې په کارولو دا د لاندې دپاره په تړلي
تابلو اړولے شو:
\[
\{\sFmla{\True}{\lnot !A},
\sFmla{\True}{!B_1}, \dots, \sFmla{\True}{!B_n}\}.
\]
\label{olpseg:OLP-0106-B010:end}

\phantomsection\label{olpseg:OLP-0106-B011:start}
بل لوري ته، که د وروستي سټ دپاره تړلے تابلو وي، د لومړي سټ
په تړلي تابلو ئې اړولے شو: هغه هر فارمول ړنګوو چې پر لومړي
فرض~\LR{$\sFmla{\True}{\lnot !A}$} د~\TRule{\True}{\lnot} له کارونې
راوتلے وي، خپله دغه فرض هم ړنګوو، او فرض
\LR{$\sFmla{\False}{!A}$} ورزياتوو۔ ځکه که کومه څانګه مخکښې د
\TRule{\True}{\lnot} د هغې پايلې له امله تړلې وه چې پر
\LR{$\sFmla{\True}{\lnot !A}$} پلې شوې، يعنې
\LR{$\sFmla{\False}{!A}$}، نو په نوي تابلو کښې اړونده څانګه هم
تړلې ده۔ که په زوړ تابلو کښې کومه څانګه د فرض
\LR{$\sFmla{\True}{\lnot !A}$} او~\LR{$\sFmla{\False}{\lnot !A}$} دواړو له
امله تړلې وه، پر~\LR{$\sFmla{\False}{\lnot !A}$} د
\LR{$\TRule{\False}{\lnot}$} په پلي کولو او د
\LR{$\sFmla{\True}{!A}$} په اخستلو ئې په تړلې څانګه اړولے شو۔ دا څانګه
ځکه تړل کېږي چې \LR{$\sFmla{\False}{!A}$} مو د فرض په توګه ورزيات کړے
دے۔
\end{proof}
\label{olpseg:OLP-0106-B011:end}

\phantomsection\label{olpseg:OLP-0106-B012:start}
\begin{prob}
ثابت کړئ چې \LR{$\Gamma \Proves \lnot !A$} هله او يوازې هله چې
\LR{$\Gamma \cup \{!A\}$} ناسازګار وي۔
\end{prob}
\label{olpseg:OLP-0106-B012:end}

\phantomsection\label{olpseg:OLP-0106-B013:start}
\begin{prop}\ollabel{prop:explicit-inc}
  که \LR{$\Gamma \Proves !A$} او~\LR{$\lnot !A \in \Gamma$} وي، نو
  \LR{$\Gamma$} ناسازګار دے۔
\end{prop}
\label{olpseg:OLP-0106-B013:end}

\phantomsection\label{olpseg:OLP-0106-B014:start}
\begin{proof}
  فرض کړئ \LR{$\Gamma \Proves !A$} او~\LR{$\lnot !A \in \Gamma$}۔ بيا داسې
  \LR{$!B_1$}، \dots، \LR{$!B_n \in \Gamma$} شته چې
  \[
  \{\sFmla{\False}{!A}, \sFmla{\True}{!B_1}, \dots, \sFmla{\True}{!B_n}\}
  \]
  تړلے تابلو لري۔ فرض~\sFmla{\False}{!A} په
  \sFmla{\True}{\lnot !A} بدل کړئ، او د فرضونو وروسته هغه
  \sFmla{\False}{!A} پايله ورننباسئ چې پر
  \sFmla{\True}{\lnot !A} د~\TRule{\True}{\lnot} له کارونې
  راځي۔ په زوړ تابلو کښې هره تړلے فارمول چې د~\LR{$1$} کرښې په
  حواله توجيه شوې وه، اوس د~\LR{$n+1$} کرښې په حواله توجيه کېږي۔ نو که
  زوړ تابلو تړلے و، نوے تابلو هم تړلے دے۔ دا د~\LR{$\Gamma$} ناسازګاري
  ښيي؛ ځکه ټول فرضونه په~\LR{$\Gamma$} کښې دي۔
  \textbf{د ژباړې د نسخې سمون (OLTAB-006):} سرچينه د «داسې چې»
  وروسته يو بې دليله تېک شاتګی ږدي؛ دلته لرې شوے دے۔
  \textbf{د ژباړې د نسخې سمون (OLTAB-007):} سرچينه وايي د رښتونې
  نفي قاعده پر دروغې جملې پلې کېږي؛ قاعده په حقيقت کښې پر رښتونې
  نفي شوې جملې پلې کېږي او دروغ مثبت فارمول راوباسي، لکه د همدې
  ثبوت په راتلونکې جمله کښې۔
\end{proof}
\label{olpseg:OLP-0106-B014:end}

\phantomsection\label{olpseg:OLP-0106-B015:start}
\begin{prop}\ollabel{prop:provability-exhaustive}
  که \LR{$\Gamma \cup \{!A\}$} او~\LR{$\Gamma \cup \{\lnot !A\}$} دواړه
  ناسازګار وي، نو \LR{$\Gamma$} ناسازګار دے۔
\end{prop}
\label{olpseg:OLP-0106-B015:end}

\phantomsection\label{olpseg:OLP-0106-B016:start}
\begin{proof}
  که داسې \LR{$!B_1$}، \dots، \LR{$!B_n \in \Gamma$} او
  \LR{$!C_1$}، \dots، \LR{$!C_m \in \Gamma$} وي چې
  \begin{align*}
    \{\sFmla{\True}{!A}, &
    \sFmla{\True}{!B_1}, \dots, \sFmla{\True}{!B_n}\} \text{\RL{ او}}\\
    \{\sFmla{\True}{\lnot !A}, &
    \sFmla{\True}{!C_1}, \dots, \sFmla{\True}{!C_m}\}
  \end{align*}
  دواړه تړلي تابلوګانې ولري، يو داسې ګډ تابلو جوړولے شو
  چې د~\LR{$\Gamma$} ناسازګاري ښيي۔ د فرضونو په توګه
  \LR{$\sFmla{\True}{!B_1}$}، \dots، \LR{$\sFmla{\True}{!B_n}$} له
  \LR{$\sFmla{\True}{!C_1}$}، \dots، \LR{$\sFmla{\True}{!C_m}$} سره يو
  ځاے کاروو، او بيا د~\Cut{} قاعده پلې کوو۔ دا دوه څانګې ورکوي:
  يوه په~\LR{$\sFmla{\True}{!A}$} او بله په~\LR{$\sFmla{\False}{!A}$} پيلېږي۔
\label{olpseg:OLP-0106-B016:end}

\phantomsection\label{olpseg:OLP-0106-B017:start}
په کيڼ لوري کښې د لومړي تابلو د فرضونو لاندې برخه ورزياته
  کړئ۔ دلته د قاعدې هره کارونه لا هم سمه ده؛ ځکه د لومړي
  تابلو هر فرض، د~\LR{$\sFmla{\True}{!A}$} په شمول، موجود دے۔
  نو د~\LR{$\sFmla{\True}{!A}$} لاندې هره څانګه تړل کېږي۔
\label{olpseg:OLP-0106-B017:end}

\phantomsection\label{olpseg:OLP-0106-B018:start}
په ښي لوري کښې د دويم تابلو د فرض لاندې برخه ورزياته کړئ،
  خو د~\LR{$\sFmla{\True}{\lnot !A}$} پر فرض د
  \LR{$\TRule{\True}{\lnot}$} د هرې کارونې پايلې ترې لرې کړئ۔ پر
  \LR{$\sFmla{\True}{\lnot !A}$} د~\LR{$\TRule{\True}{\lnot}$} پايله
  \LR{$\sFmla{\False}{!A}$} ده، چې بيا هم موجوده ده؛ ځکه د ګډ
  تابلو په ښي لوري کښې د~\Cut{} قاعدې پايله ده۔
\label{olpseg:OLP-0106-B018:end}

\phantomsection\label{olpseg:OLP-0106-B019:start}
که په دويم تابلو کښې کومه څانګه د فرض
  \LR{$\sFmla{\True}{\lnot !A}$} (چې په ګډ تابلو کښې نور فرض نۀ
  دے) او~\LR{$\sFmla{\False}{\lnot !A}$} دواړو له امله تړلې وه، پر
  \LR{$\sFmla{\False}{\lnot !A}$} د~\LR{$\TRule{\False}{\lnot}$} په پلي کولو
  \LR{$\sFmla{\True}{!A}$} اخستلے شو۔ اوس په ګډ تابلو کښې اړونده
  څانګه هم تړل کېږي؛ ځکه د~\Cut{} قاعدې ښۍ پايله،
  \LR{$\sFmla{\False}{!A}$}، پکې ده۔ که په دويم تابلو کښې يوه
  څانګه د بل کوم لامل له امله تړلې وه، په ګډ تابلو کښې
  اړونده څانګه هم تړلې ده؛ ځکه په زوړ، دويم تابلو کښې پر
  څانګه له~\LR{$\sFmla{\True}{\lnot !A}$} پرته هر بل نښه لرونکي فارمولونه
  د ګډ تابلو پر اړوندې څانګې هم راځي۔
  \end{proof}
\label{olpseg:OLP-0106-B019:end}

\phantomsection\label{olpseg:OLP-0107-B005:start}
\olfileid{fol}{tab}{ppr}
\label{olpseg:OLP-0107-B005:end}

\phantomsection\label{olpseg:OLP-0107-B006:start}
\olsection{د اشتقاق وړتيا او قضيوي رابطونه}
\label{olpseg:OLP-0107-B006:end}

\phantomsection\label{olpseg:OLP-0107-B007:start}
\begin{explain}
  ثابتوو چې د تابلوګانو د اشتقاق وړتيا اړيکه~\LR{$\Proves$} دومره
  پياوړې ده چې د قضيوي رابطونو په اړه ځينې بنسټيز حقيقتونه ثابت
  کړي؛ لکه \LR{$!A \land !B \Proves !A$} او
  \LR{$!A, !A \lif !B \Proves !B$} (مودوس پوننس)۔ دا حقيقتونه د
  بشپړتيا قضیې د ثبوت دپاره پکار دي۔
\end{explain}
\label{olpseg:OLP-0107-B007:end}

\phantomsection\label{olpseg:OLP-0107-B008:start}
\begin{prop}\ollabel{prop:provability-land}
  \begin{enumerate}
  \item \ollabel{prop:provability-land-left} دواړه
    \LR{$!A \land !B \Proves !A$} او~\LR{$!A \land !B \Proves !B$}۔
  \item \ollabel{prop:provability-land-right}
    \LR{$!A, !B \Proves !A \land !B$}۔
  \end{enumerate}
\end{prop}
\label{olpseg:OLP-0107-B008:end}

\phantomsection\label{olpseg:OLP-0107-B009:start}
\begin{proof}
  \begin{enumerate}
  \item د~\LR{$\{\sFmla{\False}{!A}, \sFmla{\True}{!A \land !B}\}$} او
    \LR{$\{\sFmla{\False}{!B}, \sFmla{\True}{!A \land !B}\}$} دواړو
    دپاره تړلي تابلوګانې شته:
    \begin{oltableau}{}
      [\sFmla{\False}{\formula{A}}, just=\TAss
        [\sFmla{\True}{\formula{A} \land \formula{B}}, just=\TAss
          [\sFmla{\True}{\formula{A}},just={\TRule{\True}{\land}[2]}
            [\sFmla{\True}{\formula{B}},just={\TRule{\True}{\land}[2]}, close
            ]
          ]
        ]
      ]
    \end{oltableau}
    \begin{oltableau}{}
      [\sFmla{\False}{\formula{B}}, just=\TAss
        [\sFmla{\True}{\formula{A} \land \formula{B}}, just=\TAss
          [\sFmla{\True}{\formula{A}},just={\TRule{\True}{\land}[2]}
            [\sFmla{\True}{\formula{B}},just={\TRule{\True}{\land}[2]}, close
            ]
          ]
        ]
      ]
    \end{oltableau}
    \item د~\LR{$\{\sFmla{\True}{!A},
      \sFmla{\True}{!B}, \sFmla{\False}{!A \land !B}\}$} دپاره
      لاندې تړلے تابلو دے:
      \begin{oltableau}
        [\sFmla{\False}{\formula{A} \land \formula{B}}, just = \TAss
          [\sFmla{\True}{\formula{A}}, just = \TAss
            [\sFmla{\True}{\formula{B}}, just=\TAss
              [\sFmla{\False}{\formula{A}}, just = {\TRule{\False}{\land}[1]}, close]
              [\sFmla{\False}{\formula{B}}, just = {\TRule{\False}{\land}[1]}, close]
            ]
          ]
        ]
      \end{oltableau}
  \end{enumerate}
  \textbf{د ژباړې د نسخې سمون (OLTAB-008):} سرچينه په لومړيو
  څلورو تابلوګانو کښې د اتو نښه لرونکو فارمولونو د نښې او فارمول
  تر منځ جلا کونکے قوس غورځوي؛ دلته د تعريف شوې دوه-ارګومېنټه
  د \texttt{\string\sFmla} دوه-ارګومېنټه بڼه بېرته راوړل شوې ده۔
\end{proof}
\label{olpseg:OLP-0107-B009:end}

\phantomsection\label{olpseg:OLP-0107-B010:start}
\begin{prop}\ollabel{prop:provability-lor}
  \begin{enumerate}
  \item \LR{$\{!A \lor !B, \lnot !A, \lnot !B\}$} ناسازګار دے۔
  \item دواړه \LR{$!A \Proves !A \lor !B$} او
    \LR{$!B \Proves !A \lor !B$}۔
  \end{enumerate}
\end{prop}
\label{olpseg:OLP-0107-B010:end}

\phantomsection\label{olpseg:OLP-0107-B011:start}
\begin{proof}
  \begin{enumerate}
  \item د~\LR{$\{\sFmla{\True}{!A \lor !B},
    \sFmla{\True}{\lnot !A}, \sFmla{\True}{\lnot !B}\}$} تړلے
    تابلو ورکوو:
    \begin{oltableau}
      [\sFmla{\True}{\formula{A} \lor \formula{B}}, just = \TAss
        [\sFmla{\True}{\lnot \formula{A}}, just = \TAss
          [\sFmla{\True}{\lnot \formula{B}}, just = \TAss
            [\sFmla{\False}{\formula{A}}, just = {\TRule{\True}{\lnot}[2]}
              [\sFmla{\False}{\formula{B}}, just = {\TRule{\True}{\lnot}[3]}
                [\sFmla{\True}{\formula{A}}, just = {\TRule{\True}{\lor}[1]}, close]
                [\sFmla{\True}{\formula{B}}, just = {\TRule{\True}{\lor}[1]}, close]
              ]
            ]
          ]
        ]
      ]
    \end{oltableau}
  \item د~\LR{$\{\sFmla{\False}{!A \lor !B}, \sFmla{\True}{!A}\}$} او
    \LR{$\{\sFmla{\False}{!A \lor !B}, \sFmla{\True}{!B}\}$} دواړو
    دپاره تړلي تابلوګانې شته:
    \begin{oltableau}{}
      [\sFmla{\False}{\formula{A} \lor \formula{B}}, just=\TAss
        [\sFmla{\True}{\formula{A}}, just=\TAss
          [\sFmla{\False}{\formula{A}},just={\TRule{\False}{\lor}[1]}
            [\sFmla{\False}{\formula{B}},just={\TRule{\False}{\lor}[1]}, close
            ]
          ]
        ]
      ]
    \end{oltableau}
    \begin{oltableau}{}
      [\sFmla{\False}{\formula{A} \lor \formula{B}}, just=\TAss
        [\sFmla{\True}{\formula{B}}, just=\TAss
          [\sFmla{\False}{\formula{A}},just={\TRule{\False}{\lor}[1]}
            [\sFmla{\False}{\formula{B}},just={\TRule{\False}{\lor}[1]}, close
            ]
          ]
        ]
      ]
    \end{oltableau}
  \end{enumerate}
\end{proof}
\label{olpseg:OLP-0107-B011:end}

\phantomsection\label{olpseg:OLP-0107-B012:start}
\begin{prop}\ollabel{prop:provability-lif}
  \begin{enumerate}
  \item \ollabel{prop:provability-lif-left}
    \LR{$!A, !A \lif !B \Proves !B$}۔
  \item \ollabel{prop:provability-lif-right}
    دواړه \LR{$\lnot !A \Proves !A \lif !B$} او
    \LR{$!B \Proves !A \lif !B$}۔
  \end{enumerate}
\end{prop}
\label{olpseg:OLP-0107-B012:end}

\phantomsection\label{olpseg:OLP-0107-B013:start}
\begin{proof}
  \begin{enumerate}
    \item د~\LR{$\{\sFmla{\False}{!B}, \sFmla{\True}{!A \lif !B},
      \sFmla{\True}{!A}\}$} يو تړلے تابلو شته:
      \begin{oltableau}
        [\sFmla{\False}{\formula{B}}, just=\TAss
          [\sFmla{\True}{\formula{A} \lif \formula{B}}, just=\TAss
            [\sFmla{\True}{\formula{A}}, just=\TAss
              [\sFmla{\False}{\formula{A}}, just = {\TRule{\True}{\lif}[2]}, close]
              [\sFmla{\True}{\formula{B}}, just = {\TRule{\True}{\lif}[2]}, close]
            ]
          ]
        ]
      \end{oltableau}
    \item د~\LR{$\{\sFmla{\False}{!A \lif !B},
      \sFmla{\True}{\lnot !A}\}$} او
      \LR{$\{\sFmla{\False}{!A \lif !B}, \sFmla{\True}{!B}\}$} دواړو
      دپاره تړلي تابلوګانې شته:
      \begin{oltableau}
        [\sFmla{\False}{\formula{A} \lif \formula{B}}, just = \TAss
          [\sFmla{\True}{\lnot \formula{A}}, just = \TAss
            [\sFmla{\True}{\formula{A}}, just = {\TRule{\False}{\lif}[1]}
              [\sFmla{\False}{\formula{B}}, just = {\TRule{\False}{\lif}[1]}
                [\sFmla{\False}{\formula{A}}, just = {\TRule{\True}{\lnot}[2]}, close]
              ]
            ]
          ]
        ]
      \end{oltableau}
      \begin{oltableau}
        [\sFmla{\False}{\formula{A} \lif \formula{B}}, just = \TAss
          [\sFmla{\True}{\formula{B}}, just = \TAss
            [\sFmla{\True}{\formula{A}}, just = {\TRule{\False}{\lif}[1]}
              [\sFmla{\False}{\formula{B}}, just = {\TRule{\False}{\lif}[1]},close
              ]
            ]
          ]
        ]
      \end{oltableau}
  \end{enumerate}
\end{proof}
\label{olpseg:OLP-0107-B013:end}

\phantomsection\label{olpseg:OLP-0108-B004:start}
\olfileid{fol}{tab}{qpr}
\olsection{د اشتقاق وړتيا او سورونه}
\label{olpseg:OLP-0108-B004:end}

\phantomsection\label{olpseg:OLP-0108-B005:start}
\begin{explain}
  د بشپړتيا قضيه دا هم غواړي چې د تابلو قاعدې د~\LR{$\Proves$} په اړه
  هغه حقيقتونه ورکړي چې په دې برخه کښې ثابتېږي۔
\end{explain}
\label{olpseg:OLP-0108-B005:end}

\phantomsection\label{olpseg:OLP-0108-B006:start}
\begin{thm}
\ollabel{thm:strong-generalization} که \LR{$c$} داسې ثابت وي چې په
\LR{$\Gamma$} يا~\LR{$!A(x)$} کښې نۀ راځي، او \LR{$\Gamma \Proves !A(c)$} وي،
نو \LR{$\Gamma \Proves \lforall[x][!A(x)]$}۔
\end{thm}
\label{olpseg:OLP-0108-B006:end}

\phantomsection\label{olpseg:OLP-0108-B007:start}
\begin{proof}
فرض کړئ \LR{$\Gamma \Proves !A(c)$}، يعنې داسې
\LR{$!B_1$}، \dots، \LR{$!B_n \in \Gamma$} او د لاندې دپاره تړلے
تابلو شته:
\begin{align*}
\{ \sFmla{\False}{!A(c)}, &
\sFmla{\True}{!B_1}, \dots, \sFmla{\True}{!B_n} \}.
\intertext{\RL{بايد وښيو چې د لاندې دپاره هم تړلے تابلو شته:}}
\{ \sFmla{\False}{\lforall[x][!A(x)]}, &
\sFmla{\True}{!B_1}, \dots, \sFmla{\True}{!B_n} \}.
\end{align*}
تړلے تابلو واخلئ، لومړے فرض ئې په
\LR{$\sFmla{\False}{\lforall[x][!A(x)]}$} بدل کړئ، او د فرضونو وروسته
\LR{$\sFmla{\False}{!A(c)}$} ورننباسئ۔
\begin{center}
\begin{tableau}{not line numbering}
  [\sFmla{\False}{\formula{A}(c)}
    [\sFmla{\True}{\formula{B}_1}
      [\vdots
      [\sFmla{\True}{\formula{B}_n}, 
        [][][]]]]]
\end{tableau}\qquad
\begin{tableau}{not line numbering}
  [\sFmla{\False}{\lforall[x][\formula{A}(x)]}
    [\sFmla{\True}{\formula{B}_1}
      [\vdots
        [\sFmla{\True}{\formula{B}_n},
          [\sFmla{\False}{\formula{A}(c)}
        [][][]]]]]]
\end{tableau}
\end{center}
تابلو لا هم تړلے دے؛ ځکه ټولې هغه تړلي فارمولونه چې مخکښې د فرضونو
په توګه موجودې وې، اوس هم د~تابلو په سر کښې موجودې دي۔
ننويستل شوې کرښه د~\LR{$\TRule{\False}{\lforall}$} قاعدې د سمې کارونې
پايله ده؛ ځکه ثابت سمبول~\LR{$c$} په~\LR{$!B_1$}، \dots، \LR{$!B_n$} يا
\LR{$\lforall[x][!A(x)]$} کښې نۀ راځي، يعنې په نوي تابلو کښې د
ننويستل شوې کرښې پاس نۀ راځي۔
\end{proof}
\label{olpseg:OLP-0108-B007:end}

\phantomsection\label{olpseg:OLP-0108-B008:start}
\begin{prop}
\ollabel{prop:provability-quantifiers}
\begin{enumerate}
\item \LR{$!A(t) \Proves \lexists[x][!A(x)]$}.
\item \LR{$\lforall[x][!A(x)] \Proves !A(t)$}.
\end{enumerate}
\end{prop}
\label{olpseg:OLP-0108-B008:end}

\phantomsection\label{olpseg:OLP-0108-B009:start}
\begin{proof}
\begin{enumerate}
\item د
  \LR{$\sFmla{\False}{\lexists[x][!A(x)]}, \sFmla{\True}{!A(t)}$} دپاره
  تړلے تابلو دا دے:
  \begin{oltableau}
    [\sFmla{\False}{\lexists[x][\formula{A}(x)]}, just = \TAss
      [\sFmla{\True}{\formula{A}(t)}, just = \TAss
        [\sFmla{\False}{\formula{A}(t)}, just = {\TRule{\False}{\lexists}[1]}, close]]]
    \end{oltableau}
\item د
  \LR{$\sFmla{\False}{\formula{A}(t)}, \sFmla{\True}{\lforall[x][!A(x)]}, $} دپاره
  تړلے تابلو دا دے:
  \begin{oltableau}
    [\sFmla{\False}{\formula{A}(t)}, just = \TAss
      [\sFmla{\True}{\lforall[x][\formula{A}(x)]}, just = \TAss
        [\sFmla{\True}{\formula{A}(t)}, just = {\TRule{\True}{\lforall}[2]}, close]]]
    \end{oltableau}
\end{enumerate}
\end{proof}
\label{olpseg:OLP-0108-B009:end}

\phantomsection\label{olpseg:OLP-0109-B004:start}
\olfileid{fol}{tab}{sou}
\label{olpseg:OLP-0109-B004:end}

\phantomsection\label{olpseg:OLP-0109-B005:start}
\olsection{سموالے}
\label{olpseg:OLP-0109-B005:end}

\phantomsection\label{olpseg:OLP-0109-B006:start}
\begin{explain}
د تابلوګانو په شان اشتقاق نظام هله \emph{سم} دے چې
هغه څيزونه اشتقاق نۀ کړي چې په حقيقت کښې نۀ وي۔ نو سموالے د
اشتقاق نظامونو د تضمين شوې خونديتوب يو ډول خاصيت دے۔
د دې له مخې چې کوم ثبوت-تيوريکي خاصيت تر بحث لاندې وي، مثلاً غواړو
پوه شو چې:
\begin{enumerate}
\item هر د اشتقاق وړ~\LR{$!A$} معتبر دے؛
\item که تړلے فارمول له ځينو نورو څخه د اشتقاق وړ وي، د هغو
  معنايي پايله هم ده؛
\item که د تړلي فارمولونه يو سټ ناسازګار وي، د صدق وړ نۀ دے۔
\end{enumerate}
دا د اشتقاق نظام مهم خاصيتونه دي۔ که له دغو څخه کوم يو
سم نۀ وي، اشتقاق نظام نيمګړے دے---له حده زيات څيزونه به
اشتقاق کوي۔ له دې امله د اشتقاق نظام سموالے ثابتول
تر ټولو زيات اهميت لري۔
\label{olpseg:OLP-0109-B006:end}

\phantomsection\label{olpseg:OLP-0109-B007:start}
ځکه دغه ټول ثبوت-تيوريکي خاصيتونه د يو يا بل ډول د تړلو
تابلوګانې له لارې تعريف شوي دي، د پورته (۱)--(۳) ثابتول غواړي
چې د تړلو تابلوګانې د معنايي خاصيتونو په اړه يو څه ثابت کړو۔
لومړے به تعريف کړو چې د يو نښه لرونکے فارمول پوره کېدل په يو
جوړښت کښې څه معنٰی لري؛ بيا به وښيو چې که يو تابلو تړلے وي،
هېڅ جوړښت ئې ټول فرضونه نۀ پوره کوي۔ وروسته (۱)--(۳) د دې پايلې
د لازمو نتيجو په توګه راځي۔
\end{explain}
\label{olpseg:OLP-0109-B007:end}

\phantomsection\label{olpseg:OLP-0109-B008:start}
\begin{defn}
جوړښت~\LR{$\Struct{M}$}
يو نښه لرونکے فارمول~\LR{$\sFmla{\True}{!A}$} هله او يوازې هله
\emph{پوره کوي} چې \LR{$\Sat{M}{!A}$}، او
\LR{$\sFmla{\False}{!A}$} هله او يوازې هله پوره کوي چې
\LR{$\Sat/{M}{!A}$}۔
\LR{$\Struct{M}$} د نښه لرونکي فارمولونه يو
سټ~\LR{$\Gamma$} هله او يوازې هله پوره کوي چې هر
\LR{$\sFmla{S}{!A} \in \Gamma$} پوره کړي۔ \LR{$\Gamma$} هله \emph{د صدق وړ}
دے چې کوم جوړښت ئې پوره کړي،
او که داسې نۀ وي \emph{د صدق وړ نۀ} دے۔
\end{defn}
\label{olpseg:OLP-0109-B008:end}

\phantomsection\label{olpseg:OLP-0109-B009:start}
\begin{thm}[سموالے]
  \ollabel{thm:tableau-soundness}
  که \LR{$\Gamma$} يو تړلے تابلو ولري، نو \LR{$\Gamma$} د صدق وړ نۀ دے۔
\end{thm}
\label{olpseg:OLP-0109-B009:end}

\phantomsection\label{olpseg:OLP-0109-B010:start}
\begin{proof}
د يو تابلو څانګه هله او يوازې هله د صدق وړ بولو چې پر هغې
د نښه لرونکي فارمولونه سټ د صدق وړ وي؛ او تابلو هله د صدق وړ
بولو چې لږ تر لږه يوه د صدق وړ څانګه ولري۔
\label{olpseg:OLP-0109-B010:end}

\phantomsection\label{olpseg:OLP-0109-B011:start}
دا خبره ښيو: د استنتاج د کومې قاعدې په وسيله د يو د صدق وړ
تابلو پراخول تل يو د صدق وړ تابلو ورکوي۔ دا به قضيه
ثابته کړي: هر تړلے تابلو د~\LR{$\Gamma$} يوازې له فرضونو څخه جوړ
تابلو ته د استنتاج د قاعدو په پلي کولو تر لاسه کېږي۔ نو که
\LR{$\Gamma$} د صدق وړ واے، د هغې هر تابلو به د صدق وړ و۔ خو تړلے
تابلو په څرګند ډول د صدق وړ نۀ دے: هره څانګه دواړه
\LR{$\sFmla{\True}{!A}$} او \LR{$\sFmla{\False}{!A}$} لري، او هېڅ جوړښت
نۀ شي کولے چې \LR{$!A$} هم پوره او هم نۀ پوره کړي۔
\label{olpseg:OLP-0109-B011:end}

\phantomsection\label{olpseg:OLP-0109-B012:start}
فرض کړئ يو د صدق وړ تابلو لرو، يعنې داسې تابلو چې لږ
تر لږه يوه د صدق وړ څانګه لري۔ د استنتاج قاعده يا يوې څانګې ته
نښه لرونکي فارمولونه ورزياتوي يا يوه څانګه په دوو وېشي۔ که په
تابلو کښې کومه د صدق وړ څانګه وي چې اړوند قاعدوي تطبيق ئې نۀ
پراخوي، هغه په پراخ شوي تابلو کښې هم د صدق وړ پاتې کېږي؛ نو
پراخ شوے تابلو د صدق وړ دے۔ ځکه يوازې هغه حالت څېړل پکار دي چې
قاعده پر يوې د صدق وړ څانګې پلي کېږي۔
\label{olpseg:OLP-0109-B012:end}

\phantomsection\label{olpseg:OLP-0109-B013:start}
\LR{$\Gamma$} پر هغې څانګې د نښه لرونکي فارمولونه سټ وبولئ، او
\LR{$\sFmla{S}{!A} \in \Gamma$} هغه نښه لرونکے فارمول وبولئ چې قاعده
پرې پلي کېږي۔ که قاعده څانګه نۀ وېشي، بايد وښيو چې پراخه شوې څانګه،
يعنې \LR{$\Gamma$} له قاعدوي پايلو سره، لا هم د صدق وړ ده۔ که قاعده
څانګه وېشي، بايد وښيو چې له دوو ترلاسه شوو څانګو لږ تر لږه يوه د
صدق وړ ده۔
\label{olpseg:OLP-0109-B013:end}

\phantomsection\label{olpseg:OLP-0109-B014:start}
لومړے هغه ممکن استنتاجونه څېړو چې څانګه نۀ وېشي۔
\begin{enumerate}
\item څانګه پر \LR{$\sFmla{\True}{\lnot !B} \in \Gamma$} د
  \LR{$\TRule{\True}{\lnot}$} په پلي کولو پراخېږي۔ بيا پراخه شوې څانګه
  د \LR{$\Gamma \cup \{\sFmla{\False}{!B}\}$} نښه لرونکي فارمولونه لري۔
  فرض کړئ \LR{$\Sat{M}{\Gamma}$}۔ په ځانګړي ډول
  \LR{$\Sat{M}{\lnot !B}$}۔ نو
  \LR{$\Sat/{M}{!B}$}، يعنې
  \LR{$\Struct{M}$}،
  \LR{$\sFmla{\False}{!B}$} پوره کوي۔
\item څانګه پر \LR{$\sFmla{\False}{\lnot !B} \in \Gamma$} د
  \LR{$\TRule{\False}{\lnot}$} په پلي کولو پراخېږي: تمرين۔
\item څانګه پر \LR{$\sFmla{\True}{!B \land !C} \in \Gamma$} د
  \LR{$\TRule{\True}{\land}$} په پلي کولو پراخېږي؛ له دې څخه پر څانګه
  دوه نوي نښه لرونکي فارمولونه ترلاسه کېږي: \LR{$\sFmla{\True}{!B}$} او
  \LR{$\sFmla{\True}{!C}$}۔ فرض کړئ
  \LR{$\Sat{M}{\Gamma}$}، او په ځانګړي ډول
  \LR{$\Sat{M}{!B \land !C}$}۔ بيا
  \LR{$\Sat{M}{!B}$} او
  \LR{$\Sat{M}{!C}$}۔ يعنې
  \LR{$\Struct{M}$} دواړه
  \LR{$\sFmla{\True}{!B}$} او \LR{$\sFmla{\True}{!C}$} پوره کوي۔
\item څانګه پر \LR{$\sFmla{\False}{!B \lor !C} \in \Gamma$} د
  \LR{$\TRule{\False}{\lor}$} په پلي کولو پراخېږي: تمرين۔
\item څانګه پر \LR{$\sFmla{\False}{!B \lif !C} \in \Gamma$} د
  \LR{$\TRule{\False}{\lif}$} په پلي کولو پراخېږي۔ له دې څخه پر څانګه
  دوه نوي نښه لرونکي فارمولونه ترلاسه کېږي: \LR{$\sFmla{\True}{!B}$} او
  \LR{$\sFmla{\False}{!C}$}۔ فرض کړئ
  \LR{$\Sat{M}{\Gamma}$}، او په ځانګړي ډول
  \LR{$\Sat/{M}{!B \lif !C}$}۔ بيا
  \LR{$\Sat{M}{!B}$} او
  \LR{$\Sat/{M}{!C}$}۔ يعنې
  \LR{$\Struct{M}$} دواړه
  \LR{$\sFmla{\True}{!B}$} او \LR{$\sFmla{\False}{!C}$} پوره کوي۔
\label{olpseg:OLP-0109-B014:end}

\phantomsection\label{olpseg:OLP-0109-B015:start}
%
\textbf{د ژباړې د نسخې سمون (OLTAB-009):} سرچينې په دواړو
کلي-کميت ټاکونکو قاعدوي حالتونو کښې د قاعدې په مقدمه کښې د بي ميتا-فارمول
ليکلے و، خو پايله او ټول ورپسې معنايي استدلال د اې ميتا-فارمول
کاروي۔ دلته اې په هر حالت کښې يو شان ساتل کېږي۔
\item څانګه پر \LR{$\sFmla{\True}{\lforall[x][!A(x)]} \in \Gamma$} د
  \LR{$\TRule{\True}{\lforall}$} په پلي کولو پراخېږي۔ له دې څخه پر څانګه
  يو نوے نښه لرونکے فارمول~\LR{$\sFmla{\True}{!A(t)}$} ترلاسه کېږي۔
  فرض کړئ \LR{$\Sat{M}{\Gamma}$}، او په ځانګړي ډول
  \LR{$\Sat{M}{\lforall[x][!A(x)]}$}۔ د
  \olref[syn][sem]{prop:quant-terms} له مخې \LR{$\Sat{M}{!A(t)}$}۔ نو
  \LR{$\Struct{M}$}، \LR{$\sFmla{\True}{!A(t)}$} پوره کوي۔
\label{olpseg:OLP-0109-B015:end}

\phantomsection\label{olpseg:OLP-0109-B016:start}
\item څانګه پر \LR{$\sFmla{\False}{\lforall[x][!A(x)]} \in \Gamma$} د
  \LR{$\TRule{\False}{\lforall}$} په پلي کولو پراخېږي۔ له دې څخه پر
  څانګه يو نوے نښه لرونکے فارمول~\LR{$\sFmla{\False}{!A(a)}$} ترلاسه
  کېږي، چې \LR{$a$} داسې ثابت سمبول دے چې په~\LR{$\Gamma$} کښې نۀ راځي۔
  ځکه \LR{$\Gamma$} د صدق وړ دے، داسې \LR{$\Struct{M}$} شته چې
  \LR{$\Sat{M}{\Gamma}$}، او په ځانګړي ډول
  \LR{$\Sat/{M}{\lforall[x][!A(x)]}$}۔ بايد وښيو چې
  \LR{$\Gamma \cup \{\sFmla{\False}{!A(a)}\}$} د صدق وړ دے۔ د دې دپاره
  مناسب~\LR{$\Struct{M'}$} داسې تعريفوو۔
\label{olpseg:OLP-0109-B016:end}

\phantomsection\label{olpseg:OLP-0109-B017:start}
د \olref[syn][ass]{prop:sat-quant} له مخې
  \LR{$\Sat/{M}{\lforall[x][!A(x)]}$} هله او يوازې هله چې د کوم~\LR{$s$}
  دپاره \LR{$\Sat/{M}{!A(x)}[s]$}۔ اوس \LR{$\Struct{M'}$} د~\LR{$\Struct{M}$} په
  شان وبولئ، خو \LR{$\Assign{a}{M'} = s(x)$}۔ د
  \olref[syn][ext]{cor:extensionality-sent} له مخې، د هر
  \LR{$\sFmla{\True}{!C} \in \Gamma$} دپاره \LR{$\Sat{M'}{!C}$}، او د هر
  \LR{$\sFmla{\False}{!C} \in \Gamma$} دپاره \LR{$\Sat/{M'}{!C}$}؛ ځکه
  \LR{$a$} په~\LR{$\Gamma$} کښې نۀ راځي۔
\label{olpseg:OLP-0109-B017:end}

\phantomsection\label{olpseg:OLP-0109-B018:start}
د \olref[syn][ext]{prop:extensionality} له مخې\linebreak
  \LR{$\Sat/{M'}{!A(x)}[s]$}۔ د \olref[syn][ext]{prop:ext-formulas} له مخې
  \LR{$\Sat/{M'}{!A(a)}[s]$}۔ ځکه \LR{$!A(a)$} يوه تړلے فارمول ده، د
  \olref[syn][ass]{prop:sentence-sat-true} له مخې
  \LR{$\Sat/{M'}{!A(a)}$}؛ يعنې \LR{$\Struct{M'}$}،
  \LR{$\sFmla{\False}{!A(a)}$} پوره کوي۔
\item څانګه پر \LR{$\sFmla{\True}{\lexists[x][!B(x)]} \in \Gamma$} د
  \LR{$\TRule{\True}{\lexists}$} په پلي کولو پراخېږي: تمرين۔
\item څانګه پر \LR{$\sFmla{\False}{\lexists[x][!B(x)]} \in \Gamma$} د
  \LR{$\TRule{\False}{\lexists}$} په پلي کولو پراخېږي: تمرين۔
\end{enumerate}
اوس هغه ممکن استنتاجونه څېړو چې څانګه وېشي۔
\begin{enumerate}
\item څانګه پر \LR{$\sFmla{\False}{!B \land !C} \in \Gamma$} د
  \LR{$\TRule{\False}{\land}$} په پلي کولو پراخېږي؛ له دې څخه دوه څانګې
  ترلاسه کېږي: کيڼه څانګه د~\LR{$\sFmla{\False}{!B}$} له لارې او ښۍ
  څانګه د~\LR{$\sFmla{\False}{!C}$} له لارې دوام کوي۔ فرض کړئ
  \LR{$\Sat{M}{\Gamma}$}، او په ځانګړي ډول
  \LR{$\Sat/{M}{!B \land !C}$}۔ بيا يا
  \LR{$\Sat/{M}{!B}$} يا
  \LR{$\Sat/{M}{!C}$}۔ په لومړي حالت کښې
  \LR{$\Struct{M}$}،
  \LR{$\sFmla{\False}{!B}$} پوره کوي؛ يعنې
  \LR{$\Struct{M}$} د کيڼې څانګې فارمولونه
  پوره کوي۔ په دويم حالت کښې
  \LR{$\Struct{M}$}،
  \LR{$\sFmla{\False}{!C}$} پوره کوي؛ يعنې
  \LR{$\Struct{M}$} د ښۍ څانګې فارمولونه
  پوره کوي۔
\item څانګه پر \LR{$\sFmla{\True}{!B \lor !C} \in \Gamma$} د
  \LR{$\TRule{\True}{\lor}$} په پلي کولو پراخېږي: تمرين۔
\item څانګه پر \LR{$\sFmla{\True}{!B \lif !C} \in \Gamma$} د
  \LR{$\TRule{\True}{\lif}$} په پلي کولو پراخېږي: تمرين۔
\item څانګه د~\Cut{} په وسيله پراخېږي۔ له دې څخه دوه څانګې ترلاسه
  کېږي: يوه \LR{$\sFmla{\True}{!B}$} او بله \LR{$\sFmla{\False}{!B}$} لري۔
  ځکه \LR{$\Sat{M}{\Gamma}$} او يا
  \LR{$\Sat{M}{!B}$} يا
  \LR{$\Sat/{M}{!B}$}،
  \LR{$\Struct{M}$} يا کيڼه يا ښۍ څانګه پوره
  کوي۔
\end{enumerate}
\end{proof}
\label{olpseg:OLP-0109-B018:end}

\phantomsection\label{olpseg:OLP-0109-B019:start}
\begin{prob}
د~\olref[fol][tab][sou]{thm:tableau-soundness} ثبوت بشپړ کړئ۔
\end{prob}
\label{olpseg:OLP-0109-B019:end}



\phantomsection\label{olpseg:OLP-0109-B021:start}
\begin{cor}
\ollabel{cor:weak-soundness}
که \LR{$\Proves !A$} وي، نو \LR{$!A$} معتبر دے۔
\end{cor}
\label{olpseg:OLP-0109-B021:end}

\phantomsection\label{olpseg:OLP-0109-B022:start}
\begin{cor}
\ollabel{cor:entailment-soundness}
که \LR{$\Gamma \Proves !A$} وي، نو \LR{$\Gamma \Entails !A$}۔
\end{cor}
\label{olpseg:OLP-0109-B022:end}

\phantomsection\label{olpseg:OLP-0109-B023:start}
\begin{proof}
  که \LR{$\Gamma \Proves !A$} وي، نو د ځينو
  \LR{$!B_1$}، \dots، \LR{$!B_n \in \Gamma$} دپاره
  \LR{$\{\sFmla{\False}{!A}, \sFmla{\True}{!B_1}, \dots,
  \sFmla{\True}{!B_n}\}$} يو تړلے تابلو لري۔ د
  \olref{thm:tableau-soundness} له مخې، هر
  جوړښت~\LR{$\Struct{M}$}
  يا کوم \LR{$!B_i$} ناسموي يا \LR{$!A$} رښتيا کوي۔ نو که
  \LR{$\Sat{M}{\Gamma}$} وي، بيا
  \LR{$\Sat{M}{!A}$} هم دے۔
\end{proof}
\label{olpseg:OLP-0109-B023:end}

\phantomsection\label{olpseg:OLP-0109-B024:start}
\begin{cor}
\ollabel{cor:consistency-soundness}
که \LR{$\Gamma$} د صدق وړ وي، نو سازګار دے۔
\end{cor}
\label{olpseg:OLP-0109-B024:end}

\phantomsection\label{olpseg:OLP-0109-B025:start}
\begin{proof}
مقابل عکس ثابتوو۔ فرض کړئ \LR{$\Gamma$} سازګار نۀ دے۔ بيا
\LR{$!B_1$}، \dots، \LR{$!B_n \in \Gamma$} شته او د
\LR{$\{\sFmla{\True}{!B_1}, \dots, \sFmla{\True}{!B_n}\}$} دپاره يو
تړلے تابلو شته۔ د \olref{thm:tableau-soundness} له مخې، داسې
\LR{$\Struct{M}$} نشته چې د ټولو
\LR{$i=1$}، \dots،~\LR{$n$} دپاره \LR{$\Sat{M}{!B_i}$} وي۔
خو بيا \LR{$\Gamma$} د صدق وړ نۀ دے۔
\end{proof}
\label{olpseg:OLP-0109-B025:end}

\phantomsection\label{olpseg:OLP-0110-B004:start}
\olfileid{fol}{tab}{ide}
\label{olpseg:OLP-0110-B004:end}

\phantomsection\label{olpseg:OLP-0110-B005:start}
\olsection{له عينيت سره تابلوګانې}
\label{olpseg:OLP-0110-B005:end}

\phantomsection\label{olpseg:OLP-0110-B006:start}
له عينيت سره تابلوګانې د استنتاج اضافي قاعدې غواړي۔
د~\LR{$\eq$} قاعدې دا دي (\LR{$t$}، \LR{$t_1$} او \LR{$t_2$} تړلي ترمونه دي):
\label{olpseg:OLP-0110-B006:end}

\phantomsection\label{olpseg:OLP-0110-B007:start}
\begin{defish}
\AxiomC{}
\RightLabel{\LR{$\eq$}}
\UnaryInfC{\sFmla{\True}{\eq[t][t]}}
\DisplayProof
\hfill
\AxiomC{\sFmla{\True}{\eq[t_1][t_2]}}
\noLine
\UnaryInfC{\sFmla{\True}{!A(t_1)}}
\RightLabel{\LR{$\TRule{\True}{\eq}$}}
\UnaryInfC{\sFmla{\True}{!A(t_2)}}
\DisplayProof
\hfill
\AxiomC{\sFmla{\True}{\eq[t_1][t_2]}}
\noLine
\UnaryInfC{\sFmla{\False}{!A(t_1)}}
\RightLabel{\LR{$\TRule{\False}{\eq}$}}
\UnaryInfC{\sFmla{\False}{!A(t_2)}}
\DisplayProof
\end{defish}
پام وکړئ چې له نورو ټولو قاعدو سره په توپير کښې،
\LR{$\TRule{\True}{\eq}$} او \LR{$\TRule{\False}{\eq}$} غواړي چې \emph{دوه}
نښه لرونکي فارمولونه له وړاندې پر څانګه وي؛ يعنې دواړه
\LR{$\sFmla{\True}{\eq[t_1][t_2]}$} او \LR{$\sFmla{S}{!A(t_1)}$}۔
\label{olpseg:OLP-0110-B007:end}

\phantomsection\label{olpseg:OLP-0110-B008:start}
\begin{ex}
که \LR{$s$} او \LR{$t$} تړلي ترمونه وي، نو
\LR{$\eq[s][t], !A(s) \Proves !A(t)$}:
\begin{oltableau}
  [\sFmla{\False}{\formula{A}(t)}, just = \TAss
    [\sFmla{\True}{\eq[s][t]}, just = \TAss
      [\sFmla{\True}{\formula{A}(s)}, just = \TAss
        [\sFmla{\True}{\formula{A}(t)}, just={\TRule{\True}{\eq}[2, 3]}, close]
      ]
    ]
  ]
\end{oltableau}
کېدے شي دا د يو شان څيزونو د تعويض د اصل يا د لايبنېڅ د قانون په
نوم درته اشنا وي۔
\label{olpseg:OLP-0110-B008:end}

\phantomsection\label{olpseg:OLP-0110-B009:start}
تابلوګانې ثابتوي چې \LR{$\eq$} متناظره ده، يعنې
\LR{$\eq[s_1][s_2] \Proves \eq[s_2][s_1]$}:
\begin{oltableau}
  [\sFmla{\False}{\eq[s_2][s_1]}, just = \TAss
    [\sFmla{\True}{\eq[s_1][s_2]}, just = \TAss
      [\sFmla{\True}{\eq[s_1][s_1]}, just = {\LR{$\eq$}}
        [\sFmla{\True}{\eq[s_2][s_1]}, just = {\TRule{\True}{\eq}[2, 3]}, close]
      ]
    ]
  ]
\end{oltableau}
دلته کرښه~\LR{$2$} د~\LR{$\TRule{\True}{\eq}$} لومړنے اړين
فارمول~\LR{$\sFmla{\True}{\eq[s_1][s_2]}$} دے۔ کرښه~\LR{$3$} دويم
اړين فارمول دے، چې د~\LR{$\sFmla{\True}{!A(s_1)}$} بڼه لري---\LR{$!A(x)$}
د~\LR{$\eq[x][s_1]$} په توګه وګڼئ؛ بيا \LR{$!A(s_1)$}،
\LR{$\eq[s_1][s_1]$} او \LR{$!A(s_2)$}، \LR{$\eq[s_2][s_1]$} کېږي۔
\textbf{د ژباړې د نسخې سمون (OLTAB-010):} سرچينې کرښه درې د اې
د دويم ترم بڼه بللې وه، خو د قاعدې په دې تطبيق کښې هغه د اې د
لومړي ترم بڼه ده، لکه ورپسې تشريح چې هم ښيي۔
\label{olpseg:OLP-0110-B009:end}

\phantomsection\label{olpseg:OLP-0110-B010:start}
تابلوګان دا هم ثابتوي چې \LR{$\eq$} انتقالي ده، يعنې
\LR{$\eq[s_1][s_2], \eq[s_2][s_3] \Proves \eq[s_1][s_3]$}:
\begin{oltableau}
  [\sFmla{\False}{\eq[s_1][s_3]}, just = \TAss
    [\sFmla{\True}{\eq[s_1][s_2]}, just = \TAss
      [\sFmla{\True}{\eq[s_2][s_3]}, just = \TAss
        [\sFmla{\True}{\eq[s_1][s_3]}, just = {\TRule{\True}{\eq}[3, 2]}, close]
      ]
    ]
  ]
\end{oltableau}
په دې تابلو کښې د~\LR{$\TRule{\True}{\eq}$} لومړنے اړين
فارمول کرښه~\LR{$3$}، \LR{$\sFmla{\True}{\eq[s_2][s_3]}$} ده
(\LR{$s_2$} د~\LR{$t_1$} او \LR{$s_3$} د~\LR{$t_2$} ونډه لري)۔ دويم اړين فارمول چې
د~\LR{$\sFmla{\True}{!A(s_2)}$} بڼه لري، کرښه~\LR{$2$} ده۔ دلته \LR{$!A(x)$}
د~\LR{$\eq[s_1][x]$} په توګه وګڼئ؛ په دې سره \LR{$!A(s_2)$}،
\LR{$\eq[s_1][s_2]$} (يعنې کرښه~\LR{$2$}) او \LR{$!A(s_3)$} په پايله کښې
فارمول~\LR{$\eq[s_1][s_3]$} کېږي۔
\textbf{د ژباړې د نسخې سمون (OLTAB-011):} سرچينې د اې د دويم
ترم د بڼې د پراخولو پر ځاے د قاعدې عمومي برابري ليکلې وه؛ خو
د همدلته ټاکل شوي اې له مخې اړوند فارمول د لومړي او دويم ترم
برابري ده۔
\end{ex}
\label{olpseg:OLP-0110-B010:end}

\phantomsection\label{olpseg:OLP-0110-B011:start}
\begin{prob}
د لاندې دپاره تړلي تابلوګانې ورکړئ:
\begin{enumerate}
\item \LR{$\sFmla{\False}{\lforall[x][\lforall[y][((x = y \land !A(x))
      \lif !A(y))]]}$}
\item \LR{$\sFmla{\False}{\lexists[x][(!A(x) \land
    \lforall[y][(!A(y) \lif y = x)])]}$},\\
  \LR{$\sFmla{\True}{\lexists[x][!A(x)] \land
    \lforall[y][\lforall[z][((!A(y) \land !A(z)) \lif y = z)]]}$}
\end{enumerate}
\end{prob}
\label{olpseg:OLP-0110-B011:end}

\phantomsection\label{olpseg:OLP-0111-B004:start}
\olfileid{fol}{tab}{sid}
\label{olpseg:OLP-0111-B004:end}

\phantomsection\label{olpseg:OLP-0111-B005:start}
\olsection{له عينيت سره سموالے}
\label{olpseg:OLP-0111-B005:end}

\phantomsection\label{olpseg:OLP-0111-B006:start}
\begin{prop}
د عينيت له قاعدو سره تابلوګانې سم دي: هېڅ تړلے تابلو
د صدق وړ نۀ دے۔
\end{prop}
\label{olpseg:OLP-0111-B006:end}

\phantomsection\label{olpseg:OLP-0111-B007:start}
\begin{proof}
يوازې بايد د پخوا په شان وښيو چې که يو تابلو د صدق وړ
څانګه ولري، نو د~\LR{$\eq$} له قاعدو څخه د يوې په پلي کولو ترلاسه شوې
څانګه هم د صدق وړ ده۔ \LR{$\Gamma$} پر څانګه د نښه لرونکو فارمولونه
سټ وبولئ، او \LR{$\Struct{M}$} داسې جوړښت وبولئ چې \LR{$\Gamma$}
پوره کوي۔
\label{olpseg:OLP-0111-B007:end}

\phantomsection\label{olpseg:OLP-0111-B008:start}
فرض کړئ څانګه د~\LR{$\eq$} په کارولو، يعنې د نښه لرونکي
فارمول~\LR{$\sFmla{\True}{\eq[t][t]}$} په زياتولو پراخېږي۔ په
څرګند ډول \LR{$\Sat{M}{\eq[t][t]}$}؛ نو \LR{$\Struct{M}$}،
\LR{$\Gamma \cup \{\sFmla{\True}{\eq[t][t]}\}$} هم پوره کوي۔
\label{olpseg:OLP-0111-B008:end}

\phantomsection\label{olpseg:OLP-0111-B009:start}
که څانګه د~\LR{$\TRule{\True}{\eq}$} په کارولو پراخېږي، موږ يو نښه
لرونکے فارمول~\LR{$\sFmla{\True}{!A(t_2)}$} ورزياتوو، خو \LR{$\Gamma$}
دواړه~\LR{$\sFmla{\True}{\eq[t_1][t_2]}$} او
\LR{$\sFmla{\True}{!A(t_1)}$} لري۔
\textbf{د ژباړې د نسخې سمون (OLTAB-012):} سرچينې د رښتوني-عينيت
قاعدې پايلې ته عمومي S نښه ورکړې وه؛ د قاعدې، دواړو مقدمو او
لاندې معنايي ثبوت له مخې دا رښتونې نښه ده۔
نو \LR{$\Sat{M}{\eq[t_1][t_2]}$} او \LR{$\Sat{M}{!A(t_1)}$} لرو۔ داسې
متغير-ټاکنه~\LR{$s$} واخلئ چې \LR{$s(x) = \Value{t_1}{M}$}۔ د
\olref[syn][ass]{prop:sentence-sat-true} له مخې
\LR{$\Sat{M}{!A(t_1)}[s]$}۔ ځکه \LR{$\varAssign{s}{s}{x}$}، د
\olref[syn][ext]{prop:ext-formulas} له مخې \LR{$\Sat{M}{!A(x)}[s]$}۔
ځکه \LR{$\Sat{M}{\eq[t_1][t_2]}$}، نو
\LR{$\Value{t_1}{M} = \Value{t_2}{M}$}، او ځکه
\LR{$s(x) = \Value{t_2}{M}$}۔ د
\olref[syn][ext]{prop:ext-formulas} په بيا پلي کولو
\LR{$\Sat{M}{!A(t_2)}[s]$} هم لرو۔ د
\olref[syn][ass]{prop:sentence-sat-true} له مخې
\LR{$\Sat{M}{!A(t_2)}$}۔ د~\LR{$\TRule{\False}{\eq}$} حالت هم په همدې ډول
حل کېږي۔
\end{proof}
\label{olpseg:OLP-0111-B009:end}

\phantomsection\label{olpseg:OLP-0112-B004:start}
\olchapter{fol}{axd}{بديهي اشتقاقونه}
\label{olpseg:OLP-0112-B004:end}

\phantomsection\label{olpseg:OLP-0112-B005:start}
\begin{editorial}
  لا تر اوسه دا هڅه نۀ ده شوې چې ډاډ ترلاسه شي چې د دې څپرکي مواد
  د هغو بېلابېلو ټاګونو درناوے کوي چې ښيي کوم نښلوونکي او کميت
  ټاکونکي بنسټيز يا تعريف شوي دي: له~\LR{$\liff$} پرته، چې تعريف شوے
  ګڼل کېږي، ټول بنسټيز ګڼل شوي دي۔ که د لومړۍ درجې منطق ټاګ رښتونے
  وي، له کميت ټاکونکو سره بڼه جوړوو؛ که نۀ وي، بې له هغو بڼه جوړوو۔
\end{editorial}
\label{olpseg:OLP-0112-B005:end}

\phantomsection\label{olpseg:OLP-0112-B008:start}
%


\label{olpseg:OLP-0112-B008:end}

\phantomsection\label{olpseg:OLP-0112-B010:start}
%


\label{olpseg:OLP-0112-B010:end}

\phantomsection\label{olpseg:OLP-0112-B013:start}
%


\label{olpseg:OLP-0112-B013:end}

\phantomsection\label{olpseg:OLP-0112-B016:start}
%


\label{olpseg:OLP-0112-B016:end}

\phantomsection\label{olpseg:OLP-0113-B004:start}
\olfileid{fol}{axd}{rul}
\label{olpseg:OLP-0113-B004:end}

\phantomsection\label{olpseg:OLP-0113-B005:start}
\olsection{قاعدې او اشتقاقونه}
\label{olpseg:OLP-0113-B005:end}

\phantomsection\label{olpseg:OLP-0113-B006:start}
\begin{explain}
  بديهي اشتقاقونه ښايي د منطق د اشتقاق تر ټولو ساده
  نظام وي۔ اشتقاق يوازې د فارمولونه يوه لړۍ ده۔ د
  اشتقاق ګڼل کېدو دپاره، په لړۍ کښې هر فارمول بايد
  يا د کوم بديهي اصل يوه بېلګه وي، يا د استنتاج د کومې قاعدې له مخې
  له يوه يا زياتو هغو فارمولونه څخه راشي چې په لړۍ کښې ترې وړاندې
  دي۔ اشتقاق خپل وروستے فارمول، اشتقاق کوي۔
\end{explain}
\label{olpseg:OLP-0113-B006:end}

\phantomsection\label{olpseg:OLP-0113-B007:start}
\begin{defn}[د اشتقاق وړتيا]
که \LR{$\Gamma$} د~\LR{$\Lang L$} د فارمولونه يو سټ وي، نو له~\LR{$\Gamma$}
څخه يو \emph{اشتقاق} د فارمولونه يوه متناهي لړۍ
\LR{$!A_1$}، \dots،~\LR{$!A_n$} ده، چې د هر \LR{$i \le n$} دپاره له لاندې څخه
يو شرط پوره کېږي:
\begin{enumerate}
\item \LR{$!A_i \in \Gamma$}؛ يا
\item \LR{$!A_i$} يو بديهي اصل دے؛ يا
\item \LR{$!A_i$} د استنتاج د کومې قاعدې له مخې له کوم \LR{$!A_j$} (او
  \LR{$!A_k$}) څخه راځي، چې \LR{$j < i$} (او \LR{$k < i$})۔
\end{enumerate}
\end{defn}
\label{olpseg:OLP-0113-B007:end}

\phantomsection\label{olpseg:OLP-0113-B008:start}
دا چې سم اشتقاق څه ګڼل کېږي، په دې پورې اړه لري چې موږ
کومې استنتاجي قاعدې منو (او البته کوم څيزونه بديهي اصول ګڼو)۔
استنتاجي قاعده يوه که-نو وينا ده چې راته وايي د ځينو شرطونو لاندې
په اشتقاق کښې يو ګام~\LR{$!A_i$} سم استنتاجي ګام دے۔
\textbf{د ژباړې د نسخې سمون (OLAX-001):} سرچينې په دې يو ځای کښې
د اې-آی د ميتا-فارمول نښه پرېښې وه؛ تعريف او د همدې پاراګراف نورې
ټولې کارونې هغه د فارمول په توګه ليکي، نو نښه دلته بېرته راوستل شوه۔
\label{olpseg:OLP-0113-B008:end}

\phantomsection\label{olpseg:OLP-0113-B009:start}
\begin{defn}[د استنتاج قاعده]
د \emph{استنتاج قاعده} يو بس شرط ورکوي چې له~\LR{$\Gamma$} څخه په
اشتقاق کښې څه شی سم استنتاجي ګام ګڼل کېږي۔
\end{defn}
\label{olpseg:OLP-0113-B009:end}

\phantomsection\label{olpseg:OLP-0113-B010:start}
د بېلګې په توګه، ځکه چې هره يو-غړيزه لړۍ~\LR{$!A$} چې
\LR{$!A \in \Gamma$} وي، په څرګند ډول اشتقاق ګڼل کېږي، لاندې
ډېره ساده استنتاجي قاعده کېدے شي:
\begin{quote}
  که \LR{$!A \in \Gamma$} وي، نو \LR{$!A$} له~\LR{$\Gamma$} څخه په هر
  اشتقاق کښې تل سم استنتاجي ګام دے۔
\end{quote}
همدارنګه که \LR{$!A$} له بديهي اصولو څخه يو وي، نو يوازې \LR{$!A$} هم
اشتقاق دے؛ ځکه دا هم يوه استنتاجي قاعده ده:
\begin{quote}
  که \LR{$!A$} بديهي اصل وي، نو \LR{$!A$} سم استنتاجي ګام دے۔
\end{quote}
خبره هغه وخت لا په زړه پورې شي چې استنتاجي قاعده هغو فارمولونه
ته رجوع کوي چې تر څېړل شوي ګام وړاندې راغلي وي۔ لاندې قاعده
\emph{مودوس پوننس} بلل کېږي:
\begin{quote}
  که \LR{$!B \lif !A$} او \LR{$!B$} په اشتقاق کښې وړاندې راغلي وي،
  نو~\LR{$!A$} سم استنتاجي ګام دے۔
\end{quote}
که همدا يوازينۍ استنتاجي قاعده وي، نو د اشتقاق پورته تعريف
مو داسې کېږي: \LR{$!A_1$}، \dots،~\LR{$!A_n$} هله او يوازې هله
اشتقاق دے چې د هر \LR{$i \le n$} دپاره له لاندې څخه يو شرط
پوره کېږي:
\begin{enumerate}
\item \LR{$!A_i \in \Gamma$}؛ يا
\item \LR{$!A_i$} يو بديهي اصل دے؛ يا
\item د کوم \LR{$j < i$} دپاره \LR{$!A_j$}، \LR{$!B \lif !A_i$} وي، او د کوم
  \LR{$k < i$} دپاره \LR{$!A_k$}،~\LR{$!B$} وي۔
\end{enumerate}
وروستے بند وايي چې \LR{$!A_i$} له~\LR{$!A_j$} (\LR{$!B \lif !A_i$}) او
\LR{$!A_k$} (\LR{$!B$}) څخه د مودوس پوننس په وسيله راځي۔ که له~\LR{$1$} څخه
تر~\LR{$n$} پورې تېرېدے شو، او هر ځل داسې فارمول~\LR{$!A_i$} ومومو
چې يا په~\LR{$\Gamma$} کښې وي، يا بديهي اصل وي، يا استنتاجي قاعده راته
وايي چې سم استنتاجي ګام دے، نو ټوله لړۍ سم اشتقاق ګڼل کېږي۔
\label{olpseg:OLP-0113-B010:end}

\phantomsection\label{olpseg:OLP-0113-B011:start}
\begin{defn}[د اشتقاق وړتيا]
يو فارمول~\LR{$!A$} هله له~\LR{$\Gamma$} څخه \emph{د اشتقاق وړ} دے،
چې \LR{$\Gamma \Proves !A$} ليکو، چې له~\LR{$\Gamma$} څخه داسې
اشتقاق وي چې په~\LR{$!A$} پاے ته رسېږي۔
\end{defn}
\label{olpseg:OLP-0113-B011:end}

\phantomsection\label{olpseg:OLP-0113-B012:start}
\begin{defn}[قضيې]
فارمول~\LR{$!A$} هله \emph{قضيه} ده چې له تش سټ څخه د~\LR{$!A$} يو
اشتقاق وي۔ که \LR{$!A$} قضيه وي \LR{$\Proves !A$}، او که نۀ وي
\LR{$\Proves/ !A$} ليکو۔
\end{defn}
\label{olpseg:OLP-0113-B012:end}

\phantomsection\label{olpseg:OLP-0114-B004:start}
\olfileid{fol}{axd}{prp}
\label{olpseg:OLP-0114-B004:end}

\phantomsection\label{olpseg:OLP-0114-B005:start}
\olsection{د قضيې نښلوونکو دپاره بديهي اصول او قاعدې}
\label{olpseg:OLP-0114-B005:end}

\phantomsection\label{olpseg:OLP-0114-B006:start}
\begin{defn}[بديهي اصول]
د قضيې نښلوونکو د بديهي اصولو سټ~\LR{$\PAx$} د لاندې ټولو بڼو
فارمولونه لري:
\begin{align}
  & (!A \land !B) \lif !A \ollabel{ax:land1}\\
  & (!A \land !B) \lif !B \ollabel{ax:land2}\\
  & !A \lif (!B \lif (!A \land !B)) \ollabel{ax:land3}\\
  & !A \lif (!A \lor !B) \ollabel{ax:lor1}\\
  & !A \lif (!B \lor !A) \ollabel{ax:lor2}\\
  & (!A \lif !C) \lif ((!B \lif !C) \lif ((!A \lor !B) \lif !C)) \ollabel{ax:lor3}\\
  & !A \lif (!B \lif !A) \ollabel{ax:lif1}\\
  & (!A \lif (!B \lif !C)) \lif ((!A \lif !B) \lif (!A \lif !C)) \ollabel{ax:lif2}\\
  & (!A \lif !B) \lif ((!A \lif \lnot !B) \lif \lnot !A) \ollabel{ax:lnot1}\\
  & \lnot !A \lif (!A \lif !B) \ollabel{ax:lnot2}\\
  & \ltrue \ollabel{ax:ltrue}\\
  & \lfalse \lif !A \ollabel{ax:lfalse1}\\
  & (!A \lif \lfalse) \lif \lnot !A \ollabel{ax:lfalse2}\\
  & \lnot\lnot !A \lif !A \ollabel{ax:dne}
\end{align}
\end{defn}
\label{olpseg:OLP-0114-B006:end}

\phantomsection\label{olpseg:OLP-0114-B007:start}
\begin{defn}[مودوس پوننس]
  که \LR{$!B$} او \LR{$!B \lif !A$} له وړاندې په اشتقاق کښې راغلي
  وي، نو \LR{$!A$} سم استنتاجي ګام دے۔
\end{defn}
\label{olpseg:OLP-0114-B007:end}

\phantomsection\label{olpseg:OLP-0114-B008:start}
د مودوس پوننس قاعده به په ``\MP'' لنډوو۔
\label{olpseg:OLP-0114-B008:end}

\phantomsection\label{olpseg:OLP-0115-B004:start}
\olfileid{fol}{axd}{qua}
\label{olpseg:OLP-0115-B004:end}

\phantomsection\label{olpseg:OLP-0115-B005:start}
\olsection{د کميت ټاکونکو دپاره بديهي اصول او قاعدې}
\label{olpseg:OLP-0115-B005:end}

\phantomsection\label{olpseg:OLP-0115-B006:start}
\begin{defn}[د کميت ټاکونکو بديهي اصول]
پر کميت ټاکونکو واکمن \emph{بديهي اصول} د لاندې ټولو بڼو بېلګې دي:
\begin{align}
\ollabel{ax:q1} & \lforall[x][!B] \lif !B(t), \\
\ollabel{ax:q2} & !B(t) \lif \lexists[x][!B].
\end{align}
د هر تړلي ترم~\LR{$t$} دپاره۔
\end{defn}
\label{olpseg:OLP-0115-B006:end}

\phantomsection\label{olpseg:OLP-0115-B007:start}
\begin{defn}[د کميت ټاکونکو قاعدې]
\item که \LR{$!B \lif !A(a)$} له وړاندې په اشتقاق کښې راغلے وي
  او \LR{$a$} په~\LR{$\Gamma$} يا~\LR{$!B$} کښې نۀ راځي، نو
  \LR{$!B \lif \lforall[x][!A(x)]$} سم استنتاجي ګام دے۔
\item که \LR{$!A(a) \lif !B$} له وړاندې په اشتقاق کښې راغلے وي
  او \LR{$a$} په~\LR{$\Gamma$} يا~\LR{$!B$} کښې نۀ راځي، نو
  \LR{$\lexists[x][!A(x)] \lif !B$} سم استنتاجي ګام دے۔
\end{defn}
\label{olpseg:OLP-0115-B007:end}

\phantomsection\label{olpseg:OLP-0115-B008:start}
له دغو دواړو هره يوه به په ``\QR'' لنډوو۔
\label{olpseg:OLP-0115-B008:end}

\phantomsection\label{olpseg:OLP-0116-B004:start}
\olfileid{fol}{axd}{pro}
\label{olpseg:OLP-0116-B004:end}

\phantomsection\label{olpseg:OLP-0116-B005:start}
\olsection{د اشتقاقونه بېلګې}
\label{olpseg:OLP-0116-B005:end}

\phantomsection\label{olpseg:OLP-0116-B006:start}
\begin{ex}
فرض کړئ غواړو \LR{$(\lnot !D \lor !E) \lif (!D \lif !E)$} ثابت کړو۔
دا په څرګند ډول زموږ د هېڅ بديهي اصل بېلګه نۀ ده، نو د هغې د
اشتقاق کولو دپاره د~\MP{} قاعده کارول پکار دي۔ زموږ يوازينۍ
قاعده مودوس پوننس ده، چې \LR{$!A$} او \LR{$!A \lif !B$} راکړل شي، د~\LR{$!B$}
توجيه راکوي۔ يوه تګلاره دا ده چې \olref[prp]{ax:lor3} داسې وکاروو
چې \LR{$!A$}، \LR{$\lnot !D$}؛ \LR{$!B$}، \LR{$!E$}؛ او \LR{$!C$}، \LR{$!D \lif !E$} وي؛ يعنې
دا بېلګه:
\[
(\lnot !D \lif (!D \lif !E)) \lif ((!E \lif (!D \lif !E)) \lif ((\lnot
!D \lor !E) \lif (!D \lif !E))).
\]
ولې؟ د مودوس پوننس دوه تطبيقونه وروستۍ برخه راکوي، او همدا مو
غوښتې ده۔ په اسانۍ وينو چې \LR{$\lnot !D \lif (!D \lif !E)$} د
\olref[prp]{ax:lnot2}، او \LR{$!E \lif (!D \lif !E)$} د
\olref[prp]{ax:lif1} بېلګه ده۔ نو زموږ اشتقاق دا دے:
\begin{derivation}
  1. &  \LR{$\lnot !D \lif (!D \lif !E)$} & \olref[prp]{ax:lnot2} \\
  2. & \LR{$(\lnot !D \lif (!D \lif !E)) \lif {}$}\\
  &\qquad \LR{$((!E \lif (!D \lif !E)) \lif ((\lnot !D \lor !E) \lif (!D \lif !E)))$} & \olref[prp]{ax:lor3}\\
  3. & \LR{$(!E \lif (!D \lif !E)) \lif ((\lnot !D \lor !E) \lif (!D \lif !E))$} & 1, 2, \MP\\
  4. & \LR{$!E \lif (!D \lif !E)$} & \olref[prp]{ax:lif1}\\
  5. & \LR{$(\lnot !D \lor !E) \lif (!D \lif !E)$} & 3, 4, \MP
\end{derivation}
\end{ex}
\label{olpseg:OLP-0116-B006:end}

\phantomsection\label{olpseg:OLP-0116-B007:start}
\begin{ex}\ollabel{ex:identity}
راځئ د~\LR{$!D \lif !D$} يو اشتقاق ولټوو۔ دا د بديهي اصل بېلګه
نۀ ده، نو په~\MP{} ئې اشتقاق کول پکار دي۔
\olref[prp]{ax:lif1} د~\LR{$!A \lif !B$} په بڼه يو بديهي اصل دے چې
\MP{} پرې پلي کېدے شي۔ د ګټې دپاره، البته، هغه~\LR{$!B$} چې \MP{} ئې په
دې حالت کښې د سم ګام په توګه توجيه کوي بايد~\LR{$!D \lif !D$} وي؛ ځکه
همدا مو غوښتل چې اشتقاق ئې کړو۔ په دې سره \LR{$!A$} هم بايد \LR{$!D$} وي؛
يعنې د \olref[prp]{ax:lif1} دې بېلګې ته کتلے شو:
\[
!D \lif (!D \lif !D)
\]
د~\MP{} د پلي کولو دپاره به د دويمې اړوندې مقدمې، يعنې~\LR{$!A$}،
توجيه هم پکار وي۔ خو زموږ په حالت کښې هغه~\LR{$!D$} دے، او \LR{$!D$} په
يوازې ځان اشتقاق کولے نۀ شو۔ نو بله تګلاره پکار ده۔
\label{olpseg:OLP-0116-B007:end}

\phantomsection\label{olpseg:OLP-0116-B008:start}
بل بديهي اصل چې يوازې~\LR{$\lif$} لري \olref[prp]{ax:lif2} دے، يعنې
\[
(!A \lif (!B \lif !C)) \lif ((!A \lif !B) \lif (!A \lif !C))
\]
د~\MP{} په دوه ځله پلي کولو وروستي دننني شرطي ته رسېدے شو۔ بيا دا
معنٰی لري چې د \olref[prp]{ax:lif2} داسې بېلګه غواړو چې
\LR{$!A \lif !C$} پکې \LR{$!D \lif !D$}، يعنې زموږ موخه فارمول وي۔
نو \LR{$!A$} او \LR{$!C$} دواړه~\LR{$!D$} دي۔ \LR{$!B$} څنګه وټاکو چې دواړه
\LR{$!A \lif (!B \lif !C)$} او \LR{$!A \lif !B$}، يعنې زموږ په حالت کښې
\LR{$!D \lif (!B \lif !D)$} او \LR{$!D \lif !B$}، هم د اشتقاق وړ وي؟
له دغو لومړنے، هر~\LR{$!B$} چې وټاکو، لا د \olref[prp]{ax:lif1} بېلګه
ده۔ او که \LR{$!B$}، \LR{$(!D \lif !D)$} وي، نو \LR{$!D \lif !B$} به هم د
\olref[prp]{ax:lif1} بله بېلګه وي۔ نو زموږ اشتقاق دا دے:
\begin{derivation}
  1. & \LR{$!D \lif ((!D \lif !D) \lif !D)$} & \olref[prp]{ax:lif1}\\
  2. & \LR{$(!D \lif ((!D \lif !D) \lif !D)) \lif {}$}\\
  & \qquad \LR{$((!D \lif (!D \lif !D)) \lif (!D \lif !D))$} & \olref[prp]{ax:lif2}\\
  3. & \LR{$(!D \lif (!D \lif !D)) \lif (!D \lif !D)$} & 1, 2, \MP\\
  4. & \LR{$!D \lif (!D \lif !D)$} & \olref[prp]{ax:lif1}\\
  5. & \LR{$!D \lif !D$} & 3, 4, \MP
\end{derivation}
\end{ex}
\label{olpseg:OLP-0116-B008:end}

\phantomsection\label{olpseg:OLP-0116-B009:start}
\begin{ex}\ollabel{ex:chain}
کله کله غواړو وښيو چې له ځينو نورو فارمولونه~\LR{$\Gamma$} څخه د کوم
فارمول يو اشتقاق شته۔ مثلاً راځئ وښيو چې له
\LR{$\Gamma = \{!A \lif !B, !B \lif !C\}$} څخه
\LR{$!A \lif !C$}، اشتقاق کولے شو۔
\begin{derivation}
  1. & \LR{$!A \lif !B$} & \Hyp\\
  2. & \LR{$!B \lif !C$} & \Hyp\\
  3. & \LR{$(!B \lif !C) \lif (!A \lif (!B \lif !C))$} & \olref[prp]{ax:lif1} \\
  4. & \LR{$!A \lif (!B \lif !C)$} & 2, 3, \MP\\
  5. & \LR{$(!A \lif (!B \lif !C)) \lif {}$}\\
  & \qquad \LR{$((!A \lif !B) \lif (!A \lif !C))$} & \olref[prp]{ax:lif2}\\
  6. & \LR{$((!A \lif !B) \lif (!A \lif !C))$} & 4, 5, \MP\\
  7. & \LR{$!A \lif !C$} & 1, 6, \MP
\end{derivation}
هغه کرښې چې ``\Hyp'' (د ``فرضیې'' دپاره) پرې ليکل شوي دي، ښيي چې
په هغې کرښه فارمول د~\LR{$\Gamma$} يو غړے دے۔
\end{ex}
\label{olpseg:OLP-0116-B009:end}

\phantomsection\label{olpseg:OLP-0116-B010:start}
\begin{prop}
  \ollabel{prop:chain} که \LR{$\Gamma \Proves !A \lif !B$} او \LR{$\Gamma
  \Proves !B \lif !C$} وي، نو \LR{$\Gamma \Proves !A \lif !C$}
\end{prop}
\label{olpseg:OLP-0116-B010:end}

\phantomsection\label{olpseg:OLP-0116-B011:start}
\begin{proof}
  فرض کړئ \LR{$\Gamma \Proves !A \lif !B$} او \LR{$\Gamma \Proves !B \lif
  !C$}۔ بيا له~\LR{$\Gamma$} څخه د~\LR{$!A \lif !B$} يو اشتقاق شته؛
  او له~\LR{$\Gamma$} څخه د~\LR{$!B \lif !C$} يو اشتقاق هم شته۔ د
  هغو په يو بل پسې نښلولو ئې په يو اشتقاق کښې يو ځاے کړئ۔
  اوس د مخکښې بېلګې د~اشتقاق ۳--۷ کرښې ورزياتې کړئ۔ دا له~\LR{$\Gamma$} څخه د
  \LR{$!A \lif !C$} يو اشتقاق دے---او همدا د نوي
  اشتقاق وروستۍ کرښه ده۔ پام وکړئ چې د ۴مې او ۷مې کرښې
  توجيه لا هم سمه پاتې کېږي که د ۱مې کرښې حواله د~\LR{$!A \lif !B$}
  د~اشتقاق د وروستۍ کرښې په حواله، او د ۲مې کرښې حواله د
  \LR{$!B \lif !C$} د~اشتقاق د وروستۍ کرښې په حواله بدله شي۔
\end{proof}
\label{olpseg:OLP-0116-B011:end}

\phantomsection\label{olpseg:OLP-0116-B012:start}
\begin{prob} 
له بديهي اصولو څخه د اشتقاقونه په وړاندې کولو وښيئ چې لاندې
خبرې سمې دي:
\begin{enumerate} 
\item \LR{$(!A \land !B) \lif (!B \land !A)$}
\item \LR{$((!A \land !B) \lif !C) \lif (!A \lif (!B \lif !C))$}
\item \LR{$\lnot(!A \lor !B) \lif \lnot !A$}
\end{enumerate} 
\end{prob}
\label{olpseg:OLP-0116-B012:end}

\phantomsection\label{olpseg:OLP-0117-B004:start}
\olfileid{fol}{axd}{prq}
\label{olpseg:OLP-0117-B004:end}

\phantomsection\label{olpseg:OLP-0117-B005:start}
\olsection{له کميت ټاکونکو سره اشتقاقونه}
\label{olpseg:OLP-0117-B005:end}

\phantomsection\label{olpseg:OLP-0117-B006:start}
\begin{ex}
راځئ د~\LR{$(\lforall[x][!A(x)] \land
\lforall[y][!B(y)]) \lif \lforall[x][(!A(x) \land !B(x))]$} يو
اشتقاق ورکړو۔
\label{olpseg:OLP-0117-B006:end}

\phantomsection\label{olpseg:OLP-0117-B007:start}
لومړے پام وکړئ چې
\begin{align*}
  (\lforall[x][!A(x)] \land \lforall[y][!B(y)]) & \lif \lforall[x][!A(x)]\\
  \intertext{\RL{د \olref[prp]{ax:land1} يوه بېلګه ده، او}}
  \lforall[x][!A(x)] & \lif !A(a) \\
  \intertext{\RL{د \olref[qua]{ax:q1} بېلګه ده۔ نو د \olref[pro]{prop:chain} له مخې پوهېږو چې}}
  (\lforall[x][!A(x)] \land \lforall[y][!B(y)]) & \lif !A(a) \\
  \intertext{\RL{د اشتقاق وړ دے۔ همدارنګه، ځکه}}
  (\lforall[x][!A(x)] \land \lforall[y][!B(y)]) & \lif \lforall[y][!B(y)] \qquad\text{\RL{او}}\\
    \lforall[y][!B(y)] & \lif !B(a)\\
    \intertext{\RL{په ترتيب د \olref[prp]{ax:land2} او \olref[qua]{ax:q1} بېلګې دي،}}
    (\lforall[x][!A(x)] \land \lforall[y][!B(y)]) & \lif !B(a)\\
    \intertext{\RL{د \olref[pro]{prop:chain} په وسيله د ثبوت وړ دے۔ د \olref[prp]{ax:land3} د يوې مناسبې بېلګې او د~\MP{} د دوو تطبيقونو په کارولو وينو چې}}
    (\lforall[x][!A(x)] \land \lforall[y][!B(y)]) & \lif (!A(a) \land !B(a))\\
    \intertext{\RL{د ثبوت وړ دے۔ اوس \QR{} پلي کولے شو او دا ترلاسه کوو:}}
(\lforall[x][!A(x)] \land \lforall[y][!B(y)]) & \lif \lforall[x][(!A(x) \land !B(x))].
\end{align*}
\end{ex}
\label{olpseg:OLP-0117-B007:end}

\phantomsection\label{olpseg:OLP-0118-B004:start}
\olfileid{fol}{axd}{ptn}
\label{olpseg:OLP-0118-B004:end}

\phantomsection\label{olpseg:OLP-0118-B005:start}
\olsection{ثبوت-تيوريکي مفاهيم}
\label{olpseg:OLP-0118-B005:end}

\phantomsection\label{olpseg:OLP-0118-B006:start}
\begin{explain}
لکه څنګه چې مو يو شمېر مهم معنايي مفاهيم
(اعتبار، معنايي لازميت او د صدق وړتيا)
تعريف کړي دي، اوس ورسره سم \emph{ثبوت-تيوريکي مفاهيم} تعريفوو۔
دا د جوړښتونه په دننه کښې د تړلي فارمولونه پوره کېدو ته په
رجوع نۀ، بلکې د ځينو فارمولونو د اشتقاق وړتيا يا
د اشتقاق نۀ وړتيا ته په رجوع تعريفېږي۔ دا يوه مهمه موندنه وه چې
دغه مفاهيم يو شان راځي۔ د هغو يو شان راتلل د \emph{سموالي} او
\emph{بشپړوالي د قضيو} منځپانګه ده۔
\end{explain}
\label{olpseg:OLP-0118-B006:end}

\phantomsection\label{olpseg:OLP-0118-B007:start}
\begin{defn}[د اشتقاق وړتيا]
فارمول~\LR{$!A$} هله له~\LR{$\Gamma$} څخه \emph{د اشتقاق وړ} دے،
چې \LR{$\Gamma \Proves !A$} ليکو، چې له~\LR{$\Gamma$} څخه داسې
اشتقاق وي چې په~\LR{$!A$} پاے ته رسېږي۔
\end{defn}
\label{olpseg:OLP-0118-B007:end}

\phantomsection\label{olpseg:OLP-0118-B008:start}
\begin{defn}[قضيې]
فارمول~\LR{$!A$} هله \emph{قضيه} ده چې له تش سټ څخه د~\LR{$!A$} يو
اشتقاق وي۔ که \LR{$!A$} قضيه وي \LR{$\Proves !A$}، او که نۀ وي
\LR{$\Proves/ !A$} ليکو۔
\end{defn}
\label{olpseg:OLP-0118-B008:end}

\phantomsection\label{olpseg:OLP-0118-B009:start}
\begin{defn}[سازګاري]
د فارمولونه يو سټ~\LR{$\Gamma$} هله او يوازې هله \emph{سازګار} دے
چې \LR{$\Gamma\Proves/ \lfalse$}؛ که داسې نۀ وي \emph{ناسازګار} دے۔
\end{defn}
\label{olpseg:OLP-0118-B009:end}

\phantomsection\label{olpseg:OLP-0118-B010:start}
\begin{prop}[انعکاس]
\ollabel{prop:reflexivity}
که \LR{$!A \in \Gamma$} وي، نو \LR{$\Gamma \Proves !A$}۔
\end{prop}
\label{olpseg:OLP-0118-B010:end}

\phantomsection\label{olpseg:OLP-0118-B011:start}
\begin{proof}
  يوازې فارمول~\LR{$!A$} له~\LR{$\Gamma$} څخه د~\LR{$!A$} يو
  اشتقاق دے۔
\end{proof}
\label{olpseg:OLP-0118-B011:end}

\phantomsection\label{olpseg:OLP-0118-B012:start}
\begin{prop}[زياتېدنه]
\ollabel{prop:monotonicity}
که \LR{$\Gamma \subseteq \Delta$} او \LR{$\Gamma \Proves !A$} وي، نو
\LR{$\Delta \Proves !A$}۔
\end{prop}
\label{olpseg:OLP-0118-B012:end}

\phantomsection\label{olpseg:OLP-0118-B013:start}
\begin{proof}
له~\LR{$\Gamma$} څخه د~\LR{$!A$} هر اشتقاق له~\LR{$\Delta$} څخه هم د~\LR{$!A$}
يو اشتقاق دے۔
\end{proof}
\label{olpseg:OLP-0118-B013:end}

\phantomsection\label{olpseg:OLP-0118-B014:start}
\begin{prop}[انتقال]
\ollabel{prop:transitivity}
که \LR{$\Gamma \Proves !A$} او \LR{$\{!A\} \cup \Delta \Proves !B$} وي،
نو \LR{$\Gamma \cup \Delta \Proves !B$}۔
\end{prop}
\label{olpseg:OLP-0118-B014:end}

\phantomsection\label{olpseg:OLP-0118-B015:start}
\begin{proof}
  فرض کړئ \LR{$\{!A\} \cup \Delta \Proves !B$}۔ بيا له
  \LR{$\{!A\} \cup \Delta$} څخه يو اشتقاق
  \LR{$!B_1$}، \dots، \LR{$!B_l = !B$} شته۔ په دې اشتقاق کښې ځينې
  ګامونه به د هغې قاعدې له مخې سم وي چې کومې مخکښې کرښې
  \LR{$!B_i = !A$} ته رجوع کوي۔ د فرض له مخې له~\LR{$\Gamma$} څخه د~\LR{$!A$}
  يو اشتقاق شته؛ يعنې داسې اشتقاق
  \LR{$!A_1$}، \dots، \LR{$!A_k = !A$} چې هر \LR{$!A_i$} يا بديهي اصل وي، يا
  د~\LR{$\Gamma$} يو غړے وي، يا د استنتاج د قاعدې له مخې سم
  وي۔ اوس دا لړۍ په پام کښې ونيسئ:
  \[
  !A_1, \dots, !A_k = !A, !B_1, \dots, !B_l = !B.
  \]
  دا له~\LR{$\Gamma \cup \Delta$} څخه د~\LR{$!B$} سم اشتقاق دے؛ ځکه
  هر \LR{$!B_i = !A$} اوس د هماغې قاعدې له مخې توجيه کېږي چې
  \LR{$!A_k = !A$} توجيه کوي۔
  \textbf{د ژباړې د نسخې سمون (OLAX-002):} سرچينې په وروستۍ جمله
  کښې يوازې له بي-آی څخه د ميتا-فارمول نښه پرېښې وه؛ هماغه کرښه
  پورته او په نښلول شوې لړۍ کښې په نښه لرلې بڼه تعريف شوې ده۔
\end{proof}
\label{olpseg:OLP-0118-B015:end}

\phantomsection\label{olpseg:OLP-0118-B016:start}
پام وکړئ چې په ځانګړي ډول دا معنٰی لري: که
\LR{$\Gamma \Proves !A$} او \LR{$!A \Proves !B$} وي، نو
\LR{$\Gamma \Proves !B$}۔ له دې څخه دا هم راځي چې که
\LR{$!A_1, \dots, !A_n \Proves !B$} او د هر~\LR{$i$} دپاره
\LR{$\Gamma \Proves !A_i$} وي، نو \LR{$\Gamma \Proves !B$}۔
\label{olpseg:OLP-0118-B016:end}

\phantomsection\label{olpseg:OLP-0118-B017:start}
\begin{prop}
\ollabel{prop:incons}
\LR{$\Gamma$} هله او يوازې هله ناسازګار دے چې د هر~\LR{$!A$} دپاره
\LR{$\Gamma \Proves !A$} وي۔
\end{prop}
\label{olpseg:OLP-0118-B017:end}

\phantomsection\label{olpseg:OLP-0118-B018:start}
\begin{proof}
تمرين۔
\end{proof}
\label{olpseg:OLP-0118-B018:end}

\phantomsection\label{olpseg:OLP-0118-B019:start}
\begin{prob}
\olref[fol][axd][ptn]{prop:incons} ثابت کړئ۔
\end{prob}
\label{olpseg:OLP-0118-B019:end}



\phantomsection\label{olpseg:OLP-0118-B021:start}
\begin{prop}[تراکم]
\ollabel{prop:proves-compact}
  \begin{enumerate}
  \item که \LR{$\Gamma \Proves !A$} وي، نو داسې متناهي فرعي سټ
    \LR{$\Gamma_0 \subseteq \Gamma$} شته چې \LR{$\Gamma_0 \Proves !A$}۔
  \item که د~\LR{$\Gamma$} هر متناهي فرعي سټ سازګار وي، نو \LR{$\Gamma$}
    سازګار دے۔
  \end{enumerate}
\end{prop}
\label{olpseg:OLP-0118-B021:end}

\phantomsection\label{olpseg:OLP-0118-B022:start}
\begin{proof}
  \begin{enumerate}
    \item که \LR{$\Gamma \Proves !A$} وي، نو د فارمولونه داسې متناهي
      لړۍ \LR{$!A_1$}، \dots،~\LR{$!A_n$} شته چې \LR{$!A \ident !A_n$}؛ او هر
      \LR{$!A_i$} يا منطقي بديهي اصل وي، يا د~\LR{$\Gamma$} يو غړے
      وي، يا د مودوس پوننس له مخې له مخکښو فارمولونه څخه راځي۔
      \LR{$\Gamma_0$} د هغو \LR{$!A_i$} سټ ونيسئ چې په~\LR{$\Gamma$} کښې دي۔ بيا
      همدا اشتقاق له~\LR{$\Gamma_0$} څخه هم يو اشتقاق
      دے، نو \LR{$\Gamma_0 \Proves !A$}۔
    \item دا د ځانګړي حالت~\LR{$!A \ident \lfalse$} دپاره د~(۱) مقابل
      عکس دے۔
\end{enumerate}
\end{proof}
\label{olpseg:OLP-0118-B022:end}

\phantomsection\label{olpseg:OLP-0119-B005:start}
\olfileid{fol}{axd}{ded}
\label{olpseg:OLP-0119-B005:end}

\phantomsection\label{olpseg:OLP-0119-B006:start}
\olsection{د استنتاج قضيه}
\label{olpseg:OLP-0119-B006:end}

\phantomsection\label{olpseg:OLP-0119-B007:start}
لکه چې مو وليدل، په بديهي نظام کښې اشتقاقونه ورکول ستونزمن
دي، او اشتقاقونه موندل هم ګرانېږي۔ د فارمولونه اوږدې لړۍ
په رښتيا ليکلو پر ځاے، زياتره دا استدلال اسانه وي چې داسې
اشتقاقونه شته؛ د دې دپاره څو ساده پايلې کاروو۔ درې داسې پايلې
مو لا ثابتې کړې دي: \olref[ptn]{prop:reflexivity} وايي کله چې
\LR{$!A \in \Gamma$} پوهېږو، تل \LR{$\Gamma \Proves !A$} ويلے شو۔
\olref[ptn]{prop:monotonicity} وايي که \LR{$\Gamma \Proves !A$} وي، نو
\LR{$\Gamma \cup \{!B\} \Proves !A$} هم دے۔ او
\olref[ptn]{prop:transitivity} ښيي که \LR{$\Gamma \Proves !A$} او
\LR{$!A \Proves !B$} وي، نو \LR{$\Gamma \Proves !B$}۔ دا د مودوس پوننس يوه
بله ساده، ``ميتا'' بڼه ده:
\label{olpseg:OLP-0119-B007:end}

\phantomsection\label{olpseg:OLP-0119-B008:start}
\begin{prop}
  \ollabel{prop:mp} که \LR{$\Gamma \Proves !A$} او
  \LR{$\Gamma \Proves !A \lif !B$} وي، نو \LR{$\Gamma \Proves !B$}۔
\end{prop}
\label{olpseg:OLP-0119-B008:end}

\phantomsection\label{olpseg:OLP-0119-B009:start}
\begin{proof}
  موږ \LR{$\{!A, !A \lif !B\} \Proves !B$} لرو:
  \begin{derivation}
    1. & \LR{$!A$} & فرض\\
    2. & \LR{$!A \lif !B$} & فرض\\
    3. & \LR{$!B$} & 1, 2, مودوس پوننس
  \end{derivation}
  د \olref[ptn]{prop:transitivity} له مخې \LR{$\Gamma \Proves !B$}۔
\end{proof}
\label{olpseg:OLP-0119-B009:end}

\phantomsection\label{olpseg:OLP-0119-B010:start}
په دې برخه کښې تر ټولو مهمه پايله چې کاروو ئې د استنتاج قضيه ده:
\label{olpseg:OLP-0119-B010:end}

\phantomsection\label{olpseg:OLP-0119-B011:start}
\begin{thm}[د استنتاج قضيه]
\ollabel{thm:deduction-thm} \LR{$\Gamma \cup \{!A\} \Proves !B$} هله
او يوازې هله چې \LR{$\Gamma \Proves !A \lif !B$}۔
\end{thm}
\label{olpseg:OLP-0119-B011:end}

\phantomsection\label{olpseg:OLP-0119-B012:start}
\begin{proof}
د ``که'' لوری سمدستي دے۔ که \LR{$\Gamma \Proves !A \lif !B$} وي، د
\olref[ptn]{prop:monotonicity} له مخې
\LR{$\Gamma \cup \{!A\} \Proves !A \lif !B$} هم دے۔ همدارنګه د
\olref[ptn]{prop:reflexivity} له مخې
\LR{$\Gamma \cup \{!A\} \Proves !A$}۔ نو د \olref{prop:mp} له مخې
\LR{$\Gamma \cup \{!A\} \Proves !B$}۔
\label{olpseg:OLP-0119-B012:end}

\phantomsection\label{olpseg:OLP-0119-B013:start}
د ``يوازې که'' لوري دپاره له~\LR{$\Gamma \cup \{!A\}$} څخه د~\LR{$!B$} د
اشتقاق پر اوږدوالي استقرا کوو۔
\label{olpseg:OLP-0119-B013:end}

\phantomsection\label{olpseg:OLP-0119-B014:start}
د استقرا د بنسټ دپاره دعوه د اوږدوالي~\LR{$1$} د هر اشتقاق
دپاره ثابتوو۔ له~\LR{$\Gamma \cup \{!A\}$} څخه د~\LR{$!B$} د اوږدوالي~\LR{$1$}
اشتقاق يوازې له~\LR{$!B$} څخه جوړ وي؛ او که سم وي، يا
\LR{$!B \in \Gamma \cup \{!A\}$} وي يا \LR{$!B$} بديهي اصل وي۔
\textbf{د ژباړې د نسخې سمون (OLAX-003):} سرچينې په دې غړيتوب کښې
له اې-متا فارمول څخه مخکښې يوازې د غړيتوب نښه ليکلې وه؛ د جملې
فاعل دلته د بشپړ فارمول په دننه کښې راوستل شو۔
که \LR{$!B \in \Gamma$} وي يا بديهي اصل وي، نو \LR{$\Gamma \Proves !B$}۔
همدارنګه د \olref[prp]{ax:lif1} له مخې
\LR{$\Gamma \Proves !B \lif (!A \lif !B)$}، او \olref{prop:mp} موږ ته
\LR{$\Gamma \Proves !A \lif !B$} راکوي۔ که \LR{$!B \in \{ !A\}$} وي، نو
\LR{$\Gamma \Proves !A \lif !B$}؛ ځکه وروستۍ تړلے فارمول
\LR{$!A \lif !B$} د~\LR{$!A \lif !A$} هماغه جمله ده، او په
\olref[pro]{ex:identity} کښې مو هغه اشتقاق کړې ده۔
\label{olpseg:OLP-0119-B014:end}

\phantomsection\label{olpseg:OLP-0119-B015:start}
د استقرا په پړاو کښې فرض کړئ له~\LR{$\Gamma \cup \{!A\}$} څخه د~\LR{$!B$}
يو اشتقاق په داسې ګام~\LR{$!B$} پاے ته رسېږي چې په مودوس پوننس
توجيه شوے دے۔ (که په مودوس پوننس توجيه شوے نۀ وي، نو
\LR{$!B \in \Gamma$}، \LR{$!B \ident !A$}، يا \LR{$!B$} بديهي اصل دے، او د استقرا
د بنسټ هماغه استدلال پرې پلي کېږي۔) بيا د کوم فارمول~\LR{$!C$}
دپاره په اشتقاق کښې ځينې مخکښې ګامونه \LR{$!C \lif !B$} او
\LR{$!C$} دي؛ يعنې \LR{$\Gamma \cup \{!A\} \Proves !C \lif !B$} او
\LR{$\Gamma \cup \{!A\} \Proves !C$}، او اړوند اشتقاقونه لنډ دي،
نو د استقرا فرضيه پرې پلي کېږي۔ ځکه دواړه دا لرو:
\begin{align*}
  & \Gamma \Proves !A \lif (!C \lif !B); \\
  & \Gamma \Proves !A \lif !C.
\end{align*}
خو د \olref[prp]{ax:lif2} له مخې دا هم لرو:
\[
\Gamma \Proves (!A \lif (!C \lif !B)) \lif
((!A\lif !C)  \lif (!A \lif !B)),
\]
او د \olref{prop:mp} دوه تطبيقونه، لکه چې پکار ده،
\LR{$\Gamma \Proves !A \lif !B$} راکوي۔
\end{proof}
\label{olpseg:OLP-0119-B015:end}

\phantomsection\label{olpseg:OLP-0119-B016:start}
پام وکړئ چې \olref[prp]{ax:lif1} او \olref[prp]{ax:lif2} په عين
دې موخه ټاکل شوي وو چې د استنتاج قضيه سمه وي۔
\label{olpseg:OLP-0119-B016:end}

\phantomsection\label{olpseg:OLP-0119-B017:start}
لاندې د~د اشتقاق وړتيا په اړه ځينې ګټور حقيقتونه دي چې د تمرينونو
په توګه ئې پرېږدو۔
\label{olpseg:OLP-0119-B017:end}

\phantomsection\label{olpseg:OLP-0119-B018:start}
\begin{prop}
  \ollabel{prop:derivfacts}
\begin{enumerate}
\item \LR{$\Proves (!A \lif !B) \lif ((!B \lif !C)
  \lif (!A \lif !C))$}؛ \ollabel{derivfacts:a}
\textbf{د ژباړې د نسخې سمون (OLAX-004):} سرچينې د لومړۍ دعوې د
بهرني وروستي قوس د تړلو نښه پرېښې وه؛ دلته د پرانيستو قوسونو سره
سمه يوه تړونکې نښه بېرته راوستل شوې ده۔
\item که \LR{$\Gamma \cup \{ \lnot !A\}
  \Proves \lnot !B$} وي، نو \LR{$\Gamma \cup \{ !B\} \Proves
  !A$} (مقابل عکس)؛ \ollabel{derivfacts:b}
\item \LR{$\{ !A, \lnot!A\} \Proves
    !B$} (له کاذبۍ هر څه، انفجار)؛ \ollabel{derivfacts:c}
\item \LR{$\{ \lnot\lnot!A\} \Proves
  !A$} (دوه ګونې نفي حذف)؛\ollabel{derivfacts:d}
\item که \LR{$\Gamma \Proves \lnot\lnot!A$} وي، نو \LR{$\Gamma \Proves
  !A$}؛\ollabel{derivfacts:e}
\end{enumerate}
\end{prop}
\label{olpseg:OLP-0119-B018:end}

\phantomsection\label{olpseg:OLP-0119-B019:start}
\begin{prob}
  \olref[fol][axd][ded]{prop:derivfacts} ثابت کړئ۔
\end{prob}
\label{olpseg:OLP-0119-B019:end}

\phantomsection\label{olpseg:OLP-0120-B004:start}
\olfileid{fol}{axd}{ddq}
\label{olpseg:OLP-0120-B004:end}

\phantomsection\label{olpseg:OLP-0120-B005:start}
\olsection{له کميت ټاکونکو سره د استنتاج قضيه}
\label{olpseg:OLP-0120-B005:end}

\phantomsection\label{olpseg:OLP-0120-B006:start}
\begin{thm}[د استنتاج قضيه]
\ollabel{thm:deduction-thm-q} که
\LR{$\Gamma \cup \{!A\} \Proves !B$} وي، نو
\LR{$\Gamma \Proves !A \lif !B$}۔
\end{thm}
\label{olpseg:OLP-0120-B006:end}

\phantomsection\label{olpseg:OLP-0120-B007:start}
\begin{proof}
يو ځل بيا له~\LR{$\Gamma \cup \{!A\}$} څخه د~\LR{$!B$} د اشتقاق
پر اوږدوالي استقرا کوو۔
\label{olpseg:OLP-0120-B007:end}

\phantomsection\label{olpseg:OLP-0120-B008:start}
د استقرا د بنسټ ثبوت د
\olref[ded]{thm:deduction-thm} په ثبوت کښې له بنسټ سره يو شان دے۔
\label{olpseg:OLP-0120-B008:end}

\phantomsection\label{olpseg:OLP-0120-B009:start}
د استقرا په پړاو کښې بيا فرض کړئ له~\LR{$\Gamma \cup \{!A\}$} څخه د
\LR{$!B$} اشتقاق په داسې ګام~\LR{$!B$} پاے ته رسېږي چې په استنتاجي
قاعده توجيه شوے دے۔ که استنتاجي قاعده مودوس پوننس وي، د
\olref[ded]{thm:deduction-thm} د ثبوت په شان مخکښې ځو۔ که قاعده
\QR{} وي، پوهېږو چې \LR{$!B \ident !C \lif \lforall[x][!D(x)]$} او د
\LR{$!C \lif !D(a)$} بڼې فارمول په اشتقاق کښې له مخکښې
راغلے دے، چې \LR{$a$} په~\LR{$!C$}، \LR{$!A$} يا~\LR{$\Gamma$} کښې نۀ راځي۔ نو دا لرو:
\begin{align*}
  \Gamma \cup \{!A\} & \Proves !C \lif !D(a),\\
  \intertext{\RL{او د استقرا فرضيه پلي کېږي؛ يعنې دا لرو:}}
    \Gamma & \Proves !A \lif (!C \lif !D(a)).\\
  \intertext{\RL{له دې څخه}}
  & \Proves (!A \lif (!C \lif !D(a))) \lif ((!A \land !C) \lif !D(a))\\
  \intertext{\RL{او مودوس پوننس سره ترلاسه کوو:}}
  \Gamma & \Proves (!A \land !C) \lif !D(a).\\
  \intertext{\RL{ځکه د ځانګړي متغير شرط لا هم پلي کېږي، دې اشتقاق ته په \QR{} توجيه شوے ګام ورزياتولے شو، او ترلاسه کوو:}}
    \Gamma & \Proves (!A \land !C) \lif \lforall[x][!D(x)].\\
    \intertext{\RL{دا هم لرو:}}
    & \Proves ((!A \land !C) \lif \lforall[x][!D(x)]) \lif (!A \lif (!C \lif \lforall[x][!D(x)])),\\
    \intertext{\RL{نو د مودوس پوننس له مخې،}}
    \Gamma & \Proves !A \lif (!C \lif \lforall[x][!D(x)]),
\end{align*}
\textbf{د ژباړې د نسخې سمون (OLAX-005):} سرچينې د پورته ماقبل
وروستۍ بديهي بېلګه کښې د بهرني شرطي يو تړونکے قوس پرېښے و؛ هغه
قوس دلته د بشپړې شېما دپاره بېرته راوستل شو۔
يعنې \LR{$\Gamma \Proves !A \lif !B$}۔
\textbf{د ژباړې د نسخې سمون (OLAX-006):} سرچينې وروستۍ پايله يوازې
د بي د ثبوت په توګه ليکلې وه؛ خو د قضیې موخه او سمدستي مخکښې کرښه
د اې-که-بي ثبوت ورکوي، نو هماغه بشپړه پايله راوستل شوې ده۔
\label{olpseg:OLP-0120-B009:end}

\phantomsection\label{olpseg:OLP-0120-B010:start}
هغه حالت د تمرين په توګه پرېږدو چې \LR{$!B$} د~\QR{} قاعدې په وسيله
توجيه کېږي، خو د~\LR{$\lexists[x][!D(x)] \lif !C$} بڼه لري۔
\end{proof}
\label{olpseg:OLP-0120-B010:end}

\phantomsection\label{olpseg:OLP-0120-B011:start}
\begin{prob}
  د~\olref[fol][axd][ddq]{thm:deduction-thm-q} ثبوت بشپړ کړئ۔
\end{prob}
\label{olpseg:OLP-0120-B011:end}

\phantomsection\label{olpseg:OLP-0121-B004:start}
\olfileid{fol}{axd}{prv}
\label{olpseg:OLP-0121-B004:end}

\phantomsection\label{olpseg:OLP-0121-B005:start}
\olsection{د اشتقاق وړتيا او سازګاري}
\label{olpseg:OLP-0121-B005:end}

\phantomsection\label{olpseg:OLP-0121-B006:start}
اوس به د~د اشتقاق وړتيا اړيکې يو شمېر خاصيتونه ثابت کړو۔ دا په
خپلواکه توګه په زړه پورې دي، خو هر يو به د بشپړوالي د قضیې په
ثبوت کښې ونډه ولري۔
\label{olpseg:OLP-0121-B006:end}

\phantomsection\label{olpseg:OLP-0121-B007:start}
\begin{prop}\ollabel{prop:provability-contr}
  که \LR{$\Gamma \Proves !A$} او \LR{$\Gamma \cup \{!A\}$} ناسازګار وي، نو
  \LR{$\Gamma$} ناسازګار دے۔
\end{prop}
\label{olpseg:OLP-0121-B007:end}

\phantomsection\label{olpseg:OLP-0121-B008:start}
\begin{proof}
  که \LR{$\Gamma \cup \{!A\}$} ناسازګار وي، نو
  \LR{$\Gamma \cup \{!A\} \Proves \lfalse$}۔ د
  \olref[ptn]{prop:reflexivity} له مخې د هر \LR{$!B \in \Gamma$} دپاره
  \LR{$\Gamma \Proves !B$}۔ د فرض له مخې \LR{$\Gamma \Proves !A$} هم دے،
  نو د هر \LR{$!B \in \Gamma \cup \{!A\}$} دپاره
  \LR{$\Gamma \Proves !B$}۔ د \olref[ptn]{prop:transitivity} له مخې
  \LR{$\Gamma \Proves \lfalse$}؛ يعنې \LR{$\Gamma$} ناسازګار دے۔
\end{proof}
\label{olpseg:OLP-0121-B008:end}

\phantomsection\label{olpseg:OLP-0121-B009:start}
\begin{prop}
\ollabel{prop:prov-incons}
\LR{$\Gamma \Proves !A$} هله او يوازې هله چې
\LR{$\Gamma \cup \{\lnot !A\}$} ناسازګار وي۔
\end{prop}
\label{olpseg:OLP-0121-B009:end}

\phantomsection\label{olpseg:OLP-0121-B010:start}
\begin{proof}
لومړے فرض کړئ \LR{$\Gamma \Proves !A$}۔ د
\olref[ptn]{prop:monotonicity} له مخې
\LR{$\Gamma \cup \{\lnot !A\} \Proves !A$}۔ د
\olref[ptn]{prop:reflexivity} له مخې
\LR{$\Gamma \cup \{\lnot !A\} \Proves \lnot !A$}۔ د
\olref[prp]{ax:lnot2} له مخې
\LR{$\Proves \lnot !A \lif (!A \lif \lfalse)$} هم لرو۔ نو د
\olref[ded]{prop:mp} په دوه ځله پلي کولو
\LR{$\Gamma \cup \{\lnot !A\} \Proves \lfalse$} لرو۔
\label{olpseg:OLP-0121-B010:end}

\phantomsection\label{olpseg:OLP-0121-B011:start}
اوس فرض کړئ \LR{$\Gamma \cup \{\lnot !A\}$} ناسازګار دے؛ يعنې
\LR{$\Gamma \cup \{\lnot !A\} \Proves \lfalse$}۔ د استنتاج د قضیې له
مخې \LR{$\Gamma \Proves \lnot !A \lif \lfalse$}۔ د
\olref[prp]{ax:lfalse2} له مخې
\LR{$\Gamma \Proves (\lnot !A \lif \lfalse) \lif \lnot\lnot !A$}، نو
د \olref[ded]{prop:mp} له مخې \LR{$\Gamma \Proves \lnot\lnot !A$}۔
ځکه \LR{$\Gamma \Proves \lnot\lnot !A \lif !A$}
(\olref[prp]{ax:dne})، د \olref[ded]{prop:mp} په بيا پلي کولو
\LR{$\Gamma \Proves !A$} لرو۔
\end{proof}
\label{olpseg:OLP-0121-B011:end}

\phantomsection\label{olpseg:OLP-0121-B012:start}
\begin{prob}
ثابته کړئ چې \LR{$\Gamma \Proves \lnot !A$} هله او يوازې هله چې
\LR{$\Gamma \cup \{!A\}$} ناسازګار وي۔
\end{prob}
\label{olpseg:OLP-0121-B012:end}

\phantomsection\label{olpseg:OLP-0121-B013:start}
\begin{prop}\ollabel{prop:explicit-inc}
  که \LR{$\Gamma \Proves !A$} او \LR{$\lnot !A \in \Gamma$} وي، نو
  \LR{$\Gamma$} ناسازګار دے۔
\end{prop}
\label{olpseg:OLP-0121-B013:end}

\phantomsection\label{olpseg:OLP-0121-B014:start}
\begin{proof}
  د \olref[prp]{ax:lnot2} له مخې
  \LR{$\Gamma \Proves \lnot !A \lif (!A \lif \lfalse)$}۔ د
  \olref[ded]{prop:mp} په دوه ځله پلي کولو
  \LR{$\Gamma \Proves \lfalse$}۔
\end{proof}
\label{olpseg:OLP-0121-B014:end}

\phantomsection\label{olpseg:OLP-0121-B015:start}
\begin{prop}\ollabel{prop:provability-exhaustive}
  که \LR{$\Gamma \cup \{!A\}$} او \LR{$\Gamma \cup \{\lnot !A\}$} دواړه
  ناسازګار وي، نو \LR{$\Gamma$} ناسازګار دے۔
\end{prop}
\label{olpseg:OLP-0121-B015:end}

\phantomsection\label{olpseg:OLP-0121-B016:start}
\begin{proof}
تمرين۔
\end{proof}
\label{olpseg:OLP-0121-B016:end}

\phantomsection\label{olpseg:OLP-0121-B017:start}
\begin{prob}
  \olref[fol][axd][prv]{prop:provability-exhaustive} ثابت کړئ۔
\end{prob}
\label{olpseg:OLP-0121-B017:end}

\phantomsection\label{olpseg:OLP-0122-B004:start}
\olfileid{fol}{axd}{ppr}
\label{olpseg:OLP-0122-B004:end}

\phantomsection\label{olpseg:OLP-0122-B005:start}
\olsection{د اشتقاق وړتيا او د قضیې نښلوونکي}
\label{olpseg:OLP-0122-B005:end}

\phantomsection\label{olpseg:OLP-0122-B006:start}
\begin{explain}
  ثابتوو چې د بديهي استنتاج د~\LR{$\Proves$} د اشتقاق وړتيا اړيکه
  دومره پياوړې ده چې د قضيې د نښلوونکو ځينې بنسټيز حقيقتونه ثابت
  کړي؛ لکه \LR{$!A \land !B \Proves !A$} او
  \LR{$!A, !A \lif !B \Proves !B$} (مودوس پوننس)۔ دغه حقيقتونه د
  بشپړوالي د قضیې د ثبوت دپاره پکار دي۔
\end{explain}
\label{olpseg:OLP-0122-B006:end}

\phantomsection\label{olpseg:OLP-0122-B007:start}
\begin{prop}\ollabel{prop:provability-land}
  \begin{enumerate}
  \item \ollabel{prop:provability-land-left} دواړه
    \LR{$!A \land !B \Proves !A$} او \LR{$!A \land !B \Proves !B$}
  \item \ollabel{prop:provability-land-right}
    \LR{$!A, !B \Proves !A \land !B$}۔
  \end{enumerate}
\end{prop}
\label{olpseg:OLP-0122-B007:end}

\phantomsection\label{olpseg:OLP-0122-B008:start}
\begin{proof}
  \begin{enumerate}
    \item د \olref[prp]{ax:land1} او \olref[prp]{ax:land2} څخه په
      مودوس پوننس۔
      \textbf{د ژباړې د نسخې سمون (OLAX-007):} سرچينې د دواړو
      جلا عطف-ايستنو دپاره لومړني بديهي اصل ته دوه ځله حواله کړې وه؛
      دويمه دعوه د دويم عطف-ايستنې بديهي اصل کاروي۔
  \item د \olref[prp]{ax:land3} څخه د مودوس پوننس په دوه تطبيقونو۔
  \end{enumerate}
\end{proof}
\label{olpseg:OLP-0122-B008:end}

\phantomsection\label{olpseg:OLP-0122-B009:start}
\begin{prop}\ollabel{prop:provability-lor}
  \begin{enumerate}
  \item \LR{$!A \lor !B, \lnot !A, \lnot !B$} ناسازګار دے۔
  \item دواړه \LR{$!A \Proves !A \lor !B$} او
    \LR{$!B \Proves !A \lor !B$}۔
  \end{enumerate}
\end{prop}
\label{olpseg:OLP-0122-B009:end}

\phantomsection\label{olpseg:OLP-0122-B010:start}
\begin{proof}
  \begin{enumerate}
  \item له \olref[prp]{ax:lnot2} څخه
    \LR{$\Proves \lnot !A \lif (!A \lif \lfalse)$} او
    \LR{$\Proves \lnot !B \lif (!B \lif \lfalse)$} ترلاسه کوو۔ نو د
    استنتاج د قضیې له مخې \LR{$\{\lnot !A\} \Proves !A \lif \lfalse$}
    او \LR{$\{\lnot !B\} \Proves !B \lif \lfalse$} لرو۔ له
    \olref[prp]{ax:lor3} څخه
    \LR{$\{\lnot !A, \lnot !B\} \Proves (!A \lor !B) \lif \lfalse$}
    ترلاسه کوو۔ د استنتاج د قضیې له مخې
    \LR{$\{!A \lor !B, \lnot !A, \lnot !B\} \Proves \lfalse$}۔
    \textbf{د ژباړې د نسخې سمون (OLAX-008):} سرچينې دلته د نفي
    لومړني بديهي اصل ته حواله کړې وه؛ ښودل شوې انفجاري شېماوې د نفي
    دويم بديهي اصل بېلګې دي۔
  \item له \olref[prp]{ax:lor1} او \olref[prp]{ax:lor2} څخه په
    مودوس پوننس۔
  \end{enumerate}
\end{proof}
\label{olpseg:OLP-0122-B010:end}

\phantomsection\label{olpseg:OLP-0122-B011:start}
\begin{prop}\ollabel{prop:provability-lif}
  \begin{enumerate}
  \item \ollabel{prop:provability-lif-left}
    \LR{$!A, !A \lif !B \Proves !B$}۔
  \item \ollabel{prop:provability-lif-right}
    دواړه \LR{$\lnot !A \Proves !A \lif !B$} او
    \LR{$!B \Proves !A \lif !B$}۔
  \end{enumerate}
\end{prop}
\label{olpseg:OLP-0122-B011:end}

\phantomsection\label{olpseg:OLP-0122-B012:start}
\begin{proof}
  \begin{enumerate}
  \item دا اشتقاق کولے شو:
    \begin{derivation}
      1. & \LR{$!A$} & \Hyp\\
      2. & \LR{$!A \lif !B$} & \Hyp\\
      3. & \LR{$!B$} & 1, 2, \MP
    \end{derivation}
\label{olpseg:OLP-0122-B012:end}

\phantomsection\label{olpseg:OLP-0122-B013:start}
\item په ترتيب د \olref[prp]{ax:lnot2} او
    \olref[prp]{ax:lif1}، او د استنتاج د قضیې له مخې۔
  \end{enumerate}
\end{proof}
\label{olpseg:OLP-0122-B013:end}

\phantomsection\label{olpseg:OLP-0123-B005:start}
\olfileid{fol}{axd}{qpr}
\label{olpseg:OLP-0123-B005:end}

\phantomsection\label{olpseg:OLP-0123-B006:start}
\olsection{د اشتقاق وړتيا او کميت ټاکونکي}
\label{olpseg:OLP-0123-B006:end}

\phantomsection\label{olpseg:OLP-0123-B007:start}
\begin{explain}
  د بشپړوالي قضيه دا هم غواړي چې بديهي استنتاجونه د~\LR{$\Proves$} په
  اړه په دې برخه کښې ثابت شوي حقيقتونه ورکړي۔
\end{explain}
\label{olpseg:OLP-0123-B007:end}

\phantomsection\label{olpseg:OLP-0123-B008:start}
\begin{thm}
\ollabel{thm:strong-generalization} که \LR{$c$} داسې ثابت سمبول وي چې
په~\LR{$\Gamma$} يا~\LR{$!A(x)$} کښې نۀ راځي، او \LR{$\Gamma \Proves !A(c)$} وي،
نو \LR{$\Gamma \Proves \lforall[x][!A(x)]$}۔
\end{thm}
\label{olpseg:OLP-0123-B008:end}

\phantomsection\label{olpseg:OLP-0123-B009:start}
\begin{proof}
د استنتاج د قضیې له مخې \LR{$\Gamma \Proves \ltrue \lif !A(c)$}۔ ځکه
\LR{$c$} په~\LR{$\Gamma$} يا~\LR{$\top$} کښې نۀ راځي، ترلاسه کوو چې
\LR{$\Gamma \Proves \ltrue \lif \lforall[x][!A(x)]$}۔ ځکه رښتيا يو
بديهي اصل دے، د ميتا-مودوس پوننس له مخې
\LR{$\Gamma \Proves \lforall[x][!A(x)]$}۔
\textbf{د ژباړې د نسخې سمون (OLAX-009):} سرچينې وروستے ګام د
استنتاج قضیې ته يو ځل بيا منسوب کړے و؛ هغه قضيه يوازې د رښتيا فرض
ورزياتوي، حال دا چې دلته رښتيا بديهي اصل او مودوس پوننس فرض بېرته
لرې کوي۔
\end{proof}
\label{olpseg:OLP-0123-B009:end}

\phantomsection\label{olpseg:OLP-0123-B010:start}
\begin{prop}
\ollabel{prop:provability-quantifiers}
\begin{enumerate}
\item \LR{$!A(t) \Proves \lexists[x][!A(x)]$}.
\label{olpseg:OLP-0123-B010:end}

\phantomsection\label{olpseg:OLP-0123-B011:start}
\item \LR{$\lforall[x][!A(x)] \Proves !A(t)$}.
\end{enumerate}
\end{prop}
\label{olpseg:OLP-0123-B011:end}

\phantomsection\label{olpseg:OLP-0123-B012:start}
\begin{proof}
\begin{enumerate}
\item د \olref[qua]{ax:q2} او د استنتاج د قضیې له مخې۔
\item د \olref[qua]{ax:q1} او د استنتاج د قضیې له مخې۔
\end{enumerate}
\end{proof}
\label{olpseg:OLP-0123-B012:end}

\phantomsection\label{olpseg:OLP-0124-B004:start}
\olfileid{fol}{axd}{sou}
\label{olpseg:OLP-0124-B004:end}

\phantomsection\label{olpseg:OLP-0124-B005:start}
\olsection{سموالے}
\label{olpseg:OLP-0124-B005:end}

\phantomsection\label{olpseg:OLP-0124-B006:start}
\begin{explain}
د بديهي استنتاج په شان اشتقاق نظام هله \emph{سم} دے چې
هغه څيزونه اشتقاق نۀ کړي چې په حقيقت کښې نۀ صدق کېږي۔ نو
سموالے د اشتقاق نظامونو د تضمين شوې خونديتوب يو ډول خاصيت
دے۔ د دې له مخې چې کوم ثبوت-تيوريکي خاصيت تر بحث لاندې وي، مثلاً
غواړو پوه شو چې:
\begin{enumerate}
\item هر د اشتقاق وړ~\LR{$!A$} معتبر دے؛
\item که \LR{$!A$} له ځينو نورو~\LR{$\Gamma$} څخه د اشتقاق وړ وي، د هغو
  معنايي پايله هم ده؛
\item که د فارمولونه يو سټ~\LR{$\Gamma$} ناسازګار وي، د صدق وړ نۀ
  دے۔
\end{enumerate}
دا د اشتقاق نظام مهم خاصيتونه دي۔ که له دغو څخه کوم يو
سم نۀ وي، اشتقاق نظام نيمګړے دے---له حده زيات څيزونه به
اشتقاق کوي۔ له دې امله د اشتقاق نظام سموالے ثابتول
تر ټولو زيات اهميت لري۔
\end{explain}
\label{olpseg:OLP-0124-B006:end}

\phantomsection\label{olpseg:OLP-0124-B007:start}
\begin{prop}
  که \LR{$!A$} يو بديهي اصل وي، نو
  د هر جوړښت~\LR{$\Struct{M}$} او متغير-ګومارنې~\LR{$s$} دپاره \LR{$\Sat{M}{!A}[s]$}۔
\end{prop}
\label{olpseg:OLP-0124-B007:end}

\phantomsection\label{olpseg:OLP-0124-B008:start}
\begin{proof}
بايد تصديق کړو چې ټول بديهي اصول معتبر دي۔ د بېلګې
  په توګه د~\olref[qua]{ax:q1} حالت دا دے: فرض کړئ \LR{$t$} په~\LR{$!A$}
  کښې د~\LR{$x$} دپاره ازاد وي، او
  \LR{$\Sat{M}{\lforall[x][!A]}[s]$} ومنئ۔ بيا د پوره کېدو د تعريف له
  مخې، د هرې \LR{$\varAssign{s'}{s}{x}$} دپاره \LR{$\Sat{M}{!A}[s']$} هم
  دے، او په ځانګړي ډول دا هله سم دے چې
  \LR{$s'(x) = \Value{t}{M}[s]$} وي۔ د
  \olref[syn][ext]{prop:ext-formulas} له مخې
  \LR{$\Sat{M}{\Subst{!A}{t}{x}}[s]$}۔ دا ښيي چې
  \LR{$\Sat{M}{(\lforall[x][!A] \lif \Subst{!A}{t}{x})}[s]$}۔
\end{proof}
\label{olpseg:OLP-0124-B008:end}

\phantomsection\label{olpseg:OLP-0124-B009:start}
\begin{thm}[سموالے]
\ollabel{thm:soundness}
که \LR{$\Gamma \Proves !A$} وي، نو \LR{$\Gamma \Entails !A$}۔
\end{thm}
\label{olpseg:OLP-0124-B009:end}

\phantomsection\label{olpseg:OLP-0124-B010:start}
\begin{proof}
له~\LR{$\Gamma$} څخه د~\LR{$!A$} د اشتقاق پر اوږدوالي استقرا کوو۔
که داسې هېڅ ګام نۀ وي چې په استنتاج توجيه شوے وي، نو په
اشتقاق کښې ټول فارمولونه يا د بديهي اصولو بېلګې دي يا
په~\LR{$\Gamma$} کښې دي۔ د مخکښې قضیې له مخې ټول بديهي اصول
معتبر دي؛ نو که \LR{$!A$} بديهي اصل وي،
\LR{$\Gamma \Entails !A$}۔ که \LR{$!A \in \Gamma$} وي، نو په څرګنده
\LR{$\Gamma \Entails !A$}۔
\label{olpseg:OLP-0124-B010:end}

\phantomsection\label{olpseg:OLP-0124-B011:start}
که د~\LR{$!A$} د اشتقاق وروستنے ګام په مودوس پوننس توجيه شوے
وي، نو په اشتقاق کښې فارمولونه~\LR{$!B$} او~\LR{$!B \lif !A$}
شته، او د استقرا فرضيه د اشتقاق پر هغو برخو پلي کېږي چې په
دغو فارمولونه پاے ته رسېږي (ځکه په هره برخه کښې لږ تر لږه يو
په استنتاج توجيه شوے ګام کم دے)۔ نو د استقرا د فرضيې له مخې
\LR{$\Gamma \Entails !B$} او \LR{$\Gamma \Entails !B \lif !A$}۔ بيا د
\olref[syn][sem]{thm:sem-deduction}
له مخې \LR{$\Gamma \Entails !A$}۔
\label{olpseg:OLP-0124-B011:end}

\phantomsection\label{olpseg:OLP-0124-B012:start}
اوس فرض کړئ وروستنے ګام په~\QR توجيه شوے دے۔ بيا هغه
ګام د~\LR{$!C \lif \lforall[x][B(x)]$} بڼه لري، او ترې مخکښې ګام
\LR{$!C \lif !B(c)$} دے، چې \LR{$c$} په~\LR{$\Gamma$}، \LR{$!C$} يا
\LR{$\lforall[x][B(x)]$} کښې نۀ راځي۔ د استقرا د فرضيې له مخې
\LR{$\Gamma \Entails !C \lif !B(c)$}۔ د
\olref[syn][sem]{thm:sem-deduction} له مخې
\LR{$\Gamma \cup \{!C\} \Entails !B(c)$}۔
\label{olpseg:OLP-0124-B012:end}

\phantomsection\label{olpseg:OLP-0124-B013:start}
داسې يو جوړښت~\LR{$\Struct{M}$} په پام کښې ونيسئ چې
\LR{$\Sat{M}{\Gamma \cup \{!C\}}$} وي۔ بايد وښيو چې
\LR{$\Sat{M}{\lforall[x][!B(x)]}$}۔ ځکه
\LR{$\lforall[x][!B(x)]$} يوه تړلے فارمول ده، د دې معنٰی دا ده چې د
هرې متغير-ګومارنې~\LR{$s$} دپاره بايد \LR{$\Sat{M}{!B(x)}[s]$} وښيو
(\olref[syn][ass]{prop:sat-quant})۔ ځکه \LR{$\Gamma \cup \{!C\}$} ټول
له جملو جوړ دے، د هر~\LR{$!D \in \Gamma$} دپاره
\LR{$\Sat{M}{!D}[s]$} د~\olref[syn][sat]{defn:satisfaction} له مخې دے۔
\LR{$\Struct{M'}$} داسې جوړښت ونيسئ چې د~\LR{$\Struct{M}$} په شان وي، خو
\LR{$\Assign{c}{M'} = s(x)$} وي۔ ځکه \LR{$c$} په~\LR{$\Gamma$} يا~\LR{$!C$} کښې نۀ
راځي، د~\olref[syn][ext]{cor:extensionality-sent} له مخې
\LR{$\Sat{M'}{\Gamma \cup \{!C\}}$}۔ ځکه
\LR{$\Gamma \cup \{!C\} \Entails !B(c)$}، نو \LR{$\Sat{M'}{!B(c)}$}۔
\textbf{د ژباړې د نسخې سمون (OLAX-010):} سرچينې په دې پوره کېدو
ادعا کښې د بي-سي د ميتا-فارمول نښه پرېښې وه؛ په مخکښې لازميت او
ورپسې ټولو استعمالونو کښې همدا نښه موجوده ده، نو دلته بېرته
راوستل شوه۔ ځکه \LR{$!B(c)$} يوه تړلے فارمول ده، د
\olref[syn][ass]{prop:sentence-sat-true} له مخې
\LR{$\Sat{M'}{!B(c)}[s]$}۔ د~\olref[syn][ext]{prop:ext-formulas} له مخې
\LR{$\Sat{M'}{!B(x)}[s]$} هله او يوازې هله چې
\LR{$\Sat{M'}{!B(c)}[s]$} (ياد ولرئ چې \LR{$!B(c)$} هماغه
\LR{$\Subst{!B(x)}{c}{x}$} دے)۔ نو \LR{$\Sat{M'}{!B(x)}[s]$}۔ ځکه \LR{$c$} په
\LR{$!B(x)$} کښې نۀ راځي، د~\olref[syn][ext]{prop:extensionality} له
مخې \LR{$\Sat{M}{!B(x)}[s]$}۔ خو \LR{$s$} اختياري متغير-ګومارنه وه؛ نو
\LR{$\Sat{M}{\lforall[x][!B(x)]}$}۔ له دې امله
\LR{$\Gamma \cup \{!C\} \Entails \lforall[x][!B(x)]$}۔ د
\olref[syn][sem]{thm:sem-deduction} له مخې
\LR{$\Gamma \Entails !C \lif \lforall[x][!B(x)]$}۔
\label{olpseg:OLP-0124-B013:end}

\phantomsection\label{olpseg:OLP-0124-B014:start}
هغه حالت د تمرين په توګه پرېږدو چې \LR{$!A$} په~\QR{} توجيه کېږي، خو
د~\LR{$\lexists[x][!B(x)] \lif !C$} بڼه لري۔
\end{proof}
\label{olpseg:OLP-0124-B014:end}

\phantomsection\label{olpseg:OLP-0124-B015:start}
\begin{prob}
  د~\olref[fol][axd][sou]{thm:soundness} ثبوت بشپړ کړئ۔
\end{prob}
\label{olpseg:OLP-0124-B015:end}

\phantomsection\label{olpseg:OLP-0124-B016:start}
\begin{cor}
\ollabel{cor:weak-soundness}
که \LR{$\Proves !A$} وي، نو \LR{$!A$} معتبر دے۔
\end{cor}
\label{olpseg:OLP-0124-B016:end}

\phantomsection\label{olpseg:OLP-0124-B017:start}
\begin{cor}
\ollabel{cor:consistency-soundness}
که \LR{$\Gamma$} د صدق وړ وي، نو سازګار دے۔
\end{cor}
\label{olpseg:OLP-0124-B017:end}

\phantomsection\label{olpseg:OLP-0124-B018:start}
\begin{proof}
نقيض عکس ثابتوو۔ فرض کړئ \LR{$\Gamma$} سازګار نۀ دے۔ بيا
\LR{$\Gamma \Proves \lfalse$}؛ يعنې له~\LR{$\Gamma$} څخه د~\LR{$\lfalse$} يو
اشتقاق شته۔ د~\olref{thm:soundness} له مخې هر
جوړښت~\LR{$\Struct{M}$}
چې \LR{$\Gamma$} پوره کوي، بايد \LR{$\lfalse$} هم پوره کړي۔ ځکه د هر
جوړښت~\LR{$\Struct{M}$} دپاره \LR{$\Sat/{M}{\lfalse}$}، نو هېڅ
\LR{$\Struct{M}$}، \LR{$\Gamma$} نۀ شي پوره
کولے؛ يعنې \LR{$\Gamma$} د صدق وړ نۀ دے۔
\end{proof}
\label{olpseg:OLP-0124-B018:end}

\phantomsection\label{olpseg:OLP-0125-B004:start}
\olfileid{fol}{axd}{ide}
\label{olpseg:OLP-0125-B004:end}

\phantomsection\label{olpseg:OLP-0125-B005:start}
\olsection{له عينيت سره اشتقاقونه}
\label{olpseg:OLP-0125-B005:end}

\phantomsection\label{olpseg:OLP-0125-B006:start}
په اشتقاقونه کښې د~\LR{$\eq$} د شاملولو دپاره يوازې نوې بديهي
شېماوې ورزياتوو۔ د~اشتقاق او~\LR{$\Proves$} تعريف هماغسې پاتې
کېږي؛ يوازې نوي بديهي اصول هم منو۔
\label{olpseg:OLP-0125-B006:end}

\phantomsection\label{olpseg:OLP-0125-B007:start}
\begin{defn}[د عينيت بديهي اصول]
\begin{align}
\ollabel{ax:id1} & \eq[t][t], \\
\ollabel{ax:id2} & \eq[t_1][t_2] \lif (!B(t_1) \lif
  !B(t_2)),
\end{align}
د هر تړلي ترم~\LR{$t$}، \LR{$t_1$} او~\LR{$t_2$} دپاره۔
\end{defn}
\label{olpseg:OLP-0125-B007:end}

\phantomsection\label{olpseg:OLP-0125-B008:start}
\begin{prop}\ollabel{prop:sound}
  بديهي اصول~\olref{ax:id1} او~\olref{ax:id2} معتبر دي۔
\end{prop}
\label{olpseg:OLP-0125-B008:end}

\phantomsection\label{olpseg:OLP-0125-B009:start}
\begin{proof}
  تمرين۔
\end{proof}
\label{olpseg:OLP-0125-B009:end}

\phantomsection\label{olpseg:OLP-0125-B010:start}
\begin{prob}
\olref[fol][axd][ide]{prop:sound} ثابت کړئ۔
\end{prob}
\label{olpseg:OLP-0125-B010:end}

\phantomsection\label{olpseg:OLP-0125-B011:start}
\begin{prop}\ollabel{prop:iden1}
 د هر تړلي ترم~\LR{$t$} او سټ~\LR{$\Gamma$} دپاره
 \LR{$\Gamma \Proves \eq[t][t]$}۔
 \textbf{د ژباړې د نسخې سمون (OLAX-011):} سرچينې دلته «هر ترم»
 ويلي وو، خو د يووالي لومړے بديهي اصل يوازې پر تړلو ترمونو شېما
 ده؛ هماغه
 لازم قيد بېرته څرګند شو۔
\end{prop}
\label{olpseg:OLP-0125-B011:end}

\phantomsection\label{olpseg:OLP-0125-B012:start}
\begin{prop}\ollabel{prop:iden2}
  که \LR{$\Gamma \Proves !A(t_1)$} او \LR{$\Gamma \Proves
  \eq[t_1][t_2]$} وي، نو \LR{$\Gamma \Proves !A(t_2)$}، په دې شرط چې
  دواړه ښودل شوي ترمونه تړلي وي۔
  \textbf{د ژباړې د نسخې سمون (OLAX-012):} سرچينې په دې قضيه کښې
  د تړلو ترمونو شرط نۀ و بيان کړے، سره له دې چې ثبوت ئې نېغ په
  نېغه د يووالي دويم بديهي اصل کاروي او هغه شېما يوازې تړلي
  ترمونه مني۔
\end{prop}
\label{olpseg:OLP-0125-B012:end}

\phantomsection\label{olpseg:OLP-0125-B013:start}
\begin{proof}
دا فارمول
\[
(\eq[t_1][t_2] \lif (!A(t_1) \lif !A(t_2)))
\]
د~\olref{ax:id2} يوه بېلګه ده۔ پايله د~\MP په دوه پلي کولو سره
ترلاسه کېږي۔
\end{proof}
\label{olpseg:OLP-0125-B013:end}

\end{document}
